% concatenated from main.tex
\documentclass[11pt]{amsart}


\usepackage{amssymb,amsfonts}
\usepackage{hyperref}
\usepackage{mdwlist}

\usepackage{amssymb,textcomp}
%%\usepackage{hyperref}
%\usepackage{showkeys}

%
\usepackage{enumerate}

\usepackage[utf8]{inputenc}
\usepackage[T1]{fontenc}

\usepackage[mathscr]{eucal}

\usepackage{hyperref}
 \hypersetup{
     colorlinks=true,
     linktocpage=true,
     linkcolor=red,
     filecolor=blue,
     citecolor = blue,
     urlcolor=cyan,
     }

\usepackage{xcolor}


% A4 Paper
\usepackage[a4paper, twoside=false, vmargin={2cm,3cm}, includehead]{geometry}

\usepackage{comment}
% Theorems
\newtheorem{lemma}{Lemma}[section]
\newtheorem{theorem}[lemma]{Theorem}
\newtheorem*{theorem*}{Theorem}
\newtheorem{claim}[lemma]{Claim}
\newtheorem{corollary}[lemma]{Corollary}
\newtheorem{question}[lemma]{Question}
\newtheorem{proposition}[lemma]{Proposition}
\newtheorem*{proposition*}{Proposition}
\newtheorem*{MCP}{Multiplicative correspondence principle}
\newtheorem{conjecture}[lemma]{Conjecture}
\newtheorem{problem}[lemma]{Problem}
\newtheorem*{problem*}{Problem}


\theoremstyle{definition}
\newtheorem*{claim*}{Claim}
\newtheorem*{convention}{Convention}
\newtheorem*{notation}{Notation}
\newtheorem{definition}[lemma]{Definition}
\newtheorem*{equivdefinition}{Equivalent Definition}
\newtheorem{example}[lemma]{Example}
\newtheorem{examples}[lemma]{Examples}
\newtheorem*{remark}{Remark}
\newtheorem*{remarks}{Remarks}



\let\ringaccent\r
\def\aa{\ringaccent a}
\def\aa{\ringaccent a}



% Enumerations
\renewcommand{\theenumi}{\roman{enumi}}
\renewcommand{\labelenumi}{{\rm (\roman{enumi})}}
\renewcommand{\theenumii}{{\rm \alph{enumii}}}

% Blackboard letters
\DeclareMathOperator*{\oE}{\overline{\mathbb{E}}}
\DeclareMathOperator*{\E}{\mathbb{E}}
\newcommand{\lE}{\mathbb{E}^{\log}}
\newcommand{\C}{{\mathbb C}}
%\newcommand{\E}{{\mathbb E}}
\newcommand{\D}{{\mathbb D}}
\newcommand{\F}{{\mathbb F}}
\newcommand{\N}{{\mathbb N}}
\renewcommand{\P}{{\mathbb P}}
\newcommand{\Q}{{\mathbb Q}}
\newcommand{\R}{{\mathbb R}}
\renewcommand{\S}{\mathbb{S}}
\newcommand{\T}{{\mathbb T}}
\newcommand{\Z}{{\mathbb Z}}
\newcommand{\U}{{\mathbb U}}
\newcommand{\X}{{\bold X}}
\newcommand{\Y}{{\bold Y}}

% Calligraphic letters
\newcommand{\CA}{{\mathcal A}}
\newcommand{\CB}{{\mathcal B}}
\newcommand{\CC}{{\mathcal C}}
\newcommand{\CD}{{\mathcal D}}
\newcommand{\CE}{{\mathcal E}}
\newcommand{\CF}{{\mathcal F}}
\newcommand{\CG}{{\mathcal G}}
\newcommand{\CH}{{\mathcal H}}
\newcommand{\CI}{{\mathcal I}}
\newcommand{\CK}{{\mathcal K}}
\newcommand{\CM}{{\mathcal M}}
\newcommand{\CN}{{\mathcal N}}
\newcommand{\CP}{{\mathcal P}}
\newcommand{\CQ}{{\mathcal Q}}
\newcommand{\CS}{{\mathcal S}}
\newcommand{\CX}{{\mathcal X}}
\newcommand{\CY}{{\mathcal Y}}
\newcommand{\CV}{{\mathcal V}}
\newcommand{\CU}{{\mathcal U}}
\newcommand{\CW}{{\mathcal W}}
\newcommand{\CZ}{{\mathcal Z}}

\newcommand{\FC}{{\mathfrak C}}
\newcommand{\FD}{{\mathfrak D}}
\newcommand{\FH}{{\mathfrak H}}
\newcommand{\FI}{{\mathfrak I}}
\newcommand{\FL}{{\mathfrak L}}
\newcommand{\FS}{{\mathfrak S}}



% Bold letters
\renewcommand{\a}{\mathbf{a}}
\newcommand{\ba}{\mathbf{a}}
\renewcommand{\b}{{\mathbf{b}}}
%\renewcommand{\b}{b}
\newcommand{\bb}{{\mathbf{b}}}
%\newcommand{\bb}{b}
\newcommand{\bc}{{\mathbf{c}}}
\newcommand{\bm}{{\mathbf{m}}}
\newcommand{\bn}{{\mathbf{n}}}
\newcommand{\be}{{\mathbf{e}}}
\newcommand{\bg}{{\mathbf{g}}}
\newcommand{\bh}{{\mathbf{h}}}
\newcommand{\bk}{{\mathbf{k}}}
\newcommand{\bp}{{\mathbf{p}}}
\newcommand{\p}{{\mathbf{p}}}
\newcommand{\bq}{{\mathbf{q}}}
\newcommand{\q}{{\mathbf{q}}}
\renewcommand{\r}{{\mathbf{r}}}
\newcommand{\br}{{\mathbf{r}}}
\newcommand{\bs}{{\mathbf{s}}}
\newcommand{\bt}{{\mathbf{t}}}
\newcommand{\bu}{{\mathbf{u}}}
\newcommand{\bv}{{\mathbf{v}}}
\newcommand{\x}{{\mathbf{x}}}
\newcommand{\bx}{{\mathbf{x}}}
\newcommand{\by}{{\mathbf{y}}}
\newcommand{\bz}{{\mathbf{z}}}
\newcommand{\bI}{{\mathbf{I}}}
\newcommand{\bJ}{{\mathbf{J}}}
\newcommand{\bM}{{\mathbf{M}}}
\newcommand{\bN}{{\mathbf{N}}}


\newcommand{\ubh}{{\underline{\mathbf{h}}}}
\newcommand{\ubk}{{\underline{\mathbf{k}}}}
\newcommand{\ubr}{{\underline{\mathbf{r}}}}


\newcommand{\balpha}{{\boldsymbol{\alpha}}}
\newcommand{\bepsilon}{{\boldsymbol{\epsilon}}}
\newcommand{\bgamma}{{\boldsymbol{\gamma}}}
\newcommand{\bdelta}{{\boldsymbol{\delta}}}
\newcommand{\bzero}{{\boldsymbol{0}}}
\newcommand{\btheta}{{\boldsymbol{\theta}}}
\newcommand{\beeta}{{\boldsymbol{\eta}}}
\newcommand{\bbeta}{{\boldsymbol{\beta}}}

\newcommand{\ub}{{\underline{b}}}
\newcommand{\uh}{{\underline{h}}}
\newcommand{\uj}{{\underline{j}}}
\newcommand{\uk}{{\underline{k}}}
\newcommand{\ul}{{\underline{l}}}
\newcommand{\un}{{\underline{n}}}
\newcommand{\um}{{\underline{m}}}
\newcommand{\ur}{{\underline{r}}}
\newcommand{\ut}{{\underline{t}}}
\newcommand{\ux}{{\underline{x}}}
\newcommand{\uu}{{\underline{u}}}
\newcommand{\uv}{{\underline{v}}}


% Various abbreciations
\newcommand{\ve}{\varepsilon}
\newcommand{\veps}{\varepsilon}
\newcommand{\wh}{\widehat}
\newcommand{\wt}{\widetilde}
\newcommand{\e}{\varepsilon}% for the function exp
\newcommand{\one}{\mathbf{1}}
\newcommand{\eps}{\epsilon}
\newcommand{\ueps}{{\underline{\epsilon}}}
\newcommand{\ue}{{\underline{\epsilon}}}
\newcommand{\supp}{\textrm{supp}}
\newcommand{\Id}{\textrm{Id}}
%\newcommand{\and}{\textrm{and}}
%\newcommand{\or}{\textrm{or}}

\renewcommand{\angle}[1]{\mathopen{}\left\langle #1\mathclose{}\right\rangle}




\newcommand{\ZN}{\Z_{\widetilde N}}
\newcommand{\tN}{{\widetilde N}}
\newcommand{\tM}{{\widetilde M}}
\newcommand{\CCM}{{\CM_1^c}}
\newcommand{\tZN}{\Z_\tN}
\newcommand{\norm}[1]{\left\Vert #1\right\Vert}
\newcommand{\nnorm}[1]{\lvert\!|\!| #1|\!|\!\rvert}
\newcommand{\Bignnorm}[1]{\Big|\!\Big|\!\Big| #1\Big|\!\Big|\!\Big|}
\newcommand{\lip}{{\text{\rm Lip}}}
\newcommand{\inv}{^{-1}}
\newcommand{\st}{{\text{\rm st}}}
\newcommand{\1}{{\mathbf{1}}}

\newcommand{\er}{{\text{\rm er}}}

\DeclareMathOperator{\sgn}{sgn}
\DeclareMathOperator{\spec}{Spec}
\DeclareMathOperator{\Spec}{Spec}
\DeclareMathOperator{\id}{id}
\DeclareMathOperator{\Span}{Span}
\DeclareMathOperator{\ExtSpan}{ExtSpan}
\DeclareMathOperator{\Extspan}{ExtSpan}
\DeclareMathOperator{\extspan}{ExtSpan}
\DeclareMathOperator{\extSpan}{ExtSpan}
%\DeclareMathOperator{\span}{Span}
\DeclareMathOperator{\Nil}{Nil}
\newcommand{\dmult}{d_\textrm{mult}}
\DeclareMathOperator{\reel}{Re}
\renewcommand{\Re}{\reel}
\DeclareMathOperator{\Imag}{Im}
\renewcommand{\Im}{\Imag}
\DeclareMathOperator{\nonpol}{nonpol}
\DeclareMathOperator{\pol}{pol}
\DeclareMathOperator{\poly}{poly}
\DeclareMathOperator{\fracdeg}{frac\; deg}
\newcommand{\Krat}{{\CK_{\text{\rm rat}}}}


\newcommand{\abs}[1]{\mathopen{}\left| #1\mathclose{}\right|}
\newcommand{\bigabs}[1]{\bigl| #1 \bigr|}
\newcommand{\Bigabs}[1]{\Bigl| #1 \Bigr|}
\newcommand{\brac}[1]{\mathopen{}\left( #1 \mathclose{}\right)}
\newcommand{\bigbrac}[1]{\bigl( #1 \bigr)}
\newcommand{\Bigbrac}[1]{\Bigl( #1 \Bigr)}
\newcommand{\floor}[1]{{\left \lfloor #1 \right \rfloor}}
\newcommand{\sfloor}[1]{{\lfloor #1 \rfloor}}
\newcommand{\Bigfloor}[1]{{\Big\lfloor #1 \Big\rfloor}}
\newcommand{\ceil}[1]{{\left \lceil #1 \right \rceil}}
\newcommand{\rem}[1]{\left \{ #1 \right \}}
\newcommand{\Bigrem}[1]{{\Big\{ #1 \Big\}}}
\newcommand{\srem}[1]{ \{ #1 \}}


\title[]{Resolving the joint ergodicity problem\\ for Hardy sequences}

\author{Sebasti\'an Donoso, Andreas Koutsogiannis, Borys Kuca,\\ Wenbo Sun, and Konstantinos Tsinas}

\address[Sebasti\'an Donoso]{Departamento de Ingenier\'{\i}a
	Matem\'atica and Centro de Modelamiento Ma\-te\-m\'a\-ti\-co, Universidad de Chile and IRL-CNRS 2807, Beauchef 851, Santiago,
	Chile.} 
\email{sdonosof@uchile.cl}

\address[Andreas Koutsogiannis]{
Department of Mathematics, Aristotle University of Thessaloniki, Thessaloniki 54124, Greece}
\email{akoutsogiannis@math.auth.gr}

\address[Borys Kuca]{Faculty of Mathematics and Computer Science, Jagiellonian University, 30-348 Krak\'ow, Poland}
\email{borys.kuca@uj.edu.pl}


\address[Wenbo Sun]{Department of Mathematics, Virginia Tech, 225 Stanger Street, Blacksburg, VA, 24061, USA}
\email{swenbo@vt.edu}

\address[Konstantinos Tsinas]{Institute of Mathematics, Ecole Polytechnique F\'{e}d\'{e}rale de Lausanne (EPFL), Lausanne 1015,
Switzerland}
\email{konstantinos.tsinas@epfl.ch}

\thanks{The first author was partially funded by Centro de Modelamiento Matemático (CMM) FB210005, BASAL funds for centers of excellence from ANID-Chile and ANID/Fondecyt/1241346. For the second author, the research project is implemented in the framework of H.F.R.I call ``3rd Call for H.F.R.I.’s
Research Projects to Support Faculty Members \& Researchers'' (H.F.R.I. Project Number: 24979).  The third author was supported by the NCN Polonez Bis 3 grant No. 2022/47/P/ST1/00854 (H2020 MSCA GA No. 945339). The fourth author was partially supported by the NSF Grant DMS-2247331. The fifth author was supported by the Swiss National Science Foundation grant TMSGI2-211214. For the purpose of Open Access, the authors have applied a CC-BY public copyright licence to any Author Accepted Manuscript (AAM) version arising from this submission.}

\subjclass[2020]{Primary: 37A44; Secondary: 11B30, 28D05}

\keywords{Ergodic averages, joint ergodicity, Host-Kra seminorms, Hardy fields.} 












\begin{document}

\begin{abstract}
    The joint ergodicity classification problem aims to characterize those sequences which are jointly ergodic along an arbitrary dynamical system if and only if they satisfy two natural, simpler-to-verify conditions on this system. These two conditions, dubbed the difference and product ergodicity conditions, naturally arise from Berend and Bergelson's pioneering work on joint ergodicity. Elaborating on our earlier work, we investigate this problem for Hardy sequences of polynomial growth, this time without making any independence assumptions on the sequences. Our main result establishes the ``difficult’’ direction of the problem: if a Hardy family satisfies the difference and product ergodicity conditions on a given system, then it is jointly ergodic for this system. We also find that, surprisingly, the converse fails for certain pathological families of Hardy sequences, even though it holds for all ``reasonable’’ Hardy families. We conclude by suggesting potential fixes to the statement of this problem. {New ideas of independent interest developed in this paper include the structure theory of a family of factors generalizing Host-Kra and box factors; a strengthening of Tao-Ziegler's concatenation results; and the most robust extension of a seminorm smoothing argument.}
\end{abstract}

\maketitle

\tableofcontents

\section{Introduction}


In 1977, Furstenberg gave his famous dynamical proof of Szemer\'edi's theorem on arithmetic progressions, initiating a highly successful program of applying ergodic tools to problems in combinatorics and number theory \cite{Fu77}. As part of his proof, he showed that whenever $(X, \CX, \mu, T)$ is a weakly mixing measure-preserving dynamical system, the averages
\begin{align*}
    \E_{n\in[N]}T^n f_1 \cdots T^{\ell n}f_\ell
\end{align*}
(where $\E_{n\in[N]}:=\frac{1}{N}\sum\limits_{n=1}^N$ {and $[N]:=\{1,\ldots,N\}$ for $N\in \N$}) converge in $L^2(\mu)$ to the product of integrals $\int f_1\; d\mu \cdots \int f_\ell\; d\mu$ for all functions $f_1, \ldots, f_\ell\in L^\infty(\mu)$. In doing so, he proved the first result on the \textit{joint ergodicity} of multiple ergodic averages {after the mean ergodic theorem}. 


\begin{definition}[Joint ergodicity]\label{D: joint ergodicity}
    Let $a_1, \ldots, a_\ell:\N\to\R$ be sequences and $(X, \CX, \mu, T_1,$ $ \ldots, T_\ell)$ be a system.\footnote{By which we mean commuting invertible measure preserving transformations $T_1, \ldots, T_\ell$ acting on a Lebesgue probability space $(X, \CX, \mu)$.} We say that $a_1, \ldots, a_\ell$ are \textit{jointly ergodic} for $(X, \CX, \mu, T_1, \ldots, T_\ell)$ (alternatively, the tuple $(T_1^{\floor{a_1(n)}}, \ldots, T_\ell^{\floor{a_\ell(n)}})_n$ is \textit{jointly ergodic} on $(X, \CX, \mu)$) if
\begin{align}\label{E: joint ergodicity}
        \lim_{N\to\infty}\norm{\E_{n\in[N]}\prod_{j=1}^\ell T_j^{\floor{a_{j}(n)}}f_j - \prod_{j=1}^\ell \int f_j\, d\mu}_{L^2(\mu)} = 0
    \end{align}
    holds for all $f_1, \ldots, f_\ell\in L^\infty(\mu)$. {For $\ell=1,$ we say that $a_1$ is \emph{ergodic} for $(X, \CX, \mu, T_1)$ (and respectively we call $(T_1^{\floor{a_1(n)}})_n$ \emph{ergodic}).}
\end{definition}

In particular, $(T^n)_n$ is ergodic if and only if $T$ is ergodic.


A few years after Furstenberg's work, Bergelson and Berend inquired about the necessary and sufficient conditions for $(T_1^n, \ldots, T_\ell^n)_n$ to be jointly ergodic, proving the following result.
\begin{theorem}[{\cite{BB84}}]
    Let $(X, \CX, \mu, T_1, \ldots, T_\ell)$ be a system. Then $(T_1^n, \ldots, T_\ell^n)_n$ is jointly ergodic if and only if the following two conditions hold:
    \begin{enumerate}
    \item $T_i T_j^{-1}$ is ergodic on $(X, \CX, \mu)$ for all distinct $i, j\in[\ell]$;
    \item $T_1\times \cdots \times T_\ell$ is ergodic on the product system $(X^\ell, \CX^{\otimes \ell}, \mu^\ell)$.
    \end{enumerate}
\end{theorem}
It is interesting to understand the extent to which the Bergelson-Berend conditions for joint ergodicity can be extended to more general multiple ergodic averages, leading to the following definitions. 
\begin{definition}\label{D: good for joint ergodicity}
    Let $a_1, \ldots, a_\ell:\N\to\R$ be sequences and $(X, \CX, \mu, T_1, \ldots, T_\ell)$ be a system. We say that on the system, the sequences $a_1, \ldots, a_\ell$ satisfy
    \begin{enumerate}
        \item\label{i: difference} the \emph{difference ergodicity condition} if the sequence $(T_i^{\floor{a_i(n)}}T_j^{-\floor{a_j(n)}})_n$ is ergodic\footnote{A sequence of transformations $(T_n)_n$ on a Lebesgue probability space $(X, \CX, \mu)$ is called \textit{ergodic} if
\[  \lim_{N\to\infty}\norm{\E_{n\in[N]} T_n f - \int f\, d\mu}_{L^2(\mu)} = 0    \]     
    for all $f\in L^\infty(\mu)$.} on $(X, \CX, \mu)$ for all distinct $i, j\in[\ell]$; and
        \item\label{i: product} the \emph{product ergodicity condition} if $(T_1^{\floor{a_1(n)}}\times \cdots \times T_\ell^{\floor{a_\ell(n)}})_n$ is ergodic on the product system $(X^\ell, \CX^{\otimes \ell}, \mu^\ell)$.
    \end{enumerate}
\end{definition}

We then have the following natural problem whose variants were previously asked in \cite{BMR20, DKS22, DKS23}.
    \begin{problem}[Joint ergodicity classification problem]\label{Pr: joint ergodicity problem}
    Classify the sequences $a_1, \ldots, a_\ell:\N\to\R$ with the following property: the sequences are jointly ergodic for a system $(X, \CX, \mu, T_1, \ldots, T_\ell)$ if and only if they simultaneously satisfy the difference and product ergodicity conditions on the system. 
    \end{problem}

Joint ergodicity has been a subject of intense scrutiny in the last decades \cite{BB84, BB86, Ber87, BeHK09, BLS16, BMR20, BS22, BS23, BFM22, CFH11, DFKS22, DKKST24, DKS22, DKS23, DKS24,  Fr10, Fr12, Fr22, Fr21, FrKr05, FrKu22b, FrKu22c, FrKu22a, KK19, Kouts21,  KoSu23, KoTs23, Ts22}. Broadly speaking, the aforementioned works circulate around three themes:
\begin{enumerate}
    \item establishing joint ergodicity and its variants for broad classes of sequences and systems;
    \item deriving easier-to-verify criteria for joint ergodicity;
    \item investigating Problem \ref{Pr: joint ergodicity problem} for various classes of sequences.
\end{enumerate}

So far, Problem \ref{Pr: joint ergodicity problem} has been examined for integer polynomials, Hardy sequences, and generalized linear functions, and the equivalence postulated by Problem \ref{Pr: joint ergodicity problem} has been verified affirmatively for the following classes of sequences:
\begin{itemize}
    \item linear functions \cite{BB84};
    \item generalized linear functions \cite{BLS16};
    \item a single integer polynomial $a_1 = \cdots = a_\ell$ \cite{DKS22} (see also \cite{DFKS22} for a slight generalization);
    \item pairwise independent integer polynomials \cite{FrKu22a};
    \item arbitrary integer polynomials \cite{FrKu22b};
    \item a single Hardy sequence $a_1 = \cdots = a_\ell$ of polynomial growth that stays logarithmically away from $\R[x]$ \cite{DKS23};
    \item pairwise independent Hardy sequence of polynomial growth \cite{DKKST24}.
\end{itemize}
Of this rich variety of results, we particularly emphasize the resolution of Problem \ref{Pr: joint ergodicity problem} for integer polynomials \cite{FrKu22b} and pairwise independent Hardy sequences of polynomial growth \cite{DKKST24}. By identifying broad classes of sequences for which Problem \ref{Pr: joint ergodicity problem} holds, the aforementioned works have given substantial credibility to why one should think of the difference and product ergodicity conditions as a right set of conditions guaranteeing joint ergodicity. 

So far, the difficult part of Problem \ref{Pr: joint ergodicity problem} has typically been to establish the converse direction, i.e. the sufficiency of the difference and product ergodicity conditions for joint ergodicity. The main result of this paper shows that the difficult direction in Problem \ref{Pr: joint ergodicity problem} holds for all Hardy sequences of polynomial growth (whose collection is denoted below by $\CH$; see Section \ref{SS: Hardy} for the formal definition). 
\begin{theorem}\label{T: main theorem}
    Suppose that $a_{1},\dots,a_{\ell}\in\mathcal{H}$ satisfy the difference and product ergodicity conditions on a system $(X, \CX, \mu, T_1, \ldots, T_\ell)$.
Then they are jointly ergodic for the system.
     \end{theorem}
We also establish the following partial converse.

\begin{definition}
    For $a\in\mathcal{H}$, we say that $a$ is an \emph{almost rational %(resp. real) 
    polynomial} if $a-p$ converges for some $p$ in $\Q[x]$. %(resp. $\R[x]$). 
    We say that $a$ \emph{stays logarithmically away from $\Q[x]$} (\emph{resp. $\R[x]$}) if $a-p\succ \log$ for all $p$ in $\Q[x]$ (resp. $\R[x]$).
\end{definition}

\begin{theorem}\label{T: easy direction}
        Suppose that $a_{1},\dots,a_{\ell}\in\mathcal{H}$ are jointly ergodic for a system $(X, \CX, \mu, T_1, $ $ \ldots, T_\ell)$. Then the following holds:
        \begin{enumerate}
            \item The sequences satisfy the product ergodicity condition on the system.
            \item Suppose additionally that for every $i\neq j$, $c_i, c_j\in\R$, the function $c_i a_i - c_j a_j$ is either an almost rational polynomial  or stays logarithmically away from $\Q[x]$. Then the sequences satisfy the difference ergodicity condition on the system.
        \end{enumerate}
\end{theorem}



As a corollary, we provide an affirmative answer to Problem \ref{Pr: joint ergodicity problem} for a large class of Hardy sequences.
\begin{corollary}\label{C:np}
    Let $a_1, \ldots, a_\ell\in\CH$ be Hardy sequences with the property that for every $i\neq j$ and $c_i, c_j\in\R$ the function $c_i a_i - c_j a_j$ is either an almost rational polynomial or it stays logarithmically away from $\Q[x]$. Then the sequences are jointly ergodic for a system $(X, \CX, \mu, T_1, \ldots, T_\ell)$ if and only  if they satisfy the product and difference ergodicity conditions on the system. 
\end{corollary}






In particular, we resolve \cite[Conjecture 6.1]{BMR20}, i.e. Problem \ref{Pr: joint ergodicity problem} {for general real and fractional polynomials}.
\begin{corollary}
    Let $\CF$ be the class of functions of the form $$a(x) = \beta_d x^{c_d} + \cdots + \beta_1 x^{c_1}$$ for $\beta_i\in\R$ and $c_i\geq 0$. Then any sequences $a_1, \ldots, a_\ell\in \CF$ are jointly ergodic for a system $(X, \CX, \mu, T_1, \ldots, T_\ell)$ if and only  if they satisfy the product and difference ergodicity conditions on the system. 
\end{corollary}

We also show that, surprisingly, joint ergodicity does not always imply the difference ergodicity condition, yielding a negative answer to Problem \ref{Pr: joint ergodicity problem} in general. The failure of the (supposedly) easy direction in Problem \ref{Pr: joint ergodicity problem} is instantiated by the following example. 
\begin{theorem}\label{T: counterexample}
    Let $(X, \CX, \mu, T)$ be a system. Then $(T^n, T^{n+\floor{\log_2 n}})_n$ is jointly ergodic if and only if $T$ is (strongly) mixing. By contrast, $(T^{\floor{\log_2 n}})_n$ is ergodic if and only the system is conjugate to a one-point system, or equivalently $\mu$ is concentrated on a single point. Consequently, if $T$ is mixing but not conjugate to a one-point system, then $(T^n, T^{n+\floor{\log_2 n}})_n$ is jointly ergodic whereas $(T^{\floor{\log_2 n}})_n$ is not ergodic.
\end{theorem}

To conclude the discussion, we remark that in this work we essentially resolve Problem~\ref{Pr: joint ergodicity problem} for polynomial and Hardy sequences of polynomial growth, i.e. classes for which we have a reduction-of-complexity argument based on van der Corput's inequality. Up until now, we have had no results at all for general real polynomials and several different Hardy sequences (the case when all of them are the same has been partially dealt with in \cite{DKS22, DKS23}).

Lastly, in Section \ref{section: repaired conjecture}, we propose several ways of amending the statement of Problem~\ref{Pr: joint ergodicity problem} for Hardy sequences that would avoid the issues posed by the counterexample for Theorem \ref{T: counterexample}. While we obtain some preliminary results towards these modified  versions of Problem \ref{Pr: joint ergodicity problem}, more research is needed in order to see if they could serve as a good fix to the statement of Problem~\ref{Pr: joint ergodicity problem}.

\subsection{Outline of the argument}
In \cite[Theorem 2.5]{FrKu22a}, Frantzikinakis and the third author derived the following simpler-to-verify criteria for joint ergodicity. 
\begin{theorem}[Joint ergodicity criteria]\label{T: joint ergodicity criteria}
Let $a_1, \ldots, a_\ell:\N\to\R$ be sequences and $(X, \CX, \mu, T_1, \ldots, T_\ell)$ be a system with each $T_1, \ldots, T_\ell$ ergodic. Then $a_1, \ldots, a_\ell$ are jointly ergodic for  $(X, \CX, \mu, T_1, \ldots, T_\ell)$ if and only if the following two conditions hold:
    \begin{enumerate}
        \item $a_1, \ldots, a_\ell$ are \emph{good for seminorm control} for the system in that there exists $s\in\N$ such that for all 1-bounded $f_1, \ldots, f_\ell\in L^\infty(\mu)$, we have
    \begin{align*}%\label{E: convergence to 0}
        \lim_{N\to\infty}\norm{\E_{n\in[N]} T_1^{\floor{a_1(n)}}f_1 \cdots T_\ell^{\floor{a_\ell(n)}}f_\ell}_{L^2(\mu)} = 0
    \end{align*}
    whenever $\nnorm{f_j}_{s, T_j} = 0$ for some $j\in[\ell]$;
        \item $a_1, \ldots, a_\ell$ are \emph{good for equidistribution} for the system
        %if for all eigenfunctions $\lambda_1, \ldots, \lambda_\ell$ of $T_1, \ldots, T_\ell$ respectively
        if for all eigenvalues $\lambda_1, \ldots, \lambda_\ell$ of $T_1, \ldots, T_\ell$ respectively, not all of which are 0, we have
        \begin{align}\label{E: good for equidistribution}
            \lim_{N\to\infty} \E_{n\in[N]}e(\lambda_1 \floor{a_1(n)} + \cdots + \lambda_\ell \floor{a_\ell(n)}) = 0.
        \end{align}
    \end{enumerate}
\end{theorem}
While the equidistribution properties of Hardy sequences are well-understood thanks to Boshernitzan \cite{Boshernitzan-equidistribution}, establishing seminorm control over multiple ergodic averages of interest is a highly challenging problem. 
Up to several years ago, few seminorm estimates were known for multiple ergodic averages considered in this paper. Then came two breakthroughs which dramatically changed the situation. On one hand, the fifth author significantly refined existing arguments and obtained seminorm estimates for essentially all possible (unweighted) multiple ergodic averages along Hardy sequences of a single transformation \cite{Ts22}.\footnote{In another interesting work that came out around the same time, Bergelson, Moreira, and Richter showed that one can get seminorm estimates for another large class of averages if one is willing to introduce a suitable weight \cite{BMR20}.} On the other hand, Frantzikinakis and the third author introduced a seminorm smoothing technique \cite{FrKu22b, FrKu22a} that gave seminorm estimates for multiple ergodic averages of \textit{commuting} transformations along integer polynomials. As a starting point, they used earlier box seminorm estimates obtained by the first, second and fourth author jointly with Ferr\'e-Moragues \cite{DFKS22}, based on a concatenation result of Tao and Ziegler \cite{TZ16}. 

Generalizing all the aforementioned arguments, we have recently shown \cite{DKKST24} that averages of commuting transformations along Hardy sequences can be controlled by certain seminorms that generalize the now classical Host-Kra and box seminorms. We then used a variant of the seminorm smoothing argument to obtain (quantitative) Host-Kra seminorm estimates for averages of commuting transformations along any family of pairwise independent Hardy sequences; from this we derived a number of corollaries of joint ergodicity, recurrence, and $L^2(\mu)$ convergence. In this work, we drop the assumption of pairwise independence and show that under some necessary conditions which are good enough for the purpose of proving Theorem \ref{T: main theorem}, we can get Host-Kra seminorm control for averages along not necessarily pairwise independent Hardy sequences.
\begin{theorem}\label{T: seminorm control}
    Let $a_1, \ldots, a_\ell\in\CH$ and $(X, \CX, \mu, T_1, \ldots, T_\ell)$ be a system with $T_1, \ldots, T_\ell$ ergodic. If the sequences satisfy the difference ergodicity condition, then they are good for seminorm control for the system. 
\end{theorem}
The proof of Theorem \ref{T: seminorm control} occupies a bulk of this paper. Starting with the generalized box seminorm estimates from \cite{DKKST24}, we use a new variant of the seminorm smoothing argument to obtain another collection of seminorm estimates for the averages in question. For instance, the estimates from \cite{DKKST24} allow us to control the limit of an average 
\begin{align}\label{E: average example intro}
    \E_{n\in[N]} T_1^{\sfloor{n^{3/2}}}f_1\cdot T_2^{\sfloor{n^{3/2}+\log n}}f_2\cdot T_3^{\sfloor{n^{3/2}+n\log n}}f_3\cdot T_4^{\sfloor{n^{3/2}+n\log n}}f_4
\end{align}
by a box seminorm of $f_1$ involving many copies of $T_1, T_1T_2\inv, T_1T_3\inv, T_1T_4\inv$. By contrast, our smoothing argument gives us control over (the limit of) \eqref{E: average example intro} by a seminorm of $f_1$ along many copies of $T_1, T_1T_2\inv, T_3T_4\inv$.
A priori, there is no reason why the new seminorm 
would be easier to work with than the original one. It is at this point where the difference ergodicity condition kicks in. It forces $T_1T_2\inv, T_3T_4\inv$ to be ergodic, hence the intermediate seminorm involving $T_1, T_1T_2\inv, T_3T_4\inv$ can be replaced by a proper Host-Kra seminorm along $T_1$. 


To carry out the argument above, we introduce a number of tools of independent interest. In Section \ref{SS: factors}, we introduce factors related to the generalized box seminorms in much the same manner in which Host-Kra factors are related to Host-Kra seminorms (Theorem \ref{T: relative concatenation}). We also obtain a weak structure theorem for these factors, showing that they can be generated via dual functions weighted by nilsequences of appropriate step. In Section \ref{S: concatenation}, we derive a Tao-Ziegler-type concatenation result for such factors. In fact, our concatenation strengthens Tao and Ziegler's results \cite{TZ16} in the following sense. While Tao and Ziegler show for example that 
\begin{align*}
    \mathcal{Z}_{H_{1},\dots,H_{d}}\cap \mathcal{Z}_{H'_{1},\dots,H'_{d'}}\subseteq \mathcal{Z}_{(H_{i}+H'_{i'})_{i\in[d], i'\in[d']}}
\end{align*}
for any finitely generated subgroups $H_1, \ldots, H_d, H'_1, \ldots, H'_{d'}\subseteq\Z^\ell$, we show that
\begin{align*}
    \mathcal{Z}_{G_{1},\dots,G_s,H_{1},\dots,H_{d}}\cap \mathcal{Z}_{G_{1},\dots,G_s,H'_{1},\dots,H'_{d'}}\subseteq \mathcal{Z}_{G_{1},\dots,G_{s},(H_{i}+H'_{i'})_{i\in[d], i'\in[d']}}
\end{align*}
for any finitely generated subgroups $G_1, \ldots, G_s, H_1, \ldots, H_d, H'_1, \ldots, H'_{d'}\subseteq\Z^\ell$. Because we are able to fix the same set of subgroups $G_1, \ldots, G_s$ in all the factors that we compare, we refer to our concatenation result as \textit{relative concatenation}.

Relative concatenation plays a key role in the seminorm smoothing argument carried out in Section \ref{S: smoothing}. The version of the smoothing argument developed in this paper (Theorem \ref{T: new estimates general}) is more robust than previous variations from \cite{DKKST24, FrKu22b, FrKu22a}. While performing smoothing for integer polynomials \cite{FrKu22b}, Frantzikinakis and the third author had to make some a priori ergodicity assumptions on the system, and then they had to deal with intermediate averages in which functions are invariant under various transformations. Not only did keeping track of all this information throughout an inductive argument require an unwieldy formalism; even worse, the argument could not be translated to the language of generalized box seminorms needed in the setting of Hardy sequences. Hence the need for a more flexible smoothing argument that we develop in this paper. Our argument is conceptually simpler than the one from \cite{FrKu22b} in that we make no a priori assumptions on the ergodicity of the transformations, and we have to track much less information throughout the proof than in \cite{FrKu22b}. The key innovation that makes it possible is the use of relative concatenation in the smoothing argument. On top of \textit{ping} and \textit{pong}, the two traditional steps of the smoothing argument, we add a third step, where we concatenate the original seminorm controlling our average (along $T_1, T_1T_2\inv, T_1T_3\inv, T_1T_4\inv$ in \eqref{E: average example intro}) with an auxiliary one obtained in the \textit{pong} step to get an upgraded seminorm (along $T_1, T_1T_2\inv, T_3T_4\inv$ in \eqref{E: average example intro}) controlling our average. The details of this innovation are explained with an explicit example in Section \ref{SS: concatenation example}.

Sections \ref{S: seminorm comparison I} and \ref{S: seminorm comparison II} are then dedicated to extracting information from the difference ergodicity assumption. The ergodicity of $(T_i^{\floor{a_i(n)}}T_j^{-\floor{a_j(n)}})_n$ imposes strict conditions on the spectral measures of functions in $L^2(\mu)$ since by the Herglotz-Bochner theorem, a spectral measure may not have a point mass on any $(t,s)$ for which the exponential sum
\begin{align}\label{E: exp sums intro}
    \E_{n\in[N]}e(t\floor{a_i(n)}-s\floor{a_j(n)})
\end{align}
does not vanish as $N\to\infty$. The bulk of these two sections is then dedicated to a case-by-case study of exponential sums \eqref{E: exp sums intro}. While the details of the computations performed in these sections are not very illuminating, we manage to extract enough information to upgrade seminorm control obtained in smoothing to a Host-Kra seminorm control over our average. Section \ref{S: proofs of main theorems} then pieces everything together and derives the theorems stated in the introduction.

\subsection{Basic notation}\label{SS:notation}
We start with presenting basic notation used throughout the paper.

The symbols $\C, \R, \R_+, \Q, \Z, \N_0, \N$ respectively stand for the sets of complex numbers, reals, positive reals, rationals, integers, nonnegative integers, and positive integers. With $\T$, we denote the one dimensional torus, and we often identify it with $\R/\Z$ or  with $[0,1)$. For a ring $R$, we denote the collection of polynomials over $R$ in variables $x_1, \ldots, x_k$ via $R[x_1, \ldots, x_k]$. If $E\subseteq G$ is a subset of a group $G$, then $\langle E \rangle$ is the subgroup generated by $E$.

We frequently use the interval notation to mean either a real interval or its restriction to $\Z$; the meaning would typically be clear from the context. If we restrict to $\Z$, then we write $[M,N] = \{M, \ldots, N\}$ and $[N]=\{1, \ldots, |N|\}$ for integers $M\leq N$.


For a finite set $E,$ we define the average of a sequence $a:E\to \C$ as 
\[\E_{n\in E}a(n) :=\frac{1}{|E|}\sum_{n\in E}a(n).\] We will mainly use finite averages along subsets of integers of the form $[N]^s$ or $[\pm N]^s$, also known as \textit{Ces\`aro averages}.

We let $\CC z := \overline{z}$ be the complex conjugate of $z\in \C$.

For a set $E$, we define its indicator function by $1_E$. Whenever convenient, we also use the notation $1(E)$ for an event $E$.


Typically, we use underlined symbols $\uh$ to denote elements of $\R^s$ (which we think as tuples of parameters), and we employ the bold notation $\bx$ to denote elements of $\R^\ell$ (which we think as points in a space); sometimes we also use bold notation to denote $s$-tuples of points in $\R^\ell$, e.g., $\bm=(\bm_1, \ldots, \bm_s)\in(\R^s)^\ell$. We denote the integer and fractional parts of $\bx$ via
\begin{align*}
    \floor{\bx} = (\floor{x_1}, \ldots, \floor{x_\ell})\quad \textrm{and}\quad \rem{\bx} =(\rem{x_1}, \ldots, \rem{x_\ell}).
\end{align*}

We often write $\ueps\in\{0,1\}^s$ for a vector of 0s and 1s of length $s$. For $\ueps\in\{0,1\}^s$ and $\uh, \uh'\in\R^s$, we set $\ueps\cdot \uh:=\eps_1 h_1+\cdots+ \eps_s h_s$, $\abs{\uh} := |h_1|+\cdots+|h_s|$ and
\begin{gather}\label{E: h^eps}
    h_j^{\eps_j} =\begin{cases}
        h_j,\; &\eps_j =0\\
        h_j',\; &\eps_j = 1
    \end{cases}
    \quad \textrm{for}\quad j\in[s].
\end{gather}

Given a system $(X, \CX, \mu, T_1, \ldots, T_\ell)$ and $\bm=(m_{1},\dots,m_{\ell})\in\R^\ell$, we define $$T^{\floor{\bm}}:=T_1^{\floor{m_1}}\cdots T_\ell^{\floor{m_\ell}}.$$
For $j\in[\ell]$, we set $\be_j$ to be the unit vector in $\R^\ell$ in the $j$-th direction, and we let $\be_0 = \mathbf{0}$, so that $T^{\be_j} = T_j$ for $j\in[\ell]$ and $T^{\be_0}$ is the identity transformation. 

For $M>0$, we say that a function $f\in L^\infty(\mu)$ is \textit{$M$-bounded} if $\norm{f}_{L^\infty(\mu)}\leq M$.

For a $\sigma$-algebra $\CA\subseteq\CX$ and $f\in L^1(\mu)$, we let $\E(f|\CA)$ be the conditional expectation of $f$ with respect to $\CA$. By abuse of notation, we can identify the $\sigma$-algebra $\CA$ with the algebra of functions $L^2(\CA)$. One algebra that we often work with is the \textit{invariant factor}
$$\CI(T) = \{A\in \CX:\; T\inv A = A\}$$
of a transformation $T$.


Given two functions $a,b:(x_0,\infty)\to\mathbb{R}$ for some $x_0\in\R$ we write
\begin{enumerate}
    \item $a\ll b$ or $a = O(b)$ if there exists $C>0$ and $x_1\geq x_0$ such that $|a(x)| \leq C |b(x)|$ for all $x\geq x_1$; 
    \item $a \asymp b$ if $a\ll b$ and $b\ll a$;
    \item $a \prec b$ or $a = o(b)$  if $\lim\limits_{x\to\infty}\frac{a(x)}{b(x)} = 0$.
    %\item $a \lll b$ if there exists $\delta > 0$ such that $|a(x)|x^\delta \prec |b(x)|$.
\end{enumerate}
If the absolute constant in (i) depends on some parameters $x_1, \ldots, x_n$, then we write $a\ll_{x_1, \ldots, x_n} b$, $a = O_{x_1, \ldots, x_n}(b)$, etc. If in (ii) we want to emphasize what parameter we are taking to infinity, we can also write $a = o_{x\to\infty; x_1, \ldots, x_n} (b)$ to emphasize that $x_1, \ldots, x_n$ are kept fixed while $x$ goes to infinity.





For a finitely generated abelian group $G$ and a bounded function $\psi:G\to\C$, we define
\begin{align*}
    \E_{\bm\in G}\psi(\bm) = \lim_{N\to\infty}\E_{\bm\in I_N}\psi(\bm)
\end{align*}
to be the limit along any F\o lner sequence $(I_N)_N$ on $G$ (if such a limit exists), and if $\psi$ is nonnegative, then we define
\begin{align*}
    \oE_{\bm\in G}\psi(\bm) = \sup_{(I_N)_N}\limsup_{N\to\infty}\E_{\bm\in I_N} \psi(\bm),
\end{align*}
where the supremum is taken along any F\o lner sequence on $G$. Likewise, if $(X,\CX,\mu)$ is a probability space and $f_\bm\in L^\infty(\mu)$ for all $\bm\in G$, then we define $\E_{\bm\in G}f_\bm$ to be the $L^2(\mu)$ limit (whenever it exists).



\subsection{Basic properties of Hardy sequences}\label{SS: Hardy}

Let $B$ be the collection of germs at $\infty$ of real valued functions defined on some halfline $(x_0,\infty)$ for $x_0\geq 0$ (thus functions that eventually agree are considered to be the same).  A \emph{Hardy field} is a subfield of the ring $(B, +, \cdot)$ that is closed under differentiation.
 
 For various technical reasons, we will consider Hardy fields $\mathcal{H}_0$ which are closed under composition and compositional inversion of functions, when defined. We say that $a\in\CH$ has \emph{polynomial growth} if it satisfies $a(x)\ll x^d$ for some $d>0$, and we then let
\begin{align*}
    \CH = \{a\in \CH_0:\; a(x) \ll x^d\;\; \textrm{for\; some}\;\; d>0\},
\end{align*}
which is still a Hardy field closed under compositions, but not closed under inverses.

A prime example is the Hardy field $\mathcal{LE}$ of \textit{logarithmico-exponential functions}, which consists of functions defined on a half-line $(x_0,+\infty)$ by a finite combination of symbols $+, -, \times, \div, \sqrt[n]{\cdot}, \exp, \log$ acting on the real variable $x$ and on real constants. While $\mathcal{LE}$ is not closed under compositional inverses, it is contained in the Hardy field of Pfaffian functions that is. See \cite{Fr10, Fr12, Hardy12} for more on Hardy fields in general and $\mathcal{LE}$ in particular. 


















\begin{theorem}[{\cite[Theorem 1.3]{Boshernitzan-equidistribution}}]\label{T: Boshernitzan}
    Let $a\in\mathcal{H}$. The sequence $(a(n))_n$ is equidistributed $\mod~\!1$ (meaning that $$\lim_{N\to\infty}\E_{n\in [N]}F(a(n)\!\!\! \mod 1)=\int_{\T} F(t)\,dt$$ for any Riemann-integrable function $F\colon\T\to\C$) if and only if for every $p\in \Q[x]$ we have \[\lim_{x\to \infty}\frac{|a(x)-p(x)|}{\log x}=\infty.\]
\end{theorem}

\begin{definition}[Independence properties of Hardy sequences]\label{D: independence}
Let $a_1, \ldots, a_\ell\in\CH$. We call them:
\begin{enumerate}
\item \textit{independent} if for any $c_1, \ldots, c_\ell\in\R$ not all 0, we have
\begin{align*}
    \lim_{x\to\infty}\frac{|c_1 a_1(x) + \cdots + c_\ell a_\ell(x)|}{\log x}=\infty,
\end{align*}
and \textit{dependent} otherwise;
\item \textit{pairwise independent} if $a_i, a_j$ are independent for any $1\leq i<j\leq \ell$, and \textit{pairwise dependent} otherwise.
 \end{enumerate}
\end{definition}

\begin{definition}[Strongly nonpolynomial Hardy sequences]
    We call $a\in\CH$ \textit{strongly nonpolynomial} if     there exists $d\in\N_0$ such that $x^d\prec a(x)\prec x^{d+1}$.
\end{definition}

 \subsection*{Acknowledgments}
 We would like to thank Nikos Frantzikinakis for helpful discussions and sharing ideas that turned into the proof of part (ii) of Lemma \ref{L: slowly growing}.

\section{Generalized box seminorms and factors}

This section discusses generalized box seminorms along finitely generated subgroups $G_1, \ldots, G_s\subseteq\R^\ell$ and associated factors, all of which extend the classical Host-Kra and box seminorms and factors defined for subgroups of $\Z^\ell$ (see e.g. \cite{DFKS22, H09, HK05a, HK18, TZ16}). The generalized box seminorms have been defined in our earlier paper \cite{DKKST24}, in which we proved their main properties; the notion of factors is new. For the duration of the entire section, we fix a system $(X, \CX, \mu, T_1, \ldots, T_\ell)$.

\subsection{Generalized box seminorms}

Let $G_1, \ldots, G_s\subseteq\R^\ell$ be finitely generated subgroups. For $f\in L^\infty(\mu)$ and $\bm,\bm'\in\R^\ell$, let $\Delta_{(\bm,\bm')}f = T^{\sfloor{\bm}}f\cdot T^{\sfloor{\bm'}}\overline{f}$. Recall from \cite{DKKST24} that we can define two seminorms of $f$ along $G_1, \ldots, G_s$. The first one is given by
\begin{align*}
    \nnorm{f}_{\substack{G_1, \ldots, G_s}} &= \brac{\E_{\bm_1, \bm_1'\in G_1}\cdots \E_{\bm_s, \bm_s'\in G_s}\int \Delta_{({\bm_1}, {\bm'_1}), \ldots, ({\bm_s}, {\bm'_s})}f\, d\mu}^{1/2^{s}}\\
    &= \brac{\E_{\bm_1, \bm_1'\in G_1}\cdots \E_{\bm_s, \bm_s'\in G_s}\int \prod_{\ueps\in\{0,1\}^s}\CC^{|\ueps|}T^{\sfloor{\bm_1^{\eps_1}}+\cdots + \sfloor{\bm_s^{\eps_s}}} f\, d\mu}^{1/2^{s}},
\end{align*}
while the second one is 
\begin{align*}
    \nnorm{f}^+_{\substack{G_1, \ldots, G_s}} &:=  \brac{\E_{\bm_1, \bm_1'\in G_1}\cdots \E_{\bm_s, \bm_s'\in G_s}\abs{\int \Delta_{({\bm_1}, {\bm_1'}), \ldots, ({\bm_s}, {\bm_s'})}f\, d\mu}^2}^{1/2^{s+1}}\\
      &= \brac{\E_{\bm_1, \bm_1'\in G_1}\cdots \E_{\bm_s, \bm_s'\in G_s}\int \Delta_{({\bm_1}, {\bm_1'}), \ldots, ({\bm_s}, {\bm_s'})}f\otimes\overline{f}\, d(\mu\times\mu)}^{1/2^{s+1}}\\
      &=\nnorm{f\otimes\overline{f}}_{G_1, \ldots, G_s}^{1/2}.
\end{align*}
See \eqref{E: h^eps} for the convention regarding expressions like $T^{\sfloor{\bm_1^{\eps_1}}}$.

If $G_i = \langle \bv_i\rangle$, we also write 
\begin{align*}
    \nnorm{f}_{\substack{G_1, \ldots, G_s}} = \nnorm{f}_{\substack{\bv_1, \ldots, \bv_s}}\quad \textrm{and}\quad \nnorm{f}^+_{\substack{G_1, \ldots, G_s}} = \nnorm{f}^+_{\substack{\bv_1, \ldots, \bv_s}},
\end{align*}
and if $G_1 = \cdots = G_s = G$, then
\begin{align*}
    \nnorm{f}_{\substack{G_1, \ldots, G_s}} = \nnorm{f}_{G^{\times s}}\quad \textrm{and}\quad \nnorm{f}^+_{\substack{G_1, \ldots, G_s}} = \nnorm{f}^+_{G^{\times s}}.
\end{align*}

For instance, if $\ell=s=2$ and $G_i = \langle \alpha_i \be_i\rangle$, then 
\begin{align*}
    \nnorm{f}_{\substack{G_1, G_2}}^4= \E\limits_{m_{1},m'_{1}\in\Z}\E\limits_{m_{2},m'_{2}\in\Z} &\int T_1^{\sfloor{\alpha_1 m_1}}T_2^{\sfloor{\alpha_2 m_2}} f \cdot T_1^{\sfloor{\alpha_1 m_1}}T_2^{\sfloor{\alpha_2 m'_2}}\overline{f}\\
        &\qquad \cdot T_1^{\sfloor{\alpha_1 m'_1}}T_2^{\sfloor{\alpha_2 m_2}}\overline{f}\cdot T_1^{\sfloor{\alpha_1 m_1'}}T_2^{\sfloor{\alpha_2 m_2'}}f\, d\mu
\end{align*}
and
\begin{align*}
    \Bigbrac{\nnorm{f}^+_{\substack{G_1, G_2}}}^8= \E\limits_{m_{1},m'_{1}\in\Z}\E\limits_{m_{2},m'_{2}\in\Z} &\left|\int T_1^{\sfloor{\alpha_1 m_1}}T_2^{\sfloor{\alpha_2 m_2}} f \cdot T_1^{\sfloor{\alpha_1 m_1}}T_2^{\sfloor{\alpha_2 m'_2}}\overline{f}\right.\\
        &\qquad \cdot\left. T_1^{\sfloor{\alpha_1 m'_1}}T_2^{\sfloor{\alpha_2 m_2}}\overline{f}\cdot T_1^{\sfloor{\alpha_1 m_1'}}T_2^{\sfloor{\alpha_2 m_2'}}f\, d\mu\right|^2.
\end{align*}

The following properties of generalized box seminorms have all been established in \cite[Section 3]{DKKST24}.
\begin{lemma}[Basic properties of seminorms]\label{L: basic properties of seminorms}
    Let $G_1, \ldots, G_s\subseteq\R^\ell$ be finitely generated subgroups and $f\in L^\infty(\mu)$. Then: 
    \begin{enumerate}
        \item (Symmetry) For any permutation $\sigma:[s]\to[s]$, we have
        \begin{align*}
            \nnorm{f}_{G_1, \ldots, G_s} = \nnorm{f}_{G_{\sigma(1)}, \ldots, G_{\sigma(s)}}\quad \textrm{and}\quad \nnorm{f}_{G_1, \ldots, G_s}^+ = \nnorm{f}_{G_{\sigma(1)}, \ldots, G_{\sigma(s)}}^+.
        \end{align*}
        \item (Monotonicity) For any finitely generated $G_{s+1}\subseteq\R^\ell$, we have
        $$\nnorm{f}_{G_1, \ldots, G_s}\leq\nnorm{f}_{G_1, \ldots, G_s}^+\leq\nnorm{f}_{G_1, \ldots, G_{s+1}}\leq\nnorm{f}_{G_1, \ldots, G_{s+1}}^+.$$
        \item (Subgroup property) If $G'_i\subseteq G_i$ for all $i\in[s]$, then $$\nnorm{f}_{G_1, \ldots, G_s}\leq\nnorm{f}_{G'_1, \ldots, G'_s} \quad \textrm{and}\quad \nnorm{f}^+_{G_1, \ldots, G_s}\leq\nnorm{f}_{G'_1, \ldots, G'_s}^+.$$ Conversely,  if the index of $G'_{i}$ in $G_{i}$ is at most $D$ for all $i\in[s]$, then %$\nnorm{f}_{G'_1, \ldots, G'_s}^+\ll_{D,s} \nnorm{f}_{G_1, \ldots, G_s}^+.$
\begin{align*}
\nnorm{f}_{G'_1, \ldots, G'_s}^+\ll_{D,s} \nnorm{f}_{G_1, \ldots, G_s}^+.
\end{align*} 
        \item (Rescaling property) Let $K>0$. If for every $i\in[s]$ we have
        \begin{align*}
            G_i = \langle \bg_{i1}, \ldots, \bg_{iK_i}\rangle\quad \textrm{and}\quad G'_i = \langle \alpha_{i1}\bg_{i1}, \ldots, \alpha_{iK_i} \bg_{iK_i}\rangle 
        \end{align*}
        for some $K_i\leq K$ and nonzero real $1/K\leq |\alpha_{ij}|\leq K$, then
    \begin{align*}
        \nnorm{f}^+_{G'_1, \ldots, G'_s}\asymp_{\ell, K, s} \nnorm{f}^+_{G_1, \ldots, G_s}
    \end{align*}
    and
    \begin{align*}
        \nnorm{f}_{G'_1, \ldots, G'_s}\asymp_{\ell, K, s} \nnorm{f}_{G_1, \ldots, G_s} \quad \textrm{(if\; additionally}\; s\geq 2).
    \end{align*}
    \item (Equality of ergodic averages) For all finitely generated subgroups $G_1, \ldots, G_s~\subseteq~\R^\ell$ for which  $\E\limits_{\bm\in G_i} T^{\floor{\bm}}f = \E\limits_{\bm\in G'_i} T^{\floor{\bm}}f$ holds for all $f\in L^\infty(\mu)$ and $i\in[s]$, we have
    \begin{align*}
        \nnorm{f}_{G_1, \ldots, G_s}=\nnorm{f}_{G'_1, \ldots, G'_s}.
    \end{align*}
    \end{enumerate}
\end{lemma}

\subsection{Generalized box factors}\label{SS: factors}

Our goal now is to construct factors $\CZ^+_{G_1, \ldots, G_s}$ with the property that
\begin{align}\label{E: seminorm vs factor}
    \nnorm{f}_{G_1, \ldots, G_s}^+ = 0 \quad \Longleftrightarrow \quad \E(f|\CZ_{G_1, \ldots, G_s}^+) = 0.
\end{align}
If $G_1, \ldots, G_s$ lie inside $\Z^\ell$, then %we could also define
there exist factors $\CZ_{G_1, \ldots, G_s}$ with the property
\begin{align}\label{E: seminorm vs factor unplussed}
    \nnorm{f}_{G_1, \ldots, G_s} = 0 \quad \Longleftrightarrow \quad \E(f|\CZ_{G_1, \ldots, G_s}) = 0
\end{align}
%(indeed, such factors have been constructed e.g. in \cite{TZ16} following \cite{H09, HK05a}). 
(see for example \cite{H09, HK05a,S21,TZ16}).
%However, it is likely that no suitable factor  exists in general if $G_1, \ldots, G_s$ do not lie inside the integer lattice (see the discussion below Definition \ref{D: factor}). 
We however do not know if there exists a suitable factor satisfying (\ref{E: seminorm vs factor unplussed}) whenever $G_1, \ldots, G_s$ do not lie inside $\Z^\ell$. By the results in \cite{DKS20}, there exists a factor corresponding to the seminorm $\nnorm{\cdot}_{\langle\alpha\rangle}$ for a system $(X,\mathcal{X},\mu,T)$ and $\alpha\in [0,1)$ (this factor is spanned by eigenfunctions with eigenvalues coming from the subgroup $(\frac{1}{\alpha}\Z)/\Z$). While this limited evidence suggests a possibility of the existence of a factor satisfying (\ref{E: seminorm vs factor unplussed}) for general finitely generated subgroups $G_1, \ldots, G_s\subseteq\R^\ell$, we do not require such factors in our arguments and hence do not pursue this question in this paper. We refer the readers to the discussion below Definition \ref{D: factor} for an explanation of difficulties inherent in constructing a factor $\CZ_{G_1, \ldots, G_s}$ even when $s=1$.

% However, if $G_1, \ldots, G_s$ do not lie inside $\Z^\ell$, then in general there exists no suitable factor satisfying (\ref{E: seminorm vs factor unplussed}). For example, \cite[Proposition 4.4]{DKS20} indicates that in general there is no factor corresponding to the seminorm $\nnorm{\cdot}_{\langle\alpha\rangle}$ for a system $(X,\mathcal{X},\mu,T)$ and $\alpha\in \mathbb{R}\backslash\mathbb{Q}$. We also refer the readers to the discussion below Definition \ref{D: factor} for this topic.

Following e.g. \cite{BHK05, FH18, GT12}, we define multiparameter nilsequences. %For the definitions of filtrations and polynomial sequences, see \cite{GT12}.

% \begin{definition}[Multiparameter nilsequences]
%     A \textit{$K$-parameter degree-$s$ nilsequence} is a bounded sequence  $\psi:\Lambda\to\C$, where: %\KT{Why do we need polynomial sequences? All error terms that we get in the proofs are "linear nilsequences". } \BK{If I recall correctly, there is one point where we consider a sequence $g\circ \eta$ for some linear isomorphism $\eta$ - this will no longer be linear.}
%     \begin{itemize}
%         \item $\Lambda\subseteq\Z^K$ is a finite-index subgroup for some $K\in\N$;
%         \item $G/\Gamma$ is a nilmanifold with a degree-$s$ filtration $G_\bullet$;
%         \item $g:\Lambda\to G$ is a polynomial sequence adapted to the filtration $G_\bullet$;
%         \item $\psi(\um) = F(g(\um)\Gamma)$ for a bounded Riemann integrable function $F:G/\Gamma\to\C$.
%     \end{itemize}
% %\AK{I will have to check the following: While I was in the group of people that agreed on the definition of nilsequences through Riemann integrable functions, we usually assume continuity. Actually, there is a difference between these notions (I remember a paper by Green plus someone answering a question of Nikos in the nil + null setting).}
    
%     A \textit{degree-$s$ nilsequence} is then a $K$-parameter degree-$s$ nilsequence for some $K\in\N$.
%     %Let $\Lambda\subseteq\Z^K$ be a finite-index subgroup for some $K\in\N$. A \textit{$K$-parameter $s$-step nilsequence} is a sequence  $\psi:\Lambda\to\C$ given by $\psi(\um) = F(\tau_1^{m_1}\cdots \tau_K^{m_K}\Gamma)$ for a bounded Riemann integrable function $F:G/\Gamma\to\C$ on an $s$-step nilmanifold $G/\Gamma$ and elements $\tau_1, \ldots, \tau_K\in G$.
% \end{definition}


% The next lemma allows us to lift a $K$-parameter $s$-step nilsequence to a sequence defined on $\R^K$ that nevertheless equals $\psi$ when restricted to $\Q^K$.
% \begin{lemma}[Nilsequence lifting]\label{L: nilsequence lifting}
%     Let $s,K\in\N$, and let $\psi:\Lambda\to\C$ be a $K$-parameter degree-$s$ nilsequence supported on a finite-index subgroup $\Lambda\subseteq\Z^K$. Let $\widetilde\psi:\R^K\to\C$ be defined via $\widetilde{\psi}=\psi\cdot 1_\Lambda$. Then $\widetilde\psi:\widetilde{\Lambda}\to\C$ is a degree-$s$ nilsequence for any subgroup $\Z^K\subseteq\widetilde{\Lambda}\subseteq\Q^K$ such that the group $\widetilde{\Lambda}/\Z^K$ is finite. 
% \end{lemma}
% \begin{proof}
% We will first lift $\psi$ to $\R^K$. Let $\psi(\um) = F(g(\um)\Gamma)$ for a bounded Riemann integrable function $F:G/\Gamma\to\C$. Since $G/\Gamma$ can be embedded into an $s$-step nilmanifold $\widetilde{G}/\widetilde{\Gamma}$ with $\widetilde{G}\supseteq G$ connected and simply-connected \BK{(cite)}, we can assume without loss of generality that $G$ itself is connected and simply-connected. Then the expression $F(g(\ut)\Gamma)$ makes sense for every $\ut\in\R^K$.

% Now, consider any subgroup $\widetilde{\Lambda}\subseteq \Q^K$ such that $\Z^K\subseteq\widetilde{\Lambda}$ is a finite-index subgroup. Then $\Lambda\subseteq\widetilde{\Lambda}$ is also a finite-index subgroup. Let $H$ be the Pontryagin dual of $\widetilde{\Lambda}/\Lambda$ and $D = \abs{H} = \abs{\widetilde{\Lambda}/\Lambda}$. Fourier expanding, we get $\widehat{1_{\Lambda}}(\chi) = \frac{1}{D}\overline{\chi}(\Lambda)$ for every $\chi\in H$, and so Fourier inversion gives us $$1_{\Lambda}(\um) = \E_{\chi\in H}\overline{\chi}(\Lambda) \chi(\um+\Lambda)$$
% for any $\um\in\widetilde{\Lambda}$. 
%     This is a finite linear combination of sequences of the form $\um\mapsto \chi(\um+\Lambda)$, which are periodic degree-1 nilsequences. Hence $\psi\cdot 1_\Lambda:\widetilde\Lambda\to\C$ is also a degree-$s$ nilsequence (we are using the assumption $s\geq 1$ here). The result follows on setting $\widetilde{\psi} = \psi \cdot 1_\Lambda$.
%     \end{proof}




% \begin{lemma}[Nilsequence lifting]\label{L: nilsequence lifting}
%     Let $s,K\in\N$, and let $\psi:\Lambda\to\C$ be a $K$-parameter degree-$s$ nilsequence supported on a finite-index subgroup $\Lambda\subseteq\Z^K$. Then there exists a function $\widetilde\psi:\R^K\to\C$, supported on $\Lambda$, such that $\widetilde{\psi}|_{\Z^K} =\psi$ and $\widetilde\psi:\Z^K\to\C$ is a $K$-parameter degree-$s$ nilsequence.
% \end{lemma}
% \begin{proof}
% We will first lift $\psi$ to $\Z^K$, and then to $\R^K$. First, let $H$ be the Pontryagin dual of $\Z^K/\Lambda$ and $D = \abs{H} = \abs{\Z^K/\Lambda}$. Fourier expanding, we get $\widehat{1_{\Lambda}}(\chi) = \frac{1}{D}\overline{\chi}(\Lambda)$ for every $\chi\in H$, and so Fourier inversion gives us $$1_{\Lambda}(\um) = \E_{\chi\in H}\overline{\chi}(\Lambda) \chi(\um+\Lambda).$$
%     This is a finite linear combination of sequences of the form $\um\mapsto \chi(\um+\Lambda)$, which are periodic $K$-parameter degree-1 nilsequences. Hence $\psi\cdot 1_\Lambda:\Z^K\to\C$ is also a $K$-parameter degree-$s$ nilsequence (we are using the assumption $s\geq 1$ here).

% Now, let $\psi(\um) = F(g(\um)\Gamma)$ for a bounded Riemann integrable function $F:G/\Gamma\to\C$. Since $G/\Gamma$ can be embedded into an $s$-step nilmanifold $\widetilde{G}/\widetilde{\Gamma}$ with $\widetilde{G}\supseteq G$ connected and simply-connected \BK{(cite)}, we can assume without loss of generality that $G$ itself is connected and simply-connected. Then the expression $F(g(\ut)\Gamma)$ makes sense for every $\ut\in\R^K$, and the result follows on setting $\widetilde{\psi}(\ut) = F(g(\ut)\Gamma)1_\Lambda(\ut)$.\end{proof}

% \begin{lemma}[Reparametrizing a nilsequence]
%     Let $s,K\in\N$, and let $\psi:\Lambda\to\C$ be a $K$-parameter degree-$s$ nilsequence for some finite-index subgroup $\Lambda\subseteq\Z^K$. Let $\eta:\Q^K\to\Q^K$ be a $\Q$-vector space isomorphism. Then $\psi\circ\eta\cdot 1_{\widetilde{\Lambda}}:\widetilde{\Lambda}\to\C$ is a degree-$s$ nilsequence for any subgroup $\Z^K\subseteq\widetilde{\Lambda}\subseteq\Q^K$ such that the group $\widetilde{\Lambda}/\Z^K$ is finite. 
% \end{lemma}
% \begin{proof}
%     By Lemma \ref{L: nilsequence lifting}, we can find a function 
%     Let $\psi(\um) = F(\tau_1^{m_1}\cdots \tau_K^{m_K}\Gamma)$ for a bounded Riemann integrable function $F:G/\Gamma\to\C$ on an $s$-step nilmanifold $G/\Gamma$ and elements $\tau_1, \ldots, \tau_K\in G$. Since $G/\Gamma$ can be embedded into an $s$-step nilmanifold $\widetilde{G}/\widetilde{\Gamma}$ with $\widetilde{G}\supseteq G$ connected and simply-connected \BK{(cite)}, we can assume without loss of generality that $G$ itself is connected and simply-connected. Letting $\eta(\um) = (\eta_1(\um), \ldots, \eta_K(\um))$, we therefore have
%     \begin{align*}
%         \psi\circ \eta(\um) = F(\tau_1^{\eta_1(\um)}\cdots\tau_K^{\eta_K(\um)}\Gamma).
%     \end{align*}
%     Using the arguments from \cite[Section 2.11]{L05c}, we can express
%     \begin{align*}
%         \psi\circ \eta(\um) = \widehat{F}(\widehat{\tau}_1^{m_1}\cdots\widehat{\tau}_K^{m_K}\widehat{\Gamma})
%     \end{align*}
%     on some other $s$-step nilmanifold $\widehat{G}/\widehat{\Gamma}\supseteq G/\Gamma$ and $\widehat{\tau}_1, \ldots, \widehat{\tau}_K\in \widehat{G}$. \BK{Please double check Leibman's arwhether the group $\widehat{G}$ is of the same nilpotency class as $G$. }
% \end{proof}










% % Following e.g. \cite[Section 2.1.2]{FH18}, we define multiparameter nilsequences. In contrast to \cite{FH18}, we allow the function defining the nilmanifold to be bounded and Riemann integrable rather than continuous. 
% % \begin{definition}[Multiparameter nilsequences]
% %     A \textit{$K$-parameter $s$-step nilsequence} is a sequence  $\psi:\Lambda\to\C$, where:
% %     \begin{itemize}
% %         \item $\Lambda\subseteq\Z^K$ is a finite-index subgroup for some $K\in\N$;
% %         \item $G/\Gamma$ is an $s$-step nilmanifold, and $\tau_1, \ldots, \tau_K\in G$;
% %         \item $\psi(\um) = F(\tau_1^{m_1}\cdots \tau_K^{m_K}\Gamma)$ for a bounded Riemann integrable function $F:G/\Gamma\to\C$.
% %     \end{itemize}
% %     %Let $\Lambda\subseteq\Z^K$ be a finite-index subgroup for some $K\in\N$. A \textit{$K$-parameter $s$-step nilsequence} is a sequence  $\psi:\Lambda\to\C$ given by $\psi(\um) = F(\tau_1^{m_1}\cdots \tau_K^{m_K}\Gamma)$ for a bounded Riemann integrable function $F:G/\Gamma\to\C$ on an $s$-step nilmanifold $G/\Gamma$ and elements $\tau_1, \ldots, \tau_K\in G$.
% % \end{definition}

% % As the following lemma shows, we can always lift a $K$-parameter $s$-step nilsequence defined on a finite index subgroup $\Lambda\subseteq\Z^K$ to a nilsequence defined on $\Z^K$.
% % \begin{lemma}[Nilsequence lifting]\label{L: nilsequence lifting}
% %     Let $\Lambda\subseteq\Z^K$ be a finite-index subgroup for some $K\in\N$ and $\psi:\Lambda\to\C$ be a $K$-parameter $s$-step nilsequence. Then $\psi\cdot 1_\Lambda:\Z^K\to\C$ is also a $K$-parameter $s$-step nilsequence.
% % \end{lemma}
% % \begin{proof}
% %     Let $H$ be the Pontryagin dual of $\Z^K/\Lambda$ and $D = \abs{H} = \abs{\Z^K/\Lambda}$. Fourier expanding, we get $\widehat{1_{\Lambda}}(\chi) = \frac{1}{D}\overline{\chi}(\Lambda)$ for every $\chi\in H$, and so Fourier inversion gives us $$1_{\Lambda}(\um) = \E_{\chi\in H}\overline{\chi}(\Lambda) \chi(\um+\Lambda).$$
% %     % \begin{align*}
% %     %     \widehat{1_{\Lambda}}(\chi) = \sum_{\um\in \Z^K/\Lambda}1_{\Lambda}(\um + \Lambda)\overline{\chi}(\um)\\
% %     %     \sum_{\chi\in H}1_{\Lambda}(\um)\chi(\um)
% %     % \end{align*}
% %     This is a finite linear combination of sequences of the form $\um\mapsto \chi(\um+\Lambda)$, which are periodic $K$-parameter 1-step nilsequences. Hence $\psi\cdot 1_\Lambda$ is also a $K$-parameter nilsequence of step $\max(s,1)$.
% % \end{proof}

% % \begin{lemma}[Reparametrizing a nilsequence]
% %     Let $\Lambda\subseteq\Z^K$ be a finite-index subgroup for some $K\in\N$ and $\psi:\Lambda\to\C$ be a $K$-parameter $s$-step nilsequence. Let $\eta:\Q^K\to\Q^K$ be a $\Q$-vector space isomorphism. Then $\psi\circ\eta|_{?}:???\to\C$ is also a $K$-parameter $s$-step nilsequence.
% % \end{lemma}
% % \begin{proof}
% %     Let $\psi(\um) = F(\tau_1^{m_1}\cdots \tau_K^{m_K}\Gamma)$ for a bounded Riemann integrable function $F:G/\Gamma\to\C$ on an $s$-step nilmanifold $G/\Gamma$ and elements $\tau_1, \ldots, \tau_K\in G$. Since $G/\Gamma$ can be embedded into an $s$-step nilmanifold $\widetilde{G}/\widetilde{\Gamma}$ with $\widetilde{G}\supseteq G$ connected and simply-connected \BK{(cite)}, we can assume without loss of generality that $G$ itself is connected and simply-connected. Letting $\eta(\um) = (\eta_1(\um), \ldots, \eta_K(\um))$, we therefore have
% %     \begin{align*}
% %         \psi\circ \eta(\um) = F(\tau_1^{\eta_1(\um)}\cdots\tau_K^{\eta_K(\um)}\Gamma).
% %     \end{align*}
% %     Using the arguments from \cite[Section 2.11]{L05c}, we can express
% %     \begin{align*}
% %         \psi\circ \eta(\um) = \widehat{F}(\widehat{\tau}_1^{m_1}\cdots\widehat{\tau}_K^{m_K}\widehat{\Gamma})
% %     \end{align*}
% %     on some other $s$-step nilmanifold $\widehat{G}/\widehat{\Gamma}\supseteq G/\Gamma$ and $\widehat{\tau}_1, \ldots, \widehat{\tau}_K\in \widehat{G}$. \BK{Please double check Leibman's arwhether the group $\widehat{G}$ is of the same nilpotency class as $G$. }
% % \end{proof}



% Our next goal is to obtain a suitable notion of nilsequence on general finitely generated subgroups $G\subseteq\R^k$. To this end, we need a way to map $G$ to a subgroup of $\Z^K$.
% \begin{definition}[Coordinate maps]
%     Let $G\subseteq\R^k$ be a finitely generated subgroup and $K=\dim(\Span_\Q G)$. We call a vector space isomorphism $\phi:\Z^K\to G$ a \textit{coordinate map} if $G$ is a finite index subgroup of $\phi(\Z^K)$.
%     %there exists a group $\widetilde{G}\supseteq G$ such that $\widetilde{G}/G$ is finite and $\phi|_{\Z^K}:\Z^K\to\widetilde{G}$ is an isomorphism.
%     %\footnote{Note that we cannot in general define a group isomorphism from $\Z^K$ to $G$ (consider e.g. $G = \langle (2,0), (0,2), (1,1)\rangle$).} \textit{for $G$} if $\phi\inv(G)\subseteq\Z^K$. %$\phi(\Z^K)\subseteq G$.
% \end{definition}

% For instance, if $G = \langle (\alpha, \beta) \rangle$, then $\phi(m) = (\alpha m, \beta m)$ is a suitable coordinate map. 

% Note that we cannot in general take $\phi(\Z^K) = G$. Indeed, if $G = \langle (2,0), (0,2), (1,1)\rangle$, then we can embed it as a finite index subgroup into $\langle (1,0), (0,1)\rangle = \Z^2$, but $G$ itself is not group-isomorphic to $\Z^K$ for any $K\in\N$.

% With the help of coordinate maps, the notion of nilsequences can be extended to general finitely generated subgroups $G\subseteq\R^k$.
% \begin{definition}[Nilsequences on $G$]
% %    Let $G\subseteq\R^k$ be a finitely generated subgroup with a coordinate map $\phi:\Z^K\to G$. Then a sequence $\psi:G\to\C$ is an \textit{$s$-step nilsequence on $G$} if the map $\psi\circ \phi: \phi\inv(G)\to\C$ is a $K$-parameter $s$-step nilsequence. For $\veps>0$, we call $\psi$ an \textit{$\veps$-nullsequence on $G$} if
% Let $G\subseteq\R^k$ be a finitely generated subgroup. Then a sequence $\psi:G\to\C$ is a \textit{basic degree-$s$ nilsequence on $G$} if there exists a coordinate map $\phi:\Z^K\to G$ such that the map $\psi\circ \phi: \phi\inv(G)\to\C$ is a $K$-parameter degree-$s$ nilsequence. We then define $\Nil_s(G)$ to be the algebra generated by all basic degree-$s$ nilsequences on $G$, and we call any $\psi\in\Nil_s(G)$ a \textit{degree-$s$ nilsequence on $G$}.
%     % Let $G\subseteq\R^k$ be finitely generated. We call a group isomorphism $\phi:G\to\Z^K$ a \textit{coordinate map for $G$}. Then a sequence $\psi:G\to\C$ is an \textit{$s$-step nilsequence on $G$} if the map $\psi\circ \phi: \Z^K\to\C$ is a $K$-parameter $s$-step nilsequence. \BK{(We should define properly what that means)} For $\veps>0$, we call $\psi$ a \textit{$\veps$-nullsequence on $G$} if
%     % \begin{align*}
%     %     \limsup_{N\to\infty}\sup_{(I_N)_N}\E_{\bm\in I_N}|\psi(\bm)|^2\leq \veps,
%     % \end{align*}
%     % where the supremum is taken along any F\o lner sequence on $G$.
% \end{definition}
% If $\psi_1, \psi_2:G\to\C$ are two basic nilsequences on $G$, then there exist coordinate maps $\phi_1, \phi_2:\Q^K\to\Span_\Q G$ such that $\psi_i \circ \phi_i: \phi_i\inv(G)\to\C$ are $K$-parameter nilsequences. Then both $\psi_1 \cdot (\psi_2\circ\phi_2\circ\phi_1\inv)$ and $(\psi_1\circ\phi_1\circ\phi_2\inv) \cdot \psi_2$ are basic $K$-parameter nilsequences with the coordinate maps $\phi_1, \phi_2$ respectively (assuming that the maps $\psi_2\circ\phi_2\circ\phi_1\inv$ and $\psi_1\circ\phi_1\circ\phi_2\inv$ are well-defined, which can be ensured\footnote{Indeed, let $\psi:\Lambda\to\C$ be a $K$-parameter, degree-$s$ nilsequence given by $\psi(\um) = F(g(\um)\Gamma)$ for a bounded Riemann integrable function $F:G/\Gamma\to\C$ on a degree-$s$ nilmanifold $G/\Gamma$. Without loss of generality, we can assume that $G/G^o$ is finitely generated, where $G^o$ is the connected component of identity, since we can always restrict to the group generated by $G^o$ and the elements  $\tau_1, \ldots, \tau_d\in G$ that are ``coefficients'' of the sequence $g$. Then $G/\Gamma$ can be embedded into an $s$-step nilmanifold $\widetilde{G}/\widetilde{\Gamma}$ with $\widetilde{G}\supseteq G$ connected and simply connected and such that any element of $G$ is represented in $\widetilde{G}$ (see e.g. \cite[Section 2.11]{L05c}).
% Hence we can assume without loss of generality that $G$ itself is connected and simply connected. Then the expression $\psi(\ut)$ makes sense for every $\ut\in\R^K$, and so $\psi\circ\eta$ is a well-defined degree-$s$ nilsequence for every vector space isomorphism $\eta:\R^K\to\R^K$ by \cite[Theorem 1.6.9]{tao_2012}.}). However, if the coordinate maps are different, then it is not immediate that the product $\psi_1 \psi_2$ is a nilsequence on $G$ with respect to some common coordinate map. {\color{red} (is it difficult to prove that the product $\psi_1 \psi_2$ is a nilsequence? if not, I think it is better to give a proof (and thus we will not need the notion of non-basic nilsequences), since having a clearer description of the factor $\mathcal{Z}^{+}$ might have applications in the future.)} \BK{(Frankly, I don't know. It was not immediate to me how to do this if we don't fix the coordinate map, but if you want, give it a try.)} To ensure that nilsequences on $G$ form an algebra, we therefore opt to define both basic nilsequences on $G$ (as the generators of the algebra) and nilsequences on $G$ (as elements of the algebra). We remark that our distinction between basic nilsequences on $G$ and nilsequences on $G$ does not coincide with a similar distinction into nilsequences and basic nilsequences common in ergodic literature (see e.g. \cite{BHK05}). 

% An alternative approach would be to assume that all nilsequences on $G$ are defined with respect to a fixed coordinate map. However, this approach seems slightly less flexible as it forces us to ensure that each nilsequence on $G$ is defined with respect to the same coordinate map in later applications.

\begin{definition}[Multiparameter nilsequences]
    A \textit{degree-$s$ $\Z^{K}$-nilsequence} is a bounded sequence  $\psi:\Z^{K}\to\C$, where: 
    \begin{itemize}
        \item $G/\Gamma$ is a nilmanifold (with $G$ connected and simply connected) with a degree-$s$ filtration\footnote{We recall that a filtration $G_\bullet = (G_i)_{i=0}^\infty$ on a group $G$ is a collection of subgroups $G = G_0=G_1\supseteq G_2 \supseteq \cdots$ satisfying $[G_i, G_j]\subseteq G_{i+j}$ for all $i,j\in\N_0$, where $[A,B]$ is the group generated by the commutators $[a,b]=a\inv b\inv a b$ for $a\in A, b\in B$.} $G_\bullet = (G_i)_{i=0}^\infty$;
        \item $g:\Z^{K}\to G$ is a polynomial sequence adapted to the filtration $G_\bullet$, meaning that $$\partial_{h_{i}}\dots\partial_{h_{1}}g(n)\in G_{i}$$
        for all $i\in\N$ and $n,h_{1},\dots,h_{i}\in\Z^{K}$ (here, $\partial_h g(n) = g(n)g(n+h)\inv$); 
        \item $\psi(\um) = F(g(\um)\Gamma)$ for a bounded Riemann integrable function $F:G/\Gamma\to\C$.
    \end{itemize}
%\AK{I will have to check the following: While I was in the group of people that agreed on the definition of nilsequences through Riemann integrable functions, we usually assume continuity. Actually, there is a difference between these notions (I remember a paper by Green plus someone answering a question of Nikos in the nil + null setting).}
    %Let $\Lambda\subseteq\Z^K$ be a finite-index subgroup for some $K\in\N$. A \textit{$K$-parameter $s$-step nilsequence} is a sequence  $\psi:\Lambda\to\C$ given by $\psi(\um) = F(\tau_1^{m_1}\cdots \tau_K^{m_K}\Gamma)$ for a bounded Riemann integrable function $F:G/\Gamma\to\C$ on an $s$-step nilmanifold $G/\Gamma$ and elements $\tau_1, \ldots, \tau_K\in G$.
\end{definition}

%Note that since the group $G$ is a connected, simply connected, nilpotent Lie group, it is divisible, and so every polynomial sequence $g:\Z^K\to G$ naturally extends to $\Z^{K}$ (in fact, to all of $\R^K$). We define nilsequences on $\Z^{K}$ as opposed to $\Z^K$ because the proof of Proposition \ref{P: same coordinate map} below uses basic linear algebra and requires us to have them defined over a vector space (rather than a group), and we choose $\Z^{K}$ rather than $\R^K$ so that the definition is compatible with the notion of coordinate maps below (which does not do the job if defined for real vector spaces).

 


Our next goal is to obtain a suitable notion of nilsequence on general finitely generated subgroups $G\subseteq\R^\ell$. To this end, we need to map $G$ to some $\Z^K$. We note that by the fundamental theorem of finitely generated abelian groups, $G$ is always group-isomorphic to $\Z^K$ for some $K\in\N_0$.
\begin{definition}[Coordinate maps]\label{D: coordinate maps}
    Let $G\subseteq\R^\ell$ be a finitely generated subgroup with rank $K$ (which we recall is the cardinality of the largest linearly independent subset of $G$). We then call a group isomorphism $\phi:\Z^K\to G$ a \textit{coordinate map}.
\end{definition}

\begin{examples}

Here are some examples of coordinate maps:
\begin{enumerate}
    \item if $G = \langle (\sqrt{2}, \sqrt{3}) \rangle\subseteq\R^2$, then $\phi(m) = (\sqrt{2} m, \sqrt{3} m)$ is a coordinate map;
    \item if $G = \langle \sqrt{2}, \sqrt{3} \rangle\subseteq\R$, then $\phi(m_1,m_2) = \sqrt{2} m_1 + \sqrt{3} m_2$ is a  coordinate map. 
\end{enumerate}    
\end{examples}
%For instance, if $G = \langle (\alpha, \beta) \rangle$, then $\phi(m) = (\alpha m, \beta m)$ is a suitable coordinate map. 

%Note that we cannot in general take $\phi(\Z^K) = G$. Indeed, if $G = \langle (2,0), (0,2), (1,1)\rangle$, then we can embed it as a finite index subgroup into $\langle (1,0), (0,1)\rangle = \Z^2$, but $G$ itself is not group-isomorphic to $\Z^K$ for any $K\in\N$.

With the help of coordinate maps, the notion of nilsequences can be extended to general finitely generated subgroups $G\subseteq\R^\ell$. In the entire ensuing discussion, $K$ denotes the rank of $G$.
\begin{definition}[Nilsequences on $G$]
%    Let $G\subseteq\R^\ell$ be a finitely generated subgroup with a coordinate map $\phi:\Z^K\to G$. Then a sequence $\psi:G\to\C$ is an \textit{$s$-step nilsequence on $G$} if the map $\psi\circ \phi: \phi\inv(G)\to\C$ is a $K$-parameter $s$-step nilsequence. For $\veps>0$, we call $\psi$ an \textit{$\veps$-nullsequence on $G$} if
Let $G\subseteq\R^\ell$ be a finitely generated subgroup. Then a sequence $\psi:G\to\C$ is a \textit{degree-$s$ nilsequence on $G$} if there exists a coordinate map $\phi:\Z^K\to G$ 
and a  degree-$s$ $\Z^{K}$-nilsequence $\theta\colon \Z^{K}\to \C$ such that $\theta=\psi\circ \phi$.
 We then define $\Nil_s(G)$ to be the collection of all degree-$s$ nilsequences on~$G$.
    % Let $G\subseteq\R^\ell$ be finitely generated. We call a group isomorphism $\phi:G\to\Z^K$ a \textit{coordinate map for $G$}. Then a sequence $\psi:G\to\C$ is an \textit{$s$-step nilsequence on $G$} if the map $\psi\circ \phi: \Z^K\to\C$ is a $K$-parameter $s$-step nilsequence. \BK{(We should define properly what that means)} For $\veps>0$, we call $\psi$ a \textit{$\veps$-nullsequence on $G$} if
    % \begin{align*}
    %     \limsup_{N\to\infty}\sup_{(I_N)_N}\E_{\bm\in I_N}|\psi(\bm)|^2\leq \veps,
    % \end{align*}
    % where the supremum is taken along any F\o lner sequence on $G$.
\end{definition}

\begin{proposition}\label{P: same coordinate map}
Let $G\subseteq\R^\ell$ be a finitely generated subgroup. If $\psi$ is a degree-$s$ nilsequence on $G$ with respect to some coordinate map, then it is a degree-$s$ nilsequence on $G$ with respect to any coordinate map.
\end{proposition}
\begin{proof}
    Let $\psi\colon G\to\C$ be a map and $\phi_{1},\phi_{2}\colon\Z^{K}\to G$ be coordinate maps. Suppose that $\psi\circ \phi_{1}=\theta_{1}$ for some degree-$s$ $\Z^{K}$-nilsequence $\theta_{1}\colon\Z^{K}\to \C$. It suffices to show that  $\psi\circ \phi_{2}=\theta_{2}$ for some  degree-$s$ $\Z^{K}$-nilsequence $\theta_{2}\colon\Z^{K}\to \C$.

    Note that for all $n\in \Z^K$, we have
    $$\psi\circ\phi_{2}(n)=\psi\circ(\phi_{1}\circ\phi_{1}^{-1})\circ\phi_{2}(n)=(\psi\circ\phi_{1})\circ(\phi_{1}^{-1}\circ\phi_{2})(n)=\theta_1\circ(\phi_{1}^{-1}\circ\phi_{2})(n).$$
  Because $\phi:=\phi_{1}^{-1}\circ\phi_{2}$ is a group automorphism of $\Z^{K}$, it suffices to show that degree-$s$ $\Z^{K}$-nilsequences are closed under compositions with group automorphisms of $\Z^k$.
    
    To see this, assume that $\theta_1=F(g(\cdot)\Gamma)$. Then $\theta_1\circ \phi=F(g\circ\phi(\cdot)\Gamma)$. Since $\phi:\Z^{K}\to\Z^{K}$ is a group isomorphism, we have that $$\partial_{h_{i}}\dots \partial_{h_{1}}(g\circ \phi)(n)=\partial_{\phi(h_{i})}\dots \partial_{\phi(h_{1})}g(\phi(n))$$ for all $i\in\N$ and $n,h_{1},\dots,h_{i}\in\Z^{K}$. This implies that $g\circ \phi$  is a polynomial sequence with the same filtration as $g$. So $\theta_2:=\theta_1\circ \phi$ is a  degree-$s$ $\Z^{K}$-nilsequence and we are done.
\end{proof}

As an immediate consequence, we have the following.

\begin{corollary}\label{C: algebra}
   Let $G\subseteq\R^\ell$ be a finitely generated subgroup. Then $\Nil_s(G)$ is a conjugation-closed $\C$-algebra.
\end{corollary}
\begin{proof}
    Let $\psi_1, \psi_2\in\Nil_s(G)$. By the definition of $\Nil_s(G)$ and Proposition \ref{P: same coordinate map}, we can find a coordinate map $\phi:\Z^K\to G$ and degree-$s$ $\Z^{K}$-nilsequences $\theta_1, \theta_2\colon \Z^{K}\to \C$ such that $\theta_i=\psi_i\circ \phi$ for $i\in\{1,2\}$. Then for any $c\in\C$, the functions $\overline{\psi_1}$, $\psi_1 + c\psi_2$, and $\psi_1\psi_2$ are elements of $\Nil_s(G)$ corresponding to degree-$s$ $\Z^{K}$-nilsequences $\overline{\theta_1}$, $\theta_1 + c\theta_2$, and $\theta_1\theta_2$. This proved the claim.
\end{proof}

The following notion will play a crucial part in the forthcoming definition of the factors. %$\CZ_{G_1, \ldots, G_s}^+$.
\begin{definition}[Dual functions]\label{D: dual functions}
    Set $G = G_1\times G_1 \times \cdots \times G_s\times G_s$, and let $\psi\in\Nil_s(G)$.
%$\psi$ be a degree-$s$ nilsequence on $G$. 
Let $\{0,1\}^s_* = \{0,1\}^s\setminus\{\underline{0}\}$ and assume that $f_\ueps\in L^\infty(\mu)$ for $\ueps\in \{0,1\}^s_*$.
We then call a function 
\begin{multline}\label{E: dual function}
    \CD_{G_1, \ldots, G_s}((f_\ueps)_{\{0,1\}^s_*}; \psi) := \E_{(\bm_1, \bm_1', \ldots, \bm_s, \bm_{s'})\in G}\; \psi(\bm_1,\bm_1',\ldots,\bm_s,\bm_s')\\
    \prod_{\ueps\in\{0,1\}^s_*}T^{\sfloor{\bm_1^{\eps_1}}-\floor{\bm_1}+\cdots + \sfloor{\bm_s^{\eps_s}}-\floor{\bm_s}} f_\ueps
\end{multline}
a \textit{level-$s$ dual function of $f$ along $G_1, \ldots, G_s$}. %{\color{red} Is this a generalization of dual functions appearing in Proposition 5.1?} \BK{Yes. In Proposition 5.1, we apply the weak structure theorem to the unplussed seminorm, in which case we don't need the nilsequence $\psi$, but if we wanted to only work with the plussed seminorm, then we could apply the weak structure theorem, we could look consider there dual functions that involve $\psi$. The choice is up to us.}

In a number of cases, we simplify the notation:
\begin{itemize}
    \item if $f_\ueps = \CC^{|\ueps|}f$ for some $f$, then we also denote the dual function as $\CD_{G_1, \ldots, G_s}(f; \psi)$;
    \item if $\psi = 1$, then we write $\CD_{G_1, \ldots, G_s}((f_\ueps)_{\{0,1\}^s_*})$;
    \item if $G_1 = \cdots = G_s$, then we use the label $\CD_{s, G_1}((f_\ueps)_{\{0,1\}^s_*}; \psi)$, and if all the groups equal $\langle \be_j \rangle$ for some $j\in[\ell]$, then we further write $\CD_{s, T_j}((f_\ueps)_{\{0,1\}^s_*}; \psi)$.
    %(and $\CD_{s, G_1}((f_\ueps)_{\{0,1\}^s_*})$ if additionally $\psi=1$).
\end{itemize}
We also combine all the following conventions freely, e.g. denoting the dual function by $\CD_{s, T_j}(f)$ when $G_1 = \cdots = G_s = \langle \be_j \rangle$, $f_\ueps = \CC^{|\ueps|}f$, and $\psi = 1$.
\end{definition}

% If $f_\ueps = \CC^{|\ueps|}f$ for some $f$, then we also denote the dual function as $\CD_{G_1, \ldots, G_s}(f; \psi)$. Similarly, if $\psi = 1$, then we write $\CD_{G_1, \ldots, G_s}((f_\ueps)_{\{0,1\}^s_*})$, and if $G_1 = \cdots = G_s$, then we use the label $\CD_{s, G_1}((f_\ueps)_{\{0,1\}^s_*}; \psi)$\AK{, and $\CD_{s, G_1}((f_\ueps)_{\{0,1\}^s_*})$ for $\psi=1$ respectively}.

\begin{lemma}
    The limit \eqref{E: dual function} exists in $L^2(\mu)$. Moreover, the single limit along $G$ can be replaced by an iterated limit along $G_1\times G_1$, ..., $G_s\times G_s$.
\end{lemma}
\begin{proof}
    Let $\phi:\Z^K\to G$ be a coordinate map and $(I_N)_N$ be a F\o lner sequence on $G$. Define
    \begin{align*}
         F(\bm) = \prod_{\ueps\in\{0,1\}^s_*}T^{\sfloor{\bm_1^{\eps_1}}-\floor{\bm_1}+\cdots + \sfloor{\bm_s^{\eps_s}}-\floor{\bm_s}} f_\ueps
     \end{align*}
     for $\bm\in G$.
    % $\widetilde{G} = \phi(\Z^K)$, and $E\subseteq \widetilde{G}$ be a finite of size $|E|=|\widetilde{G}/G|$ that contains 0 and satisfies $\widetilde{G} = G +E$. We first show that the average along $G$ can be replaced by an average along $\widetilde{G}$ so that later on, we can replace it by an average over $\Z^K$. Indeed, if $(I_N)_N$ is a F\o lner sequence on $G$, then $(I_N+E)_N$ is a F\o lner sequence on $\widetilde{G}$, and so
    % \begin{align*}
    %     \E_{{\bm}\in I_N}\; \psi(\bm)F(\bm) = |E|\cdot \E_{{\bm}\in I_N+E}\; \psi(\bm)1_G(\bm) F(\bm)
    % \end{align*}
    % for 
    % \begin{align*}
    %     F(\bm) = \prod_{\ueps\in\{0,1\}^s_*}T^{\sfloor{\bm_1^{\eps_1}}-\floor{\bm_1}+\cdots + \sfloor{\bm_s^{\eps_s}}-\floor{\bm_s}} f_\ueps.
    % \end{align*}
    % By Fourier expanding $1_G\circ\phi = 1_{\phi\inv(G)}$ along the finite group $\Z^K/\phi\inv(G)$, we see that it is a degree-1 $\Z^{K}$-nilsequence, and hence $1_G\in\Nil_1(\widetilde G)$.\footnote{To illustrate this maneuver with an example, consider $G = \langle (2\alpha,0), (0,2\beta), (\alpha,\beta)\rangle$ with the coordinate map $\phi:\Q^2\to\Span_\Q((\alpha,0),(0,\beta))$ given by $\phi(x,y) = (x/\alpha, y/\beta)$. Then on $\widetilde G=\phi(\Z^2)$, the map $1_G(x,y)$ equals the 1-step nilsequence $\frac{1}{2}\brac{1+e(\frac{x}{2\alpha}+\frac{y}{2\beta})}$.}
    % %is a degree-1 nilsequence on $\widetilde{G}$. 
    % From this and Corollary \ref{C: algebra} it follows that $\psi \cdot 1_G\in\Nil_s\widetilde{G}$.
    % %is an degree-$s$ nilsequence on $\widetilde{G}$, 
    % %and so without loss of generality we can assume that $G=\widetilde{G} = \phi(\Z^K)$. 
    Recasting the original average over $G$ as an average over $\Z^K$, it suffices to show that the average
    \begin{align}\label{E: dual over Z^K}
       \E_{{\um}\in \phi\inv(I_N)}\; \psi(\phi(\um)) \cdot  F(\phi(\um))
    \end{align}
    converges as $N\to\infty$ to the limit independent of the choice of $(I_N)_N$. Note that in \eqref{E: dual over Z^K}, $\psi(\bm)$ is replaced by a $\Z^{K}$-nilsequence $\psi\circ\phi$ and iterates $T^{\floor{\bm_i}}$ inside $F$ are replaced by expressions like $T^{\floor{\phi(\um_i)}}=T^{\sfloor{\balpha_1 m_{i,1} + \cdots + \balpha_{K_i} m_{i, K_i}}}$, where $G_i = \langle \balpha_1, \ldots, \balpha_{K_i}\rangle$. 
    
    
    From then on, the proof of the $L^2(\mu)$ convergence is completely analogous to the arguments in \cite[Section 3]{Ko18}. %for the corresponding regularity. 
    The multiparameter variant of the suspension flow trick \cite[Theorem~3.2]{Ko18} combined with the celebrated convergence result of Walsh \cite{W12} allows us to obtain the $L^2(\mu)$ convergence of \eqref{E: dual over Z^K} whenever $\psi= 1$.  
    %expressions like $T^{\floor{\alpha_1 m_1 + \cdots + \alpha_{K_i} m_{K_i}}}$ {\color{red}{Do we need here $\bm_i$ in bold too?}} via convergence results for analogous expressions for flows without the ceiling function; the latter exist by the (multiparameter) convergence result of Walsh \cite{W12}.  
    To treat the general form of \eqref{E: dual over Z^K}, it suffices to approximate the nilsequence $\psi$ in uniform density with the special weights considered in \cite[Proposition~4.2]{FH18}.\footnote{The paper \cite{FH18} defines nilsequences in terms ``linear sequences'' $\tau_1^{n_1}\cdots \tau_K^{n_K}$ rather than general polynomial sequences, and so here we also implicitly use the argument from \cite[Section 2.11]{L05c} to convert the original ``polynomial nilsequence'' into a ``linear nilsequence'' as defined in \cite{FH18}.} 
    %(which is the multiparameter variant of \cite[Proposition~2.4]{Fr15} that is used in the single variable case). 
    The claim then follows by arguments similar to those in \cite[Section~3]{Ko18}.

  % \AK{See green paragraph below; the black part is old.}  From then on, the proof of $L^2(\mu)$ convergence is completely analogous to arguments in \cite{FH18, Ko18}, and it follows from combining the suspension flow trick \cite[Theorem 4.1]{Ko18} (which allows us to remove integer parts from expressions like $T^{\floor{\alpha_1 m_1 + \cdots + \alpha_{K_i} m_{K_i}}}$) with  %\AK{Haven't checked the precise statements but we only have linear expressions here. Maybe we can be more specific on the multidimensional expression and get the result from my paper in which I show that nilsequence times the usual sequence converges in the mean (but for a single variable).}
  %    \cite[Proposition~2.1]{FH18}. Since the latter is stated in terms of ``linear nilsequences'' rather than ``polynomial sequences'', we also implicitly use the argument from \cite[Section 2.11]{L05c} to convert a ``polynomial nilsequence'' into a ``linear nilsequence''. 

    The convergence of the iterated limit can be proved analogously, and the equality of a single and iterated limits can be established using the Fubini principle for ergodic averages \cite[Lemma 1.1]{BL15} much like in \cite[Lemma 3.1]{DKKST24}.
    % \AK{From then on, the proof of $L^2(\mu)$ convergence is completely analogous to arguments in \cite{Ko18} for the corresponding regularity. The multiparameter variant of the suspension flow trick \cite[Theorem~3.2]{Ko18} allows us to obtain convergence results for expressions like $T^{\floor{\alpha_1 m_1 + \cdots + \alpha_{K_i} m_{K_i}}}$ {\color{red}{Do we need here $\bm_i$ in bold too?}} via convergence results for analogous expressions for flows without the ceiling function; the latter exist by the (multiparameter) convergence result of Walsh \cite{W12}.  Then, to treat the general, weighted by nilsequences, expressions, it suffices to deal with special weights as the ones in \cite[Proposition~4.2]{FH18} (which is the multiparameter variant of \cite[Proposition~2.4]{Fr15} that is used in the single variable case). The claim follows by analogous arguments in \cite[Section~3]{Ko18}. 
    %
    % The convergence of the iterated limit can be proved analogously, and the equality of a single and iterated limits can be established using the Fubini principle for ergodic averages \cite[Lemma~1.1]{BL15} much like in \cite[Lemma~3.1]{DKKST24}.}
\end{proof}
%By combining \cite[Lemma 3.1]{DKKST24}, the product property of limits, the uniform convergence along nilsequences \cite{???}, and the fact that a coordinate map $\phi:\Z^K\to G$ induces a correspondence between F\o lner sequences on $G$ and $\Z^K$, one gets that the single limit along $G$ can be replaced by an iterated limit along $G_1\times G_1, \ldots, G_s\times G_s$ as well as by other iterated limits. \BK{(This is not a correct explanation. We should explain why these things converge, possibly by combining the methods of \cite{FH18, Ko18})}

% \AK{This \emph{should} be exactly the same with what we did with $k$-regularity (I left a question there on how we go from $G$ to $\mathbb{Z}^k$ via $\phi^{-1}$). Via \cite[Proposition~4.2]{FH18}, we have the result we want, analogously to \cite[Proposition~2.1]{FH18} but for integer parts of real polynomials.}

Dual functions play a crucial role in the following definition of the closed space of functions ``structured'' with respect to the groups $G_1, \ldots, G_s$.
\begin{definition}\label{D: factor}
    Let $G_1, \ldots, G_s\subseteq\R^\ell$ be finitely generated subgroups.
    Given a factor $\CY\subseteq \CX$, let $Z_{G_1, \ldots, G_s}^+(\CY)\subseteq L^\infty(\mu)$ be the $L^2(\mu)$-closed linear span of functions of the form 
$\CD_{G_1, \ldots, G_s}((f_\ueps)_{\{0,1\}^s_*}; \psi)$ for $f_\ueps\in L^\infty(\CY, \mu)$ and $\psi\in\Nil_s(G)$
%degree-$s$ nilsequences $\psi: G\to\C$ 
(where we recall that $G = G_1\times G_1\times\ldots \times G_s\times G_s$). If $\CY = \CX$, then we simply write $Z_{G_1, \ldots, G_s}^+(\CX) = Z_{G_1, \ldots, G_s}^+$.



%If the functions $f_\ueps$ are measurable with respect to a sub-$\sigma$-algebra $\CY$, then we also denote
\end{definition}












This is a good place to explain the presence of a nilsequence $\psi$ in the definition \eqref{E: dual function} as well as the reason why we consider a factor associated with the seminorm $\nnorm{\cdot}_{G_1, \ldots, G_s}^+$ rather than $\nnorm{\cdot}_{G_1, \ldots, G_s}$. To this end, consider the group ${G_1} = \langle (\alpha, \beta)\rangle\subseteq\R^2$ for some irrational $\alpha,\beta\in\R$. If we wanted to find an algebra $Z_{G_1}\subseteq L^\infty(\mu)$ with the property \eqref{E: seminorm vs factor unplussed}, then it would have to be the closed linear span of functions of the form
\begin{align*}
    \CD_{G_1}(f) = \E_{m,m'\in\Z} T_1^{\sfloor{\alpha m'}-\sfloor{\alpha m}}T_2^{\sfloor{\beta m'}-\sfloor{\beta m}}f
\end{align*}
(with no $1$-step-nilsequence weight). However, it is not clear whether the resulting vector space is an algebra (i.e. closed under multiplication). The reason is that the product $\CD_{G_1}(f)\cdot\CD_{G_1}(g)$ equals
\begin{multline*}
    \CD_{G_1}(f)\cdot\CD_{G_1}(g) = \sum_{r_1,r_2=-1}^1\E_{n,n'\in\Z}\E_{m,m'\in\Z}\psi_{r_1,r_2}(m,m',n,n')\\
    T_1^{\sfloor{\alpha m'}-\sfloor{\alpha m}}T_2^{\sfloor{\beta m'}-\sfloor{\beta m}}(f\cdot T_1^{\sfloor{\alpha n'}-\sfloor{\alpha n}+r_1}T_2^{\sfloor{\beta n'}-\sfloor{\beta n}+r_2}g)
\end{multline*}
for some nontrivial weights $\psi_{r_1,r_2}$ that take care of the error terms coming from comparing expressions like $\sfloor{\alpha (m+n)}$ with $\sfloor{\alpha m} + \sfloor{\alpha n}$. It is unclear whether these weights, which can easily be seen to be degree-1 nilsequences, can always be removed. 

Once we allow for the presence of a degree-1 nilsequence weight in the definition of $\CD_{G_1}(f)$, then we get an algebra $Z_{G_1}^+$ corresponding to the seminorm $\nnorm{\cdot}_{{G_1}}^+$ rather than $\nnorm{\cdot}_{{G_1}}$ (see Lemma~\ref{L: algebra} and Proposition~\ref{P: structure theorem} below). More generally, if we want to find a factor associated with the seminorm $\nnorm{\cdot}_{G_1, \ldots, G_s}^+$, then our definition of dual functions needs to include more complicated weights that turn out to be well-approximable by degree-$s$ nilsequences.


\begin{lemma}\label{L: algebra}
    Let $G_1, \ldots, G_s\subseteq\R^\ell$ be finitely generated subgroups. Then $Z_{G_1, \ldots, G_s}^+$ is a closed $\C$-algebra invariant under $T_1, \ldots, T_\ell$ and closed under conjugation.
\end{lemma}
\begin{proof}
    Recall that $Z_{G_1, \ldots, G_s}^+\subseteq L^\infty(\mu)$ is defined to be an $L^2(\mu)$-closed linear span. It is easy to observe from the definition that it is invariant under $T_1, \ldots, T_\ell$ and closed under conjugation. The only property that remains is that it is closed under multiplication. That is, given
    \begin{align*}
D_1 &=\E_{(\bm, \bm')\in G}\; \psi_1(\bm,\bm')
     \prod_{\ueps\in\{0,1\}^s_*}T^{\sfloor{\bm_1^{\eps_1}}-\floor{\bm_1}+\cdots + \sfloor{\bm_s^{\eps_s}}-\floor{\bm_s}} f_{1,\ueps}\\
    D_2 &=\E_{(\bn, \bn')\in G}\; \psi_2(\bn, \bn')
     \prod_{\ueps\in\{0,1\}^s_*}T^{\sfloor{\bn_1^{\eps_1}}-\floor{\bn_1}+\cdots + \sfloor{\bn_s^{\eps_s}}-\floor{\bn_s}} f_{2,\ueps}, 
    %  D_1 &=\E_{(\bm_1, \bm_1', \ldots, \bm_s, \bm_{s'})\in G}\; \psi_1(\bm_1,\bm_1',\ldots,\bm_s,\bm_s')
    %  \prod_{\ueps\in\{0,1\}^s_*}T^{\sfloor{\bm_1^{\eps_1}}-\floor{\bm_1}+\cdots + \sfloor{\bm_s^{\eps_s}}-\floor{\bm_s}} f_{1,\ueps}\\
    % D_2 &=\E_{(\bn_1, \bn_1', \ldots, \bn_s, \bn_{s'})\in G}\; \psi_2(\bn_1,\bn_1',\ldots,\bn_s,\bn_s')
    %  \prod_{\ueps\in\{0,1\}^s_*}T^{\sfloor{\bn_1^{\eps_1}}-\floor{\bn_1}+\cdots + \sfloor{\bn_s^{\eps_s}}-\floor{\bn_s}} f_{2,\ueps}, 
    \end{align*}
    we want to show that $D_1\cdot D_2$ is also in $Z_{G_1, \ldots, G_s}^+$ (above, we set $G = G_1\times G_1 \times \cdots \times G_s\times G_s$, $(\bm,\bm') = (\bm_1, \bm_1',\ldots, \bm_s, \bm_s')$ and $(\bn,\bn') = (\bn_1, \bn_1',\ldots, \bn_s, \bn_s')$, and we take $f_{j,\ueps}$ and $\psi_j$ to be 1-bounded).
    %$D_i = \E_m e(\xi_i m) T_1^{\floor{\alpha m}}T_2^{\floor{-\beta m}}f_i$ for $i\in[2]$, 
    We note first that the averages converge along any F\o lner sequence, and so $D_2$ will converge to the same limit if we replace all $\bn_i$ and $\bn_i'$ by $\bm_i+\bn_i$ and $\bm_i'+\bn_i'$ respectively for any $\bn_i,\bn_i'\in G_i$. 
    Multiplying the two limits together, we therefore get
    \begin{multline*}
        %D_1 D_2 = \E_{(\bm, \bm')\in G} \psi_1(\bm,\bm') \E_{(\bn, \bn')\in G} \psi_2(\bm+\bn, \bm'+\bn')\\
        %\prod_{\ueps\in\{0,1\}^s_*}T^{\sfloor{\bm_1^{\eps_1}}-\floor{\bm_1}+\cdots + \sfloor{\bm_s^{\eps_s}}-\floor{\bm_s}} f_{1,\ueps} \cdot 
 %T^{\sfloor{\bm_1^{\eps_1}+\bn_1^{\eps_1}}-\floor{\bm_1+\bn_1}+\cdots + \sfloor{\bm_s^{\eps_s}+\bn_s^{\eps_s}}-\floor{\bm_s+\bn_s}} f_{2,\ueps}\\
        D_1\cdot D_2 = \sum_{\br\in \{-1,0,1\}^{\ell s}}\E_{(\bm, \bm')\in G}  \E_{(\bn, \bn')\in G} \psi_1(\bm,\bm')\psi_2(\bm+\bn, \bm'+\bn')1_{E_\br}(\bm,\bm',\bn,\bn')\\
        \prod_{\ueps\in\{0,1\}^s_*}T^{\sfloor{\bm_1^{\eps_1}}-\floor{\bm_1}+\cdots + \sfloor{\bm_s^{\eps_s}}-\floor{\bm_s}} (f_{1,\ueps} \cdot T^{\sfloor{\bn_1^{\eps_1}} - \sfloor{\bn_1} + \cdots + \sfloor{\bn_s^{\eps_s}} - \sfloor{\bn_s}+ \ueps\cdot \br}f_{2,\ueps}),
 %T^{\sfloor{\bm_1^{\eps_1}+\bn_1^{\eps_1}}-\floor{\bm_1+\bn_1}+\cdots + \sfloor{\bm_s^{\eps_s}+\bn_s^{\eps_s}}-\floor{\bm_s+\bn_s}} f_{2,\ueps}.
    \end{multline*}
    where %the sum over $\br\in \Z^{ks}$ is finite, 
    $\ueps\cdot \br = \eps_1\cdot \br_1+\cdots + \eps_s\cdot\br_s$, and
    \begin{multline*}
        E_{\br} = \{(\bm,\bm',\bn,\bn')\in G\times G:\; \sfloor{\bm_i^{\eps_i}+\bn_i^{\eps_i}}-\sfloor{\bm_i + \bn_i}\\
        = \sfloor{\bm_i^{\eps_i}}-\sfloor{\bm_i}+\sfloor{\bn_i^{\eps_i}} - \sfloor{\bn_i} + \eps_i\cdot \br_{i}\; \textrm{for\; all}\; i, \eps_i\}.
    \end{multline*}

    Swapping the order of the averages along $(\bm, \bm')$ and $(\bn, \bn')$ {by the Fubini principle for ergodic averages \cite[Lemma 1.1]{BL15}}, the limit will follow if we can establish that the sequence
    \begin{align*}
        \psi_{\br, \bn,\bn'}(\bm,\bm') = \psi_1(\bm,\bm')\psi_2(\bm+\bn, \bm'+\bn')1_{E_\br}(\bm,\bm',\bn,\bn')
    \end{align*}
     is in $\Nil_s(G)$.  Since $\psi_1, \psi_2\in\Nil_s(G)$ by assumption, 
     %is an degree-$s$ nilsequence on $G$.  Since both $\psi_1, \psi_2$ are degree-$s$ nilsequences by assumption, 
     the result will follow from Corollary~\ref{C: algebra} once we show that $(\bm,\bm')\mapsto 1_{E_\br}(\bm,\bm',\bn,\bn')\in\Nil_1(G)$. 
     %{\color{red}Why nilsequences are closed under multiplications, as they may have different coordinate maps?} \KT{Is there any need to consider nilsequences with different coordinate maps, or can we simply fix the coordinate map in the definition \eqref{E: dual function} of dual functions?}  {\color{red} An (almost) equivalent question is: let $\phi,\phi'\colon\mathbb{Q}^{K}\to \text{Span}_{\mathbb{Q}}G$ be two coordinate maps and let $\psi\colon G\to\C$ be a map. Show that if $\psi\circ\phi\colon\phi^{-1}(G)\to\C$ is a $K$-parameter degree-$s$ nilsequence, then so is $\psi\circ\phi'$.} \BK{(I think the issue should be resolved by now)}
     Such a sequence is a product of sequences coming from the constraints
    \begin{align*}
        \sfloor{\bm_i^{\eps_i}+\bn_i^{\eps_i}}-\sfloor{\bm_i + \bn_i}
        = \sfloor{\bm_i^{\eps_i}}-\sfloor{\bm_i}+\sfloor{\bn_i^{\eps_i}} - \sfloor{\bn_i} + \eps_i\cdot \br_{i},
    \end{align*}
    and so it suffices to show that the indicator function of the set of $(\bm_i,\bm'_i)$ satisfying such a constraint for fixed $i,\eps_i,\bn,\bn'$ is a nilsequence. A standard argument shows that such an indicator function takes the form $F_{\bn_i,\bn_i'}(\srem{\bm_i}, \srem{\bm_i'})$ for some 1-bounded Riemann integrable function $F_{\bn_i,\bn_i'}:[0,1)^{2\ell}\to\R$, and hence it is indeed a degree-1 nilsequence. 
    \end{proof}

Using a standard identification between factors of a dynamical system and closed invariant subalgebras of its Hilbert space, we obtain the following corollary.
\begin{corollary}
Let $G_1, \ldots, G_s\subseteq\R^\ell$ be finitely generated subgroups. There exists a factor $\CZ_{G_1, \ldots, G_s}^+\subseteq \CX$ such that $L^\infty(\CZ_{G_1, \ldots, G_s}^+) = Z_{G_1, \ldots, G_s}^+$.
\end{corollary}
Whenever $G_i = \langle \bv_i\rangle$, we also write $\CZ^+_{G_1, \ldots, G_s} = \CZ^+_{\bv_1, \ldots, \bv_s}$, and if $G_1 = \cdots = G_s = G$, we additionally abbreviate $\CZ^+_{G_1, \ldots, G_s} = \CZ^+_{G^{\times s}}$.

% {\color{red} A general question: although in this paper we only used the $\mathcal{Z}^{+}$ factor and seminorms, in Sections 3 and 4, should we include both the $\mathcal{Z}^{+}$ version and the $\mathcal{Z}$ version? The reason I ask this is because that the weak structure theorem (Proposition 3.10) is not yet known in the literature for $\nnorm{\cdot}_{G_{1},\dots,G_{s}}$ for subgroups of $\R^{k}$, and that the new concatenation theorem (Theorem 4.1) is also not known for subgroups of $\R^{k}$ (in fact not known even for subgroups of $\Z^{k}$ since Theorem 4.1 is stronger than the result of Tao and Ziegler). On the other hand, these results could be useful in the future (and the proofs shouldn't be too long).}
% \BK{1) I do not know if the factor $\CZ$ is well-defined in general: see the discussion below Definition \ref{D: factor}. If it were defined, then we would most likely not need the factors $\CZ^+$ at all. 2) I added a remark below Corollary 4.2 regarding the statement of relative concatenation for subgroups of $\Z^k$. }


The next result relates the factor $\CZ_{G_1, \ldots, G_s}^+$ to the seminorm $\nnorm{\cdot}_{G_1, \ldots, G_s}^+$.
\begin{proposition}[Weak structure theorem for $\nnorm{\cdot}_{G_1, \ldots, G_s}^+$]\label{P: structure theorem}
    Let $G_1, \ldots, G_s\subseteq\R^\ell$ be finitely generated subgroups. For any $f\in L^\infty(\mu)$, the relationship \eqref{E: seminorm vs factor} holds.
\end{proposition}

Before we prove Proposition \ref{P: structure theorem}, we need to understand the structure of multicorrelation sequences appearing in the definition of generalized box seminorms.

\begin{definition}[Nullsequences on $G$]
    Let $G\subseteq\R^\ell$ be a finitely generated subgroup. For $\veps>0$, we call $\psi: G\to\C$ an \textit{$\veps$-nullsequence on $G$} if
    \begin{align*}
        \oE_{\bm\in G}|\psi(\bm)|^2\leq \veps.
%        \limsup_{N\to\infty}\sup_{(I_N)_N}\E_{\bm\in I_N}|\psi(\bm)|^2\leq \veps,
    \end{align*}
\end{definition}

\begin{proposition}[Decomposition of multicorrelation sequences]\label{P: multicorrelation}
    Let $G_1, \ldots, G_s\subseteq\R^\ell$ be finitely generated subgroups and set $G = G_1\times G_1 \times \cdots \times G_s\times G_s$. For $(\bm,\bm')\in G$ and  1-bounded functions $f_\ueps\in L^\infty(\mu)$, define 
    \begin{align*}
        %\psi(\bm_1, \bm_1', \ldots, \bm_s, \bm_s')
        \psi(\bm, \bm')= \int \prod_{\ueps\in\{0,1\}^s}T^{\sfloor{\bm_1^{\eps_1}}+\cdots + \sfloor{\bm_s^{\eps_s}}} f_\ueps\; d\mu.
    \end{align*}
    Then for every $\veps>0$, we can decompose $\psi = \psi_1 + \psi_2$, where $\psi_1$ is a 1-bounded degree-$s$ nilsequence on $G$ and $\psi_2$ is an $\veps$-nullsequence on $G$.
\end{proposition}
\begin{proof}
   Proposition \ref{P: multicorrelation} can be proved via a straightforward but tedious adaptation of the methods of \cite{Fr15, FH18, Ko18}. 
Let $\phi:\Z^K\to G$ be a coordinate map.
    By a natural multiparameter extension of \cite[Theorem 2.1]{Ko18}, it suffices to show that the sequence $\psi\circ \phi:\Z^K\to\C$ satisfies two conditions: 
    \begin{enumerate}
        \item it is \textit{$(s+1)$-weak-anti-uniform} in that for every uniformly bounded $b:\Z^K\to\C$, the average
    \begin{align*}
        \oE_{\um\in\Z^K}\psi\circ \phi(\um)b(\um)
    \end{align*}
    is quantitatively controlled by a degree-$(s+1)$ uniformity seminorm on $\Z^K$ (in the sense of \cite{Ko18});
        \item it is \textit{$(s+1)$-regular} in that the uniform limit
    \begin{align*}
        \E_{\um\in\Z^K}\psi\circ \phi(\um)b(\um)
    \end{align*}
    exists
    for any $s$-step nilsequence $b:\Z^K\to\C$.
    \end{enumerate}

%We notice here that for the map $\psi\circ \phi|_{\Z^K}$ to be well-defined, $\psi$ has to be extended to $\phi(\Z^K) = \widetilde{G}$ for some subgroup $G\subseteq\widetilde{G}\subseteq\R^k$ with $\widetilde{G}/G$ finite; this can be done in the natural way.
    
The condition of $(s+1)$-weak-anti-uniformity follows by combining a suspension flow trick \cite[Theorem~4.1]{Ko18} (which allows us to remove integer parts from expressions like $T^{\floor{\bm}}$) with a standard van der Corput argument (which can be found e.g. at the end of the proof of \cite[Theorem~2.6]{FH18}). Similarly, $(s+1)$-regularity can be deduced by combining \cite[Theorem~4.1]{Ko18} with \cite[Proposition~2.1]{FH18}. %(but for integer parts of real multivariable polynomials). 
%Analogously to \cite[Theorem~2.5]{FH18}, $s$-weak-anti-uniformity follows from the multivariable variant of \cite[Theorem~4.1]{Ko18} via the $s$-anti-uniformity of \cite{FH18} (for the latter, see the end of proof of \cite[Proposition~2.6]{FH18}) and $s$-regularity  as in \cite[Proposition~2.1]{FH18} (but for integer parts of real multivariable polynomials). 
For every $\veps>0$, the natural multiparameter extension of \cite[Theorem 2.1]{Ko18} then allows us to decompose $\psi\circ \phi = \psi_1+\psi_2$, where $\psi_1$ is a degree-$s$ nilsequence on $\Z^{K}$ and $\psi_2$ is an $\veps$-nullsequence on $\Z^K$. This gives the claimed result for ${\psi}.$
% We now specialize to $G$, defining $\psi_j = \psi_j'\circ\phi\inv|_{\phi\inv(G)}:G\to\C$ for $j=1,2$ so that $\psi = \psi_1 + \psi_2$. It is immediate from the definition that the sequence $\psi_1$ is a degree-$s$ nilsequence on $G$. It is however less immediate that $\psi_2$ defines a nullsequence on $G$. What we get from our assumption is that
% \begin{align*}
%     \oE_{\bm\in\widetilde{G}}\abs{\psi_2'\circ\phi\inv(\bm)}^2 = \oE_{\um\in\Z^K}\abs{\psi_2'(\um)}^2\leq \veps, 
% \end{align*}
% whereas what we want to show is the smallness of 
% \begin{align*}
%     \oE_{\bm\in{G}}\abs{\psi_2(\bm)}^2=\oE_{\bm\in{G}}\abs{\psi_2'\circ\phi\inv(\bm)}^2.
% \end{align*}
%
% To see that the latter average is still small, let $(I_N)_N$ be a F\o lner sequence on $G$ and $E\subseteq \widetilde{G}$ be a finite of size $|E|=|\widetilde{G}/G|$ that contains 0 and satisfies $\widetilde{G} = G +E$. Then $(I_N + E)_N$ is a F\o lner sequence on $\widetilde{G}$. The nonnegativity of the summands implies that
% \begin{align*}
%     \E_{\bm\in I_N}\abs{\psi_2'\circ\phi\inv(\bm)}^2 \leq \frac{1}{|I_N|}\sum_{\bm\in I_N + E}\abs{\psi_2'\circ\phi\inv(\bm)}^2\leq |E|\cdot \E_{\bm\in I_N + E}\abs{\psi_2'\circ\phi\inv(\bm)}^2 + o(1),
% \end{align*}
% and hence taking the limit as $N\to\infty$ and supremum over all F\o lner sequences, we deduce that
% \begin{align*}
%     \oE_{\bm\in{G}}\abs{\psi_2(\bm)}^2 \leq |E| \cdot \oE_{\bm\in\widetilde{G}}\abs{\psi_2'\circ\phi\inv(\bm)}^2\leq |E|\cdot \veps.
% \end{align*}
% Hence $\psi_2$ is a $|E|\veps$-nullsequence on $G$. Since $\veps>0$ was arbitrary, the result follows.
\end{proof}
% \begin{corollary}
%     For every $\veps>0$ we can find a 1-bounded $s$-step nilsequence $\psi_\veps>0$ such that
%     \begin{align*}
%         (\nnorm{f}_{G_1, \ldots, G_s}^+)^{2^{s+1}} = \int f\cdot \CD_{G_1, \ldots, G_s}(f; \psi_\veps)\; d\mu + O(\veps).
%     \end{align*}
% \end{corollary}
% \begin{proof}
%     Expanding the definition of $\nnorm{f}_{G_1, \ldots, G_s}^+$, we get that 
%     \begin{align*}
%         (\nnorm{f}_{G_1, \ldots, G_s}^+)^{2^{s+1}} = \int f\cdot \CD_{G_1, \ldots, G_s}(f; \psi)\; d\mu,
%     \end{align*}
%     where $f_{\ueps} = \CC^{|\ueps|}f$ and 
%     \begin{align*}
%         \psi(\bm,\bm') = \int \prod_{\ueps\in\{0,1\}^s}\CC^{|\ueps|}T^{\sfloor{\bm_1^{\eps_1}}+\cdots + \sfloor{\bm_s^{\eps_s}}} \overline{f}\, d\mu.
%     \end{align*}
%     Fix $\veps>0$, and let $\psi_\veps$ be a 1-bounded $s$-step nilsequence coming from Proposition \ref{P: multicorrelation} such that
%     \begin{align*}
%         \overline{\E}_{(\bm,\bm')\in G}\abs{\psi(\bm,\bm')-\psi_\veps(\bm,\bm')} \leq \veps^{2^{s+1}}.
%     \end{align*}
%     Then
%     \begin{align*}
%         \abs{(\nnorm{f}_{G_1, \ldots, G_s}^+)^{2^{s+1}} - \int f\cdot \CD_{G_1, \ldots, G_s}(f; \psi_\veps)\; d\mu}\leq \overline{\E}_{(\bm,\bm')\in G}\abs{\psi(\bm,\bm')-\psi_\veps(\bm,\bm')}\leq \veps^{2^{s+1}},
%     \end{align*}
%     and taking the root gives the lemma.
    
% \end{proof}


\begin{proof}[Proof of Proposition \ref{P: structure theorem}]
We show  that $\nnorm{f}_{G_1, \ldots, G_s}^+ = 0$ if and only if $f$ is orthogonal to dual functions $\CD_{G_1, \ldots, G_s}((f_\ueps)_{\{0,1\}^s_*}; \psi)$ with the functions $f_\ueps$ and degree-$s$ nilsequence being 1-bounded. The invariance of the measure gives
\begin{multline*}
     \int f \cdot \CD_{G_1, \ldots, G_s}((f_\ueps)_{\{0,1\}^s_*}; \psi)\;d\mu\\
     = \E_{\bm_1, \bm_1'\in G_1}\cdots \E_{\bm_s, \bm_s'\in G_s}\psi(\bm,\bm')\cdot\int \prod_{\ueps\in\{0,1\}^s}T^{\sfloor{\bm_1^{\eps_1}}+\cdots + \sfloor{\bm_s^{\eps_s}}} f_\ueps\, d\mu,
\end{multline*}
where $f_{\underline{0}} = f$.
    Using the Cauchy-Schwarz inequality, we can remove the nilsequence and bound
    \begin{multline*}
        \abs{\int f \cdot \CD_{G_1, \ldots, G_s}((f_\ueps)_{\{0,1\}^s_*}; \psi)\; d\mu}^2\\
        \leq  \E_{\bm_1, \bm_1'\in G_1}\cdots \E_{\bm_s, \bm_s'\in G_s}\abs{\int \prod_{\ueps\in\{0,1\}^s}T^{\sfloor{\bm_1^{\eps_1}}+\cdots + \sfloor{\bm_s^{\eps_s}}} f_\ueps\; d\mu}^2.
    \end{multline*}
    By the Gowers-Cauchy-Schwarz inequality (e.g. \cite[Lemma 3.3]{DKKST24}), the last expression is bounded by $(\nnorm{f}_{G_1, \ldots, G_s}^+)^{1/2^s}$. In particular, if the seminorm is 0, then $f$ is orthogonal to the function $\CD_{G_1, \ldots, G_s}((f_\ueps)_{\{0,1\}^s_*}; \psi)$.

    In the direction covered so far, $\psi(\bm,\bm')$ could be any bounded weight; its nilsequence structure will be used in the converse direction. Expanding the definition of $\nnorm{f}_{G_1, \ldots, G_s}^+$, we get 
    \begin{align*}
        (\nnorm{f}_{G_1, \ldots, G_s}^+)^{2^{s+1}} = \int f\cdot \CD_{G_1, \ldots, G_s}(f; \psi)\; d\mu
    \end{align*}
    for
    \begin{align*}
        \psi(\bm,\bm') = \int \prod_{\ueps\in\{0,1\}^s}\CC^{|\ueps|}T^{\sfloor{\bm_1^{\eps_1}}+\cdots + \sfloor{\bm_s^{\eps_s}}} \overline{f}\, d\mu.
    \end{align*}
    Fix $\veps>0$, and let $\psi_\veps$ be a 1-bounded degree-$s$ nilsequence coming from Proposition~\ref{P: multicorrelation} such that
    \begin{align*}
        \oE_{(\bm,\bm')\in G}\abs{\psi(\bm,\bm')-\psi_\veps(\bm,\bm')} \leq \veps^{2^{s+1}}.
    \end{align*}
    Then (assuming $f$ is 1-bounded, which we may)
    \begin{align*}
        \abs{(\nnorm{f}_{G_1, \ldots, G_s}^+)^{2^{s+1}} - \int f\cdot \CD_{G_1, \ldots, G_s}(f; \psi_\veps)\; d\mu}\leq \oE_{(\bm,\bm')\in G}\abs{\psi(\bm,\bm')-\psi_\veps(\bm,\bm')}\leq \veps^{2^{s+1}}.
    \end{align*}
    The orthogonality of $f$ to $\CD_{G_1, \ldots, G_s}(f; \psi_\veps)$ gives the upper bound $\nnorm{f}_{G_1, \ldots, G_s}^+\leq \veps$, and the result follows by letting $\veps\to 0$.
\end{proof}

As a consequence of relevant properties of seminorms (see Lemma \ref{L: basic properties of seminorms}), we obtain the following properties of factors.
\begin{lemma}[Basic properties of factors]\label{L: basic properties of factors}
   Let $G_1, \ldots, G_s\subseteq\R^\ell$ be finitely generated subgroups. Then: 
    % Let $G_1, G_1', \ldots, G_s, G_s'\subseteq\R^k$ be finitely generated subgroups with $G_i\subseteq G_i$ for every $i\in[s]$. Then: \AK{The following can definitely be written in a more compact way. E.g., for the second one we have to state it for $s=2.$ Think of the optimal way to do it.} \BK{Bro, things can only "definitely" be written in a more compact way if you already have a way to do so in mind, in which case please share.}
     \begin{enumerate}
        \item (Symmetry) For any permutation $\sigma:[s]\to[s]$, we have $$\CZ_{G_{\sigma(1)}, \ldots, G_{\sigma(s)}}^+ = \CZ_{G_1, \ldots, G_s}^+;$$
        \item (Monotonicity) For any finitely generated $G_{s+1}\subseteq\R^\ell$, we have $$\CZ_{G_1, \ldots, G_s}^+\subseteq\CZ_{G_1, \ldots, G_{s+1}}^+;$$
        \item (Subgroup property) If $G'_i\subseteq G_i$ for all $i\in[s]$, then $$\CZ_{G_1, \ldots, G_s}^+\subseteq \CZ_{G_1', \ldots, G_s'}^+,$$ with equality whenever $G_i/G_i'$ has finite index for all $i\in[s]$;
        \item (Rescaling property) If for every $i\in[s]$ we have
        \begin{align*}
            G_i = \langle \bg_{i1}, \ldots, \bg_{iK_i}\rangle\quad \textrm{and}\quad G'_i = \langle \alpha_{i1}\bg_{i1}, \ldots, \alpha_{iK_i} \bg_{iK_i}\rangle 
        \end{align*}
        for some $K_i\in\N$ and $\alpha_{ij}\neq 0$, then
    \begin{align*}
       \CZ^+_{G'_1, \ldots, G'_s}=\CZ^+_{G_1, \ldots, G_s}.
    \end{align*}
    \end{enumerate}
\end{lemma}

\section{Relative concatenation}\label{S: concatenation}

Having defined generalized box factors and established their various properties, we now show that factors along smaller groups can be merged (``concatenated'') into (higher-degree) factors along larger groups. Once again, we fix a system $(X, \CX, \mu, T_1, \ldots, T_\ell)$ throughout this section. The main result of this section is as follows.

%{\color{red} I assume that this result is for the $Z^{+}$ factor not the $Z$ factor?}
\begin{theorem}[Relative concatenation]\label{T: relative concatenation}
Let $s\in\N_0$, $d,d'\in\N$, and $$G_1, \ldots, G_s, H_1, \ldots, H_d, H'_1, \ldots, H'_{d'}\subseteq\R^\ell$$ be finitely generated subgroups. Then the following statements hold:
\begin{enumerate}
    \item We have
\begin{align*}
    \mathcal{Z}^+_{G_{1},\dots,G_s,H_{1},\dots,H_{d}}\cap \mathcal{Z}^+_{G_{1},\dots,G_s,H'_{1},\dots,H'_{d'}}\subseteq \mathcal{Z}^+_{G_{1},\dots,G_{s},(H_{i}+H'_{i'})_{i\in[d], i'\in[d']}}.
\end{align*}
    \item If moreover there exists a subgroup $H\subseteq \R^\ell$ such that $H_i + H'_{i'} = H$ for all $i\in[d]$ and $i'\in[d']$, then
\begin{align*}
    \mathcal{Z}^+_{G_{1},\dots,G_s, H_{1},\dots,H_{d}}\cap \mathcal{Z}^+_{G_{1},\dots,G_s, H'_{1},\dots,H'_{d'}} \subseteq \CZ^+_{G_{1},\dots,G_s, H^{\times (d+d'-1)}},
\end{align*}
		where we recall that $H^{\times d}$ denotes $d$-copies of $H$.
\end{enumerate}
\end{theorem}
When $s=0$ and $H_i, H'_{i'}\subseteq\Z^\ell$, Theorem \ref{T: relative concatenation} has been established by Tao and Ziegler in \cite[Theorems 1.15 and 1.16]{TZ16};\footnote{Tao-Ziegler's results apply more broadly to box factors along subgroups that are sums of finitely generated and profinite groups. For the second statement, their result additionally assumes that $H_1 = \cdots = H_d$ and $H'_1 = \cdots = H_{d'}'$.} it is their influential paper that introduced the idea of concatenating factors and seminorms to ergodic theory, additive combinatorics, and analytic number theory. We extend Tao-Ziegler's result in two ways: by relativizing it (i.e. allowing $s>0$ groups that do not get concatenated), and by extending it to generalized box factors.

We shall use the following consequence of Theorem~\ref{T: relative concatenation}, which generalizes \cite[Corollaries~1.20 and 1.21]{TZ16}, and which can be deduced from Theorem~\ref{T: relative concatenation} in the same way in which \cite[Corollaries~1.20 and 1.21]{TZ16} are derived from \cite[Theorems 1.15 and 1.16]{TZ16}.

\begin{corollary}\label{C: relative concatenation}
    Let $s\in\N_0$, $d,d'\in\N$, and $$G_1, \ldots, G_s, H_1, \ldots, H_d, H'_1, \ldots, H'_{d'}\subseteq\R^\ell$$ be finitely generated subgroups. Let $f\in L^\infty(\mu)$. Then the following statements hold:
    \begin{enumerate}
        \item If $\E(f|\mathcal{Z}^+_{G_{1},\dots,G_{s},(H_{i}+H'_{i'})_{i\in[d], i'\in[d']}})=0$, then we can decompose $f = g_1 + g_2$ for some $g_i\in L^\infty(\mu)$ satisfying
        \begin{align*}
            \E(g_1|\mathcal{Z}^+_{G_{1},\dots,G_s,H_{1},\dots,H_{d}}) = 0\quad\textrm{and}\quad \E(g_2|\mathcal{Z}^+_{G_{1},\dots,G_s,H'_{1},\dots,H'_{d'}}) = 0.
        \end{align*}
        \item If moreover there exists a subgroup $H\subseteq \R^\ell$ such that $H_i + H'_{i'} = H$ for all $i\in[d]$ and $i'\in[d']$, and if $\E(f|\CZ^{+}_{G_{1},\dots,G_s, H^{\times (d+d'-1)}})=0$, then we can decompose $f = g_1 + g_2$ for some $g_i\in L^\infty(\mu)$ satisfying
        \begin{align*}
            \E(g_1|\mathcal{Z}^+_{G_{1},\dots,G_s, H_{1},\dots,H_{d}}) = 0\quad\textrm{and}\quad \E(g_2|\mathcal{Z}^+_{G_{1},\dots,G_s,H'_{1},\dots,H'_{d'}})=0.
        \end{align*}
    \end{enumerate}
\end{corollary}

%If we specialize to subgroups of $\Z^k$ in Theorem \ref{T: relative concatenation} and Corollary \ref{C: relative concatenation}, the factors $\CZ^+$ can be replaced by the more classical factors $\CZ$, defined e.g. in \cite{HK05a, TZ16} and satisfying the property \eqref{E: seminorm vs factor unplussed}. {\color{red} should we briefly explain how to prove these results for the classical factors $\CZ$? (I think the major difference is the definition of dual functions, and how this affects the proof of Theorem \ref{T: fundamental concatenation})} \BK{See if you like this line, and if it is correct: "The proof in this case is almost identical; the only difference is that since we only need to consider dual functions with nilsequences $\psi = 1$, we can control expressions like $\int g \cdot \CD_{G_1, \ldots, G_s}f\; d\mu$ by $\nnorm{g}_{G_1, \ldots, G_s}$ rather than $\nnorm{g}_{G_1, \ldots, G_s}^+$."}  For general finitely generated subgroups $G_1, \ldots, G_s\subseteq\R^k$, it is however not clear that such factors exist, see the discussion below Definition \ref{D: factor}.

 
If we specialize to subgroups of $\Z^\ell$ in Theorem \ref{T: relative concatenation} and Corollary \ref{C: relative concatenation}, the factors $\CZ^+$ can be replaced by the more classical factors $\CZ$, defined e.g. in \cite{HK05a, TZ16} and satisfying the property \eqref{E: seminorm vs factor unplussed} (recall that in this case  $Z_{G_1, \ldots, G_s}(\CY)\subseteq L^\infty(\mu)$ is defined to be the $L^2(\mu)$-closed linear span of functions of the form 
$\CD_{G_1, \ldots, G_s}((f_\ueps)_{\{0,1\}^s_*})$ for $f_\ueps\in L^\infty(\CY, \mu)$). The proof in this case is almost identical; the only difference is that since we only need to consider dual functions with nilsequences $\psi = 1$, we can control expressions like $\int g \cdot \CD_{G_1, \ldots, G_s}f\; d\mu$ by $\nnorm{g}_{G_1, \ldots, G_s}$ rather than $\nnorm{g}_{G_1, \ldots, G_s}^+$.  We leave the details to the interested readers.

For general finitely generated subgroups $G_1, \ldots, G_s\subseteq\R^\ell$, such factors in general do not exist; see the discussions at the beginning of Section \ref{SS: factors} and below Definition \ref{D: factor}.


\subsection{Fundamental concatenation theorem}

Theorem \ref{T: relative concatenation} will be deduced by iterating the following result whose proof is the main subject of this subsection.
\begin{theorem}[Fundamental concatenation theorem]\label{T: fundamental concatenation}
Let $G_1, \ldots, G_s, G_s'\subseteq\R^\ell$ be finitely generated subgroups.
Then $$\mathcal{Z}^+_{G_1,\dots,G_{s-1},G_s}\cap \mathcal{Z}^+_{G_1,\dots,G_{s-1},G'_s}=\mathcal{Z}^+_{G_1,\dots,G_{s-1},G_s+G'_s}.$$ 
\end{theorem}

Before we supply the proof of Theorem \ref{T: fundamental concatenation}, we state several lemmas that will show up in the proof. The first one is a consequence of Proposition \ref{P: structure theorem} and follows from the same proof as \cite[Corollary 2.5]{TZ16}.
\begin{lemma}[Restriction to a factor]\label{L: substituting factor}
    Let $G_1, \ldots, G_s\subseteq\R^\ell$ be finitely generated subgroups. %\KT{$G'_s$ does not appear in this statement}. 
    For any factor $\CY\subseteq\CX$, we have
    \begin{align*}
        \CY \cap \CZ^+_{G_1, \ldots, G_s}(\CX) = \CZ^+_{G_1, \ldots, G_s}(\CY).
    \end{align*}
\end{lemma}
% \begin{proof}
%     We start with the converse direction. The inclusion $\CZ^+_{G_1, \ldots, G_s}(\CY)\subseteq \CZ^+_{G_1, \ldots, G_s}(\CX)$ is immediate. For the inclusion $\CZ^+_{G_1, \ldots, G_s}(\CY)\subseteq\CY$, we observe that for any $\CD_{G_1, \ldots, G_s}((f_\ueps)_{\{0,1\}^s_*}; \psi)\in Z^+_{G_1, \ldots, G_s}(\CY)$ and $f\in L^\infty(\mu)$, we have
%     \begin{align*}
%         \int f \cdot \CD_{G_1, \ldots, G_s}((f_\ueps)_{\{0,1\}^s_*}; \psi)\; d\mu = \int \E(f|\CY) \cdot \CD_{G_1, \ldots, G_s}((f_\ueps)_{\{0,1\}^s_*}; \psi)\; d\mu
%     \end{align*}
%     by $\CY$-measurability of $f_\ueps$'s, and hence $\CD_{G_1, \ldots, G_s}((f_\ueps)_{\{0,1\}^s_*}; \psi)\in L^\infty(\CY)$.
    
% \end{proof}

Before we state the second lemma, let $\CB_{G_1, \ldots, G_s}$ be the closed convex hull of dual functions $\CD_{G_1,\dots,G_{s}}((f_{\ueps})_{\ueps\in\{0,1\}^{s}_{\ast}};\psi)$ for 1-bounded functions $f_\ueps\in L^\infty(\mu)$ and 1-bounded degree-$s$ nilsequences $\psi$ on $G = G_1\times G_1\times \cdots \times G_s\times G_s$.
\begin{lemma}[Approximate multiplicative closedness of $\CB_{G_1, \ldots, G_s}$]\label{L: convex hull}
    %There exists $A=O_{k,s}(1)$ such that 
    Let $G_1, \ldots, G_s\subseteq\R^\ell$ be finitely generated subgroups.
    If $D_1, D_2\in \CB_{G_1, \ldots, G_s}$, then $D_1\cdot D_2\in 3^{\ell s}\cdot \CB_{G_1, \ldots, G_s}$.
\end{lemma}
\begin{proof}
    By convexity, it suffices to prove the result when $D_1, D_2$ are themselves dual functions $\CD_{G_1,\dots,G_{s}}((f_{\ueps})_{\ueps\in\{0,1\}^{s}_{\ast}};\psi)$ for 1-bounded functions $f_\ueps\in L^\infty(\mu)$ and 1-bounded %degree-$s$ nilsequences 
    $\psi\in\Nil_s(G)$ for $G = G_1\times G_1\times \cdots \times G_s\times G_s$. By inspecting the proof of Lemma~\ref{L: algebra}, we see that $D_1\cdot D_2$ is a sum of at most $3^{\ell s}$ elements in $\CB_{G_1, \ldots, G_s}$, corresponding to $3^{\ell s}$ possible values of $\br\in\{-1,0,1\}^{\ell s}$.
\end{proof}

\begin{lemma}[Anti-uniformity of $\CB_{G_1, \ldots, G_s}$]\label{L: correlation with dual}
    Let $G_1, \ldots, G_s\subseteq\R^\ell$ be finitely generated subgroups. %\KT{Again, $G_s'$ does not appear in the statement}. 
    If $f\in \CB_{G_1, \ldots, G_s}$, $g\in L^\infty(\mu)$, then $\abs{\int f\cdot g\; d\mu}\leq \nnorm{g}_{G_1, \ldots, G_s}^+$.
\end{lemma}
\begin{proof}
    By convexity, it suffices to consider the case when $f$ is a dual function, in which case the result follows from the Gowers-Cauchy-Schwarz inequality much like in the proof of Proposition~\ref{P: structure theorem}. 
\end{proof}
We are now ready to prove Theorem~\ref{T: fundamental concatenation}.
\begin{proof}[Proof of Theorem~\ref{T: fundamental concatenation}]
The converse inclusion is a straightforward consequence of the monotonicity property of factors (Lemma \ref{L: basic properties of factors} (ii)). For the forward inclusion, it suffices to show via Proposition~\ref{P: structure theorem} that any
\begin{align*}
    f\in L^\infty(\mathcal{Z}^+_{G_1,\dots,G_{s-1},G_s}\cap \mathcal{Z}^+_{G_1,\dots,G_{s-1},G'_s})
\end{align*}
is orthogonal to $g\in L^{\infty}(\mu)$ with $\nnorm{g}^+_{G_1,\dots, G_{s-1}, G_s+G_s'}=0$. Indeed, if the latter statement holds, then 
\begin{align*}
    0 &= \int f\cdot (\bar{f}-\mathbb{E}(\bar{f}\vert \CZ^+_{G_1,\dots, G_{s-1}, G_s+G_s'}))\ d\mu\\
    &= \norm{f}_{L^2(\mu)}^2 - \norm{\mathbb{E}(f\vert \CZ^+_{G_1,\dots, G_{s-1}, G_s+G_s'})}_{L^2(\mu)}^2,
\end{align*}
and so $f = \mathbb{E}(f\vert \CZ^+_{G_1,\dots, G_{s-1}, G_s+G_s'})\in Z^+_{G_1,\dots, G_{s-1}, G_s+G_s'}$. 

By an $L^2(\mu)$ approximation argument, it suffices to consider $$f = \CD_{G_1,\dots,G_{s-1},G_{s}}((f_{\ueps})_{\ueps\in\{0,1\}^{s}_{\ast}};\psi)$$ for 1-bounded functions $f_\ueps\in L^\infty(\mu)$ and %a 1-bounded degree-$s$ nilsequence 
$\psi\in\Nil_s(G)$ for $G = G_1\times G_1\times \cdots \times G_s\times G_s$. Since the nilsequence $\psi$ can be approximated in  $$\norm{v}_{\ell^2(G)}:=\Bigbrac{\oE_{\bm\in G}|v(\bm)|^2}^{1/2}$$ by a linear combination of nilsequences
\begin{align*}
  \psi_s(\bm_s,\bm_s') \psi_{s-1}(\bm_1, \bm_1', \ldots, \bm_{s-1},\bm_{s-1}')  
\end{align*}
(where $\bm_i,\bm_i'\in G_i$),
we can assume that $\psi$ takes such form.
%$\psi(\bm_1,\bm'_1,\ldots, \bm_s,\bm_s')$ 

Applying Lemma \ref{L: substituting factor} with $\CY =  \mathcal{Z}^+_{G_1,\dots,G_{s-1},G'_s}$,
%$$\CY = \mathcal{Z}^+_{G_1,\dots,G_{s-1},G_s}\cap  \mathcal{Z}^+_{G_1,\dots,G_{s-1},G'_s},$$
%{\color{red} here should be $\CY =  \mathcal{Z}^+_{G_1,\dots,G_{s-1},G'_s}$?}
we can take every $f_\veps$ to be $\mathcal{Z}^+_{G_1,\dots,G_{s-1},G'_s}$-measurable. 
% with respect to
% \begin{align}\label{E: intersection of factors}
%     \mathcal{Z}^+_{G_1,\dots,G_{s-1},G_s}\cap  \mathcal{Z}^+_{G_1,\dots,G_{s-1},G'_s}.    
% \end{align}
% \BK{(I think it suffices if it is $\mathcal{Z}^+_{G_1,\dots,G_{s-1},G'_s}$-measurable?)}
A similar approximation argument as before allows us to take each $f_\ueps$ to equal
$f_{\ueps}=\mathcal{D}_{G_1,\dots,G_{s-1},G_{s}'}((f_{\ueps,\ueps'})_{\ueps'\in\{0,1\}^{s}_{\ast}}; \psi_{\ueps})$ for some $\mathcal{Z}^+_{G_1,\dots,G_{s-1},G'_s}$-measurable functions and nilsequences $\psi_\ueps\in\Nil_s(G)$, all of which can be assumed to be 1-bounded. 

Inspecting the definition of dual functions, we observe that $f$ can be recast  as
\begin{multline*}
    f = \E_{\bm_s,\bm_s'\in G_s} \psi_s(\bm_s,\bm_s')\cdot T^{\sfloor{\bm_s'}-\sfloor{\bm_s}}f_{\underline{0}1} \\ \CD_{G_1,\dots,G_{s-1}}((f_{\ueps 0}\cdot T^{\sfloor{\bm_s'}-\sfloor{\bm_s}} f_{\ueps 1})_{\ueps\in\{0,1\}^{s-1}_{\ast}}; \psi_{s-1}).
\end{multline*}
Arguing similarly, we can recast $f_{\underline{0}1}$ as 
\begin{multline*}
    f_{\underline{0}1} = \E_{\bn_s,\bn_s'\in G_s'} \psi_{\underline{0}1, s}(\bn_s,\bn_s')\cdot T^{\sfloor{\bn_s'}-\sfloor{\bn_s}}f_{\underline{0}1, \underline{0}1} \\ \CD_{G_1,\dots,G_{s-1}}((f_{\underline{0}1, \ueps 0}\cdot T^{\sfloor{\bn_s'}-\sfloor{\bn_s}} f_{\underline{0}1, \ueps 1})_{\ueps\in\{0,1\}^{s-1}_{\ast}};\psi_{\underline{0}1,s-1}).
\end{multline*}
Plugging in the formula for $f_{\underline{0}1}$ into the formula for $f$ and setting
\begin{multline*}
    D_{\bm_s,\bm_s', \bn_s,\bn_s'} = \psi_s(\bm_s,\bm_s')\cdot\CD_{G_1,\dots,G_{s-1}}((f_{\ueps 0}\cdot T^{\sfloor{\bm_s'}-\sfloor{\bm_s}} f_{\ueps 1})_{\ueps\in\{0,1\}^{s-1}_{\ast}}; \psi_{s-1})\\
    \psi_{\underline{0}1, s}(\bn_s,\bn_s')\cdot T^{\sfloor{\bm_s'}-\sfloor{\bm_s}}\CD_{G_1,\dots,G_{s-1}}((f_{\underline{0}1, \ueps 0}\cdot T^{\sfloor{\bn_s'}-\sfloor{\bn_s}} f_{\underline{0}1, \ueps 1})_{\ueps\in\{0,1\}^{s-1}_{\ast}};\psi_{\underline{0}1,s-1}),
    %     D_{\bm_s,\bm_s', \bn_s,\bn_s'} = \CD_{G_1,\dots,G_{s-1}}((f_{\ueps 0}\cdot T^{\sfloor{\bm_s'}-\sfloor{\bm_s}} f_{\ueps 1})_{\ueps\in\{0,1\}^{s-1}_{\ast}}; \psi_{s-1})\\
    % T^{\sfloor{\bm_s'}-\sfloor{\bm_s}}\CD_{G_1,\dots,G_{s-1}}((f_{\underline{0}1, \ueps 0}\cdot T^{\sfloor{\bn_s'}-\sfloor{\bn_s}} f_{\underline{0}1, \ueps 1})_{\ueps\in\{0,1\}^{s-1}_{\ast}};\psi_{\underline{0}1,s-1}),
\end{multline*}
% \begin{align*}
%     D_{1, \bm_s,\bm_s'} &= \CD_{G_1,\dots,G_{s-1}}((f_{\ueps 0}\cdot T^{\sfloor{\bm_s'}-\sfloor{\bm_s}} f_{\ueps 1})_{\ueps\in\{0,1\}^{s-1}_{\ast}}; \psi_{s-1})\\
%     D_{1, \bm_s,\bm_s'} &= T^{\sfloor{\bm_s'}-\sfloor{\bm_s}}\CD_{G_1,\dots,G_{s-1}}((f_{\underline{0}1, \ueps 0}\cdot T^{\sfloor{\bn_s'}-\sfloor{\bn_s}} f_{\underline{0}1, \ueps 1})_{\ueps\in\{0,1\}^{s-1}_{\ast}};\psi_{\underline{0}1,s-1}),
% \end{align*}
we deduce that
\begin{align*}
    f = \E_{\bm_s,\bm_s'\in G_s} \E_{\bn_s,\bn_s'\in G_s'} T^{\sfloor{\bm_s'} +\sfloor{\bn_s'}- \sfloor{\bm_s}-\sfloor{\bn_s}}f_{\underline{0}1, \underline{0}1} \cdot D_{\bm_s,\bm_s', \bn_s,\bn_s'}.
%    f = \E_{\bm_s,\bm_s'\in G_s} \E_{\bn_s,\bn_s'\in G_s'} \widetilde\psi_s(\bm_s,\bm_s',\bn_s,\bn_s')\cdot T^{\sfloor{\bm_s'} +\sfloor{\bn_s'}- \sfloor{\bm_s}-\sfloor{\bn_s}}f_{\underline{0}1, \underline{0}1} \cdot D_{\bm_s,\bm_s', \bn_s,\bn_s'}
\end{align*}
%for some degree-$s$ nilsequence $\widetilde\psi_s$ on $G_s\times G_s\times G_s'\times G_s'$. 
%for some 1-bounded dual functions $D_{\bm_s,\bm_s', \bn_s,\bn_s'}$ along $G_1, \ldots, G_{s-1}$.
By incorporating emerging error terms into $D_{\bm_s,\bm_s', \bn_s,\bn_s'}$, we can reduce to the case where $f$ equals
\begin{align*}
    f = \E_{\bm_s,\bm_s'\in G_s} \E_{\bn_s,\bn_s'\in G_s'} %\widetilde\psi_s(\bm_s,\bm_s',\bn_s,\bn_s')\cdot 
    T^{\sfloor{\bm_s'+\bn_s'}- \sfloor{\bm_s+\bn_s}}f_{\underline{0}1, \underline{0}1} \cdot D_{\bm_s,\bm_s', \bn_s,\bn_s'}.   
    %f = \E_{\bm_s,\bm_s'\in G_s} \E_{\bn_s,\bn_s'\in G_s'} \widetilde\psi_s(\bm_s,\bm_s',\bn_s,\bn_s')\cdot T^{\sfloor{\bm_s'+\bn_s'}- \sfloor{\bm_s+\bn_s}}f_{\underline{0}1, \underline{0}1} \cdot D_{\bm_s,\bm_s', \bn_s,\bn_s'}    
\end{align*}

The point of all these maneuvers is that we have by now expressed $f$ as (essentially) an ergodic average along $G_s + G_s'$ twisted by %an degree-$s$ nilsequence $\widetilde\psi_s$ and a 
1-bounded $\CZ^+_{G_1, \ldots, G_{s-1}}$-measurable functions $D_{\bm_s,\bm_s', \bn_s,\bn_s'}$. By removing the dual function in a more or less standard manner, we will be able to control $\int f\cdot g\; d\mu$ by $\nnorm{g}^+_{G_1,\dots, G_{s-1}, G_s+G_s'}$, and this will conclude the proof of the forward direction.

Using the invariance of measure and the Cauchy-Schwarz inequality, we can bound
\begin{align*}
    \abs{\int f\cdot g\; d\mu} &= \abs{\int f_{\underline{0}1, \underline{0}1} \cdot \E_{\bm_s,\bm_s'\in G_s} \E_{\bn_s,\bn_s'\in G_s'} %\widetilde\psi_s(\bm_s,\bm_s',\bn_s,\bn_s')\cdot 
    T^{\sfloor{\bm_s+\bn_s}-\sfloor{\bm_s'+\bn_s'}}(g\cdot D_{\bm_s,\bm_s', \bn_s,\bn_s'})\; d\mu}\\
    &\leq \norm{\E_{\bm_s,\bm_s'\in G_s} \E_{\bn_s,\bn_s'\in G_s'} %\widetilde\psi_s(\bm_s,\bm_s',\bn_s,\bn_s')\cdot 
    T^{\sfloor{\bm_s+\bn_s}-\sfloor{\bm_s'+\bn_s'}}(g\cdot D_{\bm_s,\bm_s', \bn_s,\bn_s'})}_{L^2(\mu)}.
\end{align*}
Applying the van der Corput lemma (e.g. \cite[Lemma 2.1]{DKS22}), composing the resulting integral with $T^{\sfloor{\bm_s'+\bn_s'}- \sfloor{\bm_s+\bn_s}}$, handling error terms, and then pigeonholing in $\bm_s,\bm_s',\bn_s,\bn_s'$, we deduce that the left-hand side vanishes whenever
\begin{align}\label{E: after vdC}
    \E_{\bh_s,\bh_s'\in G_s} \E_{\bk_s,\bk_s'\in G_s'}\abs{\int D_{\substack{\bh_s,\bh_s', \bk_s,\bk_s'}} \cdot g\cdot T^{\sfloor{\bh_s+\bk_s}-\sfloor{\bh_s'+\bk_s'}}\overline{g}\; d\mu}^2
%    \E_{\bh_s,\bh_s'\in G_s} \E_{\bk_s,\bk_s'\in G_s'}\abs{\int \E_{\bm_s,\bm_s'\in G_s} \E_{\bn_s,\bn_s'\in G_s'} D_{\substack{\bm_s,\bm_s', \bn_s,\bn_s',\\ \bh_s,\bh_s', \bk_s,\bk_s'}}  T^{\sfloor{\bh_s+\bk_s}-\sfloor{\bh_s'+\bk_s'}}g}
\end{align}
vanishes for some
%where $D_{\substack{\bm_s,\bm_s', \bn_s,\bn_s',\\ \bh_s,\bh_s', \bk_s,\bk_s'}}$ is 
1-bounded $\CZ^+_{G_1, \ldots, G_{s-1}}$-measurable $D_{\substack{\bh_s,\bh_s', \bk_s,\bk_s'}}$.

The last step is to apply Lemma \ref{L: correlation with dual} and \cite[Lemma 4.1]{DKKST24} in order to bound the expression above in terms of 
\begin{align*}
    \E_{\bh_s,\bh_s'\in G_s} \E_{\bk_s,\bk_s'\in G_s'}\brac{\nnorm{g\cdot T^{\sfloor{\bh_s+\bk_s}-\sfloor{\bh_s'+\bk_s'}}\overline{g}}_{G_1, \ldots, G_{s-1}}^+}^{2^s} = \brac{\nnorm{g}_{G_1, \ldots, G_{s-1}, G_s+G_s'}^+}^{2^{s+1}}.
\end{align*}
Analyzing how we arrived at \eqref{E: after vdC}, we see that each $D_{\substack{\bh_s,\bh_s', \bk_s,\bk_s'}}$ is a linear combination of $O_{\ell,s}(1)$ 1-bounded degree-$s$ nilsequences times 4 dual functions from $\CB_{G_1, \ldots, G_{s-1}}$. Using Lemma~\ref{L: convex hull}, we can therefore find absolute $C=C(\ell,s)>0$ such that $$D_{\substack{\bh_s,\bh_s', \bk_s,\bk_s'}}\in C\cdot \CB_{G_1, \ldots, G_{s-1}},$$ and so by Lemma~\ref{L: correlation with dual}, the average \eqref{E: after vdC} vanishes whenever $\nnorm{g}_{G_1, \ldots, G_{s-1}, G_s+G_s'}^+=0.$ This concludes the proof.
\end{proof}

\subsection{Induction scheme}

The next goal is to present an induction scheme that allows us to upgrade Theorem \ref{T: fundamental concatenation} to Theorem \ref{T: relative concatenation}. Our induction is based on the following formalism.
\begin{definition}\label{D: CIS}
   Let $G$ be an abelian group, $\text{sub}(G)$ be the set of all subgroups of $G$, and $\text{sub}(G)^{\omega}:=\sqcup_{n=1}^{\infty}\text{sub}(G)^{n}$. We say that a subset $\mathcal{C}$ of $\text{sub}(G)^{\omega}$ is a \emph{CMS-family} if the followings properties hold:
   \begin{enumerate}
 		\item (Concatenation) For all $d\in\mathbb{N}$ and subgroups $H_{1},\dots,H_{d},H'_{d}\subseteq G$,
 		$$(H_{1},\dots,H_{d-1},H_{d}), (H_{1},\dots,H_{d-1},H'_{d})\in \mathcal{C}\Rightarrow (H_{1},\dots,H_{d-1},H_{d}+H'_{d})\in \mathcal{C}.$$
 		\item (Monotonicity) For all $d\in\mathbb{N}$ and subgroups $H_{1},\dots,H_{d}\subseteq G$,  
 		$$(H_{1},\dots,H_{d-1})\in \mathcal{C}\Rightarrow (H_{1},\dots,H_{d-1},H_{d})\in \mathcal{C}.$$
 		\item (Symmetry) For all $d\in\mathbb{N}$, subgroups $H_{1},\dots,H_{d}\subseteq G$ and permutations $\tau:[d]\to[d]$, we have
        %$\tau\colon\{1,\dots,d\}\to \{1,\dots,d\}$,  
 		$$(H_{1},\dots,H_{d})\in \mathcal{C}\Rightarrow (H_{\tau(1)},\dots,H_{\tau(d)})\in \mathcal{C}.$$
 	\end{enumerate}
 %   We say that $\mathcal{C}$ is a \emph{anti-CIS-family} if the all implications $\Rightarrow$ are in the reverse order $\Leftarrow$.
\end{definition}
%{\color{red} Maybe we should define the three properties in a way consistent to Lemma 2.15. In Lemma 2.15, (ii) is called monotonicity and (iii) is called permutation invariant. We could change permutation invariance to symmetry in Lemma 2.15. For Monotonicity vs Inclusion, perhaps Monotonicity is better?} \BK{done} {\color{red} I also changed CIS to CMS throughout}

 %For example, let $(X,\mathcal{X},\mu,(T_{n})_{n\in\mathbb{Z}^{d}})$ be a system and $f\in L^{\infty}(\mu)$.  Let $\mathcal{C}_{f}$ denote the set of all tuples of subgroups $(H_{1},\dots,H_{k})$ of $G$ such that $f$ is measurable with respect to the factor $\mathcal{Z}_{H_{1},\dots,H_{k}}^+$.
For $f\in L^\infty(\mu)$, consider the following family of subgroups:
\begin{align*}
    \CC_f \colon= \{(G_1, \ldots, G_s)\in \text{sub}(\R^\ell)^\omega:\; f\in Z^+_{G_1,\ldots, G_s}\}.
\end{align*}
As it turns out, $\CC_f$ satisfies the conditions of Definition \ref{D: CIS}.
\begin{proposition}\label{P: C_f is CIS}
    For every $f\in L^\infty(\mu)$, the collection $\mathcal{C}_{f}$ is a CMS-family. 
\end{proposition}
\begin{proof}
    The Monotonicity and Symmetry properties follow from Lemma \ref{L: basic properties of factors} while the Concatenation property follows from Theorem \ref{T: fundamental concatenation}.
\end{proof}

We will derive Theorem \ref{T: relative concatenation} by combining Proposition \ref{P: C_f is CIS} (which one could see as the ``base case'' of induction) with Proposition \ref{P: CIS induction} below (which provides the inductive scheme).
\begin{proposition}\label{P: CIS induction}
	Let $s\in\N_0$, $d,d'\in\N$, and $G_1, \ldots, G_s, H_1, \ldots, H_d, H_1', \ldots, H_{d'}'$ be subgroups of an abelian group $G$.
    %$G$ be an abelian group $d,d'\in\mathbb{N}$, and $G_{i}, 1\leq i\leq \ell$, $H_{i},H'_{i'}$, $i\in[d]$, $1\leq i'\leq d'$ be subgroups of $G$. 
    Suppose that $\mathcal{C}$ is a CMS-family. Then the following statements hold:
\begin{enumerate}
    \item We have 
    \begin{gather*}
        (G_{1},\dots,G_{s},H_{1},\dots,H_{d}),(G_{1},\dots,G_{s},H'_{1},\dots,H'_{d'})\in \mathcal{C}\\
        \Rightarrow(G_{1},\dots,G_{s},(H_{i}+H'_{i'})_{i\in[d], i'\in[d']})\in \mathcal{C}.
    \end{gather*}
    \item If $H_{i}+H'_{i'}=H$ for all $i\in[d], i'\in[d']$ for some subgroup $H\subseteq G$, then 
    \begin{gather*}
        (G_1, \ldots, G_s, H_{1},\dots,H_{d}), (G_1, \ldots, G_s, H'_{1},\dots,H'_{d'})\in \mathcal{C}\\
        \Rightarrow (G_1, \ldots, G_s, H^{\times (d+d'-1)})\in \mathcal{C}.
    \end{gather*}
\end{enumerate}
		
%		
%		Similarly, if $\mathcal{C}$ is an anti-CMS-family, then the same conclusions hold with all implications $\Rightarrow$ changed to the reverse order $\Leftarrow$.
\end{proposition}

\begin{proof}[Proof of Theorem \ref{T: relative concatenation} assuming Proposition \ref{P: CIS induction}]
    Let $f\in L^\infty(\mu)$. To prove the first claim in the language of CMS-families, it suffices to show that if $\CC_f$ contains 
    \begin{align*}
        (G_{1},\dots,G_s,H_{1},\dots,H_{d})\quad \textrm{and}\quad (G_{1},\dots,G_s,H'_{1},\dots,H_{d}'),
    \end{align*}
    then it also contains 
    \begin{align*}
        (G_{1},\dots,G_{s},(H_{i}+H'_{i'})_{i\in[d], i'\in[d']}).
    \end{align*}
    This follows from the fact that $\CC_f$ is a CMS-family (Proposition \ref{P: C_f is CIS}) and the induction scheme given by the first part of Proposition \ref{P: CIS induction}. Similarly, the second claim follows from Proposition \ref{P: C_f is CIS} and the second part of Proposition \ref{P: CIS induction}.
\end{proof}

We now prove Proposition \ref{P: CIS induction}.

\begin{proof}[Proof of Proposition \ref{P: CIS induction}]

To simplify the notation, we write $[H_{1},\dots,H_{d}]$ if $(H_{1},\dots,H_{d})$ belongs to $\mathcal{C}$, and we write $A+B$ to indicate that claims $A$ and $B$ both hold.

We start with the first part of the proposition.  Our claim can then be rephrased in this language as
\begin{align*}
    [G_1, \ldots, G_s, H_1, \ldots, H_d] + [G_1, \ldots, G_s, H'_1, \ldots, H_d'] \Rightarrow   [G_1, \ldots, G_s, (H_{i}+H'_{i'})_{i\in[d], i'\in[d']}].
\end{align*}

For any $(d,d')\in\N^2$, we say that $(d,d')$ is \textit{good} if the implication
	\begin{equation}\label{E: (d,d')}
	\begin{split}
	[\bold{J},H_{1},\dots,H_d]+[\bold{J},H'_{1},\dots,H'_{d'}]
	\Rightarrow[\bold{J},(H_{i}+H'_{i'})_{i\in[d], i'\in[d']}]
	\end{split}
	\end{equation} 
	holds for any finite multifamily of subgroups $\bold{J}=(K_{j})_{j\in J}$ of $G$ (that is, $\vert J\vert<\infty$, elements are counted with multiplicity, and the family may be empty).
	Then our goal is to show that every $(d,d')\in\N^2$ is good. We will prove this by double induction, first showing that all pairs $(d,1)$ and $(1,d')$ are good and then establishing the statement for arbitrary $(d,d')\in\N^2$.
    
    That $(1,1)$ is good is equivalent to the Concatenation property. Suppose that $(d,1)$ is good for some $d\in\N$, i.e. 
    \begin{align}\label{E: (d,1)}
        [\bJ, H_1, \ldots, H_d] + [\bJ, H_1'] \Rightarrow [\bJ, (H_i+H_1')_{i\in[d]}]
    \end{align}
    holds for all $\bJ$. We want to show that $(d+1,1)$ is good as well. Applying Monotonicity, and then Symmetry as well as the implication \eqref{E: (d,1)} for the family $\bJ\cup\{H_{d+1}\}$, we obtain the implication
    %we have $[\bJ, H_1']\Rightarrow [\bj, H_1', H_{d+1}]$, and hence
    \begin{align*}
        [\bJ, H_1, \ldots, H_{d+1}] + [\bJ, H_1'] &\Rightarrow [\bJ, H_1, \ldots, H_{d+1}] + [\bJ, H_1', H_{d+1}]\\
        &\Rightarrow [\bJ, H_{d+1}, (H_i+H_1')_{i\in[d]}].
    \end{align*}
    It remains to concatenate $H_1'$ with $H_{d+1}$. By Monotonicity, followed by Symmetry and Concatenation,
    %\eqref{E: (d,1)} for the family $\bJ\cup\{H_i+H_1':\; i\in[d]\}$, 
    we deduce that
    \begin{align*}
        [\bJ, H_{d+1}, (H_i+H_1')_{i\in[d]}] + [\bJ, H_1'] &\Rightarrow [\bJ, H_{d+1}, (H_i+H_1')_{i\in[d]}] + [\bJ, H_1', (H_i+H_1')_{i\in[d]}]\\
        &\Rightarrow [\bJ, (H_i+H_1')_{i\in[d+1]}].
    \end{align*}
    Hence $(d+1,1)$ is good. By induction and symmetry, $(d,1)$ and $(1,d')$ are good for any $d,d'\in\N$.

       Now, suppose that $(d,d')\in\N^2$ is good; we want to show that $(d+1,d')$ is good. In doing so, we will inductively apply the goodness of $(d,d')$ and $(1,d')$. Applying Monotonicity first and then Symmetry and \eqref{E: (d,d')} for the family $\bJ\cup\{H_{d+1}\}$ (this is where we apply the case $(d,d')$), we get
       \begin{align*}
           [\bold{J},H_{1},\dots,H_{d+1}]+[\bold{J},H'_{1},\dots,H'_{d'}] &\Rightarrow [\bold{J},H_{1},\dots,H_{d+1}]+[\bold{J},H'_{1},\dots,H'_{d'}, H_{d+1}]\\
           &\Rightarrow [\bold{J}, H_{d+1}, (H_{i}+H'_{i'})_{i\in[d], i'\in[d']}].
       \end{align*}
        It remains to concatenate $H_{d+1}$ with the subgroups $H'_{i'}$. By Monotonicity, followed by Symmetry and  \eqref{E: (d,d')} for the family $\bJ\cup\{H_{i}+H'_{i'}:\; i\in[d], i'\in[d']\}$ (this is where we apply the case $(1,d')$), we obtain
        \begin{multline*}
            [\bold{J}, H_{d+1}, (H_{i}+H'_{i'})_{i\in[d], i'\in[d']}] + [\bold{J},H'_{1},\dots,H'_{d'}]\\ \Rightarrow [\bold{J}, H_{d+1}, (H_{i}+H'_{i'})_{i\in[d], i'\in[d']}] + [\bold{J},H'_{1},\dots,H'_{d'}, (H_{i}+H'_{i'})_{i\in[d], i'\in[d']}]\\
            \Rightarrow [\bold{J}, (H_{i}+H'_{i'})_{i\in[d+1], i'\in[d']}].
        \end{multline*}
        This establishes the case $(d+1,d')$. The case $(d,d'+1)$ follows by an identical argument, and hence the claim holds by induction.

        \

	We move on to the second part of the proposition. %This time, we write $[H_{1},\dots,H_{d}]$ if $(H_{1},\dots,H_{d})$ belongs to $\mathcal{C}$.
	Our starting assumptions can then be rephrased %in this language 
    as $[\bJ, H_{1},\dots,H_{d}]$ and $[\bJ, H'_{1},\dots,H'_{d'}]$, where this time we fix $\bJ = \{G_1, \ldots, G_s\}$.
	For $(i,i')\in([0,d]\times[0,d'])\backslash\{(0,0)\}$, we say that $(i,i')$ is \emph{good} if 
    $$[\bJ, H_{1},\dots,H_{d-i},H'_{1},\dots,H'_{d'-i'},H^{\times (i+i'-1)}].$$ 
    Then our goal is to show that $(d,d')$  is good, which will follow from the following subclaims:
    \begin{enumerate}
        \item if $(i+1,i')$ and $(i,i'+1)$ are good for some  $i\in[0,d-1], i'\in[0,d'-1]$, then so is $(i+1,i'+1)$;
        \item $(i,0)$ and $(0,i')$ are good for all $i\in[d], i'\in[d']$.
    \end{enumerate}
      It is not hard to see that both claims imply together that $(d,d')$ is good; to see this, visualize $(i,i')$ as lattice points on the rectangle $([0,d]\times[0,d'])\backslash\{(0,0)\}$.
	
	% If $(i+1,i')$ and $(i,i'+1)$ are good, then we have
 %    \begin{align*}
 %        &[\bJ, H_{1},\dots,H_{d-(i+1)},H'_{1},\dots,H'_{d'-(i'+1)},H'_{d'-i'},H^{\times (i+i')}]\\
 %        \textrm{and}\quad &[\bJ, H_{1},\dots,H_{d-(i+1)},H'_{1},\dots,H'_{d'-(i'+1)},H_{d-i},H^{\times (i+i')}].    
 %    \end{align*}
 %    {\color{red} maybe these two lines can be removed} \AK{I agree}
%    \BK{(I corrected $H^{\times (i+i'-1)}$ to $H^{\times (i+i')}$ in both statements above)}
     For the first subclaim, we infer from the Symmetry and Concatenation properties that
	\begin{equation}\nonumber
	\begin{split}
	&\quad [\bJ, H_{1},\dots,H_{d-(i+1)},H'_{1},\dots,H'_{d'-(i'+1)},H'_{d'-i'},H^{\times (i+i')}]
	\\&+[\bJ, H_{1},\dots,H_{d-(i+1)},H'_{1},\dots,H'_{d'-(i'+1)},H_{d-i},H^{\times (i+i')}]
	\\& \Rightarrow[\bJ, H_{1},\dots,H_{d-(i+1)},H'_{1},\dots,H'_{d'-(i'+1)},H^{\times(i+i'+1)}],
	\end{split}
	\end{equation} 
	which implies that $(i+1,i'+1)$ is good whenever $(i+1,i')$ and $(i,i'+1)$ are. 
	
	To show that $(0,i')$ is good, we use the Monotonicity property to deduce that
	$$[\bJ, H_{1},\dots,H_{d}]\Rightarrow [\bJ, H_{1},\dots,H_{d},H'_{1},\dots,H'_{d'-i'},H^{\times(i'-1)}]$$
	for all $i'\in[d']$. Since $[\bJ, H_{1},\dots,H_{d}]$ by assumption, we deduce that $(0,i')$ is good for all $i'\in[d']$. By an analogous argument, $(i,0)$ is good for all $i\in[d]$. This proves the second subclaim and thus finishes the proof.
%
 %        We now prove (1). This time, we write $[H_{1},\dots,H_{d}]$ if $(G_{1},\dots,G_{\ell},H_{1},\dots,H_{d})$ belongs to $\mathcal{C}$, and we say that $(u,u')$ is good if the implication
	% \begin{equation}\label{temp:1}
	% \begin{split}
	% [\bold{J},H_{1},\dots,H_{u}]+[\bold{J}',H'_{1},\dots,H'_{u'}]
	% \Rightarrow[\bold{J},\bold{J}',(H_{i}+H'_{i'})_{1\leq i\leq u, 1\leq i'\leq u'}]
	% \end{split}
	% \end{equation} 
	% holds for any finite multifamilies of subgroups $\bold{J}=(K_{j})_{j\in J},$ $\bold{J}'=(K'_{j})_{j\in J'}$ of $G$ (that is, $\max(\vert J\vert,\vert J'\vert)<\infty$, elements are counted with multiplicity, and the families may be empty)
	% Then our goal is to show that $(d,d')$ is good.
    %
	% First, observe that by applying first Inclusion and Symmetry first and then Concatenation, we get
	% \begin{equation}\nonumber
	% \begin{split}
	% [\bold{J},H_{1}]+[\bold{J}',H'_{1}]\Rightarrow  [\bold{J},\bold{J}',H_{1}]+[\bold{J},\bold{J}',H'_{1}]
	% \Rightarrow[\bold{J},\bold{J}',H_{1}+H'_{1}].
	% \end{split}
	% \end{equation} 
	% Thus, (1,1) is good.
%	
	% Suppose now that $(u,u')$ is good for some $u,u'\geq 1$.
	% We wish to prove that $(u,u'+1)$ is good as well. From now on, fix  two families $\bold{J}$ and $\bold{J}'$, and suppose that  $[\bold{J},H_{1},\dots,H_{u}]$ and $[\bold{J}',H_{1},\dots,H'_{u'+1}]$. Since $(u,u')$ is good, applying  (\ref{temp:1}) to $\bold{J},\bold{J}'\cup \{H'_{u'+1}\}$ and then Inclusion and Symmetry gives us
	% \begin{equation}\label{temp:2}
	% \begin{split}
	% &\quad [\bold{J},H_{1},\dots,H_{u}]+[\bold{J}',H'_{u'+1},H'_{1},\dots,H'_{u'}]
	% \\& \Rightarrow[\bold{J},\bold{J}',H'_{u'+1},(H_{i}+H'_{i'})_{1\leq i\leq u, 1\leq i'\leq u'}]
 %        \\& \Rightarrow [\bold{J},\bold{J}',H_{1},\dots,H_{u-1},H'_{u'+1},(H_{i}+H'_{i'})_{1\leq i\leq u, 1\leq i'\leq u'}].
	% \end{split}
	% \end{equation}
	% On the other hand, many applications of Inclusion and Symmetry give
 %    \begin{align}\label{temp:3}
 %        [\bold{J},H_{1},\dots,H_{u}]\Rightarrow [\bold{J},\bold{J}',H_{1},\dots,H_{u},(H_{i}+H'_{i'})_{1\leq i\leq u, 1\leq i'\leq u'}],
 %    \end{align}	
 %        Combining \eqref{temp:2} and \eqref{temp:3} via Concatenation, we get
	% \begin{equation}\nonumber
	% \begin{split}
	% %&\quad [\bold{J},\bold{J}',H_{1},\dots,H_{u},(H_{i}+H'_{i'})_{1\leq i\leq u, 1\leq i'\leq u'}]
	% %\\&+[\bold{J},\bold{J}',H_{1},\dots,H_{u-1},H'_{u'+1},(H_{i}+H'_{i'})_{1\leq i\leq u, 1\leq i'\leq u'}]\\
	% & \Rightarrow[\bold{J},\bold{J}',H_{1},\dots,H_{u-1},H_{u}+H'_{u'+1},(H_{i}+H'_{i'})_{1\leq i\leq u, 1\leq i'\leq u'}].
	% \end{split}
	% \end{equation}
	% Now suppose that $[\bold{J},\bold{J}',H_{1},\dots,H_{k},H_{k+1}+H'_{u'+1},\dots,,H_{u}+H'_{u'+1},(H_{i}+H'_{i'})_{1\leq i\leq u, 1\leq i'\leq u'}]$
	% for some $1\leq k\leq u-1$. Since 
	% \begin{equation}\nonumber
	% \begin{split}
	% &\quad [\bold{J},\bold{J}',H'_{u'+1},(H_{i}+H'_{i'})_{1\leq i\leq u, 1\leq i'\leq u'}]
	% \\&\Rightarrow[\bold{J},\bold{J}',H_{1},\dots,H_{k-1},H'_{u'+1},H_{k+1}+H'_{u'+1},\dots,H_{u}+H'_{u'+1},(H_{i}+H'_{i'})_{1\leq i\leq u, 1\leq i'\leq u'}],
	% \end{split}
	% \end{equation}
	% we have
	% \begin{equation}\nonumber
	% \begin{split}
	% &\quad [\bold{J},\bold{J}',H_{1},\dots,H_{k},H_{k+1}+H'_{u'+1},\dots,,H_{u}+H'_{u'+1},(H_{i}+H'_{i'})_{1\leq i\leq u, 1\leq i'\leq u'}]
	% \\&+[\bold{J},\bold{J}',H_{1},\dots,H_{k-1},H'_{u'+1},H_{k+1}+H'_{u'+1},\dots,H_{u}+H'_{u'+1},(H_{i}+H'_{i'})_{1\leq i\leq u, 1\leq i'\leq u'}]
	% \\&\Rightarrow[\bold{J},\bold{J}',H_{1},\dots,H_{k-1},H_{k}+H'_{u'+1},\dots,,H_{u}+H'_{u'+1},(H_{i}+H'_{i'})_{1\leq i\leq u, 1\leq i'\leq u'}].
	% \end{split}
	% \end{equation} 
	% By induction, we have  
	% \begin{equation}\nonumber
	% \begin{split}
	% &\quad [\bold{J},\bold{J}',H_{1}+H'_{u'+1},\dots,,H_{u}+H'_{u'+1},(H_{i}+H'_{i'})_{1\leq i\leq u, 1\leq i'\leq u'}]
	% \\&\Leftrightarrow[\bold{J},\bold{J}',(H_{i}+H'_{i'})_{1\leq i\leq u, 1\leq i'\leq u'+1}].
	% \end{split}
	% \end{equation} 	
	% So  Property-$(u,u'+1)$ holds. Similarly,  Property-$(u+1,u')$ holds. This finishes the proof.
\end{proof}


\section{Seminorm smoothing}\label{S: smoothing}
Our next goal is to derive improved seminorm estimates on arbitrary multiple ergodic averages of commuting transformations along Hardy sequences of polynomial growth. In \cite[Theorem 10.1]{DKKST24}, we obtained very general estimates on such averages; in this section, we upgrade them to ones that are more useful for our applications via a new variant of the seminorm smoothing method from \cite{DKKST24, FrKu22b, FrKu22a}. This is the first step in proving that the difference and product conditions in Definition \ref{D: good for joint ergodicity} are \textit{sufficient} for joint ergodicity.

To state our main results, we need the following definitions. %{\color{red} A general question: in this paper, for families/order families, do we write $\{q_1, \ldots, q_m\}$ or $(q_1, \ldots, q_m)$? I think we used both notations} \AK{For general families we can use the first, for the ordered we should use the second} \BK{I will put the basis in the ordered way throughout this section to stay consistent even if the order is not necessary.}

\begin{definition}[Ordered Hardy family]\label{D: ordered Hardy family}
    Let $\CQ = (q_1, \ldots, q_m)\subseteq\CH$ be a family of Hardy sequences such that $q_i\succ 1$ for all $i\in[m]$. We say that $\CQ$ is \textit{an ordered Hardy family} (or simply \textit{ordered}) if there exist integers $1\leq m_1\leq m_2\leq m_3\leq m$ such that the following holds:
        \begin{enumerate}
            \item $1\prec q_1(x) \prec \cdots \prec q_{m_1}(x)\ll\log x$;
            \item $\log x\prec q_{m_1+1}(x) \prec \cdots \prec q_{m_2}(x) \prec x^\veps$ for every $\veps>0$;
            \item for $m_2 < i \leq m_3$, we have $q_i(x) = x^{l_i}$ for positive integers $l_{m_2+1}< \cdots~<~l_{m_3}$;
            \item $x^\veps \prec q_{m_3+1}(x) \prec \cdots \prec q_m(x)$ for some $\veps>0$, and for each $m_3 < i\leq m$, the function $q_i$ is strongly nonpolynomial.
            %in that $t^{d_i}\prec q_i(t)\prec t^{d_i+1}$ for some $d_i\in\N$. \BK{(State this definition earlier on)}
        \end{enumerate}    
\end{definition}
For instance, the family $(\log x, (\log x)^2, x, x^2, x^{2/3}, x\log x, x^{3/2})$ is ordered. It is a priori strange to put monomials $x, x^2$ before slower growing strongly nonpolynomial functions like $x^{2/3}$ of positive fractional degree, but this ordering turns out to come out naturally in the seminorm bounds from \cite{DKKST24} stated below. More specifically, strongly nonpolynomial functions of positive fractional degree like $x^{2/3}$ can be Taylor-expanded as polynomials of arbitrarily high degree on short intervals, and in this sense they attain higher degrees than polynomials $x, x^2, \ldots$ whose Taylor expansions are polynomials of the same degree.
\begin{definition}[Extended span]
    Let $\CQ = (q_1, \ldots, q_m)\subseteq\CH$ be an ordered Hardy family. Then the \textit{extended span} of $\CQ$, denoted by $\extspan_\R\CQ$, is the collection of sequences $a\in\CH$ satisfying
    $$\lim_{x\to\infty}\abs{a(x) - \sum\limits_{i=1}^m \beta_i q_i(x) - \beta_0} = {0}$$ for some $\beta_0, \ldots,\beta_m\in\R$.
\end{definition}

%{\color{red} The family in Definition 4.2 needs not to be ordered, although Definition 4.3 needs} \BK{You're right, but to be consistent, I put the ordered notation.}
 
\begin{definition}[Leading coefficients]\label{d43}
    Let $\CQ = (q_1, \ldots, q_m)\subseteq\CH$ be an ordered Hardy family, and suppose that $\ba\in(\extspan_\R\CQ)^\ell$ satisfies $$\lim_{x\to\infty}\abs{\ba(x) - \sum\limits_{i=1}^m \bbeta_i q_i(x) - \bbeta_0} = {0}$$ for some $\bbeta_0, \ldots,\bbeta_m\in\R^\ell$. Then we call $\bbeta_i$ the \textit{leading coefficient} of $\ba$ if $\bbeta_i\neq \mathbf{0}$ and $\bbeta_{i'} = \mathbf{0}$ for all $i<i'\leq m$ {(note that the choices of $\bbeta_{i}$ are unique)}. If $\bbeta_0 = \cdots = \bbeta_m = \bf{0}$, then we say that $\ba$ has leading coefficient $\bf{0}.$
\end{definition}

All the concepts in this section involving degrees, leading coefficients and independence are defined with respect to a fixed ordered family. We will not specify the family as it is always clear from the context.

Then, \cite[Theorem 10.1]{DKKST24} gives the following seminorm bounds. 
     \begin{theorem}\label{T: old estimates}
        %Let $m, r, k_{1},\dots,k_{r}\in\N$ and $k=k_{1}+\dots+k_{r}$. Let $\CQ = \{q_1, \ldots, q_m\}\subseteq\CH$ be an ordered family of Hardy sequences, and for $i\in[r]$ and $j\in[k_i]$, let $a_{i,j}\in\CH$ be given by $$\lim\limits_{t\to\infty} \abs{a_{i,j}(t)-\sum\limits_{l=1}^m \beta_{i,j,l} q_l(t) + \beta_{i,j,0}} = 0$$ for some $\beta_{i,j,l}\in\R$.

        Let $\ell\in\N$. Let $\CQ\subseteq\CH$ be an ordered Hardy family, and let $a_1, \ldots, a_\ell\in\extspan_\R\CQ$ with $a_1\succ \log$. Then there exists a positive integer $s=O_{\ell, \CQ}(1)$ and vectors $\bv_{1},\dots,\bv_{s}$ coming from the set of the leading coefficients of $$\{a_1 \be_1,\; a_1 \be_1 - a_2 \be_2,\; \ldots,\; a_1 \be_1 - a_\ell \be_\ell\},$$
         such that for all systems $(X, \CX, \mu, T_1, \ldots, T_\ell)$ and functions $f_1, \ldots, f_\ell\in L^\infty(\mu)$, we have 
\begin{align*}%\label{goal}
         \lim_{N\to\infty} %\sup_{|c_n|\leq 1}
         \norm{\E_{n\in [N]} 
         %c_{n}\cdot
         \prod_{j\in[\ell]}  T_{j}^{\sfloor{a_{j}(n)}} f_{j}}_{L^2(\mu)}=0
    \end{align*}
    whenever $\nnorm{f_1}_{\bv_{1},\dots,\bv_{s}}=0$.
    \end{theorem}

The upgraded seminorm estimates that we prove in this paper can then be restated as follows. We recall that various notions of independence of Hardy sequences are defined in Definition \ref{D: independence}. 
     \begin{theorem}\label{T: new estimates}
        %Let $m, r, k_{1},\dots,k_{r}\in\N$ and $k=k_{1}+\dots+k_{r}$. Let $\CQ = \{q_1, \ldots, q_m\}\subseteq\CH$ be an ordered family of Hardy sequences, and for $i\in[r]$ and $j\in[k_i]$, let $a_{i,j}\in\CH$ be given by $$\lim\limits_{t\to\infty} \abs{a_{i,j}(t)-\sum\limits_{l=1}^m \alpha_{i,j,l} q_l(t) + \alpha_{i,j,0}} = 0$$ for some $\alpha_{i,j,l}\in\R$.

        Let $\ell\in\N$. Let $\CQ\subseteq\CH$ be an ordered Hardy family, and let $a_1, \ldots, a_\ell\in\extspan_\R\CQ$ with $a_1\succ \log$.  Then there exists a positive integer $s=O_{\ell, \CQ}(1)$ and vectors $\bv_{1},\dots,\bv_{s}$ coming from the set of the leading coefficients of 
        \begin{align*}
            \{a_1 \be_1\}\cup\{a_i \be_i - a_j \be_j:\; i\neq j,\; a_i, a_j\; \textrm{are\; dependent},\; a_i, a_j\succ \log\},
        \end{align*}
        %$a_1 \be_1$ and those $a_i \be_i - a_j \be_j$ for which $i\neq j$ and $a_i$, $a_j$ are pairwise dependent,
        %$$a_1 \be_1\qquad \textrm{and}\qquad a_i \be_i - a_j \be_j\quad \textrm{as\; long\; as}\; i\neq j\; \textrm{and}\; a_i,\; a_j\; \textrm{are\; pairwise\; dependent},$$
         such that for all systems $(X, \CX, \mu, T_1, \ldots, T_\ell)$ and functions $f_1, \ldots, f_\ell\in L^\infty(\mu)$, we have 
\begin{align}\label{goal00}
        \lim_{N\to\infty} %\sup_{|c_n|\leq 1}
         \norm{\E_{n\in [N]} 
         %c_{n}\cdot
         \prod_{j\in[\ell]}  T_{j}^{\sfloor{a_{j}(n)}} f_{j}}_{L^2(\mu)}=0
    \end{align}
    whenever $\nnorm{f_1}_{\bv_{1},\dots,\bv_{s}}=0$. 
    \end{theorem}
    In fact, a close inspection of Theorem~\ref{T: new estimates general}, from which this result follows, shows that we only need to consider the leading coefficient of $a_i \be_i - a_j \be_j$ for those dependent $a_i, a_j$ of the same or higher degree than $a_1$ in the sense of Definition~\ref{D: degree and independence}.


A priori, there seems to be no good reason why the vectors featuring in Theorem~\ref{T: new estimates} should be any better than those from Theorem~\ref{T: old estimates}. However, using the difference ergodicity condition, we will later be able to further replace all $\bv_i$'s in Theorem~\ref{T: new estimates} by copies of $\be_1$. This will give us Host-Kra seminorm control over the average in Theorem~\ref{T: new estimates} assuming the difference ergodicity condition.


 
% {\color{red} 

% One can of course use the concatenation theorem to merge Theorems~\ref{T: old estimates} and ~\ref{T: new estimates} into a stronger result. However, this result is neither artistic \AK{Aesthetic?} nor necessary for this paper. This leads to the following natural open question: what is the optimal box seminorm that controls the average (\ref{goal00})? \AK{If we end up keeping this question, maybe we can formulate it as a ``Question''}
% } \BK{I'd leave it out, tbh. It is not clear that the class of possible seminorms controlling a given expression has a minimal element (i.e. if seminorms along the vectors $\CA$ and $\CB$ control an expression, it is not obvious that the seminorm along $\CA\cap\CB$ would control it as well), in which case the question may not be well-defined (and if you wanted to use concatenation, you would get yet a different family of vectors/subgroups controlling an expression). Sounds to me a bit too technical and vague to deserve a mention.} {\color{red} the question can be formulated precisely, but it will require some more space and I am not sure if we want to state it. Let $\mathcal{A}$ and $\mathcal{B}$ be two families of subgroups of $\Z^{d}$. We define a partial order $\mathcal{A}\preceq\mathcal{B}$ if for every $H\in\mathcal{A}$, there exists $H'\in\mathcal{B}$ with $H\leq H'$. The question is then to find the maximum family of subgroups $\mathcal{B}$ of $\Z^{d}$ such that the multiple ergodic average can be controlled by $\nnorm{\cdot}_{H_{1},\dots,H_{s}}$ for some $s$ and some $H_{1},\dots,H_{s}\in \mathcal{B}$. The concatenation theorem says that if you have two such families, then the concatenated family is larger (in the sense of $\preceq$) than both of them.  So whenever you get a new seminorm estimate, the concatenation theorem will give you a larger family $\mathcal{B}$.
% }
% \BK{Yeah, but when we speak of optimal seminorm control, we should not only want the groups to be as large as possible, but also the degree to be as small as possible. The degree is relevant since some inverse theorems are much easier for small-degree box norms (e.g. the inverse theorem for the degree-2 seminorm along $\be_1, \be_2$ is trivial on magic systems, but for the degree-4 seminorm $\be_1, \be_1,\be_2,\be_2$, it seems very difficult). Concatenation makes one better but the other worse. For that reason, concatenation would ultimately be insufficient by itself to give optimal seminorm, and therefore I would leave this discussion out.}



\begin{example}
    For comparison, consider the tuple $$(T_1^{\sfloor{n^{3/2}}}, T_2^{\sfloor{n^{3/2}+\log n}}, T_3^{\sfloor{n^{3/2}+n\log n}}, T_4^{\sfloor{n^{3/2}+n\log n}})_n.$$ Then Theorem \ref{T: old estimates} gives that the corresponding multiple ergodic average would be controlled by a seminorm of $f_1$ involving many copies of $\be_1,\; \be_1 - \be_2,\; \be_1-\be_3,\; \be_1 - \be_4$. The iterates can be divided into two families of pairwise dependent functions, $\{n^{3/2},\; n^{3/2} + \log n\}$ and $\{n^{3/2}+ n\log n,\; n^{3/2} + n\log n\}$,  hence Theorem~\ref{T: new estimates} gives us control in terms of a seminorm of $f_1$ involving many copies of $\be_1,\; \be_1 - \be_2,\; \be_3-\be_4$. If the difference ergodicity condition holds, i.e. the sequences
    \begin{equation}\label{E: ergodic actions}
        (T_1^{\sfloor{n^{3/2}}} T_2^{-\sfloor{n^{3/2}+\log n}})_n\quad \textrm{and} \quad (T_3^{\sfloor{n^{3/2} + n\log n}}T_4^{-\sfloor{n^{3/2} + n\log n}})_n
    \end{equation}
    %$(T_1^{\sfloor{n^{3/2}}} T_2^{-\sfloor{n^{3/2}+\log n}})_n$ and $(T_3^{\sfloor{n^{3/2} + n\log n}}T_4^{-\sfloor{n^{3/2} + n\log n}})_n$ 
    are ergodic, then we can replace each of the vectors $\be_1 - \be_2,\; \be_3-\be_4$ in the seminorm by many copies of $\be_1$. For $\be_3-\be_4$, this is immediate: the ergodicity of the second action in \eqref{E: ergodic actions}
    %$(T_3^{\sfloor{n^{3/2} + n\log n}}T_4^{-\sfloor{n^{3/2} + n\log n}})_n$ 
    implies that $T_3T_4\inv$ is ergodic, hence $\nnorm{f}_{\be_3-\be_4}\leq \nnorm{f}_{\be_1}$ by e.g. \cite[Lemma 2.2]{FrKu22a}. For $\be_1-\be_2$, this requires extra work, and the general method for unpacking the ergodicity of sequences like $(T_1^{\sfloor{n^{3/2}}} T_2^{-\sfloor{n^{3/2}+\log n}})_n$ will be presented in the next sections.
\end{example}

It is worth noting that one can use the concatenation theorem to merge Theorems~\ref{T: old estimates} and ~\ref{T: new estimates} into a stronger result. For instance, the previous example gives one estimate involving the vectors $\be_{1},\be_{1}-\be_{2},\be_{1}-\be_{3},\be_{1}-\be_{4}$ and another one involving $\be_{1},\be_{1}-\be_{2},\be_{3}-\be_{4}$. Concatenating them, we get an estimate involving $\be_{1},\be_{1}-\be_{2},\langle\be_{1}-\be_{3},\be_{1}-\be_{4}\rangle$. However, we do not require strengthenings like this in the current paper.



\subsection{General formalism}

The proof of Theorem \ref{T: new estimates} uses no properties of Hardy functions except the initial seminorm estimates provided by Theorem \ref{T: old estimates}. We therefore phrase its proof in a more general language that makes no reference to Hardy functions.

Our basic setup is captured by the following definition.
\begin{definition}\label{d47}[Ordered family, degree, and independence]\label{D: degree and independence}
   An \textit{ordered family of sequences} is an ordered collection $\CQ =(q_1, \ldots, q_m)$ of sequences $q_i: \N\to\R$ whose germs are affinely independent.\footnote{That means that for every  $N\in\N$, the restrictions of $1, q_1, \ldots, q_m$ to $[N, \infty)$ are $\R$-independent.} Additionally, we assume that each ordered family of sequences comes with a fixed integer $0\leq m_1<m$.

    Suppose that $a_1, \ldots, a_\ell:\N\to\R$ satisfy
        $$\lim_{x\to\infty}\abs{a_j(x) - \sum\limits_{i=1}^m \beta_{j,i} q_i(x) - \beta_{j,0}} = {0}$$ for some $\beta_{ji}\in\R$; like before, we call the collection of such sequences the \textit{extended span} of $\CQ$ and denote it via $\extspan_\R\CQ$.
    %$a_j(n) = \sum\limits_{i=1}^m \beta_{ji} q_i(n)$. 
    We define the \textit{degree} of $a_j$ via $$\deg a_j := \max\{i\in[0, m]:\; \beta_{j,i}\neq 0\}$$
    (the notion is well-defined since the choices of $\beta_{j,i}$ are unique),
    letting it be 0 if $\beta_{j,m} = \cdots = \beta_{j,0} = 0$. Lastly, we call $\beta_{j, \deg a_j}$ the \textit{leading coefficient} of $a_j$.
    %and its \textit{complexity} by $\max\limits_{i\in[m]}\{|\beta_{ji}|,\; 1/|\beta_{ji}|:\; \beta_{ji}\neq 0\}.$ 
    We call the sequences:
    \begin{itemize}
        \item $a_j, a_j'$  \textit{independent}  if $(\beta_{j, m_1+1}, \ldots, \beta_{j, m})$ and $(\beta_{j', m_1+1}, \ldots, \beta_{j',m})$ are $\R$-linearly independent, otherwise we call them \textit{dependent};
        \item $a_1, \ldots, a_m$ \textit{pairwise independent} if for every $1\leq j < j' \leq \ell$, the sequences $a_j, a_{j'}$ are independent. 
    \end{itemize}
     %We also denote $d_{jj'} = \max\{i\in[m]:\; \alpha_{ji}\neq\alpha_{j'i}\}$ to be the \textit{degree} of $a_j-a_{j'}$, setting $d_{jj'} = 0$ if $a_j = a_{j'}$.
\end{definition}
In other words, pairwise independence means that no nontrivial linear combination $c_j a_j + c_{j'}a_{j'}$ for $1\leq j < j' \leq \ell$ lies in $\extspan_\R\{a_1, \ldots, a_{m_1}\}$. For Hardy sequences defined with respect to an ordered Hardy basis, the notion of pairwise independence above coincides with the analogous notion from Definition \ref{D: independence}.

The definitions of degree, leading coefficients and pairwise (in)dependence naturally extend to sequences $\N^\ell\to\R$ defined as $\R^\ell$-linear combinations of elements of the ordered basis $\CQ$.

The averages studied in this section will typically be twisted by dual sequences defined below. For the definition of dual functions $\CD_{s, T_j}(f)$, see Definition \ref{D: dual functions}. 
\begin{definition}[Dual sequences]\label{D: dual sequences}
    Let $d,\ell\in\N$ and $(X, \CX, \mu, T_1, \ldots, T_\ell)$ be a system. Then $\FD_d$ denotes the class of sequences $\CD:\Z\to L^\infty(\mu)$ given by $\CD(n) = T_j^n\CD_{s, T_j}(f)$ for some $s\in[d]$, $j\in[\ell]$, and $f\in L^\infty(\mu)$. We refer to elements of $\FD_d$ as \textit{level-$d$ dual sequences}.
\end{definition}

The initial seminorm estimates instantiated by Theorem \ref{T: old estimates} are a special case of the following definition. 
\begin{definition}[Families good for smoothing]
    We call an ordered family of sequences $\CQ$ \textit{good for smoothing} if for all $d,\ell\in\N$, $J\in\N_0$, sequences $a_1, \ldots, a_\ell$, $b_1, \ldots, b_J\in \extspan_\R\CQ$  
    % with $a_1, \ldots, a_\ell$ being pairwise independent {\color{red} I think we only need $a_{1},a_{i}$ to be independent for $2\leq i\leq \ell$. This comment also applies to the definition in paper 1} \BK{(This is equivalent after swapping the role of the sequences (which we do throughout smoothing))}
    % {\color{red} actually why do we need the functions to be independent, since the transformations $T_{j}$ are already different? Should $a_{i}$ be vector valued?}
    % %and of complexity at most $M$, 
    % %and $a_1$ having maximum degree among $a_1, \ldots, a_\ell$, 
     satisfying $$\deg a_1 = \max_j\deg a_j > m_1,$$
    systems $(X, \CX, \mu, T_1, \ldots, T_\ell)$, %indexing tuples $\eta\in[k]^\ell$, 
    and functions $f_1, \ldots, f_\ell\in L^\infty(\mu)$, we have 
%    \begin{multline}\label{E: gen box bound}
 %        \limsup_{N\to\infty}\sup_{\substack{\CD_1, \ldots, \CD_J \in \FD_d}} \sup_{|c_n|\leq 1}\norm{\E_{n\in [N]} c_{n}\cdot \prod_{j\in[\ell]}  T_{\eta_j}^{\floor{a_j(n)}} f_j \cdot \prod_{j\in[J]}\CD_j(\floor{b_j(n)})}_{L^2(\mu)}^{{O_{d,J,\ell, \CQ}(1)}}\\ 
 %        \ll_{d, J, k, \ell, M, \CQ} \nnorm{f_1}_{(\alpha_{1d_{10}}\be_{\eta_1})^s,\; (\alpha_{1d_{12}}\be_{\eta_1}-\alpha_{2d_{12}}\be_{\eta_2})^s,\; \ldots,\;  (\alpha_{1d_{1\ell}}\be_{\eta_1}-\alpha_{\ell d_{1\ell}}\be_{\eta_\ell})^s}^+
%    \end{multline}
     \begin{align}\label{E: gen box bound}
         \lim_{N\to\infty}\sup_{\substack{\CD_1, \ldots, \CD_J \in \FD_d}} \sup_{|c_n|\leq 1}\norm{\E_{n\in [N]} c_{n}\cdot \prod_{j\in[\ell]}  T_{j}^{\floor{a_j(n)}} f_j \cdot \prod_{j\in[J]}\CD_j(\floor{b_j(n)})}_{L^2(\mu)} = 0
%         \limsup_{N\to\infty}\sup_{\substack{\CD_1, \ldots, \CD_J \in \FD_d}} \sup_{|c_n|\leq 1}\norm{\E_{n\in [N]} c_{n}\cdot \prod_{j\in[\ell]}  T_{j}^{\floor{a_j(n)}} f_j \cdot \prod_{j\in[J]}\CD_j(\floor{b_j(n)})}_{L^2(\mu)}^{{O_{d,J,\ell, \CQ}(1)}}\\ 
         %\ll_{d, J, k, \ell, M, \CQ} \nnorm{f_1}_{(\alpha_{1d_{10}}\be_{1})^s,\; (\alpha_{1d_{12}}\be_{1}-\alpha_{2d_{12}}\be_{2})^s,\; \ldots,\;  (\alpha_{1d_{1\ell}}\be_{1}-\alpha_{\ell d_{1\ell}}\be_{\ell})^s}
    \end{align}
    %{\color{red} check if $\FD$ and $\CD_{j}$ are defined} \BK{check Definition \ref{D: dual sequences}}
    whenever $\nnorm{f_1}_{\bv_{1},\dots,\bv_{s}}=0$
    % \begin{align*}
    %     \nnorm{f_1}_{(\alpha_{1d_{10}}\be_{1})^s,\; (\alpha_{1d_{12}}\be_{1}-\alpha_{2d_{12}}\be_{2})^s,\; \ldots,\;  (\alpha_{1d_{1\ell}}\be_{1}-\alpha_{\ell d_{1\ell}}\be_{\ell})^s} = 0
    % \end{align*}
    for some $s = O_{d,J,\ell, \CQ}(1)$ and vectors $\bv_1, \ldots, \bv_s$ from the set of the leading coefficients of %{\color{red} not defined in Definition 4.7} \BK{(added to the definition plus a remark below that we extend to $\R^\ell$ combinations)}
    $$\{a_1 \be_1,\; a_1 \be_1 - a_2 \be_2,\; \ldots,\; a_1 \be_1 - a_\ell \be_\ell\}.$$
        %\footnote{If we wanted, we could replace the seminorm $\nnorm{f_1}_{(\alpha_{1d_{10}}\be_{1})^s,\; (\alpha_{1d_{12}}\be_{1}-\alpha_{2d_{12}}\be_{2})^s,\; \ldots,\;  (\alpha_{1d_{1\ell}}\be_{1}-\alpha_{\ell d_{1\ell}}\be_{\ell})^s}$ in the definition of being good for smoothing by $\nnorm{f_1}_{(\alpha_{1d_{10}}\be_{1})^s,\; (\alpha_{1d_{12}}\be_{1}-\alpha_{2d_{12}}\be_{2})^s,\; \ldots,\;  (\alpha_{1d_{1\ell}}\be_{1}-\alpha_{\ell d_{1\ell}}\be_{\ell})^s}^+$, as one is related to the other via \cite[Lemma 3.9]{DKKST24}. However, it is somewhat more convenient to use the former throughout the argument.}
    \end{definition}

    Any ordered Hardy family $\CQ$ is good for smoothing via \cite[Corollary~6.4 \& Theorem~10.2]{DKKST24},  where $m_1$ in Definition~\ref{D: degree and independence} is the same as $m_1$ in Definition~\ref{D: ordered Hardy family}. Theorem~\ref{T: new estimates} is then a consequence of the following general result.
         \begin{theorem}\label{T: new estimates general}
        %Let $m, r, k_{1},\dots,k_{r}\in\N$ and $k=k_{1}+\dots+k_{r}$. Let $\CQ = \{q_1, \ldots, q_m\}\subseteq\CH$ be an ordered family of Hardy sequences, and for $i\in[r]$ and $j\in[k_i]$, let $a_{i,j}\in\CH$ be given by $$\lim\limits_{t\to\infty} \abs{a_{i,j}(t)-\sum\limits_{l=1}^m \alpha_{i,j,l} q_l(t) + \alpha_{i,j,0}} = 0$$ for some $\alpha_{i,j,l}\in\R$.
        Let $\ell\in\N$. Suppose that an ordered family of sequences $\CQ$ is {good for smoothing}, and let $a_1, \ldots, a_\ell\in\extSpan_\R \CQ$ with $\deg(a_{1})>m_{1}$.
        %$\min\limits_j\deg a_j > m_1$.
        %Let $\CQ =\{q_1, \ldots, q_m\}$ be an ordered collection of sequences $q_i: \N\to\R$ whose germs are affinely independent, and let $a_1, \ldots, a_\ell\in\Span_\R\CQ$. 
        Then there exists a positive integer $s=O_{\ell, \CQ}(1)$ and vectors $\bv_{1},\dots,\bv_{s}$ coming from the set of the leading coefficients of 
        \begin{align*}
%            \{a_1 \be_1\}\cup\{a_i \be_i - a_j \be_j:\; i\neq j,\; a_i, a_j\; \textrm{are\; dependent}\},
            \{a_1 \be_1\}\cup\{a_i \be_i - a_j \be_j:\; i\neq j,\; a_i, a_j\; \textrm{are\; dependent and } \deg(a_{i})=\deg(a_{j})\geq \deg(a_{1})\},
        \end{align*}
        %$a_1 \be_1$ and those $a_i \be_i - a_j \be_j$ for which $i\neq j$ and $a_i$, $a_j$ are pairwise dependent,
        %$$a_1 \be_1\qquad \textrm{and}\qquad a_i \be_i - a_j \be_j\quad \textrm{as\; long\; as}\; i\neq j\; \textrm{and}\; a_i,\; a_j\; \textrm{are\; pairwise\; dependent},$$
         such that for all systems $(X, \CX, \mu, T_1, \ldots, T_\ell)$ and functions $f_1, \ldots, f_\ell\in L^\infty(\mu)$, we have 
\begin{align*}%\label{goal}
                  \lim_{N\to\infty} %\sup_{|c_n|\leq 1}
         \norm{\E_{n\in [N]} 
         %c_{n}\cdot
         \prod_{j\in[\ell]}  T_{j}^{\sfloor{a_{j}(n)}} f_{j}}_{L^2(\mu)}=0
    \end{align*}
    whenever $\nnorm{f_1}_{\bv_{1},\dots,\bv_{s}}=0$. 
    \end{theorem}
    % {\color{red} As a consequence of the modification of condition (ii) of the next proposition, in this proposition, (i) the assumption can be changed to  $\deg(a_{1})>m_{1};$ (ii) the set of functions can be reduced to 
    %  \begin{align*}
    %         \{a_1 \be_1\}\cup\{a_i \be_i - a_j \be_j:\; i\neq j,\; a_i, a_j\; \textrm{are\; dependent and } \deg(a_{i})=\deg(a_{j})\geq \deg(a_{1})\},
    %     \end{align*}
    % }

In the proof of Theorem \ref{T: new estimates general}, it will matter a lot whether pairs of sequences $a_i, a_j$ are dependent or not. We therefore state a slight generalization of Theorem \ref{T: new estimates general} that makes it easier to read off the dependence or independence of particular sequences and better matches with the way the proof goes. In what follows, let $\hat{a}_{i}$ denote the leading coefficient of $a_{i}$. %Our goal is to prove the following. %generalization of Proposition \ref{newestimate}:
     
     \begin{proposition}\label{newestimate3}
        Let $d, r, \ell_{1},\dots,\ell_{r}\in\N$, $\ell=\ell_{1}+\dots+\ell_{r}$, and $J\in\N_0$. Suppose that an ordered family of sequences $\CQ$
        is {good for smoothing}, and let $a_{i,j}\in\extSpan_\R \CQ$ (for $i\in[r], j\in[\ell_i]$) be sequences
        %the sequences $a_{i,j}, 1\leq i\leq r, 1\leq j\leq \ell_{i}, a'_{i'}, 1\leq i'\leq r'\in\Span_\R\CQ$ are 
        %of complexity at most $M$
        with the following properties:
        \begin{enumerate}
%             \item all of $a_{i,j}$ are of the same degree;
            % \item the degree of $a_{1,1}$ is higher than any of $a'_{i'}$;
             \item $a_{i,j}$ is independent from $a_{i',j'}$ if and only if $i\neq i'$;
              \item %there exists $r'\in[r]$ such that $\deg(a_{1,j})>\deg(a_{i,j'})$ for all $i>r'$, $j\in[\ell_1]$, and $j'\in[\ell_i]$; 
              there exists $r'\in[r]$ such that for all $i\in[r]$, $j\in[\ell_1]$, and $j'\in[\ell_i]$, we have that $\deg(a_{1,j})>\deg(a_{i,j'})$ if and only if $i>r'$;
              %$\deg(a_{1,j})>\deg(a_{i,j'})$ if $i>r'$ and $\deg(a_{1,j})\leq\deg(a_{i,j'})$ if $i\leq r'$;
             \item $\deg(a_{1,1})>m_{1}$.
             %$a_{1,1}\neq 0$. \BK{I think that instead, we want $\min\limits_{i\in[r], j\in[\ell_i]}\deg a_{i,j} > m_1$.}
             %\BK{(I think that it is fine if $a_{1,1}=a_{1,j}$: consider e.g. the tuple $(T_{1,1}^n, T_{1,2}^n$)}
        \end{enumerate}

         
                 Then there exists a positive integer $s=O_{d,J,\ell, \CQ}(1)$\footnote{Note that $r,r'\leq \ell$, and so the dependence of $s$ and other parameters on $r, r'$ is implicit in their dependence on $\ell$.} and vectors $\bv_{1},\dots,\bv_{s}$ coming from the set %of following vectors 
                 \begin{align*}
                     \CG:=\{\hat{a}_{1,1}\be_{1,1}, \hat{a}_{i,j}\be_{i,j}-\hat{a}_{i,j'}\be_{i,j'}:\; i\in[r'],\; j,j'\in[\ell_{i}],\; j\neq j'\}
                 \end{align*}
                 %{\color{red} $\be_{i,j}$ does not matter when $i>r'$}
         such that for all systems $(X, \CX, \mu, (T_{i,j})_{\substack{i\in[r], j\in[\ell_i]}})$,  functions $f_{i,j}\in L^\infty(\mu)$ and sequences $b_1, \ldots, b_J\in\Span_\R\CQ$, we have 
\begin{multline}\label{goal}
         \limsup_{N\to\infty}\sup_{\substack{
         \CD_1, \ldots, \CD_J \in \FD_d}} \sup_{|c_n|\leq 1}\norm{\E_{n\in [N]} c_{n}\cdot \prod_{i\in[r]}\prod_{j\in[\ell_i]}  T_{i,j}^{\floor{a_{i,j}(n)}} f_{i,j} \cdot \prod_{j\in[J]}\CD_j(\floor{b_j(n)})}_{L^2(\mu)}%^{{O_{d,J,\ell,r', \CQ}(1)}}\\ 
         %\ll_{d, J, k, M, r', \CQ} \nnorm{f_1}_{S_{1},\dots,S_{s}}.
    \end{multline}
    is zero whenever $\nnorm{f_1}_{\bv_{1},\dots,\bv_{s}}=0$. 
    %\BK{(I don't think there is much advantage in having a separate product over $i'\in[r']$ for iterates that have lower degree than $a_{1,1}$ since throughout the almost entire argument, you assume that $a_{1,1}$ has maximum degree, in which case other $a_{i,j}$'s are implicitly assumed to also have maximum degree. I would suggest that we kill the product over $i'\in[r']$ altogether - this will help simplify the formulas a bit. Instead, throughout Case 1, we could introduce a parameter $r_0\in[r]$ to keep track of those $a_{i,j}$'s that have the same degree as $a_{1,1}.$ )} {\color{red} I made the changes accordingly}
    \end{proposition}
We observe that the set $\CH$ does not depend on the ``lower-degree'' sequences $a_{i,j}$ with $r'<i\leq r$ and $j\in[\ell_i]$. This is one way in which Proposition \ref{newestimate3} is stronger than Theorem \ref{T: new estimates general}; the other is the presence of dual functions in \eqref{goal}.

\subsection{A motivating example}\label{SS: concatenation example}

Before we give the proof of Proposition \ref{newestimate3} in full, we discuss its strategy for the average
\begin{align}\label{E: smoothing example average}
    \limsup_{N\to\infty}\norm{\E_{n\in [N]} T_{1,1}^{a_{1}(n)} f_{1,1}\cdot T_{1,2}^{a_{1}(n)} f_{1,2}\cdot T_{2,1}^{a_{2}(n)} f_{2,1}\cdot T_{2,2}^{a_{2}(n)} f_{2,2}}_{L^2(\mu)},
\end{align}
where, say, $a_1(n) = \sfloor{n^{3/2}}$ and $a_2(n) = \sfloor{n^{3/2}+n\log n}$ (what matters is that $a_1, a_2$ are independent and have the same leading term $n^{3/2}$). Throughout, all the functions are 1-bounded. While we strive to make this section as self-contained as possible, we strongly encourage the reader to first familiarize themselves with the examples of smoothing carried out in \cite[Section 8.1]{FrKu22a} and \cite[Section 4]{FrKu22b} for more details on various technical steps (e.g. dual-difference interchange) on which we will not elaborate in this expository section.


Theorem \ref{T: old estimates} gives us the following control over \eqref{E: smoothing example average}, which serves as a starting point of our argument.
\begin{proposition}\label{P: old estimates example}
    There exists $s\in\N$ and vectors
    \begin{align}\label{E: old vectors example}
        \bv_1, \ldots, \bv_s\in \{\be_{1,1},\; \be_{1,1}-\be_{1,2},\; \be_{1,1}-\be_{2,1},\; \be_{1,1}-\be_{2,2}\}.
    \end{align}
    such that the average \eqref{E: smoothing example average} is $0$ whenever $\nnorm{f_{1,1}}_{\bv_1, \ldots, \bv_s}=0$.
\end{proposition}
To get the conclusion of Theorem \ref{T: new estimates} (or Proposition \ref{newestimate3}) we want the following.
\begin{proposition}\label{P: new estimates example}
    There exists $s\in\N$ and vectors 
    \begin{align}\label{E: new vectors example}
    \bv_1, \ldots, \bv_s\in\{\be_{1,1},\; \be_{1,1}-\be_{1,2},\; \be_{2,1}-\be_{2,2}\}.
\end{align}
    such that the average \eqref{E: smoothing example average} is $0$ whenever $\nnorm{f_{1,1}}_{\bv_1, \ldots, \bv_s}=0$.
\end{proposition}
    
% Theorem \ref{T: old estimates} allows us to control this average via $\nnorm{f_{1,1}}_{\bv_1, \ldots, \bv_s}$ for some vectors $\bv_1, \ldots, \bv_s$ lying in $$\{\be_{1,1},\; \be_{1,1}-\be_{1,2},\; \be_{1,2}-\be_{2,1},\; \be_{1,2}-\be_{2,2}\},$$ in the sense that \eqref{E: smoothing example average} vanishes whenever $\nnorm{f_{1,1}}_{\bv_1, \ldots, \bv_s}=0$, whereas to get the conclusion of Theorem \ref{T: new estimates} (or Proposition \ref{newestimate3}) we want control by a seminorm with vectors from  
% \begin{align}\label{E: new vectors example}
%     \{\be_{1,1},\; \be_{1,1}-\be_{1,2},\; \be_{2,1}-\be_{2,2}\}.
% \end{align}
The goal thus is to show that each vector $\be_{1,1}-\be_{2,j}$ among \eqref{E: old vectors example} can be replaced by a bounded number of vectors from \eqref{E: new vectors example} so that the resulting seminorm still controls \eqref{E: smoothing example average}. In doing so, we will pass to auxiliary averages of the form
\begin{align}\label{E: smoothing example average eta}
        \limsup_{N\to\infty}\norm{\E_{n\in [N]} T_{\eta(1,1)}^{a_{1}(n)} f_{1,1}\cdot T_{\eta(1,2)}^{a_{1}(n)} f_{1,2}\cdot T_{\eta(2,1)}^{a_{2}(n)} f_{2,1}\cdot T_{\eta(2,2)}^{a_{2}(n)} f_{2,2}}_{L^2(\mu)}
\end{align}
for various choices $\eta: I\to I$ (here, $I = \{(1,1), (1,2), (2,1), (2,2)\}$), possibly twisted by weights $c_n$ or dual terms $\CD_j$. For convenience, we say that \eqref{E: smoothing example average eta} has \textit{type} $\eta$
or $$(\eta(1,1),\eta(1,2);\eta(2,1),\eta(2,2)).$$

%Suppose that $r=2$, $\ell_{1}=\ell_{2}=2$ $a_{i,1}=a_{i,2}=a_{i}$ for $i=1,2$ for some independent $a_{1},a_{2}$ with the same degree. For simplicity we assume that all of $a_{i,j}$ are integer coefficient polynomials with leading coefficient equal to 1, and assume that the terms $c_{n}$ and $\mathcal{D}_{j}$ does not appear. For convenience we call $(\eta(1,1),\eta(1,2);\eta(2,1),\eta(2,2))$  the \emph{type} of $\eta$. 

Much like in previous smoothing arguments \cite{DKKST24, FrKu22b, FrKu22a}, in order to prove Proposition~\ref{P: new estimates example}, we will induct on two parameters: 
\begin{itemize}
    \item a ``complexity'' of the type $\eta$, which can be measured in terms of the variance $\sum\limits_{i,j\in[2]}|\eta\inv(i,j)|^2$ (the higher the variance, the lower the complexity, hence averages in which all transformations are distinct are most complex);
    \item the length of the average, measured by the number of functions appearing in the average (weights $c_n$ or dual terms $\CD_j$ do not count).
\end{itemize}
Hence Proposition \ref{P: new estimates example} can be deduced inductively from the following two statements, one providing seminorm control for averages of types $((2,j),(1,2);(2,1),(2,2))$ (for $j\in[2]$), and the other dealing with averages of length 3.
\begin{claim}\label{C: lower type example}
    Fix $j\in[2]$. There exists $s\in\N$ such that 
    \begin{align*}%\label{E: auxiliary average 1}
    \limsup_{N\to\infty}\norm{\E_{n\in [N]} T_{2,j}^{a_{1}(n)} f_{1,1}\cdot T_{1,2}^{a_{1}(n)} f_{1,2}\cdot T_{2,1}^{a_{2}(n)} f_{2,1}\cdot T_{2,2}^{a_{2}(n)} f_{2,2}}_{L^2(\mu)} = 0
    \end{align*}
    whenever $\nnorm{f_{1,2}}_{\be_{1,2}^{\times s}, (\be_{1,2}-\be_{2,j})^{\times s}, (\be_{2,1}-\be_{2,2})^{\times s}}=0$.
\end{claim}
\begin{claim}\label{C: lower length example}
    Fix $j\in[2]$ and $d\in\N$. There exists $s\in\N$ such that for every level-$d$ dual function $\CD$, we have
    \begin{align*}%\label{E: auxiliary average 1}
    \limsup_{N\to\infty}\norm{\E_{n\in [N]} T_{1,1}^{a_{1}(n)} f_{1,1}\cdot \CD({a_{1}(n)})\cdot T_{2,1}^{a_{2}(n)} f_{2,1}\cdot T_{2,2}^{a_{2}(n)} f_{2,2}}_{L^2(\mu)} = 0
    \end{align*}
    whenever $\nnorm{f_{1,1}}_{\be_{1,1}^{\times s}, (\be_{2,1}-\be_{2,2})^{\times s}}=0$.
\end{claim}
Note that the seminorm appearing in both claims above are of the form provided by Proposition \ref{newestimate3}: we can therefore think of both claims as the induction hypothesis.
%Claim \ref{C: lower type example} is of the form provided by Theorem \ref{T: new estimates general}.

%We note that to control the original average \eqref{E: smoothing example average} in terms of a seminorm of $f_{1,1}$ as in Proposition \ref{P: new estimates example}, as want to control the auxiliary average \eqref{E: auxiliary average 1} in terms of a seminorm of $f_{1,2}$. \BK{Comment on that.}

\begin{proof}[Sketch of proof of Proposition \ref{P: new estimates example} assuming Claims \ref{C: lower type example} and \ref{C: lower length example}]
   We now outline the proof of Proposition \ref{P: new estimates example} assuming the preceding auxiliary results. Our focus will be on the sequence of necessary reductions to averages of lower complexity or length rather than on the exact way in which these reductions are carried out. To get the sense of the technicalities involved in the latter, the reader should consult the examples from \cite[Section 8.1]{FrKu22a} and \cite[Section 4]{FrKu22b}.

   Our goal is to show that \eqref{E: smoothing example average} is 0 if $\nnorm{f_{1,1}}_{\bv_{1},\dots,\bv_{s}}=0$ for some $s\in\N$ and $$\bv_{1},\dots,\bv_{s}\in \mathcal{G}:=\{\be_{1,1},\; \be_{1,1}-\be_{1,2},\; \be_{2,1}-\be_{2,2}\}.$$
   %(for convenience in this case we say that $\nnorm{f_{1,1}}_{\bv_{1},\dots,\bv_{t}}$ \emph{controls} (\ref{fwfwefwe})).

By Proposition \ref{P: old estimates example}, we have a preliminary control which asserts that \eqref{E: smoothing example average} is controlled by
  $\nnorm{f_{1,1}}_{\bv_{1},\dots,\bv_{t},\bu_{1},\dots,\bu_k}$ for some $k,t\in\N$ and vectors $\bv_{1},\dots,\bv_{t}\in \mathcal{G}$ (in fact, they lie in $\{\be_{1,1},\; \be_{1,1}-\be_{1,2}\}$) and $\bu_{1},\dots,\bu_k\in\{\be_{1,1}-\be_{2,1},\be_{1,1}-\be_{2,2}\}$. 
  %(for convenience in this case we say that $\nnorm{f_{1,1}}_{\bv_{1},\dots,\bv_{t},\bu_{1},\dots,\bu_k}$ \emph{controls} (\ref{fwfwefwe})). 
  Our aim is to remove $\bu_k$ at the cost of enlarging the number $t$, and then iterate this procedure with other $\bu_1, \ldots, \bu_{k-1}$.

 \smallskip
   \textbf{Step 1: Ping.} 
   \smallskip
   
   This step emulates the Ping step from the earlier smoothing arguments \cite{DKKST24, FrKu22b, FrKu22a}. Assume for concreteness that $\bu_k=\be_{1,1}-\be_{2,1}$ (the case $\bu_k=\be_{1,1}-\be_{2,2}$ being handled analogously). Then the Ping step essentially asserts that if some $\nnorm{f_{1,1}}_{\bold{w}_{1},\dots,\bold{w}_{r},\bu_k}$ controls (\ref{E: smoothing example average}) and if some $\nnorm{f_{i,j}}_{\bold{w}'_{1},\dots,\bold{w}'_{r'}}$ with $(i,j)\neq (1,1)$ controls
   \begin{equation}\label{E: smoothing example average ping}
    \limsup_{N\to\infty}\norm{\E_{n\in [N]} T_{2,1}^{a_{1}(n)} f_{1,1}\cdot T_{1,2}^{a_{1}(n)} f_{1,2}\cdot T_{2,1}^{a_{2}(n)} f_{2,1}\cdot T_{2,2}^{a_{2}(n)} f_{2,2}}_{L^2(\mu)}
\end{equation}
(which is an expression derived from (\ref{E: smoothing example average}) by replacing the transformation $T_{1,1}$ with $T_{2,1}$), then (\ref{E: smoothing example average}) is controlled by $\nnorm{f_{i,j}}_{\bold{w}_{1},\dots,\bold{w}_{r}, \bold{w}'_{1},\dots,\bold{w}'_{r'}}$.

In our case, Claim \ref{C: lower type example} gives us that \eqref{E: smoothing example average ping} is controlled by $\nnorm{f_{1,2}}_{\bu'_{1},\dots,\bu'_{k'},\be_{1,2}^{\times d}}$ for some $d, k'\in\N$ and $\bu'_{1},\dots,\bu'_{k'}\in \{\be_{1,2}-\be_{2,1},\be_{2,1}-\be_{2,2}\}$. %\BK{(Previously, the marked term was $\be_{1,2}-\be_{2,2}$ - please check)}. 
Therefore, we conclude that (\ref{E: smoothing example average}) is controlled by 
\begin{align}\label{E: auxiliary ping control}
    \nnorm{f_{1,2}}_{\bv_{1},\dots,\bv_{t},\bu_{1},\dots,\bu_{k-1},\bu'_{1},\dots,\bu'_{k'}, \be_{1,2}^{\times d}}.
\end{align}

 \smallskip
   \textbf{Step 2: Pong.} 
   \smallskip 
   
This step again emulates the Pong step from \cite{DKKST24, FrKu22b, FrKu22a}, and it essentially asserts that if some $\nnorm{f_{1,2}}_{\bold{w}_{1},\dots,\bold{w}_{r},\be_{i,j}^{\times d}}$ controls (\ref{E: smoothing example average}), and if $\nnorm{f_{1,1}}_{\bold{w}'_{1},\dots,\bold{w}'_{r'}}$ controls
   \begin{equation}\label{E: smoothing example average pong}
    \limsup_{N\to\infty}\norm{\E_{n\in [N]} T_{1,1}^{a_{1}(n)}f_{1,1}\cdot \CD(a_1(n)) \cdot T_{2,1}^{a_{2}(n)} f_{2,1}\cdot T_{2,2}^{a_{2}(n)} f_{2,2}}_{L^2(\mu)},
\end{equation}
%\BK{(previously, the marked terms were $T_{2,1}^{a_{1}(n)}$ and $T_{1,2}^{a_{1}(n)}$ respectively - please check)}
where $\CD$ is a level-$d$ dual sequence of $T_{1,2}$,
%twisted by a weight of dual functions (which is an expression derived from (\ref{fwfwefwe}) by removing the term $f_{1,2}$), 
then $\nnorm{f_{1,1}}_{\bold{w}_{1},\dots,\bold{w}_{r},\bold{w}'_{1},\dots,\bold{w}'_{r'}}$ controls (\ref{E: smoothing example average}). 

In our case, it follows from the Ping step that (\ref{E: smoothing example average}) is controlled by \eqref{E: auxiliary ping control}. %$$\nnorm{f_{1,2}}_{\bv_{1},\dots,\bv_{t},\bu_{1},\dots,\bu_{k-1},\bu'_{1},\dots,\bu'_{k'}, \be_{1,2}^{s}}.$$ 
On the other hand, Claim \ref{C: lower length example} gives control of (\ref{E: smoothing example average pong}) by $\nnorm{f_{1,1}}_{\bu''_{1},\dots,\bu''_{k''}}$ for some $k''\in\N$ and $\bu''_{1},\dots,\bu''_{k''}\in \mathcal{G}$. So we have that (\ref{E: smoothing example average}) is controlled by 
\begin{align}\label{E: auxiliary pong control}
    \nnorm{f_{1,1}}_{\bv_{1},\dots,\bv_{t},\bu_{1},\dots,\bu_{k-1},\bu'_{1},\dots,\bu'_{k'},\bu''_{1},\dots,\bu''_{k''}}.
\end{align}

 \smallskip
\textbf{Step 3: Relative concatenation}.
 \smallskip

This step is a completely new contribution to the seminorm smoothing machinery, and it crucially relies on the relative concatenation results from Section \ref{S: concatenation}. The preliminary estimate and the conclusion from the Pong step asserts that (\ref{E: smoothing example average}) is zero when $f_{1,1}$ is orthogonal to one of
\begin{align*}
    \mathcal{Z}^+_{\bv_{1},\dots,\bv_{t},\bu_{1},\dots,\bu_k}\quad \textrm{or}\quad  \mathcal{Z}^+_{\bv_{1},\dots,\bv_{t},\bu_{1},\dots,\bu_{k-1},\bu'_{1},\dots,\bu'_{k'},\bu''_{1},\dots,\bu''_{k''}}. 
\end{align*}
%\BK{(To be consistent with this step, we should probably run the entire argument in terms of plussed seminorms)}
It then follows from Corollary \ref{C: relative concatenation} that (\ref{E: smoothing example average}) is zero if $f_{1,1}$ is orthogonal to $$\mathcal{Z}^+_{\bv_{1},\dots,\bv_{t},\bu_{1},\dots,\bu_{k-1},\{\langle \bu_k,\bu \rangle:\; u\in\mathcal{R}\}},$$ or equivalently if $$\nnorm{f_{1,1}}^+_{\bv_{1},\dots,\bv_{t},\bu_{1},\dots,\bu_{k-1}, \{\langle \bu_k,\bu \rangle:\; u\in\mathcal{R}\}}=0,$$ where $\mathcal{R}$ is the multi-set $\{\bu'_{1},\dots,\bu'_{k'},\bu''_{1},\dots,\bu''_{k''}\}$.
However, for any $\bu\in\mathcal{R}$,  either $\bu\in\mathcal{G}$ or $\bu=\be_{1,2}-\be_{2,1}$. Now, for any such $\bu$, the group $\langle \bu_k,\bu \rangle$ contains an element of $\CG$. In the former case, this is immediate, while in the latter case, the group $\langle \bu_k,\bu \rangle=\langle \be_{1,1}-\be_{2,1}, \be_{1,2}-\be_{2,1} \rangle$ contains $\be_{1,1}-\be_{1,2}$. By Lemma \ref{L: basic properties of seminorms}, this implies that  (\ref{E: smoothing example average}) is controlled by $\nnorm{f_{1,1}}_{\bv_{1},\dots,\bv_{t'},\bu_{1},\dots,\bu_{k-1}}=0$ for some $t'\in\N$ and $\bv_{1},\dots,\bv_{t'}\in\mathcal{G}$.

We may then repeat this process to remove all of $\bu_{1},\dots,\bu_{k-1}$ and end up with a control of (\ref{E: smoothing example average}) by $\nnorm{f_{1,1}}_{\bv_{1},\dots,\bv_{t''}}=0$ for some $t''\in\N$ and $\bv_{1},\dots,\bv_{t''}\in\mathcal{G}$,  as desired.
    \end{proof}



\subsection{Proof of Proposition \ref{newestimate3}}
%We prove this result in the rest of this section. 
We move on to prove Proposition \ref{newestimate3}. From now on, we fix a system $(X, \CX, \mu, (T_{i,j})_{\substack{i\in[r], j\in[\ell_i]}})$.

%It suffices to show that  for some $Q\in\N$, we have that $$\lim_{N\to\infty}\Bigl\Vert\E_{n\in[N]}\prod_{i=1}^{r}\prod_{j=1}^{k_{i}}T_{i,j}^{a_{i,j}p_{i}(Qn+q)}f_{i,j}\prod_{v=1}^{L}\mathcal{D}_{v}(q_{v}(n))\Bigr\Vert_{2}=0$$
%    whenever $\nnorm{f_{1,1}}_{S_{1},\dots,S_{\ell}}=0$ for all $0\leq q\leq Q-1$. It is clear that there exists $Q\in\N$ such that $p_{i}(Q\cdot+q)$ is an integer coefficient polynomial for all $q\in\Z$ and $1\leq i\leq r$. So replacing $T_{i,j}$ by $T_{i,j}^{a_{i,j}}$ if necessary, it suffices to consider the case when $a_{i,j}=1$ and $p_{i}$ is an integer coefficient polynomial for all $i,j$.



%Replacing $T_{i,j}$ by $T_{i,j}^{a_{i,j}}$ if necessary, it suffices to consider the case when $a_{i,j}=1$ for all $i,j$. 


%We may first assume without loss of generality that  $a_{i,1}, a_{i,j}-a_{i,j'}\neq 0$ for all $i\in[r]$ and distinct $ j,j'\in[k_{i}]$. \BK{(Where is $a_{i,j}-a_{i,j'}\neq 0$ used?)}

\smallskip
\textbf{Case 1: $a_{1,1}$ has the maximum degree.}
\smallskip

We first prove this result under the assumption that $a_{1,1}$ has the maximum degree. 
%In this case, we may assume without loss of generality that $a_{i,j}$'s have the same degree for all $i\in[r']$ and $j\in[k_{i}]$, for otherwise we can relegate such $a_{i,j}$ to the portion of sequences indexed by $i>r'$.
%{\color{red} Thanks to the new condition (ii), we may replace this sentence with "In this case, all of $a_{i,j}$'s have the same degree greater than $m_{1}$ for all $i\in[r']$ and $j\in[k_{i}]$."}
In this case, all of $a_{i,j}$'s (for $i\in[r']$ and $j\in[\ell_{i}]$) have the same degree greater than $m_{1}$.
 We first consider the case when $\ell_{1}+\dots+\ell_{r'}=1$. Then $a_{1,1}$ has a higher degree than all other $a_{i,j}$'s, and so the maps $a_{1,1}$ and $a_{1,1}-a_{i,j}$ have the same leading coefficients. So the fact that $\CQ$ is {good for smoothing} implies that (\ref{goal}) is zero if $\nnorm{f_1}_{\be_{1,1}^{\times s}}=0$ for some $s=O_{d,J,\ell, \CQ}(1)$.

Now suppose that the conclusion holds when $\ell_{1}+\dots+\ell_{r'}=\ell'-1$ for some $\ell'\geq 2$; our aim is to show that it also holds when $\ell_{1}+\dots+\ell_{r'}=\ell'$. The argument below will make use of the following definitions.
\begin{definition}[Types, minority images and pre-images, and good terms]
    Let $$I:=\{(i,j)\in\N^{2}\colon\; i\in[r], j\in[\ell_{i}]\}\quad \text{and}\quad I':=\{(i,j)\in\N^{2}\colon\; i\in[r'], j\in[\ell_{i}]\}.$$
    We call a map $\eta\colon I\to I$ a \textit{type}, and we call it \emph{proper} if it satisfies the following conditions:
    \begin{enumerate}
        \item $\eta(I') \subseteq I'$;
        \item $\eta|_{I\backslash I'}$ is an identity;
        \item $\eta(i,j)\neq \eta(i,j')$ for all $i\in[r']$ and distinct $ j,j'\in[\ell_{i}]$.
    \end{enumerate}
    We say that $\tilde{i}\in I'$ is a \emph{minority image} of $\eta$ if
   $$\vert \eta^{-1}(\tilde{i})\vert=\min_{\tilde{i}'\in I'}\max(\vert\eta^{-1}(\tilde{i}')\vert,1).$$
   In other words, $\be_{\tilde{i}}$ is a vector  that appears least frequently in the multi-set $$\{\be_{\eta(i,j)}\colon\; i\in[r'],\; j\in [\ell_{i}]\}$$
%   $T_{\tilde{i}}$ is a transformation that appears least frequently in $\mathcal{F}_{\eta}$ 
(but it does appear at least once). We call every element in $\eta^{-1}(\tilde{i})$ a \emph{minority pre-image} of $\eta$. %and every term of $\mathcal{F}_{\eta}$ associated to a minority pre-image of $\eta$ a \emph{minority term} of $\mathcal{F}_{\eta}$.
   Lastly, we say that $(i_{0},j_{0})$ is a \emph{good term of} $\eta$  if
   there exists a positive integer $s=O_{d,J,\ell,\CQ}(1)$ and vectors $\bv_{1},\dots,\bv_{s}$ coming from the set
   %\footnote{Note that the vector $\hat{a}_{\eta(i,j)}\be_{\eta(i,j)}-\hat{a}_{\eta(i,j')}\be_{\eta(i,j')}$ might be the $\mathbf{0}$ vector if $\eta(i,j)$ and $\hat{a}_{\eta(1,1)} = \hat{a}_{\eta(1,1)}$ (this is not prevented by the properness of $\eta$, since we might still have $\eta(1,1) = \eta(i,j)$ if $i\neq 1$). To fix this and ensure that the leading coefficient is nonzero, we should replace $\hat{a}_{\eta(1,1)}\be_{\eta(1,1)}-\hat{a}_{\eta(i,j)}\be_{\eta(i,j)}$, which is the difference of the leading coefficients of some two terms, by the leading coefficient of the difference of these two terms. The former may be 0, the latter would not)}
            %   $$\be_{i_{0},j_{0}},\be_{i_{0},j_{0}}-\be_{i_{0},j}, j\in [k_{i_{0}}]\backslash\{j_{0}\},\be_{i,j}-\be_{i,j'}, i\in[r]\backslash\{i_{0}\}, j,j'\in[k_{i}], j\neq j'$$
             $$
             \{\hat{a}_{\eta(i_{0},j_{0})}\be_{\eta(i_{0},j_{0})},\hat{a}_{\eta(i,j)}\be_{\eta(i,j)}-\hat{a}_{\eta(i,j')}\be_{\eta(i,j')}:\; i\in[r'],\; j,j'\in[\ell_{i}],\; j\neq j'\}$$
         such that for all 1-bounded functions $f_{i,j}\in L^\infty(\mu)$ and sequences $b_1, \ldots, b_J\in\Span_\R\CQ$, the average 
\begin{multline}\label{E: average to be controled}
         \limsup_{N\to\infty}\sup_{\substack{
         \CD_1, \ldots, \CD_J \in \FD_d}} \sup_{|c_n|\leq 1}\norm{\E_{n\in [N]} c_{n}\cdot \prod_{i\in[r]}\prod_{j\in[\ell_{r}]}  T_{\eta(i,j)}^{\floor{\frac{\hat{a}_{\eta(i,j)}}{\hat{a}_{i,j}}a_{i,j}(n)}} f_{i,j}\cdot \prod_{j\in[J]}\CD_j(\floor{b_j(n)})}_{L^2(\mu)}%^{{O_{d,J,\ell,r', \CQ}(1)}}\\ 
         %\ll_{d, J, k, M, r', \CQ} \nnorm{f_1}_{S_{1},\dots,S_{s}}.
    \end{multline}
    is zero whenever $\nnorm{f_{i_{0},j_{0}}}_{\bv_{1},\dots,\bv_{s}}=0$.
    %, where for convenience we write $\eta(i,j):=(i,j)$ for $i>r'$ in the rest of the proof.

\end{definition}
    
   
   %{\color{red} Note there was a mistake in the previous version, where I wrote all of $\frac{\hat{a}_{\eta(i,j)}}{\hat{a}_{i,j}}$ as $\frac{\hat{a}_{i,j}}{\hat{a}_{\eta(i,j)}}$}
   To prove the statement, it suffices to show that any minority pre-image of any proper $\eta$  is good, and then apply this to the identity map $\eta(\tilde{i})=\tilde{i}$, which is proper and in which every term is minority.

    The next definition explains how we pass from more to less complex types.
    \begin{definition}[Descendance] %\BK{($\eta$ is defined on $I$, not $I'$, but $\tilde{i},\tilde{i}'$ take values in $I'$. Fixed that.)}
   For $\tilde{i},\tilde{i}'\in I'$, let $\sigma_{\tilde{i}\to \tilde{i}'}(\eta)\colon I\to I$ be the map given by $\sigma_{\tilde{i}\to \tilde{i}'}(\eta)(\tilde{i})=\eta(\tilde{i}')$ and $\sigma_{\tilde{i}\to \tilde{i}'}(\eta)(\tilde{j})=\eta(\tilde{j})$ for all  $\tilde{j}\in I\backslash\{\tilde{i}\}$. For proper maps $\eta$ and $\eta'$, we say that $\eta$ \emph{descends to} $\eta'$ (write as $\eta\searrow \eta'$) if $\eta'=\sigma_{\tilde{i}\to \tilde{i}'}(\eta)$ for some $\tilde{i},\tilde{i}'\in I'$ such that $\eta(\tilde{i})\neq \eta(\tilde{i}')$ and $\vert\eta^{-1}(\eta(\tilde{i}))\vert\leq \vert\eta^{-1}(\eta(\tilde{i}'))\vert$.     
    \end{definition}
   It is not hard to see that
   \begin{equation}\nonumber
      \eta\searrow \eta'\Rightarrow \sum_{\tilde{i}\in I}\vert\eta^{-1}(\tilde{i})\vert^{2}<\sum_{\tilde{i}\in I}\vert{\eta'}^{-1}(\tilde{i})\vert^{2},
   \end{equation}
 as an immediate consequence of which we have that the relation $\searrow$ is \emph{loopless}, meaning that there exist no $n\in\N$ and no proper types $\eta_{1},\dots,\eta_{n}$ with $\eta_{1}\searrow \eta_{2}\searrow\dots\searrow\eta_{n}\searrow\eta_{1}$.

 The base case of induction will come from averages of the following type.\footnote{A reader familiar with the earlier seminorm smoothing argument might inquire about the connection between extreme and basic types, the latter of which were playing the role of the base case in the smoothing arguments from \cite{DKKST24, FrKu22b, FrKu22a}. Despite superficial similarity, these two notions are rather different. For instance, the smoothing argument for pairwise dependent integer polynomials from \cite{FrKu22b} would allow a descendance $((1,1),(1,2);(2,1),(2,2))\searrow ((2,1),(1,2);(2,1),(2,2))\searrow((2,1),(2,1);(2,1),(2,2))$. The latter reduction is not permissible in this smoothing argument even though it is perfectly allowed in \cite{FrKu22b}. The type $((2,1),(2,1);(2,1),(2,2))$ reached this way is basic in the sense of \cite{FrKu22b}, but neither proper not extreme according to this paper's convention. Conversely, the arguments of this paper allow us to reduce types via the descendance $((1,1),(1,2), (1,3);(2,1),(2,2))\searrow ((1,1),(1,2), (2,1);(2,1),(2,2))\searrow ((1,1),(1,2), (2,1);(2,1),(1,2))$. The type reached at the end is extreme but not basic as defined in \cite{FrKu22b}.}
 \begin{definition}[Extreme types]
 We say that a proper map $\eta$ is \emph{extreme} if there is no proper map $\eta'$ with $\eta\searrow\eta'$.      
 \end{definition}
% {\color{red} extreme is not the same as basic! for example, $(T_{1}^{a_{1}(n)},T_{2}^{a_{1}(n)};T_{1}^{a_{2}(n)},T_{2}^{a_{2}(n)})$, where $a_{1},a_{2}$ are independent but with the same degree, is extreme but not basic.} \BK{Well, this particular tuple is basic (recall, basic means in your language that there exists $i_0\in[r']$ such that for all $i\in[r], j\in[k_i]$, we have $\eta(i,j) = (i_0, j')$ for some $j'\in[k_{i_0}]$. Wouldn't extreme types necessarily take the same form? If $\eta(i,j) = (i_1,j)$ for some $i_1\neq i_0$, then we should be able to descend further. It seems to me that basic = extreme + proper. In this case, the main improvement in your version of smoothing is that we don't need to consider non-proper types, which were the main source of issues in \cite{FrKu22b}.}
 
 % For two different proper maps $\eta,\eta'$, let $d(\eta,\eta')$ denote the largest number $n\in\N$ such that $\eta=\eta_{0}\searrow \eta_{1}\searrow\dots\searrow\eta_{n}=\eta'$ for some proper maps $\eta_{0},\dots,\eta_{n}$. 
 % If such an $n$ does not exist, then we write $d(\eta,\eta')=\infty$. \BK{(I don't think this option is even theoretically possible, is it? You can always take a minority preimage and descend until you hit an extreme type.)}  Clearly, if such an $n$ exists, the largest such $n$ also exists and is at most $k-1$ since $\searrow$ is loopless. \BK{Can you explain in a footnote where this bound comes from?}
 % %\BK{(I changed $k-1$ to $O_k(1)$ since I don't think the bound $k-1$ is correct. Consider the case when $k_1 = k_2 = k_3 = k_4 = 3$. Then you need 3 moves to get rid of all $(4,j)$. In the worst case, you will end up with an average of type corresponding to $k_1'=k'_2=k'_3 = 4$. Then you need 4 moves to get rid of all $(3,j)$. In the worst case, you end up with an average with $k''_1 = k''_2 = 6$, and then you need 6 extra moves to get rid of all $(2,j)$.)}. 
 % For any proper map $\eta$, denote $w(\eta)=0$ if $\eta$ is extreme and denote
 % $$w(\eta):=\max_{\eta' \text{ is extreme}\colon d(\eta,\eta')<\infty} d(\eta,\eta')$$
 % otherwise. Since $\searrow$ is loopless, there is always at least one extreme $\eta'$ with $d(\eta,\eta')<\infty$, and so $w(\eta)$ is well defined and $w(\eta)\leq k-1$. 
 % Informally speaking, $w(\eta)$ is the length of the longest chain of $\searrow$ which connects $\eta$ and an extreme map.

 
 For any proper map $\eta$, denote $w(\eta)=0$ if $\eta$ is extreme. Otherwise, 
 let $w(\eta)$ denote the largest number $n\in\N$ such that $\eta=\eta_{0}\searrow \eta_{1}\searrow\dots\searrow\eta_{n}$ for some proper maps $\eta_{0},\dots,\eta_{n}$ with $\eta_{n}$ being extreme. Since the relation $\searrow$ is loopless, $w(\eta)$ is well defined and is at most $O_{\ell}(1)$.  
    
   Using the fact that $\eta\searrow\eta'\Rightarrow w(\eta')<w(\eta)$, 
   to complete the proof, it suffices to prove the following claim (and then apply it repeatedly $O_\ell(1)$ times).

\begin{claim}[Minority pre-images are good]\label{C: smoothing claim}
    Let $\eta$ be a proper type. Suppose that for any proper type $\eta'$ with %$w(\eta')<w(\eta)$, 
   $\eta\searrow\eta'$, any minority pre-image of $\eta'$ is good. Then any minority pre-image of $\eta$ is good (in particular, any minority pre-image of $\eta$ is good when $\eta$ is extreme).
\end{claim}

\begin{example}
    This example shows how Claim \ref{C: smoothing claim} is used iteratively in the proof of Proposition \ref{P: new estimates example}. Phrased in the language of this section, the purpose of this proposition was to show that $(1,1)$ is a good term for $\eta = ((1,1),(1,2);(2,1),(2,2))$. Note that $\eta\searrow\eta'$ by changing the first entry of $\eta$ if and only if $\eta'$ is of type $((2,1),(1,2);(2,1),(2,2))$ or $((2,2),(1,2);(2,1),(2,2))$ (the other possible type $((1,2),(1,2);(2,1),(2,2))$ is not proper). Claim \ref{C: lower type example}, which showed that $(1,2)$ is good for either of these two proper types, can be seen as a special case of Claim \ref{C: smoothing claim}. 
    
    Recall from the proof of Proposition \ref{P: new estimates example} that Claim \ref{C: lower type example} was applied in the Ping step of the proof. This will remain true more generally: Claim \ref{C: smoothing claim} will always turn up in the Ping step of the smoothing argument.

   To prove Claim \ref{C: lower type example}, we would similarly want to show that the second term of $\eta' = ((2,1),(1,2);(2,1),(2,2))$ is good (and analogously with $((2,2),(1,2);(2,1),(2,2))$). Note that  $\eta'\searrow\eta''$ by changing the second entry of $\eta$ if and only if $\eta''=((2,1),(2,2);(2,1),$ $(2,2))$  (the other possible type $((2,1),(2,1);(2,1),(2,2))$ is not proper). This is an extreme type in which all four terms are minorities. Seminorm estimates for this type that are in line with Proposition \ref{newestimate3} will follow directly from the assumption of being good for smoothing.
   Indeed, the expression we wish to control is of the form
   $$\limsup_{N\to\infty}\norm{\E_{n\in [N]} T_{2,1}^{a_{1}(n)} f_{1,1}\cdot T_{2,2}^{a_{1}(n)} f_{1,2}\cdot T_{2,1}^{a_{2}(n)} f_{2,1}\cdot T_{2,2}^{a_{2}(n)} f_{2,2}}_{L^2(\mu)},$$
and since $\CQ$ is good for smoothing, this average is 0 if $\nnorm{f_{1,1}}_{\be_{2,1}^{\times s}, (\be_{2,1}-\be_{2,2})^{\times s}}=0$ for some $s\in\N$. This implies that the first term of $\eta''$ is good, and an analogous argument would give desired seminorm control in terms of seminorms of all the other terms.
\end{example}

   % {\color{red} We give an example to illustrate the use of the Claim (this is essentially the example I wrote in the proof of Proposition B.11). Suppose that $r=2$ and $k_{1}=k_{2}=2$ and all of $a_{i,j}$ have the same degree. For convenience we call $(\eta(1,1),\eta(1,2);\eta(2,1),\eta(2,2))$  the \emph{type} of $\eta$. We wish to show that $(1,1)$ is a good term for $\eta$ of the type $(1,2;3,4)$. Note that $\eta\searrow\eta'$ if and only if $\eta'$ is of type $(3,2;3,4)$ or $(4,2;3,4)$ (the type $(2,2;3,4)$ is not proper). 

   % By symmetry and by the Claim, it suffices to show that the minority term (i.e. the 2nd and 4th terms) of $(3,2;3,4)$ are good. By symmetry, it suffices to show that the 2nd term of $\eta$ of type $(3,2;3,4)$ is good. Note that  $\eta\searrow\eta'$ if and only if $\eta'$ is of type $(3,4;3,4)$  (the type $(3,3;3,4)$ is not proper). 

   % In this case, all the 4 terms of $(3,4;3,4)$ are minorities. By the Claim, it suffices to show that all the terms of $(3,4;3,4)$ are good. However, this is an extreme case, and all the terms are good because of the leading coefficient comparison result (in this case the expression is $T_{3}^{\frac{\hat{a}_{2,1}}{\hat{a}_{1,1}}a_{1,1}(n)}f_{1,1}\cdot T_{4}^{\frac{\hat{a}_{2,2}}{\hat{a}_{1,2}}a_{1,2}(n)}f_{1,2}\cdot T_{3}^{a_{2,1}(n)}f_{2,1}\cdot T_{4}^{a_{2,2}(n)}f_{2,2}$).
   % }
   % \BK{The example is good. I think we should have a separate subsection (before the main argument) that works out this example at the same level of detail as in \cite[Section 8.1]{FrKu22a} and \cite[Section 4]{FrKu22a}, with an emphasis on both how the induction differs from the arguments in \cite{DKKST24, FrKu22a, FrKu22b}, but also how the seminorm control obtained in the Pong step is lifted through relative concatenation. Such examples tend to facilitate understanding much more than full proofs.}
   % {\color{red} the examples are in the next subsection}

  
We proceed to prove Claim \ref{C: smoothing claim}. Assume without loss of generality that $(1,1)$ is a minority pre-image of $\eta$. 
%, and moreover $(1,1)\in \eta^{-1}(1,1)$. \BK{(I don't think we should assume $(1,1)\in \eta^{-1}(1,1)$, anyway we're not using it.)} {\color{red} see the comment on the next page on where this is used. Of course we could rewrite everything without this assumption, if it looks better.}
Let %$\CG_{\eta}$ be the set of the following vectors: 
\begin{align*}
    \CG_\eta :=\{\hat{a}_{\eta(1,1)}\be_{\eta(1,1)}, \hat{a}_{\eta(i,j)}\be_{\eta(i,j)}-\hat{a}_{\eta(i,j')}\be_{\eta(i,j')}:\; i\in[r'],\; j,j'\in[\ell_{i}],\; j\neq j'\}.
\end{align*}
The properness of $\eta$ implies in particular that the elements of $\CG_\eta$ are nonzero.
             % $$\be_{\eta(1,1)},\be_{\eta(1,1)}-\be_{\eta(1,j)}, 2\leq j\leq \ell_{1}, \be_{\eta(i,j)}-\be_{\eta(i,j')}, 2\leq i\leq r, 1\leq j<j'\leq \ell_{i}.$$
It suffices to show that (\ref{E: average to be controled})  can be controlled by a box seminorm of $f_{1,1}$ involving only transformations from $\CG_{\eta}$, which we will do by establishing the following property for $k = 0$.

\begin{definition}[Property $k$]
    We say that \emph{Property $k$} holds if there exist $t\in\N$ and vectors  $\bv_1, \ldots, \bv_t, \bu_1, \ldots, \bu_k\in\R^\ell$ (where $\bv_i\in \CG_{\eta}$ while $\bu_i$'s are the leading coefficients of $\frac{\hat{a}_{\eta(1,1)}}{\hat{a}_{1,1}}a_{1,1}\be_{\eta(1,1)}-\frac{\hat{a}_{\eta(i,j)}}{\hat{a}_{i,j}}a_{i,j}\be_{\eta(i,j)}$ for some $(i,j)\in I'\backslash\{(1,1)\}$), such that for all functions $f_{i,j}\in L^\infty(\mu)$ and sequences $b_1, \ldots, b_J\in\Span_\R\CQ$, the average \eqref{E: average to be controled} is 0 whenever $$\nnorm{f_{1,1}}_{\bv_{1},\dots,\bv_{t},\bu_{1},\dots,\bu_k}=0.$$
\end{definition}
Since $\CQ$ is good for smoothing %and $\eta$ is proper,
and since $$\deg(a_{1,1})=\dots=\deg(a_{r',1})>\max_{r'< i\leq r}\max_{j\in[\ell_i]}\deg(a_{i,j}),$$  Property $k$ holds for some $k\in\N$.     
 Our goal is to show that Property 0 holds. To achieve this, it suffices to show that if Property $k$ holds for some $k\in\N$, then Property $k-1$ also holds.
     
 %    Assume that
  %    $\bu_k=\hat{a}_{1,1}\be_{1,1}-\hat{a}_{i_{0},j_{0}}\be_{\eta(i_{0},j_{0})}$ for some $(i_{0},j_{0})\in I\backslash\{(1,1)\}$.
%      If $\hat{a}_{i_{0},j_{0}}\be_{\eta(i_{0},j_{0})}=\hat{a}_{1,j}\be_{\eta(1,j)}$ for some $2\leq j\leq \ell_{1}$, then $\nu_{\ell}$  belongs to $\CG_{\eta}$
 %    and so Property $z-1$ holds by absorbing $\bu_{\ell_{0}+1},\dots,\bu_k$ into $\bv_{1},\dots,\bv_{t}$ (and by enlarging $t$). Now assume that $\hat{a}_{i_{0},j_{0}}\be_{\eta(i_{0},j_{0})}\neq\hat{a}_{1,j}\be_{\eta(1,j)}$ for some $2\leq j\leq \ell_{1}$.
      

          Assume that 
$\bu_k$ is the leading coefficient of $$\frac{\hat{a}_{\eta(1,1)}}{\hat{a}_{1,1}}a_{1,1}\be_{\eta(1,1)}-\frac{\hat{a}_{\eta(i_{0},j_{0})}}{\hat{a}_{i_{0},j_{0}}}a_{i_{0},j_{0}}\be_{\eta(i_{0},j_{0})}$$ for some $(i_{0},j_{0})\in I'\backslash\{(1,1)\}$. We consider several possible scenarios.

If $\eta(i_{0},j_{0})=\eta(1,1)$, then since $\eta$ is proper, we have that $i_{0}\geq 2$ and thus $a_{1,1}$ is independent from $a_{i_{0},j_{0}}$. This implies that $\bu_k$ is a nonzero multiple of $\be_{\eta(1,1)}$. Since $\CG_{\eta}$ contains $\hat{a}_{\eta(1,1)}\be_{\eta(1,1)}$, Property $k-1$ holds by Lemma \ref{L: basic properties of seminorms} after relabeling $\bu_k$ as $\bv_{t+1}$. 

%So we now assume that $\eta(i_{0},j_{0})\neq\eta(1,j)$. In this case $\bu_k=\hat{a}_{\eta(1,1)}\be_{\eta(1,1)}-\hat{a}_{\eta(i_{0},j_{0})}\be_{\eta(i_{0},j_{0})}$. 
If instead $\eta(i_{0},j_{0})=\eta(1,j)$ for some $2\leq j\leq \ell_1$, then $\bu_k$ belongs to $\CG_{\eta}$ and again Property $k-1$ holds by by relabeling $\bu_k$ as $\bv_{t+1}$.
       
       It remains to consider the case when $\eta(i_{0},j_{0})\neq \eta(1,j)$ for all $1\leq j\leq \ell_1$. In this case,  %$\bu_k=\hat{a}_{\eta(1,1)}\be_{\eta(1,1)}-\hat{a}_{\eta(i_{0},j_{0})}\be_{\eta(i_{0},j_{0})}$ and 
       $i_{0}\neq 1$. 
       %\BK{(It is not clear why we can assume $i_{0}\neq 1$. I think that this is: (i) not true, (ii) unnecessary. Consider the descendance from $((2,1),(1,2);(2,1),(2,2))$ to $((2,1),(2,2);(2,1),(2,2))$ in the example above. If $\bu_k =\be_{1,2} - \be_{2,1}$, then $\be_{2,1}$ might be seen as coming from both the first and the third term of the original type.)} {\color{red} could you explain more about the question? The reason $i_{0}\neq 1$ is, if $i_{0}=1$, then $\eta(i_{0},j_{0})=\eta(1,j)$ for some $j$ (in fact for $j=j_{0}$). But this case is covered by the previous two paragraphs}
       Denote $\eta':=\sigma_{(1,1)\to (i_{0},j_{0})}(\eta)$.
  %  
 %     $\bu_k=\hat{a}_{\eta(1,1)}\be_{\eta(1,1)}-\hat{a}_{\eta(i_{0},j_{0})}\be_{\eta(i_{0},j_{0})}$ for some $(i_{0},j_{0})\in I'\backslash\{(1,1)\}$. If $\eta(i_{0},j_{0})=\eta(1,j)$ for some $2\leq j\leq k_{1}$, then
 %     $\bu_k=\hat{a}_{\eta(1,1)}\be_{\eta(1,1)}-\hat{a}_{\eta(1,j)}\be_{\eta(1,j)}$
 %      belongs to $\CG_{\eta}$ and so Property $k-1$ holds by absorbing $\bu_k$ into $\bv_{1},\dots,\bv_{t}$ (and enlarging $t$). So we can assume that %%$\eta(1,j)\neq \tilde{i}$ for all $j\in[k_{1}]$. Pick any $(i_{0},j_{0})\in I'$ with $\eta(i_{0},j_{0})=\tilde{i}$; 
  %     this is not the case. Then $i_{0}\neq 1$. Denote $\eta':=\sigma_{(1,1)\to (i_{0},j_{0})}(\eta)$.
     Since 
  $$\eta'(1,1)=\eta(i_{0},j_{0})\neq \eta(1,j)=\eta'(1,j)$$ for all $2\leq j\leq \ell_1$, 
     the map $\eta'$ is proper. Moreover, $\eta\searrow \eta'$ %(and hence also $w(\eta')<w(\eta)$) 
     as a consequence of $(1,1)$ being a minority pre-image of $\eta$.
      %Since (1,1) is a minority pre-image of $\eta$, we have that $\eta\searrow \eta'$. %Since $\vert(\sigma_{1\to j}(\eta))^{-1}(i)\vert\geq 2$ (this is because $\eta^{-1}(i)$ is at least 1, and $\sigma_{1\to j}$ adds in another term). 
     %Since $\eta\searrow \eta'$, it is clear that $w(\eta')<w(\eta)$.

 
     
      In particular, if $\eta$ is extreme (i.e. $w(\eta) = 0$), then the non-existence of such $\eta'$ forces $\eta(i_{0},j_{0})=\eta(1,j)$ for some $j\in[\ell_1]$, and we are back in one of the previous cases. 
      
      So, in what follows we assume that $\eta(i_{0},j_{0})\neq \eta(1,j)$ for all $1\leq j\leq \ell_1$ since all the remaining cases can be covered, and we take $\eta'=\sigma_{(1,1)\to (i_{0},j_{0})}(\eta)$.
      
      %\BK{I think that for emphasis, we should have a separate claiming that extreme types are all of the form $\eta(i,j) = (1,*)$ for $i\in[r']$, and that we get the claimed bounds in this case more or less automatically from the property of being good for seminorm estimates. The latter requires us to also ensure that if $\eta$ is extreme, then all the possible values of $\eta(i,j)$ with $i\in[r']$ can be realized as $\eta(1,j')$ for some $j'$, which in general might not be the case. 
     
     % (which we can, Nikos and I did sth like this in \cite[Proposition 6.8](vi){FrKu22b}) that if $\eta$ is extreme, then $\eta(1,j) = (1,j)$ for all $j\in [k_1]$. To get this, we might need to impose more stringent conditions on what we consider minority preimage. Here is one example. Consider the type $(1,2; 3,4,5)$. The current framework allows us to pass from this to, say, $(1,2; 1,4,5)$, then to $(1,4; 1, 4, 5)$, and then we're stuck
     
     % }%Denote $\eta':=\sigma_{(1,1)\to (i_{0},j_{0})}(\eta)$.

     \smallskip
     \textbf{Step 1: Ping.}
     \smallskip
     
     %\BK{(I think that everywhere below, we want $\floor{b_j(n)}$ in place of $q_{j}(n)$?)} 
     Suppose  that (\ref{E: average to be controled}) is positive.
   %  \begin{multline}\label{E: average to be controled}
  %    \limsup_{N\to\infty}\sup_{\substack{
  %       \CD_1, \ldots, \CD_J \in \FD_d}} \sup_{|c_n|\leq 1}\Bigl\Vert\E_{n\in[N]}c_{n}\cdot\prod_{(i,j)\in I}  T_{\eta_{(i,j)}}^{\floor{\frac{\hat{a}_{i,j}}{\hat{a}_{\eta(i,j)}}a_{i,j}(n)}} f_{i,j}\cdot\prod_{j\in[J]}\mathcal{D}_{j}(\floor{b_{j}(n)})\Bigr\Vert_{2}>0.
   %  \end{multline}
     %is positive.
     It follows from  \cite[Lemma 11.7]{DKKST24} that 
     % \begin{equation}\nonumber
     % \begin{split}
     % &\qquad\limsup_{N\to\infty}\sup_{\substack{
     %     \CD_1, \ldots, \CD_J \in \FD_d}} \sup_{|c_n|\leq 1}\lim_{N\to\infty}\Bigl\Vert\E_{n\in[N]}c_{n}\cdot T_{\eta(1,1)}^{\floor{a_{1,1}(n)}}\tilde{f}_{1,1}\prod_{(i,j)\in I\backslash\{(1,1)\}}T_{\eta(i,j)}^{\floor{\frac{\hat{a}_{i,j}}{\hat{a}_{\eta(i,j)}}a_{i,j}(n)}}f_{i,j}
     %     \\&\prod_{i'\in [r']}{T'_{i'}}^{\floor{a'_{i'}(n)}}f'_{i}\prod_{j\in[J]}\mathcal{D}_{j}(q_{j}(n))\Bigr\Vert_{2}>0
     % \end{split}
     % \end{equation}
    % \BK{(I think that everywhere below, we want $\frac{\hat{a}_{1,1}}{\hat{a}_{\eta(1,1)}}a_{1,1}(n)$ in place of $a_{1,1}(n)$?)} {\color{red} this is because I assumed that $\eta(1,1)=(1,1)$!}
     \begin{multline*}
         \limsup_{N\to\infty}\sup_{\substack{
         \CD_1, \ldots, \CD_J \in \FD_d}} \sup_{|c_n|\leq 1}\lim_{N\to\infty}\Bigl\Vert\E_{n\in[N]}c_{n}\cdot T_{\eta(1,1)}^{\floor{\frac{\hat{a}_{\eta(1,1)}}{\hat{a}_{1,1}}a_{1,1}(n)}}\tilde{f}_{1,1}\cdot\prod_{(i,j)\in I\backslash\{(1,1)\}}T_{\eta(i,j)}^{\floor{\frac{\hat{a}_{\eta(i,j)}}{\hat{a}_{i,j}}a_{i,j}(n)}}f_{i,j}\\
         \prod_{j\in[J]}\mathcal{D}_{j}(\floor{b_{j}(n)})\Bigr\Vert_{L^2(\mu)}>0
     \end{multline*}
     for some $\tilde{f}_{1,1}$ which is ``structured'' in that it is the weak limit $\lim\limits_{w\to\infty} A_{N_{w}}$ for 
% \begin{equation}\label{E: tilded}
%     \begin{split}
%         &\qquad A_{N_{w}}:=\E_{n\in[N_{w}]} \overline{c}'_{n}T_{\eta(1,1)}^{-\floor{a_{1,1}(n)}}\overline{g}_{w}\prod_{(i,j)\in I\backslash\{(1,1)\}}T_{\eta(i,j)}^{\floor{\frac{\hat{a}_{i,j}}{\hat{a}_{\eta(i,j)}}a_{i,j}(n)}}T_{\eta(1,1)}^{-\floor{a_{1,1}(n)}}\overline{f}_{i,j}
%         \\&\prod_{i'\in[r']}{T'_{i'}}^{\floor{a'_{i'}(n)}}T_{\eta(1,1)}^{-\floor{a_{1,1}(n)}}\overline{f_i'}\prod_{j\in [J]}T_{\eta(1,1)}^{-\floor{a_{1,1}(n)}}\overline{\mathcal{D}}'_{j}(q_{j}(n))
%     \end{split}
% \end{equation}
\begin{multline}\label{E: tilded}
             A_{N_{w}}:=\E_{n\in[N_{w}]} \overline{c}'_{n}T_{\eta(1,1)}^{-\floor{\frac{\hat{a}_{\eta(1,1)}}{\hat{a}_{1,1}}a_{1,1}(n)}}\overline{g}_{w}\prod_{(i,j)\in I\backslash\{(1,1)\}}T_{\eta(i,j)}^{\floor{\frac{\hat{a}_{\eta(i,j)}}{\hat{a}_{i,j}}a_{i,j}(n)}}T_{\eta(1,1)}^{-\floor{\frac{\hat{a}_{\eta(1,1)}}{\hat{a}_{1,1}}a_{1,1}(n)}}\overline{f}_{i,j}\\
         \prod_{j\in [J]}T_{\eta(1,1)}^{-\floor{\frac{\hat{a}_{\eta(1,1)}}{\hat{a}_{1,1}}a_{1,1}(n)}}\overline{\mathcal{D}}'_{j}(\floor{b_{j}(n)})
\end{multline}
     for some increasing sequence $(N_{w})_{w}$,  1-bounded weights $(c'_{n})_n$,  1-bounded functions $g_{w}$, and dual sequences $\mathcal{D}'_{1},\dots,\mathcal{D}'_{J}\in\mathfrak{D}_{d}$. %In addition, we have that $\tilde{f}_{1,1}=\lim_{w\to\infty} A_{N'_{w}}$ strongly in $L^{2}(\mu)$ for every subsequence $(N'_{w})_{w}$ of $(N_{w})_{w}$.
     By assumption, we have $$\nnorm{\tilde{f}_{1,1}}_{\bv_{1},\dots,\bv_{t},\bu_{1},\dots,\bu_k}>0.$$
      We then unpack this seminorm using \cite[Lemma 11.11]{DKKST24} and the formula \eqref{E: tilded} for $\tilde{f}_{1,1}$; after composing the resulting expression with $T_{\eta(1,1)}^{\floor{\frac{\hat{a}_{\eta(1,1)}}{\hat{a}_{1,1}}a_{1,1}(n)}}$ and applying the Cauchy-Schwarz inequality, we have  
%     \begin{equation}\nonumber
%    \begin{split}
  %      &\qquad   \liminf_{M\to\infty}\E_{\um, \um',\um'',\um'''\in [\pm M]^{t+k-1}}\limsup_{w\to\infty} \E_{n\in[N_w]}
 %          \\&  \int\Delta_{\substack{(\floor{\bv_1 m_1'} - \floor{\bv_1 m_1}, \floor{\bv_1 m_1'''} - \floor{\bv_1 m_1''}), \ldots,\\ (\floor{\bu_{k-1} m_{t+k-1}'} - \floor{\bu_{k-1} m_{t+k-1}}, \floor{\bu_{k-1} m_{t+k-1}'''} - \floor{\bu_{k-1} m_{t+k-1}''})}}\;
 %            \Bigl(\overline{c}'_{n}T_{\eta(1,1)}^{-\floor{a_{1,1}(n)}}\overline{g}_{w}
 %            \\&\prod_{(i,j)\in I\backslash\{(1,1)\}}T_{\eta(i,j)}^{\floor{a_{i,j}(n)}}T_{\eta(1,1)}^{-\floor{a_{i,j}(n)}}\overline{f}_{i,j}
  %      \prod_{i'\in[r']}{T'_{i'}}^{\floor{a'_{i'}(n)}}T_{\eta(1,1)}^{-\floor{a_{1,1}(n)}}\overline{f_i'}\prod_{j\in [J]}T_{\eta(1,1)}^{-\floor{a_{1,1}(n)}}\overline{\mathcal{D}}'_{j}(q_{j}(n))\Bigr) \cdot u_{\um,\um',\um'',\um'''} \, d\mu
 %   \end{split}
%    \end{equation} 
  %     \begin{equation}\nonumber
 %   \begin{split}
 %       &\qquad   \liminf_{M\to\infty}\E_{\um, \um',\um'',\um'''\in [\pm M]^{t+k-1}}\limsup_{w\to\infty} \E_{n\in[N_w]}
 %          \\&  \int\Delta_{\substack{(\floor{\bv_1 m_1'} - \floor{\bv_1 m_1}, \floor{\bv_1 m_1'''} - \floor{\bv_1 m_1''}), \ldots,\\ (\floor{\bu_{k-1} m_{t+k-1}'} - \floor{\bu_{k-1} m_{t+k-1}}, \floor{\bu_{k-1} m_{t+k-1}'''} - \floor{\bu_{k-1} m_{t+k-1}''})}}\;
 %            \Bigl(\overline{c}'_{n}\overline{g}_{w}
 %            \\&\prod_{(i,j)\in I\backslash\{(1,1)\}}T_{\eta(i,j)}^{\floor{a_{i,j}(n)}}\overline{f}_{i,j}
 %       \prod_{i'\in[r']}{T'_{i'}}^{\floor{a'_{i'}(n)}}\overline{f_i'}\prod_{j\in [J]}\overline{\mathcal{D}}'_{j}(q_{j}(n))\Bigr) \cdot T_{\eta(1,1)}^{\floor{a_{1,1}(n)}}u_{\um,\um',\um'',\um'''} \, d\mu
%    \end{split}
%    \end{equation}
% \begin{equation}\nonumber
% \begin{split}
%    &\qquad\liminf_{M\to\infty}\E_{\um, \um',\um'',\um'''\in [\pm M]^{t+k-1}}\limsup_{N\to\infty} \sup_{|c_n|\leq 1}\Bigl\Vert
%    \E_{n\in[N]}c_n\cdot T_{\eta(1,1)}^{\floor{a_{1,1}(n)}} u_{\um,\um',\um'',\um'''}
%    \\&\prod_{(i,j)\in I\backslash\{(1,1)\}}T_{\eta(i,j)}^{\floor{\frac{\hat{a}_{i,j}}{\hat{a}_{\eta(i,j)}}a_{i,j}(n)}}f_{i,j,\um,\um',\um'',\um'''}\cdot\prod_{i'\in[r']}{T'_{i'}}^{\floor{a'_{i'}(n)}}f'_{i',\um,\um',\um'',\um'''}\\ \prod_{j\in[J]} \CD'_{j,\um,\um',\um'',\um'''}(\floor{q_j(n)})\Bigr\Vert_{L^2(\mu)}
% \end{split}
% \end{equation}
\begin{multline*}
    \liminf_{M\to\infty}\E_{\um, \um',\um'',\um'''\in [\pm M]^{t+k-1}}\limsup_{N\to\infty} \sup_{|c_n|\leq 1}\Bigl\Vert
   \E_{n\in[N]}c_n\cdot T_{\eta(1,1)}^{\floor{\frac{\hat{a}_{\eta(1,1)}}{\hat{a}_{1,1}}a_{1,1}(n)}} h_{\um,\um',\um'',\um'''}
   \\
   \prod_{(i,j)\in I\backslash\{(1,1)\}}T_{\eta(i,j)}^{\floor{\frac{\hat{a}_{\eta(i,j)}}{\hat{a}_{i,j}}a_{i,j}(n)}}f_{i,j,\um,\um',\um'',\um'''}\cdot \prod_{j\in[J]} \CD'_{j,\um,\um',\um'',\um'''}(\floor{b_j(n)})\Bigr\Vert_{L^2(\mu)}>0
\end{multline*}
  %  for 1-bounded functions $u_{\um,\um',\um'',\um'''}\in L^\infty(\mu)$ which are product of $2^{t+k-1}$ 1-bounded approximately ${\bu_k}$-invariant functions.
        where
        \begin{align*}
            f_{i,j,\um,\um',\um'',\um'''} &= \Delta_{\substack{(\floor{\bv_1 m_1'} - \floor{\bv_1 m_1}, \floor{\bv_1 m_1'''} - \floor{\bv_1 m_1''}), \ldots,\\ (\floor{\bu_{k-1} m_{t+k-1}'} - \floor{\bu_{k-1} m_{t+k-1}}, \floor{\bu_{k-1} m_{t+k-1}'''} - \floor{\bu_{k-1} m_{t+k-1}''})}}\overline{f_{i,j}}\\
           %  f'_{i',\um,\um',\um'',\um'''} &= \Delta_{\substack{(\floor{\bv_1 m_1'} - \floor{\bv_1 m_1}, \floor{\bv_1 m_1'''} - \floor{\bv_1 m_1''}), \ldots,\\ (\floor{\bu_{k-1} m_{t+k-1}'} - \floor{\bu_{k-1} m_{t+k-1}}, \floor{\bu_{k-1} m_{t+k-1}'''} - \floor{\bu_{k-1} m_{t+k-1}''})}}\overline{f'_{i}}\\
            \textrm{and}\quad\CD'_{j,\um,\um',\um'',\um'''} &= \Delta_{\substack{(\floor{\bv_1 m_1'} - \floor{\bv_1 m_1}, \floor{\bv_1 m_1'''} - \floor{\bv_1 m_1''}), \ldots,\\ (\floor{\bu_{k-1} m_{t+k-1}'} - \floor{\bu_{k-1} m_{t+k-1}}, \floor{\bu_{k-1} m_{t+k-1}'''} - \floor{\bu_{k-1} m_{t+k-1}''})}}\overline{\CD'_{j}}
        \end{align*}
    and each $h_{\um,\um',\um'',\um'''}$ is a product of $2^{t+k-1}$ 1-bounded \emph{approximately $\bu_k$-invariant functions}, meaning functions of the form $\E\limits_{n\in\Z}T^{\floor{\bu_kn}}f$ for some $f\in L^\infty(\mu)$.
   % \BK{(define)}. 
    Since $\bu_k=\hat{a}_{\eta(1,1)}\be_{\eta(1,1)}-\hat{a}_{\eta'(1,1)}\be_{\eta'(1,1)}$, it follows from \cite[Lemma 11.10]{DKKST24} that
%     \begin{equation}\label{qtrgg1}
% \begin{split}
%    &\qquad\liminf_{M\to\infty}\E_{\um, \um',\um'',\um'''\in [\pm M]^{t+k-1}}\limsup_{N\to\infty} \sup_{|c_n|\leq 1}\Bigl\Vert
%    \prod_{j\in[J]}\E_{n\in[N]}c_n\cdot T_{\eta'(1,1)}^{\floor{\frac{\hat{a}_{1,1}}{\hat{a}_{\eta'(1,1)}}a_{1,1}(n)}} u'_{\um,\um',\um'',\um'''}
%    \\&\prod_{(i,j)\in I\backslash\{(1,1)\}}T_{\eta'(i,j)}^{\floor{\frac{\hat{a}_{i,j}}{\hat{a}_{\eta'(i,j)}}a_{i,j}(n)}}f_{i,j,\um,\um',\um'',\um'''}\prod_{i'\in[r']}{T'_{i'}}^{\floor{a'_{i'}(n)}}f'_{i',\um,\um',\um'',\um'''} \cdot \prod_{j\in[J]} \CD'_{j,\um,\um',\um'',\um'''}(\floor{q_j(n)})\Bigr\Vert_{L^2(\mu)}
% \end{split}
% \end{equation}
\begin{multline}\label{qtrgg1}
       \liminf_{M\to\infty}\E_{\um, \um',\um'',\um'''\in [\pm M]^{t+k-1}}\limsup_{N\to\infty} \sup_{|c_n|\leq 1}\Bigl\Vert
   \prod_{j\in[J]}\E_{n\in[N]}c_n\cdot T_{\eta'(1,1)}^{\floor{\frac{\hat{a}_{\eta'(1,1)}}{\hat{a}_{1,1}}a_{1,1}(n)}} h'_{\um,\um',\um'',\um'''}
   \\\prod_{(i,j)\in I\backslash\{(1,1)\}}T_{\eta'(i,j)}^{\floor{\frac{\hat{a}_{\eta'(i,j)}}{\hat{a}_{i,j}}a_{i,j}(n)}}f_{i,j,\um,\um',\um'',\um'''}\cdot  \prod_{j\in[J]} \CD'_{j,\um,\um',\um'',\um'''}(\floor{b_j(n)})\Bigr\Vert_{L^2(\mu)}>0
\end{multline}
 for some 1-bounded functions $h'_{\um,\um',\um'',\um'''}$. It is at this step that we have passed from an average of type $\eta$ to one of type $\eta'$, and this key reduction was enabled by the $\bu_k$-almost invariance of $h_{\um,\um',\um'',\um'''}$'s.
      
     
      Note that we may find some $\tilde{u}=(u',u'')\in I'\backslash\{(1,1)\}$ which is a minority pre-image of $\eta'$ (to see this, let $\tilde{v}$ be a minority pre-image of $\eta'$; if $\tilde{v}\neq (1,1)$, then we pick $\tilde{u}=\tilde{v}$ whereas if $\tilde{v}=(1,1)$, then we pick $\tilde{u}=(i_{0},j_{0})$ since   $\eta'(1,1)=\eta'(i_{0},j_{0})$). By assumption,
      the $\tilde{u}$-th term is good for $\eta'$.
      %the term of $\mathcal{F}_{\eta'}$ indexed by $\tilde{u}$ %can be controlled by transformations in the set 
 % $\CG_{\eta'}$.   
   %  $$\CG_{\eta'}=\{{\color{green} T_{\eta(\tilde{u})}}, T_{\eta(\tilde{u})}T_{\eta'(u',j)}^{-1}, T_{\eta'(j',j'')}T_{\eta'(j,j''')}^{-1}\colon 1\leq j\leq k_{u'}, j\neq u'', 1\leq j'\leq r, j'\neq u', 1\leq j''<j'''\leq k_{j'}\}.$$
           Therefore, it follows %from (\ref{qtrgg1})
           from our inductive assumption\footnote{And specifically by its soft quantitative strengthening, which is a consequence of \cite[Proposition~A.2]{FrKu22a}.}
           that for a set of $(\um,\um',\um'',\um'')$ of positive lower density, the seminorm
%               \begin{equation}\nonumber
% \begin{split}
%    &\qquad\nnorm{f_{\tilde{u},\um,\um',\um'',\um'''}}_{\bu'_{1},\dots,\bu'_{t'},(\hat{a}_{\eta(\tilde{u})}\be_{\eta(\tilde{u})})^{s'}}
%    \\&=\nnorm{\Delta_{\substack{(\floor{\bv_1 m_1'} - \floor{\bv_1 m_1}, \floor{\bv_1 m_1'''} - \floor{\bv_1 m_1''}), \ldots,\\ (\floor{\bu_{k-1} m_{t+k-1}'} - \floor{\bu_{k-1} m_{t+k-1}}, \floor{\bu_{k-1} m_{t+k-1}'''} - \floor{\bu_{k-1} m_{t+k-1}''})}}\overline{f_{\tilde{u}}}}_{\bu'_{1},\dots,\bu'_{t'},(\hat{a}_{\eta(\tilde{u})}\be_{\eta(\tilde{u})})^{s'}}
% \end{split}
% \end{equation}      
\begin{multline*}
    \nnorm{f_{\tilde{u},\um,\um',\um'',\um'''}}_{\bu'_{1},\dots,\bu'_{t'},(\hat{a}_{\eta(\tilde{u})}\be_{\eta(\tilde{u})})^{\times s'}}
   \\=\Bignnorm{\Delta_{\substack{(\floor{\bv_1 m_1'} - \floor{\bv_1 m_1}, \floor{\bv_1 m_1'''} - \floor{\bv_1 m_1''}), \ldots,\\ (\floor{\bu_{k-1} m_{t+k-1}'} - \floor{\bu_{k-1} m_{t+k-1}}, \floor{\bu_{k-1} m_{t+k-1}'''} - \floor{\bu_{k-1} m_{t+k-1}''})}}\overline{f_{\tilde{u}}}}_{\bu'_{1},\dots,\bu'_{t'},(\hat{a}_{\eta(\tilde{u})}\be_{\eta(\tilde{u})})^{\times s'}}
\end{multline*}
           is (uniformly) bounded away from 0
         for some $s',t'\in\N$ (depending only on $d,J,\ell,\CQ$) and $\bu'_{1},\dots,\bu'_{t'}\in \CG_{\eta'}\backslash\{\hat{a}_{\eta(\tilde{u})}\be_{\eta(\tilde{u})}\}$.
 %     So we have that the $x$-th term of $\mathcal{F}_{\eta}$ is controlled by transformations in $\CG_{\eta'}$.
   Averaging over all $\um,\um',\um'',\um''$, raising it to the power $2^{s'+t'}$ using the H{\"o}lder inequality and applying the inductive formula for generalized box seminorms, we get the auxiliary norm control
      $$\nnorm{f_{\tilde{u}}}_{(\hat{a}_{\eta(\tilde{u})}\be_{\eta(\tilde{u})})^{\times s'}, \bu'_{1},\dots,\bu'_{t'},\bv_{1},\dots,\bv_{t},\bu_{1},\dots,\bu_{k-1}}>0.$$
     
%      Therefore, it follows from (\ref{qtrgg1}) that $$\liminf_{H\to\infty}\E_{\uh,\uh'\in[N]^{t+k-1}}\nnorm{\Delta_{S_{1},\dots,S_{t},R_{1},\dots,R_{\ell-1};\uh-\uh'}f_{x}}_{T_{\eta(x)}^{\times s},R'_{1},\dots,R'_{t'}}>0$$
%      for some $s,t'\in\N$ and $R'_{1},\dots,R'_{t'}\in \CG_{\eta'}\backslash\{T_{\eta(x)}\}$,
 %%     So we have that the $x$-th term of $\mathcal{F}_{\eta}$ is controlled by transformations in $\CG_{\eta'}$.
%      which implies that
%      $$\nnorm{f_{x}}_{T_{\eta(x)}^{\times s},R'_{1},\dots,R'_{t'},S_{1},\dots,S_{t},R_{1},\dots,R_{\ell-1}}>0.$$

      
    %   \textbf{Step 2 (pong).} %Since (\ref{qtrgg0}) implies that $\nnorm{f_{x}}_{T_{\eta(x)}^{\times s},R'_{1},\dots,R'_{t'}}>0$, 
    %  {\color{blue} the pong step is unnecessary} \BK{(I think you still need a pong step since the original estimate is for $f_{1,1}$ while the new one obtained in ping is for some $f_{(i,j)}$ with $(i,j)\neq (1,1)$. The price you pay in the ping step is that the original function $f_{1,1}$ is replaced by sth arbitrary, and so one of the points of the pong step is to go back to an average that contains $f_{1,1}$.
      
    %   If so, then maybe the need for $T_{i,j}'s$ other than $T_{1,1}$ in the conclusion will disappear (although these transformations will reappear after applying ergodicity assumptions).)}
    %   {\color{red} I included the Pong step below.  However, I think other $T_{i,j}'s$ are still needed at least for this method (although I think they are removable by some better method). This is because if $\deg(p_{i})>\deg(p_{1})$, then $T_{i,j}$ would appear in the DKS estimate. This $T_{i,j}$ can not be removed by the Ping step.
    %   The Pong step (i.e. next step) allows you to remove $T_{i,j}$, but not other $T_{i',j'}$. We need to do it in a more clever way to get rid of all of $T_{i,j}$.}
    % \BK{(Fair enough)}

  \smallskip
         \textbf{Step 2: Pong.} 
        \smallskip
         
         %Since (\ref{qtrgg0}) implies that $\nnorm{f_{x}}_{T_{\eta(x)}^{\times s},R'_{1},\dots,R'_{t'}}>0$, 
     Suppose that  (\ref{E: average to be controled}) is positive. 
      It follows from \cite[Lemma 11.7]{DKKST24} that 
     % \begin{equation}\nonumber
     % \begin{split}
     % &\qquad\limsup_{N\to\infty}\sup_{\substack{
     %     \CD_1, \ldots, \CD_J \in \FD_d}} \sup_{|c_n|\leq 1}\lim_{N\to\infty}\Bigl\Vert\E_{n\in[N]}c_{n}\cdot T_{\eta(\tilde{u})}^{\floor{\frac{\hat{a}_{\tilde{u}}}{\hat{a}_{\eta(\tilde{u})}}a_{\tilde{u}}(n)}}\tilde{f}_{\tilde{u}}\prod_{(i,j)\in I\backslash\{\tilde{u}\}}T_{\eta(i,j)}^{\floor{\frac{\hat{a}_{i,j}}{\hat{a}_{\eta(i,j)}}a_{i,j}(n)}}f_{i,j}
     %     \\&\prod_{i'\in [r']}{T'_{i'}}^{\floor{a'_{i'}(n)}}f'_{i}\prod_{j\in[J]}\mathcal{D}_{j}(q_{j}(n))\Bigr\Vert_{2}>0
     % \end{split}
     % \end{equation}
     \begin{multline*}
         \limsup_{N\to\infty}\sup_{\substack{
         \CD_1, \ldots, \CD_J \in \FD_d}} \sup_{|c_n|\leq 1}\lim_{N\to\infty}\Bigl\Vert\E_{n\in[N]}c_{n}\cdot T_{\eta(\tilde{u})}^{\floor{\frac{\hat{a}_{\eta(\tilde{u})}}{\hat{a}_{\tilde{u}}}a_{\tilde{u}}(n)}}\tilde{f}_{\tilde{u}}\cdot\prod_{(i,j)\in I\backslash\{\tilde{u}\}}T_{\eta(i,j)}^{\floor{\frac{\hat{a}_{\eta(i,j)}}{\hat{a}_{i,j}}a_{i,j}(n)}}f_{i,j}\\
         \prod_{j\in[J]}\mathcal{D}_{j}(\floor{b_{j}(n)})\Bigr\Vert_{2}>0
     \end{multline*}
     for some $\tilde{f}_{\tilde{u}}$ which is the weak limit $\lim\limits_{w\to\infty} A_{N_{w}}$ for 
% \begin{equation}\label{E: tilded2}
%     \begin{split}
%         &\qquad A_{N_{w}}:=\E_{n\in[N_{w}]} \overline{c}'_{n}T_{\eta(\tilde{u})}^{-\floor{\frac{\hat{a}_{\tilde{u}}}{\hat{a}_{\eta(\tilde{u})}}a_{\tilde{u}}(n)}}\overline{g}_{w}\prod_{(i,j)\in I\backslash\{\tilde{u}\}}T_{\eta(i,j)}^{\floor{\frac{\hat{a}_{i,j}}{\hat{a}_{\eta(i,j)}}a_{i,j}(n)}}T_{\eta(\tilde{u})}^{-\floor{\frac{\hat{a}_{\tilde{u}}}{\hat{a}_{\eta(\tilde{u})}}a_{\tilde{u}}(n)}}\overline{f}_{i,j}
%         \\&\prod_{i'\in[r']}{T'_{i'}}^{\floor{a'_{i'}(n)}}T_{\eta(\tilde{u})}^{-\floor{\frac{\hat{a}_{\tilde{u}}}{\hat{a}_{\eta(\tilde{u})}}a_{\tilde{u}}(n)}}\overline{f_i'}\prod_{j\in [J]}T_{\eta(\tilde{u})}^{-\floor{\frac{\hat{a}_{\tilde{u}}}{\hat{a}_{\eta(\tilde{u})}}a_{\tilde{u}}(n)}}\overline{\mathcal{D}}'_{j}(q_{j}(n))
%     \end{split}
% \end{equation}
\begin{multline}\label{E: tilded2}
    A_{N_{w}}:=\E_{n\in[N_{w}]} \overline{c}'_{n}\cdot T_{\eta(\tilde{u})}^{-\floor{\frac{\hat{a}_{\eta(\tilde{u})}}{\hat{a}_{\tilde{u}}}a_{\tilde{u}}(n)}}\overline{g}_{w}\cdot\prod_{(i,j)\in I\backslash\{\tilde{u}\}}T_{\eta(i,j)}^{\floor{\frac{\hat{a}_{\eta(i,j)}}{\hat{a}_{i,j}}a_{i,j}(n)}}T_{\eta(\tilde{u})}^{-\floor{\frac{\hat{a}_{\eta(\tilde{u})}}{\hat{a}_{\tilde{u}}}a_{\tilde{u}}(n)}}\overline{f}_{i,j}
        \\
       \prod_{j\in [J]}T_{\eta(\tilde{u})}^{-\floor{\frac{\hat{a}_{\eta(\tilde{u})}}{\hat{a}_{\tilde{u}}}a_{\tilde{u}}(n)}}\overline{\mathcal{D}}'_{j}(\floor{b_{j}(n)})
\end{multline}
     for some increasing sequence $(N_{w})_{w}$, 1-bounded weights $(c'_{n})_n$, 1-bounded function $g_{w}$, and dual sequences $\mathcal{D}'_{1},\dots,\mathcal{D}'_{J}\in\mathfrak{D}_{d}$.
    The seminorm control established in the Ping step then gives $$\nnorm{\tilde{f}_{\tilde{u}}}_{(\hat{a}_{\eta(\tilde{u})}\be_{\eta(\tilde{u})})^{\times s'},\bu'_{1},\dots,\bu'_{t'},\bv_{1},\dots,\bv_{t},\bu_{1},\dots,\bu_{k-1}}>0.$$
      We unpack this seminorm using \cite[Lemma 11.11]{DKKST24} and the formula \eqref{E: tilded2} for $\tilde{f}_{\tilde{u}}$; after composing the resulting expression with $T_{\eta(\tilde{u})}^{\floor{\frac{\hat{a}_{\eta(\tilde{u})}}{\hat{a}_{\tilde{u}}}a_{\tilde{u}}(n)}}$ and applying the Cauchy-Schwarz inequality, we have that 
% \begin{equation}\nonumber
% \begin{split}
%    &\qquad\liminf_{M\to\infty}\E_{\um, \um',\um'',\um'''\in [\pm M]^{t+t'+k-1}}\limsup_{N\to\infty} \sup_{|c_n|\leq 1}\Bigl\Vert
%    \prod_{j\in[J]}\E_{n\in[N]}c_n\cdot T_{\eta(\tilde{u})}^{\floor{\frac{\hat{a}_{\tilde{u}}}{\hat{a}_{\eta(\tilde{u})}}a_{\tilde{u}}(n)}} u_{\um,\um',\um'',\um'''}
%    \\&\prod_{(i,j)\in I\backslash\{\tilde{u}\}}T_{\eta(i,j)}^{\floor{\frac{\hat{a}_{i,j}}{\hat{a}_{\eta(i,j)}}a_{i,j}(n)}}f_{i,j,\um,\um',\um'',\um'''}\cdot\prod_{i'\in[r']}{T'_{i'}}^{\floor{a'_{i'}(n)}}f'_{i',\um,\um',\um'',\um'''} \\
%    \prod_{j\in[J]} \CD'_{j,\um,\um',\um'',\um'''}(\floor{q_j(n)})\Bigr\Vert_{L^2(\mu)}
% \end{split}
% \end{equation}
\begin{multline*}
    \liminf_{M\to\infty}\E_{\um, \um',\um'',\um'''\in [\pm M]^{t+t'+k-1}}\limsup_{N\to\infty} \sup_{|c_n|\leq 1}\Bigl\Vert
  \E_{n\in[N]}c_n\cdot T_{\eta(\tilde{u})}^{\floor{\frac{\hat{a}_{\eta(\tilde{u})}}{\hat{a}_{\tilde{u}}}a_{\tilde{u}}(n)}} h_{\um,\um',\um'',\um'''}
   \\
   \prod_{(i,j)\in I\backslash\{\tilde{u}\}}T_{\eta(i,j)}^{\floor{\frac{\hat{a}_{\eta(i,j)}}{\hat{a}_{i,j}}a_{i,j}(n)}}f_{i,j,\um,\um',\um'',\um'''}\cdot
   \prod_{j\in[J]} \CD'_{j,\um,\um',\um'',\um'''}(\floor{b_j(n)})\Bigr\Vert_{L^2(\mu)} > 0,
\end{multline*}
  %  for 1-bounded functions $u_{\um,\um',\um'',\um'''}\in L^\infty(\mu)$ which are product of $2^{t+k-1}$ 1-bounded approximately ${\bu_k}$-invariant functions.
        where
        \begin{align*}
            f_{i,j,\um,\um',\um'',\um'''} &= \Delta_{\substack{(\floor{\bv_1 m_1'} - \floor{\bv_1 m_1}, \floor{\bv_1 m_1'''} - \floor{\bv_1 m_1''}), \ldots,\\ (\floor{\bu_{k-1} m_{t+t'+k-1}'} - \floor{\bu_{k-1} m_{t+t'+k-1}}, \floor{\bu_{k-1} m_{t+t'+k-1}'''} - \floor{\bu_{k-1} m_{t+t'+k-1}''})}}\overline{f_{i,j}},\\
           %  f'_{i',\um,\um',\um'',\um'''} &= \Delta_{\substack{(\floor{\bv_1 m_1'} - \floor{\bv_1 m_1}, \floor{\bv_1 m_1'''} - \floor{\bv_1 m_1''}), \ldots,\\ (\floor{\bu_{k-1} m_{t+t'+k-1}'} - \floor{\bu_{k-1} m_{t+t'+k-1}}, \floor{\bu_{k-1} m_{t+t'+k-1}'''} - \floor{\bu_{k-1} m_{t+t'+k-1}''})}}\overline{f'_{i}}\\
            \CD'_{j,\um,\um',\um'',\um'''} &= \Delta_{\substack{(\floor{\bv_1 m_1'} - \floor{\bv_1 m_1}, \floor{\bv_1 m_1'''} - \floor{\bv_1 m_1''}), \ldots,\\ (\floor{\bu_{k-1} m_{t+t'+k-1}'} - \floor{\bu_{k-1} m_{t+t'+k-1}}, \floor{\bu_{k-1} m_{t+t'+k-1}'''} - \floor{\bu_{k-1} m_{t+t'+k-1}''})}}\overline{\CD'_{j}},
        \end{align*}
    and each $h_{\um,\um',\um'',\um'''}$ is a product of $2^{t+t'+k-1}$ 1-bounded dual functions of level $s'$ with respect to $T_{\eta(\tilde{u})}$.
       By induction on $k'$,  for a set of $(\um,\um',\um'',\um'')$ of positive lower density, we have that
       \begin{multline*}
           \nnorm{f_{1,1,\um,\um',\um'',\um'''}}_{\bu''_{1},\dots,\bu''_{t''}}\\
           =\Bignnorm{\Delta_{\substack{(\floor{\bv_1 m_1'} - \floor{\bv_1 m_1}, \floor{\bv_1 m_1'''} - \floor{\bv_1 m_1''}), \ldots,\\ (\floor{\bu_{k-1} m_{t+t'+k-1}'} - \floor{\bu_{k-1} m_{t+t'+k-1}}, \floor{\bu_{k-1} m_{t+t'+k-1}'''} - \floor{\bu_{k-1} m_{t+t'+k-1}''})}}\overline{f_{1,1}}}_{\bu''_{1},\dots,\bu''_{t''}}
       \end{multline*}
           %$$\nnorm{f_{1,1,\um,\um',\um'',\um'''}}_{\bu''_{1},\dots,\bu''_{t''}}=\nnorm{\Delta_{\substack{(\floor{\bv_1 m_1'} - \floor{\bv_1 m_1}, \floor{\bv_1 m_1'''} - \floor{\bv_1 m_1''}), \ldots,\\ (\floor{\bu_{k-1} m_{t+t'+k-1}'} - \floor{\bu_{k-1} m_{t+t'+k-1}}, \floor{\bu_{k-1} m_{t+t'+k-1}'''} - \floor{\bu_{k-1} m_{t+t'+k-1}''})}}\overline{f_{1,1}}}_{\bu''_{1},\dots,\bu''_{t''}}$$
           is (uniformly) bounded away from 0
         for some $t''\in\N$ (depending only on $d,J,\ell,r',\CQ$) and $\bu''_{1},\dots,\bu''_{t''}$ belonging to the set $\CG_{\eta}$.\footnote{Once again, we use here \cite[Proposition A.2]{FrKu22a} to upgrade qualitative seminorm control to a soft quantitative estimate.}
 %     So we have that the $x$-th term of $\mathcal{F}_{\eta}$ is controlled by transformations in $\CG_{\eta'}$.
   Averaging over all $\um,\um',\um'',\um''$, raising it to the power $2^{t''}$, using the H{\"o}lder inequality and applying the inductive formula for generalized box seminorms, we get the auxiliary norm control
       $$\nnorm{f_{1,1}}_{\bv_{1},\dots,\bv_{t},\bu_{1},\dots,\bu_{k-1},\bu'_{1},\dots,\bu'_{t'},\bu''_{1},\dots,\bu''_{t''}}>0.$$
     In conclusion, we have that  (\ref{E: average to be controled}) is zero if $\nnorm{f_{1,1}}_{\bv_{1},\dots,\bv_{t},\bu_{1},\dots,\bu_{k-1},\bu'_{1},\dots,\bu'_{t'},\bu''_{1},\dots,\bu''_{t''}}=0$.

    \smallskip
      \textbf{Step 3: Relative concatenation.} 
      \smallskip

      Combining our original seminorm control with the conclusion of the Pong step and Lemma \ref{L: basic properties of seminorms}, we deduce that     
       (\ref{E: average to be controled}) is zero if
       \begin{align*}
           \min(\nnorm{f_{1,1}}^{+}_{\bv_{1},\dots,\bv_{t},\bu_{1},\dots,\bu_k}, \nnorm{f_{1,1}}^{+}_{\bv_{1},\dots,\bv_{t},\bu_{1},\dots,\bu_{k-1},\bu'_{1},\dots,\bu'_{t'},\bu''_{1},\dots,\bu''_{t''}}) = 0.
       \end{align*}
    We want to combine the previous estimates to show that \eqref{E: average to be controled} is 0 if 
       \begin{equation}\label{pojj}
           \nnorm{f_{1,1}}^{+}_{\bv_{1},\dots,\bv_{t},\bu_{1},\dots,\bu_{k-1},\{\langle \bu_k, \bu\rangle: \bu\in\mathcal{R}\}} = 0
       \end{equation}
 where $\mathcal{R}$ is the multi-set $\{\bu'_{1},\dots,\bu'_{t'},\bu''_{1},\dots,\bu''_{t''}\}$. Using Proposition \ref{P: structure theorem}
 % , it suffices to show the vanishing of \eqref{E: average to be controled} under the assumption that
 % \begin{align*}
 %     \E(f_{1,1}|\CZ^+_{\bv_{1},\dots,\bv_{t},\bu_{1},\dots,\bu_{k-1},\{\langle \bu_k, \bu\rangle: \bu\in\mathcal{R}\}})=0.
 % \end{align*}
 and Corollary \ref{C: relative concatenation}, we decompose $f = g_1 + g_2$ where $g_1, g_2\in L^\infty(\mu)$ are 1-bounded, and 
     \begin{align*}
     \nnorm{g_1}^+_{\bv_{1},\dots,\bv_{t},\bu_{1},\dots,\bu_{k-1},\bu'_{1},\dots,\bu'_{t'},\bu''_{1},\dots,\bu''_{t''}} =  \nnorm{g_2}^+_{\bv_{1},\dots,\bv_{t},\bu_{1},\dots,\bu_k} = 0.
     %\E(g_1\vert \mathcal{Z}^{+}_{\bv_{1},\dots,\bv_{t},\bu_{1},\dots,\bu_{k-1},\bu'_{1},\dots,\bu'_{t'},\bu''_{1},\dots,\bu''_{t''}})=\E(g_2\vert \mathcal{Z}^{+}_{\bv_{1},\dots,\bv_{t},\bu_{1},\dots,\bu_k})=0.    
     \end{align*}
      Hence the contribution of $g_1$ to \eqref{E: average to be controled} vanishes because of the seminorm control established in the Pong step while the contribution of $g_2$ is 0 due to the original seminorm control. It follows by the triangle inequality that \eqref{E: average to be controled} is 0 as well, and so the seminorm in \eqref{pojj} controls \eqref{E: average to be controled}.
      

      
%       Hence it also remains unchanged if we replace $f_{1,1}$ by
% \begin{equation}\nonumber
%     \begin{split}
%         &\E(\E(f_{1,1}\vert \mathcal{Z}^{+}_{\bv_{1},\dots,\bv_{t},\bu_{1},\dots,\bu_{k-1},\bu'_{1},\dots,\bu'_{t'},\bu''_{1},\dots,\bu''_{t''}})\vert \mathcal{Z}^{+}_{\bv_{1},\dots,\bv_{t},\bu_{1},\dots,\bu_k})
%         \\&=\E(f\vert \mathcal{Z}^{+}_{\bv_{1},\dots,\bv_{t},\bu_{1},\dots,\bu_{k-1},\bu'_{1},\dots,\bu'_{t'},\bu''_{1},\dots,\bu''_{t''}}\cap \mathcal{Z}^{+}_{\bv_{1},\dots,\bv_{t},\bu_{1},\dots,\bu_k}).
%     \end{split}
% \end{equation}
% \BK{(Can you justify this identity? For general sub-$\sigma$-algebras $\CA, \CB$, we have $\E(f|\CA\cap\CB)\neq \E(\E(f|\CA)|\CB)$.}
%      By Proposition \ref{T: relative concatenation}, this means that (\ref{E: average to be controled}) remains unchanged if we replace $f_{1,1}$ by $$\E(f_{1,1}\vert \mathcal{Z}^{+}_{\bv_{1},\dots,\bv_{t},\bu_{1},\dots,\bu_{k-1},\{\langle \bu_k, \bu\rangle: \bu\in\mathcal{R}\}}),$$
%   where $\mathcal{R}$ is the multi-set $\{\bu'_{1},\dots,\bu'_{t'},\bu''_{1},\dots,\bu''_{t''}\}$.
%  Therefore, (\ref{E: average to be controled}) is 0 if
%      $$\nnorm{f_{1,1}}^{+}_{\bv_{1},\dots,\bv_{t},\bu_{1},\dots,\bu_{k-1},\{\langle \bu_k, \bu\rangle: \bu\in\mathcal{R}\}}=0.$$


     

%       {\color{red} Below is the new proof, remove everything from Step 3 to here. I just copied Tao--Zielger's proof of Corollary 1.20 of their concatenation paper, check if everything makes sense:
      
%       We first combine the previous estimates to show that
%        Therefore, 
%        \begin{equation}\label{pojj}
%            \text{(\ref{E: average to be controled}) is 0 if } \nnorm{f_{1,1}}^{+}_{\bv_{1},\dots,\bv_{t},\bu_{1},\dots,\bu_{k-1},\{\langle \bu_k, \bu\rangle: \bu\in\mathcal{R}\}}=0,
%        \end{equation}
%  where $\mathcal{R}$ is the multi-set $\{\bu'_{1},\dots,\bu'_{t'},\bu''_{1},\dots,\bu''_{t''}\}$.

%  For convenience denote $\mathcal{I}:=\{\bv_{1},\dots,\bv_{t},\bu_{1},\dots,\bu_{k-1},\{\langle \bu_k, \bu\rangle: \bu\in\mathcal{R}\}\}$, $\mathcal{I}_{1}:=\{\bv_{1},\dots,\bv_{t},\bu_{1},\dots,\bu_k\}$ and $\mathcal{I}_{2}:=\{\bv_{1},\dots,\bv_{t},\bu_{1},\dots,\bu_{k-1},\bu'_{1},\dots,\bu'_{t'},\bu''_{1},\dots,\bu''_{t''}\}$, all of which are understood as multi-sets. 
%      By Proposition \ref{P: structure theorem}, we have that $f_{1,1}\in L^{\infty}(\mathcal{Z}^{+}_{\mathcal{I}})$.
%       By Propositions \ref{P: structure theorem} and \ref{T: relative concatenation}, any function in $L^{\infty}(\mathcal{Z}^{+}_{\mathcal{I}_{1}})$ which is orthogonal to $L^{\infty}(\mathcal{Z}^{+}_{\mathcal{I}})$ will have vanishing $\nnorm{\cdot}_{\mathcal{I}_{2}}^{+}$ seminorm with respect to the factor system whose $\sigma$-algebra is $\mathcal{Z}^{+}_{\mathcal{I}_{1}}$. This senimorm remains unchanged when we pass from the factor to the original system. So such a function must be orthogonal to $L^{\infty}(\mathcal{Z}^{+}_{\mathcal{I}_{2}})$ by Proposition \ref{P: structure theorem}. In conclusion, the restrictions of the spaces $L^{\infty}(\mathcal{Z}^{+}_{\mathcal{I}_{1}})$ and  $L^{\infty}(\mathcal{Z}^{+}_{\mathcal{I}_{2}})$ to $L^{\infty}(\mathcal{Z}^{+}_{\mathcal{I}})$
% are orthogonal. Therefore, we may write $f_{1,1}$ as $f_{1,1}=g_{1}+g_{2}$ with $g_{i}$ being orthogonal to $\mathcal{Z}^{+}_{\mathcal{I}_{i}}$ for $i=1,2$. By assumption and by the conclusion of the Pong step, (\ref{E: average to be controled}) is 0 if $f_{1,1}$ is replaced by $g_{1}$ or $g_{2}$. So (\ref{pojj}) holds.
%       }

 Recall that $\bu'_{1},\dots,\bu'_{t'}\in \CG_{\eta'}\backslash\{\hat{a}_{\eta(\tilde{u})}\be_{\eta(\tilde{u})}\}$  and $\bu''_{1},\dots,\bu''_{t''}\in\CG_\eta$. Hence  every element $\bu\in\mathcal{R}$ lies either in $\CG_{\eta}$ or in $\CG_{\eta'}\backslash(\CG_{\eta}\cup\{\hat{a}_{\eta(\tilde{u})}\be_{\eta(\tilde{u})}\})$. In the latter case, it takes the form 
 $$\hat{a}_{\eta'(1,1)}\be_{\eta'(1,1)}-\hat{a}_{\eta'(1,j)}\be_{\eta'(1,j)}=\hat{a}_{\eta(i_{0},j_{0})}\be_{\eta(i_{0},j_{0})}-\hat{a}_{\eta(1,j)}\be_{\eta(1,j)}$$ 
 for some $2\leq j\leq \ell_1$. Recalling that $\bu_k=\hat{a}_{\eta(1,1)}\be_{\eta(1,1)}-\hat{a}_{\eta(i_{0},j_{0})}\be_{\eta(i_{0},j_{0})}$, we deduce that the group $\langle \bu_k, \bu\rangle$ contains the element $\hat{a}_{\eta(1,1)}\be_{\eta(1,1)}-\hat{a}_{\eta(1,j)}\be_{\eta(1,j)}$ which belongs to $\CG_{\eta}$.
     In conclusion, it follows from (\ref{pojj}) that (\ref{E: average to be controled}) is zero if 
     $$\nnorm{f_{1,1}}^{+}_{{\bv}_{1},\dots,{\bv}_{t''},\bu_{1},\dots,\bu_{k-1}}=0$$
     for some $t''\in\N$ and $\bv_1, \ldots, \bv_{t''}\in\CG_{\eta}$. 
   By Lemma \ref{L: basic properties of seminorms}, we have that  (\ref{E: average to be controled}) is zero if 
     $$\nnorm{f_{1,1}}_{\bv_{1},\dots,{\bv}_{t''},\bv,\bu_{1},\dots,\bu_{k-1}}=0,$$
     for any $\bv\in\CG_{\eta}$.
     So Property $k-1$ holds and we are done.
    %\end{proof}
    
     
 %  Finally, we may use Proposition \ref{newestimate3} to provide a generalization of Proposition \ref{newestimate}: 
    
 %    \begin{proposition}\label{newestimate4}
  %      Let $d, m, M,r,r'\in\N, k_{1},\dots,k_{r}\in\N, k=k_{1}+\dots+k_{r}$, $J\in\N_0$. Suppose that the family of sequences $\CQ$
 %       is {good for smoothing}, and the sequences $a_{i,j}, 1\leq i\leq r, 1\leq j\leq k_{i}, a'_{i'}, 1\leq i'\leq r'\in\Span_\R\CQ$ are of complexity at most $M$ with the following properties:
%        \begin{itemize}
%             \item {\color{green} the degree of $a_{1,1}$ does not exceed any of $a_{i,j}$};
%\item {\color{green} we no longer require that} all of $a_{i,j}$ are of the same degree;
%             \item the degree of $a_{1,1}$ is higher than any of $a'_{i'}$;
 %            \item $a_{i,j}$ is independent of $a_{i',j'}$ if and only if $i\neq i'$;
 %            \item $a_{1,1}, a_{1,1}-a_{1,j}\neq 0$ for $2\leq j\leq k_{1}$.
 %       \end{itemize}
 
  %                Then there exists a positive integer $s=O_{d,J,\ell,r', \CQ}(1)$ and vectors $\bv_{1},\dots,\bv_{s}$ coming from the following set 
 %               % $$\{{\color{green}\hat{a}_{i,j}\be_{i,j}, 1\leq i\leq r, 1\leq j\leq k_{i},} \hat{a}_{1,1}\be_{1,1}-\hat{a}_{1,j}\be_{1,j}, 2\leq j\leq k_{1}, \hat{a}_{i,j}\be_{i,j}-\hat{a}_{i,j'}\be_{i,j'}, 2\leq i\leq r, 1\leq j<j'\leq k_{i}\}$$
    %            $$\{{\color{green}\hat{a}_{i,j}\be_{i,j},} \hat{a}_{i,j}\be_{i,j}-\hat{a}_{i,j'}\be_{i,j'}, i\in[r], j,j'\in[k_{i}], j\neq j'\}$$
    %     such that for all systems $(X, \CX, \mu, T_{i,j}, 1\leq i\leq r, 1\leq j\leq k_{i}, T'_{i'}, 1\leq i'\leq r')$,  1-bounded functions $f_{i,j},f'_{i'}\in L^\infty(\mu)$ and sequences $b_1, \ldots, b_J\in\Span_\R\CQ$, we have 
%\begin{multline}\label{goal}
 %        \limsup_{N\to\infty}\sup_{\substack{
 %        \CD_1, \ldots, \CD_J \in \FD_d}} \sup_{|c_n|\leq 1}\norm{\E_{n\in [N]} c_{n}\cdot \prod_{i\in[r]}\prod_{j\in[k_{r}]}  T_{i,j}^{\floor{a_{i,j}(n)}} f_{i,j} \prod_{i'\in[r']}  {T'_{i'}}^{\floor{a'_{i'}(n)}} f'_{i'}\cdot \prod_{j\in[J]}\CD_j(\floor{b_j(n)})}_{L^2(\mu)}%^{{O_{d,J,\ell,r', \CQ}(1)}}\\ 
  %       %\ll_{d, J, k, M, r', \CQ} \nnorm{f_1}_{S_{1},\dots,S_{s}}.
  %  \end{multline}
  %  is zero whenever $\nnorm{f_1}_{\bv_{1},\dots,\bv_{s}}=0$.
 %   \end{proposition}
%\begin{proof}

\

\smallskip
\textbf{Case 2: $a_{1,1}$ does not have the maximum degree.}
\smallskip

We now drop the assumption that $a_{1,1}$ has maximum degree.
Suppose that (\ref{goal}) is positive. By the previous case, Proposition \ref{newestimate3} holds when $\ell_1+\dots+\ell_{r'}=1$. We now induct on $\ell'$. Suppose that Proposition \ref{newestimate3} holds when $\ell_1+\dots+\ell_{r'}=\ell'-1$ for some $\ell'\geq 2$. We show that it also holds when $\ell_1+\dots+\ell_{r'}=\ell'$. 
% Let 
% \begin{align*}
%     \veps = \limsup_{N\to\infty}\sup_{\substack{\CD_1, \ldots, \CD_J \in \FD_d}} \sup_{|c_n|\leq 1}\norm{\E_{n\in [N]} c_{n}\cdot \prod_{j\in[\ell]}  T_{\eta_j}^{\floor{a_j(n)}} f_j \cdot \prod_{j\in[J]}\CD_j(\floor{b_j(n)})}_{L^2(\mu)}.
% \end{align*}
If $a_{1,1}$ has maximum degree, then the conclusion follows from the previous case. If not, then we may assume without loss of generality that $a_{2,1}$ has maximum degree. Denote $\tilde{u}:=(2,1)$.
   It follows from  \cite[Lemma 11.7]{DKKST24} that 
     % \begin{equation}\nonumber
     % \begin{split}
     % &\qquad\limsup_{N\to\infty}\sup_{\substack{
     %     \CD_1, \ldots, \CD_J \in \FD_d}} \sup_{|c_n|\leq 1}\lim_{N\to\infty}\Bigl\Vert\E_{n\in[N]}c_{n}\cdot T_{\tilde{u}}^{\floor{a_{\tilde{u}}(n)}}\tilde{f}_{\tilde{u}}\prod_{(i,j)\in I\backslash\{\tilde{u}\}}T_{i,j}^{\floor{a_{i,j}(n)}}f_{i,j}
     %     \\&\prod_{i'\in [r']}{T'_{i'}}^{\floor{a'_{i'}(n)}}f'_{i}\prod_{j\in[J]}\mathcal{D}_{j}(q_{j}(n))\Bigr\Vert_{2}>0
     % \end{split}
     % \end{equation}
     \begin{multline*}
         \limsup_{N\to\infty}\sup_{\substack{
         \CD_1, \ldots, \CD_J \in \FD_d}} \sup_{|c_n|\leq 1}\lim_{N\to\infty}\Bigl\Vert\E_{n\in[N]}c_{n}\cdot T_{\tilde{u}}^{\floor{a_{\tilde{u}}(n)}}\tilde{f}_{\tilde{u}}\cdot\prod_{(i,j)\in I\backslash\{\tilde{u}\}}T_{i,j}^{\floor{a_{i,j}(n)}}f_{i,j}
         \\
         \prod_{j\in[J]}\mathcal{D}_{j}(\floor{b_{j}(n)})\Bigr\Vert_{2}>0
     \end{multline*}
     for some $\tilde{f}_{\tilde{u}}$ which is the weak limit $\lim\limits_{w\to\infty} A_{N_{w}}$ for 
% \begin{equation}\label{E: tilded22}
%     \begin{split}
%         &\qquad A_{N_{w}}:=\E_{n\in[N_{w}]} \overline{c}'_{n}T_{\tilde{u}}^{-\floor{a_{\tilde{u}}(n)}}\overline{g}_{w}\prod_{(i,j)\in I\backslash\{\tilde{u}\}}T_{i,j}^{\floor{a_{i,j}(n)}}T_{\tilde{u}}^{-\floor{a_{\tilde{u}}(n)}}\overline{f}_{i,j}
%         \\&\prod_{i'\in[r']}{T'_{i'}}^{\floor{a'_{i'}(n)}}T_{\tilde{u}}^{-\floor{a_{\tilde{u}}(n)}}\overline{f_i'}\prod_{j\in [J]}T_{\tilde{u}}^{-\floor{a_{\tilde{u}}(n)}}\overline{\mathcal{D}}'_{j}(q_{j}(n))
%     \end{split}
% \end{equation}
\begin{multline}\label{E: tilded22}
    A_{N_{w}}:=\E_{n\in[N_{w}]} \overline{c}'_{n}\cdot T_{\tilde{u}}^{-\floor{a_{\tilde{u}}(n)}}\overline{g}_{w}\cdot\prod_{(i,j)\in I\backslash\{\tilde{u}\}}T_{i,j}^{\floor{a_{i,j}(n)}}T_{\tilde{u}}^{-\floor{a_{\tilde{u}}(n)}}\overline{f}_{i,j}
        \\
        \prod_{j\in [J]}T_{\tilde{u}}^{-\floor{a_{\tilde{u}}(n)}}\overline{\mathcal{D}}'_{j}(\floor{b_{j}(n)})
\end{multline}
     for some increasing sequence $(N_{w})_{w}$, 1-bounded weights $(c'_{n})_n$, 1-bounded function $g_{w}$, and dual sequences $\mathcal{D}'_{1},\dots,\mathcal{D}'_{J}\in\mathfrak{D}_{d}$.
     Since $a_{2,1}$ has maximum degree, we can apply the previous case to conclude that $\nnorm{\tilde{f}_{\tilde{u}}}_{(\hat{a}_{\tilde{u}}\be_{\tilde{u}})^{\times s'},\bu_{1},\dots,\bu_{t}}>0$ for some $t\in\N$ (depending only on $d,J,\ell,\CQ$) and some $\bu_{1},\dots,\bu_{t}\in\CG$.\footnote{In fact, the previous case implies that $\bu_{1},\dots,\bu_{t}$ necessarily take the form $\hat{a}_{i,j}\be_{i,j}-\hat{a}_{i,j'}\be_{i,j'}$ with ($1\leq j<j'\leq \ell_i$) for those $i\in[r']$ for which $a_{i,j}, a_{i,j'}$  have the highest degree; this is a subset of $\CG$.}
      We then unpack this seminorm using \cite[Lemma~11.11]{DKKST24} and the formula \eqref{E: tilded2} for $\tilde{f}_{\tilde{u}}$; after composing the resulting expression with $T_{\tilde{u}}^{\floor{a_{\tilde{u}}(n)}}$ and applying the Cauchy-Schwarz inequality, we have that 
      \begin{multline*}
          \liminf_{M\to\infty}\E_{\um, \um',\um'',\um'''\in [\pm M]^{t}}\limsup_{N\to\infty} \sup_{|c_n|\leq 1}\Bigl\Vert\E_{n\in[N]}c_n\cdot T_{\tilde{u}}^{\floor{a_{\tilde{u}}(n)}} h_{\um,\um',\um'',\um'''}
   \\\prod_{(i,j)\in I\backslash\{\tilde{u}\}}T_{i,j}^{\floor{a_{i,j}(n)}}f_{i,j,\um,\um',\um'',\um'''}\cdot \prod_{j\in[J]} \CD'_{j,\um,\um',\um'',\um'''}(\floor{b_{j}(n)})\Bigr\Vert_{L^2(\mu)} >0,
      \end{multline*}
% \begin{equation}\nonumber
% \begin{split}
%    &\qquad\liminf_{M\to\infty}\E_{\um, \um',\um'',\um'''\in [\pm M]^{t+t'+k-1}}\limsup_{N\to\infty} \sup_{|c_n|\leq 1}\Bigl\Vert
%    \prod_{j\in[J]}\E_{n\in[N]}c_n\cdot T_{\tilde{u}}^{\floor{a_{\tilde{u}}(n)}} u_{\um,\um',\um'',\um'''}
%    \\&\prod_{(i,j)\in I\backslash\{\tilde{u}\}}T_{i,j}^{\floor{a_{i,j}(n)}}f_{i,j,\um,\um',\um'',\um'''}\prod_{i'\in[r']}{T'_{i'}}^{\floor{a'_{i'}(n)}}f'_{i',\um,\um',\um'',\um'''} \cdot \prod_{j\in[J]} \CD'_{j,\um,\um',\um'',\um'''}(\floor{q_j(n)})\Bigr\Vert_{L^2(\mu)}
% \end{split}
% \end{equation}
%     is positive,
  %  for 1-bounded functions $u_{\um,\um',\um'',\um'''}\in L^\infty(\mu)$ which are product of $2^{t+k-1}$ 1-bounded approximately ${\bu_k}$-invariant functions.
        where
        \begin{align*}
            f_{i,j,\um,\um',\um'',\um'''} &= \Delta_{\substack{(\floor{\bu_1 m_1'} - \floor{\bu_1 m_1}, \floor{\bu_1 m_1'''} - \floor{\bu_1 m_1''}), \ldots,\\ (\floor{\bu_{t} m_{t}'} - \floor{\bu_{t} m_{t}}, \floor{\bu_{t} m_{t}'''} - \floor{\bu_{t} m_{t}''})}}\overline{f_{i,j}},\\
           %  f'_{i',\um,\um',\um'',\um'''} &= \Delta_{\substack{(\floor{\bu_1 m_1'} - \floor{\bu_1 m_1}, \floor{\bu_1 m_1'''} - \floor{\bu_1 m_1''}), \ldots,\\ (\floor{\bu_{t} m_{t}} - \floor{\bu_{t} m_{t}}, \floor{\bu_{t} m_{t}'''} - \floor{\bu_{t} m_{t}''})}}\overline{f'_{i}}\\
            \CD'_{j,\um,\um',\um'',\um'''} &= \Delta_{\substack{(\floor{\bu_1 m_1'} - \floor{\bu_1 m_1}, \floor{\bu_1 m_1'''} - \floor{\bu_1 m_1''}), \ldots,\\ (\floor{\bu_{t} m_{t}'} - \floor{\bu_{t} m_{t}}, \floor{\bu_{t} m_{t}'''} - \floor{\bu_{t} m_{t}''})}}\overline{\CD'_{j}},
        \end{align*}
    and each $h_{\um,\um',\um'',\um'''}$ is a product of $2^{t}$ 1-bounded dual functions of level $s'$ with respect to $T_{\tilde{u}}$.
     Since this average involves only $\ell'-1$ terms of degree greater than or equal to $\deg a_{1,1}$, we can apply (the soft quantitative version of) the induction hypothesis to conclude that for a set of $(\um,\um',\um'',\um'')$ of positive lower density, the seminorm
$$\nnorm{f_{1,1,\um,\um',\um'',\um'''}}_{\bu'_{1},\dots,\bu'_{t'}}=\Bignnorm{\Delta_{\substack{(\floor{\bu_1 m_1'} - \floor{\bu_1 m_1}, \floor{\bu_1 m_1'''} - \floor{\bu_1 m_1''}), \ldots,\\ (\floor{\bu_{t} m_{t}'} - \floor{\bu_{t} m_{t}}, \floor{\bu_{t} m_{t}'''} - \floor{\bu_{t} m_{t}''})}}\overline{f_{1,1}}}_{\bu'_{1},\dots,\bu'_{t'}}$$
           is (uniformly) bounded away from 0
         for some $t'\in\N$ (depending only on $d,J,\ell,\CQ$) and $\bu'_{1},\dots,\bu'_{t'}$ belonging to the set $\CG$.
 %     So we have that the $x$-th term of $\mathcal{F}_{\eta}$ is controlled by transformations in $\CG_{\eta'}$.
   Averaging over all $\um,\um',\um'',\um''$, raising it to the power $2^{t'}$ using the H{\"o}lder inequality and applying the inductive formula for generalized box seminorms, we get the claimed norm control
       $$\nnorm{f_{1,1}}_{\bu_{1},\dots,\bu_{t},\bu'_{1},\dots,\bu'_{t'}}>0$$
     %In conclusion, we have that  (\ref{qtrgg00}) is zero if $\nnorm{f_{1,1}}_{\bv_{1},\dots,\bv_{t},\bu_{1},\dots,\bu_{k-1},\bu'_{1},\dots,\bu'_{t'},\bu''_{1},\dots,\bu''_{t''}}=0$.
by a seminorm with all vectors lying in $\CG$.
%This concludes the proof {\color{red} of what?}. \AK{of 4.10?}

  
\section{Seminorm comparison under ergodicity assumptions I: preliminaries}\label{S: seminorm comparison I}
%\section{Unpacking ergodicity properties: preliminaries}
From now on, our goal is to use the difference ergodicity condition to infer the following result, which is crucial for ensuring the Host-Kra seminorm control of our averages.
\begin{proposition}[Seminorm comparison]\label{P: seminorm comparison}
    Let $\ell\in\N, s\in\N_0$, and let $(X, \CX, \mu, T_1, \ldots, T_\ell)$ be a system. 
    Let $a_1, a_2\in\CH$
    %with $a_1, a_2\succ\log$ \BK{(needed or not?)}, 
    and suppose that there exist nonzero $\beta_1, \beta_2\in\R$ for which $\beta_2 a_1 - \beta_1 a_2 \ll \log$. Suppose moreover that $(T_1^{\floor{a_1(n)}}T_2^{-\floor{a_2(n)}})_n$ is ergodic for $(X, \CX, \mu)$.    Then for any $f\in L^\infty(\mu)$, the following holds:
    \begin{enumerate}
        \item $\nnorm{f}_{\langle\beta_1 \be_1 - \beta_2 \be_2\rangle} = 0$ whenever $\nnorm{f}_{\Z^\ell, \langle\be_1, \be_2\rangle}^+ = 0$;
        \item if $s\geq 1$ and $G_1, \ldots, G_s\subseteq\R^\ell$ are finitely generated, then 
    \end{enumerate}
    \begin{align*}%\label{E: seminorm reduction}
        \nnorm{f}_{G_1, \ldots, G_s, \langle\beta_1 \be_1 - \beta_2 \be_2\rangle} = 0 \quad \textrm{whenever} \quad \nnorm{f}_{G_1, \ldots, G_s, \langle\be_1, \be_2\rangle}^+ = 0.
    \end{align*}
\end{proposition}

Proposition \ref{P: seminorm comparison} needs no lower bound on the growth of $a_1, a_2$. In the degenerate case when they are both slower than log, the ergodicity assumption forces $\mu$ to be the Dirac measure $\delta_x$ at some $x\in X$ and $T_1 x = T_2 x = x$ (see Lemma \ref{L: slowly growing} for a related computation).

The proof of Proposition \ref{P: seminorm comparison} will be broken down into several cases depending on the rationality or irrationality of $\beta_1/\beta_2$ as well as the properties of $a_1, a_2$. Its first step is the following reduction which allows us to simplify the iterates appearing in the difference ergodicity condition.
\begin{proposition}\label{P: reduction of ergodicity}
    Let $(X, \CX, \mu, T_1, T_2)$ be a system. Let $a_1, a_2\in\CH$,
    %with $a_1, a_2\succ\log$ \BK{(needed or not?)}, 
    and suppose that there exist nonzero $\beta_1, \beta_2\in\R$ for which $\beta_2 a_1 - \beta_1 a_2 \ll \log$. Suppose moreover that $(T_1^{\floor{a_1(n)}}T_2^{-\floor{a_2(n)}})_n$ is ergodic for $(X, \CX, \mu)$. 
    Then there exist $u\in\R$ and $a\in \CH$, which is either in $\R[x]$ or stays logarithmically away from $\R[x]$, such that $(T_1^{\floor{{a}(n)}}T_2^{-\floor{\beta{a}(n)+u}})_n$ is also ergodic for $(X, \CX, \mu)$, where $\beta = \beta_2/\beta_1$. 
\end{proposition}

In other words, Proposition \ref{P: reduction of ergodicity} allows us to cast away several (partially overlapping) cases: 
\begin{enumerate}
    \item $1\prec \beta_2 a_1 - \beta_1 a_2 \ll \log$, i.e. $a_2 = (\beta_2/\beta_1)a_1 + g$ for some $1\prec g \ll \log$;
    \item $\beta_2 a_1 - \beta_1 a_2$ is not constant but converges to a constant, i.e. $a_2 = (\beta_2/\beta_1)a_1 + u + g$, where $g$ is nonconstant but goes to 0;
    \item $a_1 = a + g_1$ and $a_2 = \beta a + u + g_2$ for some $a\in\R[x]$, $u\in\R$, and nonconstant $g_i$, each of which either goes to 0 or satisfies $1\prec g_i\ll \log$.
    %any of $a_1, a_2$ is a sum of a real polynomial and a nonconstant function $g\ll\log$.
\end{enumerate}


Proposition \ref{P: reduction of ergodicity} will be proved in Section \ref{SS: preliminary reduction}.

The gist of the proof of Proposition \ref{P: seminorm comparison} consists in obtaining the following result, which one should (morally) think of as the $s=0$ case of Proposition \ref{P: seminorm comparison}.
\begin{proposition}\label{P: seminorm comparison, s = 0}
    Let $(X, \CX, \mu, T_1, T_2)$ be a system.
    Let $u\in\R$ and $a\in\CH$, where $a$ is either in $\R[x]$ or stays logarithmically away from $\R[x]$.
    Suppose that $(T_1^{\floor{a(n)}}T_2^{-\floor{\beta a(n)+u}})_n$ is ergodic for $(X, \CX, \mu)$. Then for any $f\in L^\infty(\mu)$, we have
    \begin{align*}
        \nnorm{f}_{\be_1 - \beta \be_2} = 0 \quad \textrm{whenever} \quad \nnorm{f}_{\langle\be_1, \be_2\rangle}^+ = 0.
    \end{align*}
    %For any $f\in L^\infty(\mu)$, we have $\nnorm{f}_{\be_1 - \beta \be_2} = 0$ whenever $\nnorm{f}_{\langle \be_1, \be_2\rangle}^+=0$.
\end{proposition}

The proof of Proposition \ref{P: seminorm comparison} proceeds rather differently depending on whether $\beta$ is rational or irrational. These two cases are covered separately in Corollary \ref{C: seminorm comparison, beta rational 2} and Proposition \ref{P: beta irrational}, and their proofs occupy Sections \ref{SS: beta rational} and \ref{S: seminorm comparison II} respectively.

\begin{proof}[Proof of Proposition \ref{P: seminorm comparison} assuming Proposition \ref{P: seminorm comparison, s = 0}]
    %In this section, we assemble all the pieces to complete the proof of Proposition \ref{P: seminorm comparison}. 
    % Let $G_1, \ldots, G_s\subseteq\R^\ell$ be finitely generated subgroups and $a_1, a_2\in\CH$ be two sequences such that $\beta_2 a_1 - \beta_1 a_2 \ll \log$ for some nonzero $\beta_1, \beta_2\in\R$. Suppose that $(T_1^{\floor{a_1(n)}}T_2^{-\floor{a_2(n)}})_n$ is ergodic for $(X, \CX, \mu)$. We recall that the objective is to show that for any $f\in L^\infty(\mu)$, we have
    % \begin{align}\label{E: seminorm comparison; implication}
    %     \nnorm{f}_{G_1, \ldots, G_s, \langle\beta_1 \be_1 - \beta_2 \be_2\rangle} = 0 \quad \textrm{whenever} \quad \nnorm{f}_{G_1, \ldots, G_s, \Z^\ell, \langle\be_1, \be_2\rangle}^+ = 0.
    % \end{align}

Let $a_1, a_2\in\CH$ be two sequences such that $\beta_2 a_1 - \beta_1 a_2 \ll \log$ for some nonzero $\beta_1, \beta_2\in\R$, and suppose that $(T_1^{\floor{a_1(n)}}T_2^{-\floor{a_2(n)}})_n$ is ergodic for $(X, \CX, \mu)$. Let $\beta = \beta_2/\beta_1$. By Proposition \ref{P: reduction of ergodicity}, there exist a number $u\in\R$ and a sequence $a\in\CH$, which is either in $\R[x]$ or stays logarithmically away from $\R[x]$, such that $(T_1^{\floor{a(n)}}T_2^{-\floor{\beta a(n)+u}})_n$ is ergodic for $(X, \CX, \mu)$. By \cite[Proposition A.2]{FrKu22a}, we obtain the following soft quantitative strengthening of Proposition \ref{P: seminorm comparison, s = 0}: for every $\veps>0$, we can find $\delta>0$ such that for any $f\in L^\infty(\mu)$,
\begin{align}\label{E: soft quantitative seminorm comparison}
    \nnorm{f}_{\be_1 - \beta \be_2} \leq \veps \quad \textrm{whenever} \quad \nnorm{f}_{\langle \be_1, \be_2\rangle}^+\leq \delta.
\end{align}
%$\nnorm{f}_{\be_1 - \beta \be_2} \leq \veps$ whenever $\nnorm{f}_{\langle \be_1, \be_2\rangle}^+\leq \delta$ for $f\in L^\infty(\mu)$. 
Importantly, the parameter $\delta$ is allowed to depend on the system and $\beta$, but is uniform in the choice of $f\in L^\infty(\mu)$. We will use \eqref{E: soft quantitative seminorm comparison} in conjunction with the basic properties of seminorms in order to derive Proposition \ref{P: seminorm comparison}.

 
Suppose first that $s\geq 1$. Lemma \ref{L: basic properties of seminorms} gives the upper bound 
\begin{align*}
    \nnorm{f}_{G_1, \ldots, G_s, \langle\beta_1 \be_1 - \beta_2 \be_2\rangle}\ll_{\beta} \nnorm{f}_{G_1, \ldots, G_s, \langle \be_1 - \beta \be_2\rangle}.
\end{align*}
To show Proposition \ref{P: seminorm comparison}, it suffices to show that the rightmost seminorm vanishes whenever $\nnorm{f}_{G_1, \ldots, G_s, \langle\be_1, \be_2\rangle}^+ = 0.$

Suppose then that $\nnorm{f}_{G_1, \ldots, G_s, \langle\be_1, \be_2\rangle}^+=0$; setting
\begin{align*}
 G = G_1\times \cdots\times G_{s},\quad \textrm{and} \quad\bm = (\bm_1, \ldots, \bm_{s})\quad \textrm{for} \quad\bm_i\in G_i,   
\end{align*}
%$G_{s+1} = \Z^\ell$, $G = G_1\times \cdots\times G_{s+1}$ and $\bm = (\bm_1, \ldots, \bm_{s+1})$ for $\bm_i\in G_i$, 
we then have
\begin{align*}
    \E_{(\bm, \bm')\in G\times G}\brac{\nnorm{\Delta_{(\bm_1, \bm'_1), \ldots, (\bm_{s}, \bm'_{s})}f}^+_{\langle \be_1, \be_2\rangle}}^{2^{s+1}} =0.
\end{align*}
Let $\veps>0$, and let $\delta>0$ be the parameter satisfying \eqref{E: soft quantitative seminorm comparison}. The set of $(\bm, \bm')\in G\times G$ satisfying
\begin{align*}
    \nnorm{\Delta_{(\bm_1, \bm'_1), \ldots, (\bm_{s}, \bm'_{s})}f}^+_{\langle \be_1, \be_2\rangle}> \delta
\end{align*}
has 0 density with respect to any F\o lner sequence on $G\times G$. 
For all the rest, \eqref{E: soft quantitative seminorm comparison} gives
%the soft quantitative strengthening of Proposition \ref{P: seminorm comparison, s = 0} gives 
\begin{align*}
    \nnorm{\Delta_{(\bm_1, \bm'_1), \ldots, (\bm_{s}, \bm'_{s})}f}_{\be_1 -\beta\be_2}\leq \veps.
\end{align*}
Hence
\begin{align*}
    \nnorm{f}_{G_1, \ldots, G_s, \langle\be_1 - \beta \be_2\rangle}^{2^{s+1}} = \E_{(\bm, \bm')\in G\times G}\nnorm{\Delta_{(\bm_1, \bm'_1), \ldots, (\bm_{s}, \bm'_{s})}f}_{\be_1 -\beta\be_2}^{2^{s}} \leq\veps^{2^{s}},
\end{align*}
and the claim follows on taking $\veps\to 0$.

For the case $s=0$, Lemma \ref{L: basic properties of seminorms} gives
\begin{align*}
    \nnorm{f}_{\langle\beta_1 \be_1 - \beta_2 \be_2\rangle}\ll_{\beta} \nnorm{f}^+_{\langle \be_1 - \beta \be_2\rangle}\leq \nnorm{f}_{\Z^\ell, \langle \be_1 - \beta \be_2\rangle}.
\end{align*}
The previous argument shows that the rightmost seminorm vanishes if $\nnorm{f}_{\Z^\ell, \langle\be_1, \be_2\rangle}^+ = 0$ (by substituting $G$ with $\Z^{\ell}$).
\end{proof}


\subsection{Preliminary reduction}\label{SS: preliminary reduction}

This section is devoted to the proof Proposition \ref{P: reduction of ergodicity}, a preliminary reduction in the proof of Proposition \ref{P: seminorm comparison}. We begin with an auxiliary result on the density of some Hardy orbits on certain subtori.
\begin{lemma}\label{L: density}
    Let $\ell, m\in\N$. Let $g_1, \ldots, g_m\in \CH$ with $1\prec g_1(x)\prec \cdots \prec g_m(x)\prec x$, and define
    \begin{align*}
        a_j(x) &= \sum_{i=1}^m \alpha_{j,i}g_i(x) + \alpha_{j0}+r_j(x),\\ %\quad \textrm{for}\quad j\in[\ell].
        b_j(x_1, \ldots, x_m) &= \sum_{i=1}^m \alpha_{j,i}x_i + \alpha_{j0}
    \end{align*}
    for all $j\in[\ell]$, some $\alpha_{j,i}\in\R$, and some $r_j\in\CH$ with $\lim\limits_{x\to\infty}r_j(x)=0$.
    Then $$(\{a_1(n)\}, \ldots, \{a_\ell(n)\})_n$$ is dense in
    \begin{align*}
        Z = \overline{\{(\{b_1(x_1, \ldots, x_m)\}, \ldots, \{b_\ell(x_1, \ldots, x_m)\}):\; x_1, \ldots, x_m\in\R\}},
    \end{align*}
    which is an affine subtorus of $\T^\ell$.\footnote{I.e. a set $\bu + H\subseteq \T^\ell$ for some compact (not necessarily connected) subgroup $H\subseteq \T^\ell$.}
\end{lemma}
Note that because of possible rational dependencies between the coefficients $\alpha_{j,i}$, it is easier to describe the subtorus less explicitly (as we did above) than try to explicitly list its continuous generators.
\begin{proof} Let $\bu = (\alpha_{10}, \ldots, \alpha_{\ell 0})$ and $\Lambda\subseteq\Z^\ell$ be the subgroup of frequencies defining $Z$ via
\begin{align*}
    Z = \{\bz\in \T^\ell:\; \bk\cdot(\bz-\bu)\in\Z\;\; \textrm{for\; all}\;\; \bk\in \Z^\ell\}
    %\Lambda = \{k\in\Z^\ell:\; k\cdot(z-z_0) \in\Z\;\; \textrm{for\; all}\;\; z\in \T^\ell\}.
\end{align*}
(note that $\bk\cdot\bz := k_1z_1 + \cdots + k_\ell z_\ell$ is well-defined for any $\bk\in\Z^\ell$ and $\bz\in\T^\ell$).
    Take a function $W\in\CH$ with growth $1\prec W(x) \ll x$ satisfying
    \begin{equation*}
        1\prec \log W(x)\prec a_j(x)\prec x
    \end{equation*}
    for all $j\neq \ell$ for which $a_j\succ 1$ (the simplest such example would be $W(x) = g_1(x)$). %$W(t) = \exp(\sqrt{g_1(t)})$ \SD{and $W=g_1$?}).
    Then, using~\cite[Theorem 5.1]{R22}, these growth inequalities yield 
    \begin{equation*}
     \lim\limits_{N\to\infty}   \frac{1}{W(N)}\sum_{n=1}^N \left(W(n+1)-W(n)\right)e\left(k_1a_1(n)+\cdots + k_\ell a_\ell(n)\right)=0
    \end{equation*}
    for all $\bk\in\Z^\ell\setminus{\Lambda}$. In the language of \cite{R22}, this means that $(\{a_1(n)\}, \ldots, \{a_\ell(n)\})_n$ is equidistributed on $Z$ with respect to the $W$-averaging scheme above. 

    A standard variation of Weyl's equidistribution theorem implies that for any Riemann integrable function $f: Z\to\C$, we get 
    $$\lim\limits_{N\to\infty}   \frac{1}{W(N)}\sum_{n=1}^N \left(W(n+1)-W(n)\right)F\left(\{a_1(n)\},\ldots, \{a_\ell(n)\}\right)=\int F(x)\,d\mu_Z(x),$$
    where $\mu_Z$ is the Haar measure on $Z$.
    For any (nonempty) open ball $U \subseteq Z$, the previous asymptotic implies that 
    \begin{equation*}
        \frac{1}{W(N)}\sum_{n=1}^N \left(W(n+1)-W(n)\right)1_U\left(\{a_1(n)\},\ldots, \{a_\ell(n)\}\right)=\mu_Z(U)>0,
    \end{equation*}
    and the result follows.
\end{proof}

All the information that we actually need is captured in the following corollary, which is an easy consequence of Lemma \ref{L: density} and a simple decomposition result for Hardy sequences \cite[Lemma A.3]{R22}. %{\color{red} to get this we need to cite Richter's decomposition theorem} \BK{added}
\begin{corollary}\label{C: density}
    Let $\ell\in\N$, and let $a_1, \ldots, a_\ell\in\CH$ with $a_j(x)\prec x$ for each $j\in[\ell]$. Then $$(\{a_1(n)\}, \ldots, \{a_\ell(n)\})_n$$ is dense in an affine subtorus of $\T^\ell$.
\end{corollary}

We also need one more convergence result, an easy corollary of the convergence results of Boshernitzan, Kolesnik, Quas, and Wierdl \cite{BKQW05}.
\begin{lemma}\label{L: convergence of exp sum}
    Let $\ell\in\N$. Let $a\in\CH$, and suppose that $a$ is either in $\R[x]$ or stays logarithmically away from $\R[x]$. Then:
    \begin{enumerate}
        \item For every $\beta_1, s_1, u_1, \ldots, \beta_\ell, s_\ell, u_\ell\in\R$, the exponential sum
        \begin{align}\label{E: convergence of exponential sum}
        \E_{n\in[N]}e\brac{s_1\sfloor{\beta_1 a(n)+u_1} + \cdots + s_\ell\sfloor{\beta_\ell a(n)+u_\ell}}    
        \end{align}
         converges as $N\to\infty$.
        \item For every system $(X, \CX, \mu, T_1, \ldots, T_\ell)$ and every $f\in L^2(\mu)$, the average
        \begin{align*}
            \E_{n\in[N]}T_1^{\sfloor{\beta_1 a(n)+u_1}} \cdots T_\ell^{\sfloor{\beta_\ell a(n)+u_\ell}}f
        \end{align*}
        converges in $L^2(\mu)$ as $N\to\infty$.
    \end{enumerate}
\end{lemma}
\begin{proof}
    For the first statement, we first decompose \eqref{E: convergence of exponential sum} as
    \begin{align*}
        \E_{n\in[N]}e\brac{s a(n) + u - s_1\srem{\beta_1 a(n)+u_1} - \cdots - s_\ell\srem{\beta_\ell a(n)+u_\ell}}
    \end{align*}
    for $s = s_1 \beta_1 + \cdots + s_\ell\beta_\ell$ and $u = s_1 u_1 + \cdots + s_\ell u_\ell$. Since the functions $f_i(x) = e(s_i\srem{x})$ are Riemann integrable, we can approximate them arbitrarily well by trigonometric polynomials, and so it suffices to prove the convergence of 
    \begin{align*}
        \E_{n\in[N]}e\brac{\tilde{s}a(n) + \tilde{u}},
    \end{align*}
    where $\tilde{s} = s + m_1 \beta_1 + \cdots + m_\ell \beta_\ell$ and $\tilde{u}= u + m_1 u_1 + \cdots + m_\ell u_\ell$ for all choices of $m_1, \ldots, m_\ell\in\Z$. This in turn is a straightforward consequence of \cite[Theorem A]{BKQW05}. The second part of the lemma follows from the first part and the spectral theorem.
\end{proof}

Lastly, we require the following averaging lemma, which shows that Ces\`aro convergence in Hilbert space implies convergence along dyadic intervals. 
\begin{lemma}\label{L: dyadic intervals}
Let $(a_N)_N$ be a sequence whose sequence of ratios $(a_{N+1}/a_N)_N$ converges to some number in $(1, \infty]$, and let $(v_n)_n$ be a bounded sequence in a Hilbert space. Suppose that $\E\limits_{n\in[N]}v_n$ converges strongly (resp. weakly) to some limit $L$. Then $\E\limits_{n\in(a_N, a_{N+1}]} v_n$ converges strongly (resp. weakly) to $L$ as well.
\end{lemma}
\begin{proof}
We assume that all the limits are strong; the case when they are all weak follows identically.
    Let $S_N=v_1+\dots+v_N$. {For all sufficiently large $N$, the ratio condition implies that $a_{N+1}>a_N$,} in which case the average in question equals
    \begin{align*}
        \E\limits_{n\in(a_N, a_{N+1}]} v_n =\frac{1}{a_{N+1}-a_N} \big(S_{a_{N+1}} -S_{a_N}\big) = \frac{a_{N+1}}{a_{N+1}-a_N}\brac{\E\limits_{n\in[a_{N+1}]} v_n - \frac{a_N}{a_{N+1}}\E\limits_{n\in[a_{N}]} v_n}.
    \end{align*}
    By assumption, we have $\lim\limits_{N\to\infty}\frac{a_{N+1}}{a_N}>1$. If the limit is infinite, then $\frac{a_{N+1}}{a_{N+1}-a_N}\to 1$, and so
\begin{align*}
    \lim_{N\to\infty} \E\limits_{n\in(a_N, a_{N+1}]} v_n = \lim_{N\to\infty} \E\limits_{n\in[a_{N+1}]} v_n = L.
\end{align*}
If $\lim\limits_{N\to\infty}\frac{a_{N+1}}{a_N} = 1+c$ for some $c>0$, then $\frac{a_{N+1}}{a_{N+1}-a_N}\to \frac{1+c}{c}$, in which case
\begin{align*}
    \lim_{N\to\infty}\E\limits_{n\in(a_N, a_{N+1}]} v_n = \frac{1+c}{c}\Bigbrac{1-\frac{1}{1+c}}L = L.
\end{align*}
\end{proof}


We now proceed to prove Proposition \ref{P: reduction of ergodicity}. 
\begin{proof}[Proof of Proposition \ref{P: reduction of ergodicity}]
If  $\beta_2 a_1 - \beta_1 a_2 \ll \log$ and $\beta_1,\beta_2\neq 0$, then $(\beta_2/\beta_1) a_1 - a_2 \ll \log$ as well. To facilitate the notation, we reparametrize $a_1 = \tilde a$, $\beta = \beta_2/\beta_1$, and $a_2 = \beta \tilde a + \tilde g$ for some $\tilde g\ll\log$. Then our assumption reduces to $(T_1^{\floor{\tilde a(n)}}T_2^{-\floor{\beta \tilde a(n) + \tilde g(n)}})_n$ being ergodic.

We now split $\tilde a ={a} + {g}$ as follows:
    \begin{enumerate}
        \item if $\tilde a$ stays logarithmically away from $\R[x]$, then $a = \tilde{a} $ and ${g} = 0$;
        \item if $\tilde a = p + {g}$ for some $p\in\R[x]$ and ${g} \ll \log$, then ${a} = p$ (the representation is unique up to shifts by constants). %Without the loss of generality, we can assume that if $\tilde{g}$ converges, then it converges to 0, for otherwise we can swap the limit into $p$.
    \end{enumerate}
    Then we can recast
    \begin{align*}
        T_1^{\floor{\tilde a(n)}}T_2^{-\floor{\beta \tilde a(n) + \tilde g(n)}} = T_1^{\floor{{a}(n)+g_1(n)}}T_2^{-\floor{\beta {a}(n) + g_2(n)}}
    \end{align*}
    for $g_1 = {g}$ and $g_2 = \beta {g} +\tilde g$.
    %some $g_1, g_2\ll\log $, where we can assume that either $g_1\succ 0$ or $g_1\to 0$ (since $g_1 = \tilde{g}$ above). 
    The statement of the proposition then amounts to showing that the ergodicity of $$(T_1^{\floor{{a}(n)+g_1(n)}}T_2^{-\floor{\beta {a}(n) + g_2(n)}})_n$$ implies the ergodicity of $(T_1^{\floor{{a}(n)}}T_2^{-\floor{\beta {a}(n) + u}})_n$ for some $u\in\R$.

    We now take a step back and discuss the distribution of $(\{g_1(n)\}, \{g_2(n)\})_n$ in $\T^2$, which we will need shortly. By Corollary~\ref{C: density}, $(\{g_1(n)\}, \{g_2(n)\})_n$ is dense in an {affine subtorus} $Z = (u_1,u_2)+H$ of $\T^2$. By absorbing $u_1$ in the definition of ${a}$, we can replace the shift $(u_1, u_2)$ by $(0,u)$ for some $u\in\R$. Then $(\{g_1(n)\}, \{g_2(n)-u\})_n$ is dense in $H$. Now, one of the two things happens: either $\{g_1(n)\} = \{g_2(n)-u\} = 0$ for all sufficiently large $n\in\N$, or $Z$ is 1- or 2-dimensional. In both cases, Corollary~\ref{C: density} implies that for all $\delta>0$, we can find an increasing sequence $(N_k)_k$ such that $(\{g_1(N_k)\}, \{g_2(N_k)-u\})_k$ lies in $[0, \delta)^2$ for all $k\in\N$.
    
    % We split into four cases:
    % \begin{enumerate}
    %     \item both $g_1, g_2$ converge: then the result is immediate after shifting $\tilde{a}$ by the limit of $g_1$ (which we can do since the definitions of $\tilde{a}$ and $\tilde{g}$ are only unique up to scalar shifts);
    %     \item exactly one of $g_1, g_2$ converges: then by symmetry we can assume that $g_1$ is the convergent one, and by shifting $\tilde{a}$ by the limit of $g_1$ we can once again assume that $g_1$ in fact converges to 0;
    %     \item $g_1$ and $g_2$ are affinely dependent and go to $\infty$;
    %     \item $g_1$ and $g_2$ are affinely independent and go to $\infty$.
    % \end{enumerate}
    %, where $u$ is the limit of $g_2$ if $g_2$ converges, and it is 0 otherwise. The point is that we can remove $g_1, g_2$ if they are slowly growing function going to infinity.
    %The point of this decomposition is to ensure later on that certain exponential sums involving $\tilde{a}$ converge.
    
Let us see how to complete the proof of the proposition using this newly acquired information about the distribution of $(\{g_1(n)\}, \{g_2(n)\})_n$ in $\T^2$. We fix a small parameter $\delta>0$.
By the ergodicity assumption, we have
\begin{align*}
    \lim_{N\to\infty}\E_{n\in[N]} T_1^{\floor{{a}(n)+g_1(n)}}T_2^{-\floor{\beta {a}(n) + g_2(n)}}f = \int f\; d\mu,
\end{align*}
and so 
\begin{align*}
    \lim_{N\to\infty}\E_{n\in [N, (1+\delta)N)} T_1^{\floor{{a}(n)+g_1(n)}}T_2^{-\floor{\beta {a}(n) + g_2(n)}}f = \int f\; d\mu
\end{align*}
as well by Lemma \ref{L: dyadic intervals}. Now, while integer parts are not additive, we can split
\begin{align*}
    \floor{{a}(n)+g_1(n)} &= \floor{{a}(n)} + \floor{g_1(n)}\\
    \floor{\beta {a}(n) + g_2(n)} &= \floor{\beta {a}(n)+u} + \floor{g_2(n)-u}
\end{align*}
unless
\begin{align}\label{E: boundary condition}
    \rem{{a}(n)} > 1 - \rem{g_1(n)} \quad \textrm{or}\quad \rem{\beta {a}(n)+u} > 1 - \rem{g_2(n)-u}.
\end{align}
Our goal is to show that \eqref{E: boundary condition} happens very rarely. 

Because $g_j\ll\log$, we have $g_j(n) = g_j(N) + O(\delta)$ for $n\in[N, (1+\delta)N)$, so if the boundary condition \eqref{E: boundary condition} is satisfied, then we must have
 \begin{align*}
     \rem{{a}(n)} > 1 - \rem{g_1(N)}-O(\delta) \quad \textrm{or}\quad \rem{\beta{a}(n)+u} > 1 - \rem{g_2(N)-u}-O(\delta).
 \end{align*}
 By the density considerations above, we can find an increasing sequence $(N_k)_k$ such that
\begin{align}\label{E: a dense subsequence}
    \rem{g_1(N_k)}, \rem{g_2(N_k)-u}\leq \delta,
\end{align}
so that \eqref{E: boundary condition} holds for $n\in[N_k, (1+\delta)N_k)$ only if
 \begin{align}\label{E: boundary condition 1}
     \min(\rem{{a}(n)}, \srem{\beta {a}(n)+u}) > 1 -O(\delta). % \quad \textrm{or}\quad \rem{\beta \tilde{a}(n)} > 1 -O(\delta+\veps).
 \end{align}
Now, since ${a}$ is either in $\R[x]$ or stays logarithmically away from $\R[x]$, we can deduce that at most a $O(\delta)$-proportion of $n\in[N_k, (1+\delta)N_k)$ can satisfy \eqref{E: boundary condition 1} for all sufficiently small $\delta>0$.\footnote{
%If ${a}\in\Q[t]+\R$, then $\rem{{a}(n)}$ is bounded away from 1, so \eqref{E: boundary condition 1} fails for all $n\in\Z$ if $\delta>0$ is sufficiently small. If $a\in\R[t]$ has at least one irrational nonconstant coefficient, or if $a$ stays logarithmically away from $\R[t]$, then $(\rem{{a}(n)})_n$ is equidistributed, and so this claim follows from equidistribution and Lemma~\ref{L: dyadic intervals}. \BK{(Is this a correct explanation?)} \AK{How about giving more details? 
If $a$ is in $\R[x],$ then each of the polynomials $a, \beta a+u$ is either in $\Q[x]+\R$ or in its complement. Hence the claim follows by periodicity (for sufficiently $\delta>0$) or by well distribution respectively. If $a$ stays logarithmically away from $\R[x],$ so does $\beta a +u,$ and the claim follows by equidistribution and Lemma~\ref{L: dyadic intervals}.}
 % In the second case, where, we recall, $g_1(t)\to 0$ and $1 \prec g_2(t) \ll \log t$, this follows from the density of $(\{g(N)\})_N$ in $[0,1)$ (this was proved in \cite{Boshernitzan-equidistribution} but can also be deduced from Lemma \ref{L: density}). In the third case, where $1 \prec g_1(t), g_2(t) \ll \log t$, the argument differs slightly depending on whether $g_1, g_2$ are affinely dependent or not. If $g_2=\gamma g_1$ for some $\gamma\in\R$, then again such a sequence exists by density of $g$. \BK{(Does it matter whether $\gamma\in\Q$ or not?) STOPPED SOMEWHERE HERE}
 % If $g_2=\gamma g_1 + \xi$ for some $\gamma, \xi\in\R$, then we can incorporate $\xi$ in the definition of $\tilde{a}$ and we land in the previous case. Lastly, if $\tilde{g}, g$ are not affinely independent, then this follows from Lemma \ref{L: density}.
%It therefore remains to consider the case when $\tilde{g}, g$ are not affinely independent. Suppose without loss of generality that $1\prec \tilde{g} \prec g$, with the other case following similarly. We will set $N_k = b(M_k)$ for some sequence $b\in\CH$ such that $\tilde{g}\circ b$ and $g\circ b$ are both strongly nonpolynomial sequences growing faster than log. Then $(\{\alpha\tilde{g}\circ b(M)\}, \{\beta\tilde{g}\circ b(M) + g\circ b(M)\})_M$ is equidistributed on $\T^2$, and so we can once again choose a sequence $M_k\to\infty$ such that \eqref{E: a dense subsequence} holds. \BK{Hmm, I think this is still difficult e.g. when $\tilde{g}(N) = \log\log N$ and $g(N) = \log N$.}
Restricting to the sequence $(N_k)_k$, we then have
\begin{multline}\label{E: slow 1}
        \E_{n\in [N_k, (1+\delta)N_k)} T_1^{\floor{{a}(n)+g_1(n)}}T_2^{-\floor{\beta {a}(n) + g_2(n)}}f \\
        = T_1^{\floor{g_1(N_k)}}T_2^{-\floor{g_2(N_k)-u}}\E_{n\in [N_k, (1+\delta)N_k)} T_1^{\floor{{a}(n)}}T_2^{-\floor{\beta {a}(n)+u}}f +O(\delta).
%&= T_2^{\floor{\log_2 N_k}}\E_{n\in [N_k, (1+\delta)N_k)} T_1^{-\floor{\alpha n}}T_2^{\floor{\beta n}}f +O(\delta + \veps)
\end{multline}
%(if $\beta\in\Q$, we can choose $\veps, \delta$ sufficiently small for the equality to be exact; if $\beta\notin\Q$, this follows from the equidistribution of $\beta n$). 

For reasons that will become clear later on, we want the average on $[N_k, (1+\delta)N_k)$ to converge along the usual interval $[N]$. Fortunately this is the case, and the $L^2(\mu)$ limit
\begin{align}\label{E: slow 2}
    L:= \lim_{N\to\infty}\E\limits_{n\in[N]} T_1^{\floor{{a}(n)}}T_2^{-\floor{\beta {a}(n)+u}}f
\end{align}
exists by Lemma \ref{L: convergence of exp sum}.
% the reason is that because $\tilde{a}$ is either in $\R[t]$ or stays logarithmically away from $\R[t]$, the exponential sums $\lim\limits_{N\to\infty}\E\limits_{n\in[N]} e(t\floor{ \tilde{a}(n)} + s{\floor{\beta \tilde{a}(n)+u}})$ converge, \KT{(do we have a justification of this later on?)} \BK{(cite some relevant result)} and so \eqref{E: slow 2} converges as well by the spectral theorem. 
Therefore 
\begin{align*}
    \E_{n\in [N_k, (1+\delta)N_k)} T_1^{\floor{{a}(n)}}T_2^{-\floor{\beta {a}(n)+u}}f \to L
\end{align*}
in $L^2(\mu)$ as $k\to\infty$. If we can show that $L$ is constant, then $L = \int L\; d\mu = \int f\; d\mu$, and Proposition~\ref{P: reduction of ergodicity} will follow.



Comparing both sides of \eqref{E: slow 1} with the fact that the left-hand side converges to $\int f\; d\mu = \int L\; d\mu$, we get that
\begin{align*}
    \limsup_{k\to\infty}\norm{T_1^{\floor{g_1(N_k)}}T_2^{-\floor{g_2(N_k)-u}}\E_{n\in [N_k, (1+\delta)N_k)} T_1^{\floor{{a}(n)}}T_2^{-\floor{\beta {a}(n)+u}}f - \int L\; d\mu}_{L^2(\mu)}\ll \delta. 
\end{align*}
Using the invariance of measures, we can compose the integral with $T_1^{-\floor{g_1(N_k)+}}T_2^{\floor{g_2(N_k)-u}}$ and then use \eqref{E: slow 2} to get
\begin{align*}
    \norm{L - \int L\; d\mu}_{L^2(\mu)}\ll \delta.
\end{align*}
Since $\delta>0$ was arbitrary, we get that $L = \int L\; d\mu$, implying the result.
\end{proof}

% \AK{Older comments of Borys and Wenbo:}
% \BK{For Wenbo's convergence to work, we need to assume that $(T_1^{\floor{a_1(n)}}T_2^{-\floor{a_2(n)}})_n$ is weakly ergodic for $(X, \CX, \mu)$. Then we want to conclude that $(T_1^{\floor{{a}(n)}}T_2^{-\floor{\beta{a}(n)+u}})_n$ is strongly ergodic for $(X, \CX, \mu)$ (if the latter is weakly ergodic, then it is strongly ergodic by Lemma 6.5). }
% {\color{red}

% If we only assume that weak ergodicity, then comparing both sides of \eqref{E: slow 1} gives us 
% $$ \limsup_{k\to\infty}\Bigl\vert \langle \E_{n\in [N_k, (1+\delta)N_k)} T_1^{\floor{{a}(n)}}T_2^{-\floor{\beta {a}(n)+u}}f, T_1^{-\floor{g_1(N_k)}}T_2^{\floor{g_2(N_k)-u}}g\rangle-(\int f)(\int g)\Bigr\vert\ll \delta$$
% for all $f,g\in L^{2}(\mu)$. Since $\E_{n\in [N_k, (1+\delta)N_k)} T_1^{\floor{{a}(n)}}T_2^{-\floor{\beta {a}(n)+u}}f$ converges to $L$ strongly, by Cauchy-Schwartz, we have
% $$ \limsup_{k\to\infty}\Bigl\vert \langle L, T_1^{-\floor{g_1(N_k)}}T_2^{\floor{g_2(N_k)-u}}g\rangle-(\int f)(\int g)\Bigr\vert\ll \delta,$$
% or 
% $$ \limsup_{k\to\infty}\Bigl\vert \langle T_1^{\floor{g_1(N_k)}}T_2^{-\floor{g_2(N_k)-u}} L, g\rangle-(\int f)(\int g)\Bigr\vert\ll \delta.$$
% So $T_1^{\floor{g_1(N_k)}}T_2^{-\floor{g_2(N_k)-u}} L$ converges weakly to $\int f$.
% %Question: let $L\in L^{2}(\mu)$ and $S_{k}$ be measure preserving transformations. If $S_{k}L$ converges to a constant weakly, does it imply that $L$ is a constant?
% }

\subsection{The case of rational $\beta$}\label{SS: beta rational}

% \AK{Are we working under the assumption that $(T_{1}^{[ka(n)]}T_{2}^{-\floor{la(n)+u}})_{n}$ is ergodic, or we are dealing with $(T_{1}^{\floor{ka(n)}}T_{2}^{-\floor{la(n)+u}})_{n}$ working under the assumption that $(T_{1}^{\floor{ka(n)}}T_{2}^{-[la(n)]})_{n}$ is ergodic?} \BK{The former}

The goal now is to prove Proposition \ref{P: seminorm comparison, s = 0} under the assumption that $\beta\in\Q$. This will follow from the more general result presented below.

\begin{proposition}\label{P: beta rational}
    Let $k,l\in\mathbb{Z}\backslash\{0\}$ and $u\in\R$. Then there exists a finite index subgroup $H\subseteq\mathbb{Z}^{2}$, which contains $(k,l)$ and depends only on $k$ and $l$, such that for any system $(X,\CX,\mu, T_{1},T_{2})$ and any sequence $a\colon\mathbb{N}\to\mathbb{R}$, the ergodicity of $(T_{1}^{\floor{ka(n)}}T_{2}^{-\floor{la(n)+u}})_{n}$ implies that $\CI(H)=\CI(T_{1}^{k}T_{2}^{-l})$.   
\end{proposition}
We emphasize the level of generality of Proposition \ref{P: beta rational}: it works for \textit{any} sequence $a:\N\to\R$, not necessarily a Hardy one. %In degenerate cases, its statement becomes trivial, e.g. if $a$ is constant, then the ergodicity assumption forces $\mu$ to be the Dirac measure $\delta_x$ at some $x\in X$ and $T_1 x = T_2 x = x$.

As an immediate consequence of Proposition \ref{P: beta rational} and Lemma \ref{L: basic properties of seminorms} (v), we get the following corollary.
\begin{corollary}\label{C: seminorm comparison, beta rational}
    Under the assumptions of Proposition \ref{P: beta rational}, we have $\nnorm{f}_{k\be_1 - l\be_2} = \nnorm{f}_H$ for any $f\in L^\infty(\mu)$.
\end{corollary}
Corollary \ref{C: seminorm comparison, beta rational} combined with Proposition \ref{P: reduction of ergodicity} quickly yield Proposition \ref{P: seminorm comparison, s = 0} whenever $\beta\in\Q$, which is restated in the corollary below.
\begin{corollary}\label{C: seminorm comparison, beta rational 2}
    Proposition \ref{P: seminorm comparison, s = 0} holds for nonzero $\beta\in\Q$. 
    
    More specifically, let $(X, \CX, \mu, T_1, T_2)$ be a system, let $u\in\R$ and $a\in\CH$, where $a$ is either in $\R[x]$ or stays logarithmically away from $\R[x]$, and suppose that $(T_1^{\floor{a(n)}}T_2^{-\floor{\beta a(n)+u}})_n$ is ergodic for $(X, \CX, \mu)$. Then for any $f\in L^\infty(\mu)$, we have
    \begin{align*}
        \nnorm{f}_{\be_1 - \beta \be_2} = 0 \quad \textrm{whenever} \quad \nnorm{f}_{\langle\be_1, \be_2\rangle}^+ = 0.
    \end{align*}
\end{corollary}
\begin{proof}[Proof of Corollary \ref{C: seminorm comparison, beta rational 2} assuming Corollary \ref{C: seminorm comparison, beta rational}]
Let $\beta = l/k\in\Q$ for nonzero $k,l\in\Z$. Then Proposition \ref{P: beta rational} applied to the sequence $(a(n)/k)_n$ gives a subgroup $H\subseteq \langle \be_1, \be_2\rangle$ of index $O_{\beta}(1)$ for which $\CI(T_1^k T_2^{-l}) = \CI(H)$. Then
\begin{align*}
    \nnorm{f}_{\langle \be_1, \be_2\rangle}^+ \asymp_{\beta} \nnorm{f}_{H}^+ = \nnorm{f}_{k\be_1 - l\be_2}^+ \asymp_{\beta}\nnorm{f}_{\be_1 - \beta\be_2}^+\geq \nnorm{f}_{\be_1 - \beta\be_2}
\end{align*}
by (various parts of) Lemma \ref{L: basic properties of seminorms}, from which the claimed result follows.
\end{proof}

%however, we will give the full proof of Proposition \ref{P: seminorm comparison, s = 0} in the next section, once we handle the case $\beta\notin\Q$.
We finally proceed to prove Proposition \ref{P: beta rational}.
\begin{proof}[Proof of Proposition \ref{P: beta rational}]
Rescaling $k,l$ and $a$ if necessary, we may assume without loss of generality that $k$ is coprime to $l$. Let $I_1, \ldots, I_K$ be the finite partition of $[0,1)$ into all the intervals with endpoints from the set
    \begin{align}\label{E: endpoints}
        \Bigl\{\frac{i}{\vert k\vert},\frac{j}{\vert l\vert}, \frac{j-\{u\}}{\vert l\vert}\colon 0\leq i\leq \vert k\vert,\; 0\leq j\leq \vert l\vert\Bigr\}.
    \end{align}
% \BK{(This set is slightly different from the one that Wenbo originally wrote because of the presence of $u$. Please check if this makes sense.)} {\color{red} should be
% \begin{align}%\label{E: endpoints}
%         \Bigl\{\frac{i}{\vert k\vert},\frac{j}{\vert l\vert}, \Bigl\{\frac{j}{\vert l\vert}-u\Bigr\}\colon 0\leq i\leq \vert k\vert,\; 0\leq j\leq \vert l\vert\Bigr\}.
%     \end{align}
% } \BK{Are you sure? My computations give
% \begin{align*}
%     \floor{la(n) + u} = \floor{la(n)}+\floor{u} + 1_{\{la(n)\}+\{u\}\geq 1}.
% \end{align*}
% Now, if $\{a(n)\}\in[\frac{j}{l},\frac{j+1}{l})$, then $\{la(n)\} = l\{a(n)\}-j$ and $\floor{l a(n)}=l\floor{a(n)}+j$, and so we have
% \begin{align*}
%     \floor{la(n) + u} = l\floor{a(n)}+\floor{u} + j + 1_{l\{a(n)\}-j+\{u\}\geq 1}.
% \end{align*}
% Then the condition can be rewritten as $\{a(n)\}\geq \frac{j+1-\{u\}}{l}$.
% }
% {\color{red} I see. But then it should be $\{\frac{j-\{u\}}{\vert l\vert}\}$ since it belongs to [0,1)}
By the Bolzano-Weierstrass theorem, there exists a subsequence $(N_j)_j$ of $\mathbb{N}$ such that
the limit 
\begin{align}\label{E: weights}
    c_i =\lim_{j\to\infty}\frac{\vert\{n\in[N_j]\colon\; \{a(n)\}\in I_i\}\vert}{N_{j}}
\end{align}
    exists for all $i\in[K]$.


    % $$\lim_{j\to\infty}\frac{\vert\{n\in[N_j]\colon\; \{a(n)\}\in [x,y)\}\vert}{N_{j}}$$
    % exists for all $x,y$ in
    % \begin{align}\label{E: endpoints}
    %     \Bigl\{\frac{i}{\vert k\vert},\frac{j}{\vert l\vert}\colon 0\leq i\leq \vert k\vert, 0\leq j\leq \vert l\vert\Bigr\}.
    % \end{align}

%Denote $m:=(k,-l)$. 
From now on, assume that $f$ has zero integral and is invariant under the transformation $R:=T_1^kT_2^{-l}$. Our objective is to show that $f$ is $\CI(H)$-measurable for some fixed finite index subgroup $H\subseteq \Z^2$ {which contains $(k,l)$}.

Since $(T_{1}^{\floor{ka(n)}}T_{2}^{-\floor{la(n)+u}})_{n}$ is ergodic, we have
\begin{equation}\label{kl1}
    \lim_{j\to\infty}\E_{n\in[N_j]}T_{1}^{\floor{ka(n)}}T_{2}^{-\floor{la(n)+u}}f=0
\end{equation}
in $L^2(\mu)$.
%for any zero-integral function $f$.
%For convenience write $T_{(x,y)}:=T_{1}^{x}T_{2}^{y}$ for all $x,y\in\mathbb{Z}$.
Using the identities
\begin{gather*}
 \floor{n x} = n\floor{x} + i \quad\textrm{whenever}\quad n\in\N,\; x\in\left[\frac{i}{n}, \frac{i+1}{n}\right)\\
 \textrm{and}\quad  \floor{x + y}=\floor{x}+\floor{y}+1_{\rem{x}+\rem{y}\geq 1},
\end{gather*}
% $\floor{n x} = n\floor{x} + i$ whenever $n\in\N$ and $x\in\left[\frac{i}{n}, \frac{i+1}{n}\right)$ as well as $\floor{x + y}=\floor{x}+\floor{y}+1_{\rem{x}+\rem{y}\geq 1}$, 
we can find a function $r:[0,1)\to\Z^2$, constant on each $I_1, \ldots, I_K$ and satisfying the bound
\begin{align}\label{E: bound on r}
\norm{r}_\infty \leq  \max(|k|-1, |l|),    
\end{align}
 such that 
\begin{align*}
    T_{1}^{\floor{ka(n)}}T_{2}^{-\floor{la(n)+u}}=T_{1}^{k\floor{a(n)} + r_1(n)}T_{2}^{-l\floor{a(n)}-\floor{u}-r_2(n)} = R^{\floor{a(n)}}T_1^{r_1(n)}T_2^{-\floor{u}-r_2(n)},
\end{align*}
where $r(n)=(r_1(n), r_2(n))$.
Composing (\ref{kl1}) with $T_2^{\floor{u}}$ and using the $R$-invariance of $f$, we obtain that
\begin{equation}\label{kl2}
    0=\lim_{j\to\infty}\E_{n\in[N_{j}]}T_1^{r_1(n)}T_2^{-r_2(n)}f=\sum_{i=1}^{K}c_{i}\cdot T_1^{r_{1,i}}T_2^{-r_{2,i}} f,
\end{equation}
%for all zero-integral $R$-invariant functions $f\in L^\infty(\mu)$, 
where %$R:=T_1^kT_2^{-l}$, 
the pair $(r_{1,i}, r_{2,i})$ is the value of $r$ on $I_i$ and $c_i$ is defined in \eqref{E: weights}.

Since $k$ is coprime to $l$, there exist $w,v\in\mathbb{Z}$ with $\max(\vert w\vert, |v|)=O_{k,l}(1)$ such that $wl-vk=1$. For all $i\in[K]$, we can then write $(r_{1,i}, r_{2,i}) = d_{i}(w,v)+d'_{i}(k,-l)$ for a unique choice of $d_i, d'_i\in\Z$, which are of size $O_{k,l}(1)$ by \eqref{E: bound on r}. Denote $S:=T_1^w T_2^v$. Using the $R$-invariance of $f$ again, we can recast
$T_1^{r_{1,i}}T_2^{-r_{2,i}} f = S^{d_i} f$ for all $i\in[K]$. % and all $R$-invariant functions $f$.
So it follows from (\ref{kl2}) that 
\begin{equation}\label{kl3}
    0=\sum_{i=1}^{K}c_{i}S^{d_{i}}f. %\text{ for all zero-integral $T_{m}$-invariant function $f$.}
\end{equation}
Denote 
$$L(f):=\sum_{i=1}^{K}c_{i}S^{d_{i}+d_{\ast}}f$$
for all $f$, where $d_{\ast}=\max\{-d_{1},\dots,-d_{K},0\}$; clearly, $L(f) = 0$. 
% Then (\ref{kl3}) implies that 
% \begin{equation}\label{kl4}
%     L(f)=0 \text{ for all zero-integral $T_{m}$-invariant function $f$.}
% \end{equation}
By the definition of $d_*$ and the bounds on $d_i$, we have $0\leq d_{i}+d_{\ast} \ll_{k,l}1$, and hence
$$g(z):=\sum_{i=1}^{K}c_{i}z^{d_{i}+d_{\ast}}$$ is a real polynomial of degree $O_{k,l}(1)$. Since $c_{i}\geq 0$ and at least one $c_{i}$ is nonzero, the polynomial $g$ is nonzero. 

Thus, the study of the limit $L(f)$ can be reduced to the examination of the polynomial $g$. The point of the next few claims is that we can throw away all the roots of $g$ that are not roots of unity, and that we can combine them all together to find a finite-index subgroup $H\subseteq\Z^2$ containing $(k,-l)$ such that $f$ is $H$-invariant.
%Our goal now is to establish more properties of the polynomial $g$ that will be crucial in completing the proof of the proposition.
    
\begin{claim}\label{Claim 1}
% Let $f$ be a nonzero $T_{m}$-invariant function with
Suppose that $f\neq 0$ and $Sf=\lambda f$ for some $\lambda\in\mathbb{C}$.
Then $\lambda$ must be a root~of~$g$ and a root of unity. 
\end{claim}
%{\color{red} in the claims, it is better to say that it works for all zero-integral $R$-invariant function, not just for a fixed $f$. (or another way is to say clearly that \eqref{kl3} holds for all zero-integral $R$-invariant function)}

Indeed, assume that $\lambda$ is not a root of unity. Then for all $n\in \N$, $$\int f^n\; d\mu = \int S(f^n)\; d\mu= \int (Sf)^n\; d\mu=\lambda^n\int f^n\; d\mu,$$ so $f^n$ has zero integral.  
Since $f^n$ is $R$-invariant by assumption, we get $L(f^n)=0$ by \eqref{kl3}. Because $f^n\neq 0$, we further deduce that $\lambda\neq 0$ and $g(\lambda^n)=0$. By the fundamental theorem of algebra, the polynomial $g$ has only finitely many roots, and so the only way in which $g(\lambda^n)=0$ can happen for all $n\in\N$ is if $\lambda$ is a root of unity.

\

As an immediate corollary, we get the following.

\begin{claim}\label{Claim 2}
   Suppose that $\lambda\in\mathbb{C}$ is not a root of unity. If $(S-\lambda)f=0$, then $f=0$.
\end{claim}






 %We are now ready to complete the proof the proposition. Let $f$ be any zero-integral $T_{m}$-invariant function.
 Next, we move on to show that we can remove all the roots of $g$ that are not roots of unity. Since $g\neq 0$, we may factorize $g$ as 
 $$g(z)=c\prod_{j=1}^{M}(z-z_{j})$$
 for some $c\neq 0, M=O_{k,l}(1)$ and $z_{1},\dots,z_{M}\in\mathbb{C}$. Then it follows from (\ref{kl3}) that
 $$L(f)=c\prod_{j=1}^{M}(S-z_{j})f=0.$$

If $M=0$, then since $c\neq 0$, we must have $f=0$, in which case we are done. So we assume that $M>0$. Rearranging the roots, we can assume that  $z_{1},\dots,z_{M'}$ are roots of unity while $z_{M'+1},\dots,z_{M}$ are not. By repeatedly applying Claim \ref{Claim 2} to $\prod_{j=1}^{M'+t}(S-z_{j})f$ in place of $f$, we deduce that
$$\prod_{j=1}^{M'+t}(S-z_{j})f=0$$
for all $0\leq t\leq M-M'-1$. Hence
 $$\prod_{j=1}^{M'}(S-z_{j})f=0.$$
 %Since $Sf=0$ implies that $f=0$, we may further assume that all of $z_{1},\dots,z_{M'}$ are  roots of unity. 
 
 Assume that $z_j$ is a primitive $q_j$-th root of unity for all $1\leq j\leq M'$.
Then  
\begin{equation}\label{kl5}
    \prod_{j=1}^{M'}(S^{q_{j}}-1)f=0
\end{equation}
since $z-z_{j}$ divides $z^{q_{j}}-1$.

The following claim allows us to bound the size of $q_j$'s, and hence the index of the subgroup $H$ later on.
\begin{claim}\label{Claim 3}
    We have $q_{j}=O_{k,l}(1)$.
\end{claim}
 
% \BK{(This is essentially what Wenbo wrote modulo the proof of $\lim\limits_{t\to\infty}\phi(t)=\infty$, which I commented out. However, I am uncertain about what's going on in the argument below. Why does $\phi$ matter? Are you trying to show that the polynomial $\prod_{j=1}^{M'}(z^{q_{j}}-1)$ divides $g$, and hence it must have a degree at most $O_{k,l}(1)$? I don't see why this would be true given that a priori, $g$ may have arbitrary nonnegative real coefficients.)}{\color{red} Since  $z_j$ is a primitive $q_j$-th root of unity, its minimal polynomial $h$ is of degree $\phi(q_{j})$. Since $z_{j}$ is also a root of $g$, we must have that $h\vert g$ and so $\phi(q_{j})=\deg(h)\leq \deg(g)=O_{k,l}(1)$ unless $g\equiv 0$. Since $\lim\limits_{t\to\infty}\phi(t)=\infty$, $\phi(q_{j})=O_{k,l}(1)$ also implies that $q_{j}=O_{k,l}(1)$.} 
 To see this, recall that since  $z_j$ is a primitive $q_j$-th root of unity, its minimal polynomial $h$ is of degree $\phi(q_{j})$, where $\phi$ is the Euler totient function. Since $z_{j}$ is also a root of $g$, the polynomial $h$ must divide $g$, and so $\phi(q_{j})=\deg(h)\leq \deg(g)=O_{k,l}(1)$. Since $\lim\limits_{t\to\infty}\phi(t)=\infty$, we have $\phi(q_{j})=O_{k,l}(1)$ only if $q_{j}=O_{k,l}(1)$.
 
 
 %let $\phi(t)$ be the Euler totient function, i.e. the number of $j\in[t]$ coprime to $t$. It is known that the minimal polynomial for the $t$-th root of unity is of degree $\phi(t)$, and that $\lim\limits_{t\to\infty}\phi(t)=\infty$. The claim follows since  $z_{j}$ is both a root of unity and a root of a polynomial of degree $O_{k,l}(1)$.

%  Indeed, let $t=p_{1}^{m_{1}}\dots p_{r}^{m_{r}}$ be the prime factorization of $t$ with $p_{1}<\dots<p_{r}$. Then 
%  $$t-\phi(t)=(\sum_{i=1}^{r}\frac{1}{p_{i}}-\sum_{1\leq i_{1}<i_{2}\leq r}\frac{1}{p_{i_{1}}p_{i_{2}}}+\dots+(-1)^{r}\frac{1}{p_{1}\dots p_{r}})t$$
% and so
%  $$\phi(t)=t\prod_{i=1}^{r}(1-\frac{1}{p_{i}})=\prod_{i=1}^{r}(p_{i}-1)p_{i}^{m_{i}-1}.$$
%  Since $(p_{i}-1)p_{i}^{m_{i}-1}\geq p_{i}^{m_{i}/2}$ if $m_{i}\geq 2$ or if $m_{i}=1$ and $p_{i}\geq 3$. We have that $\phi(t)\geq \frac{1}{100}\sqrt{t}$. This completes the proof of Claim 3.
 
% \

The next claim enables us to combine the contributions of all the roots of unity $z_1, \ldots, z_{M'}$.

\begin{claim}\label{Claim 4}
Let $m,n\in\mathbb{N}$. If $(S^{m}-1)(S^{n}-1)f=0$, then $(S^{2mn}-1)f=0$.    
\end{claim}

To see this, note that $(S^{m}-1)(S^{n}-1)f=0$ implies that $(S^{mn}-1)^{2}f=0$ since $z^{n}-1$ and $z^{m}-1$ both divide $z^{mn}-1$. Then 
$(S^{2mn}+1)f=2S^{mn}f.$
     So $$4\int \vert f\vert^{2}=4\int \vert S^{mn}f\vert^{2}=\int \vert (S^{2mn}+1)f\vert^{2}=2\int \vert f\vert^{2}+\langle S^{2mn}f,f\rangle+\langle f,S^{2mn}f\rangle.$$
     By the Cauchy-Schwarz inequality, $\langle S^{2mn}f,f\rangle\leq \vert f\vert^{2}$ and the equality holds only if $S^{2mn}f=f$. Hence $(S^{2mn}-1)f=0$, and the claim follows. 

     \

     %\BK{(I changed the quantities below a bit. Note that the choice of $q$ depends only on $g$ (and hence on $a,k,l,u$) and not on $f$)} {\color{red} If we take $q=2^{M'-1}(O_{k,l}(1)!)^{M'}$, then $q$ depends only on $k,l$}
     Recall that $M'\leq M=O_{k,l}(1)$ and $q_{j}=O_{k,l}(1)$, the latter of which follows from Claim \ref{Claim 3}. A repeated application of Claim \ref{Claim 4} and (\ref{kl5}) allows us to find $W=O_{k,l}(1)$ and $q = 2^{M'-1} q_1 \cdots q_{M'}\leq W$ such that $S^{q}f=f$.
     %there exist $q, W=O_{k,l}(1)\in \N$ with $q\leq W$ such that 
     Let $H = \langle W! (w,v), (k,-l)\rangle$: this is a finite index subgroup of $\mathbb{Z}^{2}$ that depends only on $k$ and $l$. 
     Since $q|W!$, the fact that $f$ is invariant under both $S^{q}$ and $R$ implies that $f$ is $H$-invariant. This means that $\CI(R)=\CI(H)$. 
     %\BK{(In our statement, we have that $H$ can be made independent of $a$. While not particularly important, we should give reason why this is the case. Presumably this would follow if we know that $\Z^2$ admits $O_q(1)$ distinct index-$q$ subgroups - is this true?)}
     %Since $M'=O_{k,l}(1)$ and since $q_{j}=O_{k,l}(1)$ by Claim 3, it follows from Claim 4 and (\ref{kl5}) that there exists $W=O_{k,l}(1)\in \N$ and $q\in \mathbb{N}_{+}$ with $q\leq W$ such that $S^{q}f=f$. Let $H$ be the $\mathbb{Z}$-span of $W!n$ and $m$, which is certainly a finite index subgroup of $\mathbb{Z}^{2}$ depending only on $k$ and $l$. The fact that $f$ is $S^{q}$ and $T_{m}$ invariant implies that $f$ is $H$ invariant. This means that $I(T_{m})=I(H)$. 
     \end{proof}

% \begin{corollary} \BK{(Didn't check this one yet since maybe we would combine all these results together)}
%      Let $k,l\in\mathbb{Z}\backslash\{0\}$ and $u\in\R$. %and $a\colon\mathbb{Z}\to\mathbb{R}$ be a function such that the limit 
%     %$$\lim_{N\to\infty}\frac{\vert [N]\cap\{n\in\mathbb{N}\colon \{a(n)\}\in [x,y)\}\vert}{N}$$
%    % exists for all $$x,y\in \Bigl\{\frac{i}{\vert k\vert},\frac{j}{\vert l\vert}\colon 0\leq i\leq \vert k\vert, 0\leq j\leq \vert l\vert\Bigr\}.$$
%     Then there exists a finite index subgroup $H$ of $\mathbb{Z}^{2}$ which contains $(k,l)$ and depends only on $k$ and $l$ such that
%      for any ergodic $\mathbb{Z}^{d}$-system $(X,\mathcal{X},\mu, T_{1},\dots,T_{d})$ with $d\geq 2$ and any function $a\colon\mathbb{Z}\to\mathbb{R}$, if $(T_{1}^{\floor{ka(n)}}T_{2}^{-\floor{la(n)+u}})_{n}$ is ergodic for $\mu$, then for all function $f$, $K\in\mathbb{N}\cup\{0\}$ and subgroups $G_{1},\dots,G_{K}\subseteq \mathbb{Z}^{d}$, we have that
%      $$\nnorm{f}_{T_{1}^{k}T_{2}^{-l},G_{1},\dots,G_{K}}\leq [\mathbb{Z}^{2}:H]\nnorm{f}_{\mathbb{Z}^{d},G_{1},\dots,G_{K}}.$$
%      in particular, 
%      $$\nnorm{f}_{T_{1}^{k}T_{2}^{-l},G_{1},\dots,G_{K}}\ll_{k,l}\nnorm{f}_{\mathbb{Z}^{d},G_{1},\dots,G_{K}}.$$
% \end{corollary}
% \begin{proof} 
% Let $H$ denote the group given by Proposition \ref{P: beta rational}.  
% Then for all function $g$,
% $$\mathbb{E}_{n\in\mathbb{Z}}(T_{1}^{k}T_{2}^{-l}))^{n}g=\mathbb{E}(g\vert I(T_{1}^{k}T_{2}^{-l}))=\mathbb{E}(g\vert I(H))=\mathbb{E}_{h\in H}T_{h}g.$$
%  By Lemma 1.5 in "RealPolys(V2)", we have 
%   $$\nnorm{f}_{T_{1}^{k}T_{2}^{-l},G_{1},\dots,G_{K}}=\nnorm{f}_{H,G_{1},\dots,G_{K}}.$$
%   By ???,
%  $$\nnorm{f}_{H,G_{1},\dots,G_{K}}\leq [\mathbb{Z}^{2}:H]\nnorm{f}_{\langle T_{1},T_{2}\rangle,G_{1},\dots,G_{K}}.$$
%  On the other hand, since $(T_{1}^{\floor{ka(n)}}T_{2}^{-\floor{la(n)+u}})_{n}$ is ergodic, we must have that $\langle T_{1},T_{2}\rangle$ is ergodic. So we have that
%  $$\nnorm{f}_{\langle T_{1},T_{2}\rangle,G_{1},\dots,G_{K}}=\nnorm{f}_{\mathbb{Z}^{d},G_{1},\dots,G_{K}}.$$
%  This completes the proof.
% \end{proof}


\section{Seminorm comparison under ergodicity assumptions II: irrational $\beta$}\label{S: seminorm comparison II}

We continue our exploration of consequences that can be drawn from the ergodicity of $(T_1^{\floor{a_1(n)}}T_2^{-\floor{a_2(n)}})_n$ whenever $\beta_2 a_1 -\beta_1 a_2\ll \log$ for some nonzero $\beta_1,\beta_2\in\R$. 
% So far, we have shown two things:
% \begin{enumerate}
%     \item the ergodicity of this action implies the ergodicity of $(T_1^{\floor{{a}(n)}}T_2^{-\floor{\beta{a}(n)+u}})_n$ for $\beta = \beta_2/\beta_1$, some $u\in\R$, and $a\in\CH$, the latter of which is either in $\R[x]$ or stays logarithmically away from $\R[x]$;
%     \item a control of $\nnorm{f}_{\be_1 - \beta \be_2}$ by a seminorm along $\langle \be_1, \be_2\rangle$ whenever $(T_1^{\floor{ {a}(n)}}T_2^{-\floor{\beta {a}(n)+u}})_n$ is ergodic and $\beta\in\Q$.
% \end{enumerate}
In the preceding section, we have proved Corollary \ref{C: seminorm comparison, beta rational}, a special case of Proposition \ref{P: seminorm comparison, s = 0} for rational $\beta$. Our objective now is to complete the proof of Proposition \ref{P: seminorm comparison, s = 0} (and hence also of Proposition \ref{P: seminorm comparison}) by covering the case of $\beta\notin\Q$. This turns out to require significantly more ad hoc computations than the case of rational $\beta$. For convenience, we state the main output of this section as a separate proposition.

 %Our goal now is to use the ergodicity of $(T_1^{\floor{\alpha a(n)}}T_2^{-\floor{\beta a(n)+u}})_n$, now with $\beta/\alpha\notin\Q$, to control the seminorm along $\alpha \be_1 - \beta \be_2$ by a (higher-degree) seminorm along $\langle \be_1, \be_2\rangle$. 
\begin{proposition}\label{P: beta irrational}
    Proposition \ref{P: seminorm comparison, s = 0} holds for $\beta\notin\Q$. 
    
    More specifically, let $(X, \CX, \mu, T_1, T_2)$ be a system, let $u\in\R$ and $a\in\CH$, where $a$ is either in $\R[x]$ or stays logarithmically away from $\R[x]$, and suppose that $(T_1^{\floor{a(n)}}T_2^{-\floor{\beta a(n)+u}})_n$ is ergodic for $(X, \CX, \mu)$. Then for any $f\in L^\infty(\mu)$, we have
    \begin{align*}
        \nnorm{f}_{\be_1 - \beta \be_2} = 0 \quad \textrm{whenever} \quad \nnorm{f}_{\langle\be_1, \be_2\rangle}^+ = 0.
    \end{align*}
\end{proposition}
%Proposition \ref{P: beta irrational} will allow us to replace the instances of $ \be_1 - \beta \be_2$ by any of $\be_1$ or $\be_2$ in the seminorm, controlling an average that involves $T_1^{\floor{a(n)}}, T_2^{\floor{\beta a(n)+u}}$.

% Throughout this section, let $a\in\CH$ with $a\succ \log$. Since
% \begin{align}\label{E: alpha e_1 - beta e_2 vs. e_1, e_2}
% \nnorm{f}_{\be_1 - \beta \be_2}\leq \nnorm{f}_{\alpha \be_1 - \beta \be_2}^+ \asymp_\alpha \nnorm{f}_{\be_1 - (\beta/\alpha) \be_2}^+\leq \nnorm{f}_{\langle\be_1, \be_2\rangle,\; \be_1 - (\beta/\alpha) \be_2},
% \end{align}
% and similarly we can redefine $a(n)$ as $\alpha a(n)$, we can (and will) assume without loss of generality that $\alpha =1$ while $\beta$ is irrational and greater than 1. 
 

% extract as much information as possible from the assumption that $(T_1^{\floor{ a(n)}}T_2^{-\floor{\beta a(n)}})_n$ is ergodic. We assume throughout that $\beta\notin\Q$, as the case $\beta\in\Q$ has been handled separately. Without loss of generality, we also assume that $\beta>1$. Specifically, we want to show that if $(T_1^{\floor{ a(n)}}T_2^{-\floor{\beta a(n)}})_n$ is ergodic, then
% \begin{align*}
%     \norm{\E_n T_1^{n}T_2^{\floor{-\beta n}} f}_{L^2(\mu)} = \nnorm{f}_{ e_1-\beta e_2} \ll \nnorm{f}_{\langle e_1, e_2\rangle} \leq \min_{i\in[2]}\nnorm{f}_{e_i},
% \end{align*}
% as this would allow us to replace the instances of $e_1 - \beta e_2$ by any of $e_1$ or $e_2$ in the box seminorm that controls our average. 
% If $T_1, T_2$ are ergodic, than the bound above would follow from
% \[\norm{\E_n T_1^{n}T_2^{-\floor{\beta n}}f}_2\ll \norm{\E_n T_1^{\floor{ a(n)}}T_2^{-\floor{\beta a(n)}}f}_2,\]
% which is the opposite of what we usually do.

From now on, all the way up to the end of Section \ref{S: seminorm comparison II}, we fix all the data in the statement of Proposition \ref{P: beta irrational}: the system $(X, \CX, \mu, T_1, T_2)$, the numbers $u\in\R$ and $\beta\notin\Q$, and the  {function} $a\in\CH$. {We also make the following standing assumption:

\smallskip
($\clubsuit$)\quad \emph{$a$ is either in $\R[x]$ or stays logarithmically away from $\R[x]$.}
\smallskip}

For the discussion below, we also fix $f\in L^\infty(\mu)$.
Our main tool to analyze the ergodic averages induced by the action $(T_1^{\floor{ a(n)}}T_2^{-\floor{\beta a(n)+u}})_n$ is the classical Herglotz-Bochner theorem. It shows that
\begin{align}\label{E: L^2 for A}
    \lim_{N\to\infty}\norm{\E_{n\in[N]} T_1^{\floor{ a(n)}}T_2^{-\floor{\beta a(n)+u}} f}_{L^2(\mu)} %\lim_{N\to\infty}\norm{\E_{n\in[N]} e(\floor{ a(n)}t -\floor{\beta a(n)+u}s)}_{L^2(\sigma_f(t,s))}\\
    = \norm{A}_{L^2(\sigma_f)},
\end{align}
where $\sigma_f$ is the spectral measure on $[0,1)^2$ associated with $f$ and
\begin{multline*}
    A(t,s):= e(us)\cdot \lim_{N\to\infty}\E_{n\in[N]} e(\floor{ a(n)}t -\floor{\beta a(n)+u}s)\\
    =  \lim_{N\to\infty}\E_{n\in[N]} e((t - \beta s) a(n)- \{ a(n)\}t +\{\beta a(n)+u\}s).
\end{multline*}
We recall that the limits in \eqref{E: L^2 for A} exist by Lemma \ref{L: convergence of exp sum}, just like the limits below.

The exponential sum that we have to compare $A(t,s)$ with %\footnote{Technically, we should look at $\E_n e(t n +\floor{-\beta n}s)$, but since $\beta\notin\Q$, we have $\floor{-\beta n} = -\floor{\beta n} - 1$ for all $n\neq 0$, and so $\E_n e(t n +\floor{-\beta n}s) = e(-s)\E_n e(t n -\floor{\beta n}s)$.}
 is
\begin{multline*}
    B(t,s):= e(s) \cdot \lim_{N\to\infty}\E_{n\in[N]} e(t n +\floor{-\beta n}s) = \lim_{N\to\infty}\E_{n\in[N]} e(t n -\floor{\beta n}s)\\
    =  \lim_{N\to\infty}\E_{n\in[N]} e((t - \beta s) n+\{\beta n\}s)
\end{multline*}
(in the second equality, we use the irrationality of $\beta$ to write $\floor{-\beta n} = -\floor{\beta n} - 1$ for $n\in\Z\backslash\{0\}$)
 because 
% \AK{careful in the following about this $+u$ term. What we actually get now is that $B(t,s)$ equals 
% \[e(-\floor{u}s)\E_n \left(1_{\{n:\floor{\beta n+u}=\floor{\beta n}+\floor{u}\}}+e(-s)1_{\{n:\floor{\beta n+u}=\floor{\beta n}+\floor{u}+1\}}\right)e(t n -\floor{\beta n}s).\]
% The issue here is that any approximation will potentially introduce error terms that we don't want to have. So, we might need to find explicitly the value of $B(t,s).$ Most likely (this needs checking) this new $B$ will have the same value as the other one.}\BK{The $u$ term doesn't appear in the definition of $B$, only in the definition of $A$.}
\begin{align}\label{E: L^2 for B}
    \nnorm{f}_{\be_1 - \beta \be_2} = \lim_{N\to\infty}\norm{\E_{n\in[N]} T_1^n T_2^{\floor{-\beta n}} f}_{L^2(\mu)} 
    %= \norm{\E_n e(t n +\floor{-\beta n}s)}_{L^2(\sigma_f(t,s))}
    = \norm{B}_{L^2(\sigma_f)}
\end{align}
by the Herglotz-Bochner theorem again.

Our main goal now is to obtain a precise description of the values
%\footnote{We identify $\T^2$ with $[0,1)^2$ so that whenever $t,s$ appear in expressions that are not invariant under shifts $t\mapsto t\pm 1$ or $s\mapsto s\pm1$, we think of $(t,s)$ as lying in $[0,1)^2$.} 
$(t,s)\in[0,1)^2$ for which any of $A(t,s), B(t,s)$ vanishes, so that we can compare \eqref{E: L^2 for A} with \eqref{E: L^2 for B} in a way that allows us to control \eqref{E: L^2 for B} by a seminorm along $\langle \be_1, \be_2\rangle$ as in Proposition \ref{P: beta irrational}.
To this end, we split 
\begin{align*}
    [0,1)^2 = E_1 \cup E_2 \cup E_3,
\end{align*}
where 
\begin{align*}
    E_1 &= \{(t,s)\in[0,1)^2:\; A(t,s) = B(t,s) = 0\},\\
    E_2 &= \{(t,s)\in[0,1)^2:\; A(t,s) \neq 0\},\\
    E_3 &= \{(t,s)\in[0,1)^2:\; A(t,s) = 0,\; B(t,s) \neq 0\}.
\end{align*}

Naturally, the contribution of $E_1$ to \eqref{E: L^2 for A} and \eqref{E: L^2 for B} is zero. To evaluate the contribution of $E_2$, we use the ergodicity assumption on our action to observe that 
\begin{align*}
    \abs{\int f\, d\mu}^2 = \sigma_f(\{(0,0)\})
    %\norm{\E_n T_1^{\floor{ a(n)}}T_2^{-\floor{\beta a(n)}} f}_{L^2(\mu)} = \norm{\E_n e(\floor{ a(n)}t -\floor{\beta a(n)}s)}_{L^2(\sigma_f(t,s))}
    = \norm{A}_{L^2(\sigma_f)}.
\end{align*}
Since $(0,0)\in E_2$ and $|A(t,s)|>0$ for every $(t,s)\in E_2$, we must have $\sigma_f(E_2) = \sigma_f(\{0,0\})$ from the continuity of the spectral measure. In particular, it follows that 
\begin{align*}
    \int_{E_2}|A(t,s)|^2\, d\sigma_f(t,s) = \int_{E_2}|B(t,s)|^2\, d\sigma_f(t,s) = \sigma_f(\{0,0\}) = \abs{\int_X f\, d\mu}^2,
\end{align*}
and so the contribution of $E_2$ to \eqref{E: L^2 for A} and \eqref{E: L^2 for B} is equal. Our objective therefore is to evaluate the contribution of $E_3$, and for that we need to understand when $A$ vanishes while $B$ does not. 
% \BK{(Needed?)} In doing so, it will be beneficial to split $E_3$ into the following two sets:
%          \begin{align}
%                 \label{E: E_3'}
%                 {E}_3'&=\{(t,s)\in[0,1)^2:\; A(t,s) = 0,\; B(t,s)\neq 0,\; A(Mt, Ms)\neq 0\\
%                 \nonumber 
%                 &\qquad\qquad\qquad\qquad\qquad\qquad\textrm{for\; some}\; M\in\Z\setminus\{0\} {\text{ with } (Mt,Ms)\notin\Z^{2}}\}\\  
%                 \label{E: E_3''}
%                 {E}_3''&=\{(t,s)\in[0,1)^2:\; A(Mt,Ms) = 0,\; B(t,s)\neq 0\\
%                 \nonumber 
%                 &\qquad\qquad\qquad\qquad\qquad\qquad\textrm{for\; all}\; M\in\Z\setminus\{0\} {\text{ with } (Mt,Ms)\notin\Z^{2}}\}.
%          \end{align}
%This will be achieved by explicit computations of various exponential sums, and 
After doing so, we arrive at the following description of $E_3$.
\begin{proposition}\label{P: E_3}
The set $E_3$ is countable, and we have
\begin{align*}
E_3 \subseteq \{\brac{\beta \gamma_i, \gamma_i}:\; i\in\N\} \mod\Z^2
%+ \Z^2
\end{align*}
for some sequence $(\gamma_n)_n$ of nonzero real numbers. %Moreover, we have  $\sigma_f({E}_3')=0$.
\end{proposition}
For a set $E\subseteq\R^2$, we interpret $E$ mod $\Z^2$ as $\{(\rem{t}, \rem{s}):\; (t,s)\in E\}$.

The proof of Proposition \ref{P: E_3} will proceed by explicit computations of various exponential sums and can be broken down into the following four cases:
\begin{enumerate}
    \item $a(x)$ does not equal
    \begin{align}\label{E: special form of a}
        \frac{p(x)-ku}{k\beta + l} \quad \textrm{for\; any}\;\; (k, l)\in\Z^2\setminus{\{(0,0)\}},\; p\in\Z[x]{+\R}; %\; g \ll \log;
    \end{align}
    \item $a(x)$ equals \eqref{E: special form of a} with $k\neq 0, l=0$;
    \item $a(x)$ equals \eqref{E: special form of a} with $k=0, l\neq 0$;
    \item $a(x)$ equals \eqref{E: special form of a} with $k, l\neq 0$.
\end{enumerate}
% In all cases, we will establish several facts about the set $E_3$:
% \begin{enumerate}
%     \item the set $E_3$ is countable;
%     \item (Lemma \ref{L: E_3'}) for any $(t,s)\in E_3'$, we have $\sigma_f(t,s) = 0$ (hence, combined with the previous point, we get $\sigma_f({E}_3')=0$);\footnote{Here and elsewhere, we abbreviate $\sigma_f(t,s)$ for $\sigma_f(\{(t,s)\})$.}
%     \item the set $E_3''$ will give the \textit{structured} part of $E_3$, i.e. the finitely many groups $\left\langle\brac{\beta \gamma_i, \gamma_i} \right\rangle_{[0,1)^2}$.
% \end{enumerate}
% Unfortunately, we have not found a way of establishing the first and third facts other than through a tedious case-by-case analysis. By contrast, the second fact can be established regardless of the form of $a$.
% \begin{lemma}\label{L: E_3'}
%     If $(t,s)\in E_3'$, then $\sigma_f(t,s) = 0$.
% \end{lemma}
% \begin{proof}
%     By the ergodicity of $(T_{1}^{\floor{a(n)}}T_{2}^{-\floor{\beta a(n)+u}})_{n}$, the only point $(t,s)\in[0,1)^2$ with $|A(t,s)|~\cdot~\sigma_f(t,s) > 0$ is $(t,s) = (0,0)$. In particular, if $(t,s)\neq (0,0)$ and $A(Mt,Ms)\neq 0$ for some nonzero $M\in\mathbb{Z}$ with $(Mt,Ms)\neq (0,0)$, then $\sigma_f(Mt,Ms) = 0$. Thus, $(Mt, Ms)$ is
%     not an eigenvalue of $f$. Since eigenvalues form a group, it follows that $(t,s)$ is not an eigenvalue of $f$ either, and so $\sigma_f(t,s) = 0$.
% \end{proof}

The analysis of $A(t,s)$ and $B(t,s)$ differs markedly depending on whether $t-\beta s, \beta, 1$ are $\Q$-independent or $\Q$-dependent. The first case can be excluded almost immediately by the lemma below. 
\begin{lemma}\label{L: Q-independence}
    Let $(t,s)\in[0,1)^2$ so that $t-\beta s, \beta, 1$  are $\Q$-independent.
     Then $B(t,s) = 0$.  As a consequence, $$E_3\subseteq \{(t,s)\in[0,1)^2:\; t - \beta s = q + \beta r\; \textrm{for\; some}\; q,r\in\Q\}.$$   

\end{lemma}

\begin{proof}
By the assumption on $(t,s)$ and Weyl's equidistribution theorem, we have
\begin{align*}
    B(t,s) = \lim_{N\to\infty}\E_{n\in[N]} e((t - \beta s) n+\{\beta n\}s) =  \int_{[0,1)^2}e(x + zs)\; dxdz = 0.
\end{align*}
In particular, such $(t,s)$ cannot lie in $E_3$.
\end{proof}


\subsection{Exponential sum computations for $\Q$-dependent $t-\beta s, \beta, 1$}\label{SS: Q-dependent}
The computations become much more intricate when $ t-\beta s, \beta, 1$ are $\Q$-dependent, in which case $t-\beta s = q+\beta r$ for some $q,r\in\Q$.
The following lemmas give sufficient and necessary conditions for $A(t,s) = 0$ to hold for different cases of $a$ under the assumption of $\Q$-dependence of $ t-\beta s, \beta, 1$.
\begin{lemma}\label{L: a, beta a eq}
Let $q, r\in\Q$ and suppose that $a\in\CH$ does not equal $\frac{p - ku}{k\beta + l}$ for any nonconstant $p\in\Z[x]{+\R}$ and $(k, l)\in\Z^2\setminus{\{(0,0)\}}$. Then
\begin{multline*}
    \lim_{N\to\infty}\E_{n\in[N]} e(q a(n)+r \beta a(n)-\{ a(n)\}t +\{\beta a(n)+u\}s)\\
=    e(-ru)1_\Z(q) \brac{1_{q=t} + 1_{\R\setminus\Z}(q-t)\frac{1}{2\pi i(q-t)}(e(q-t)-1)}\\
1_\Z(r)\brac{1_{r=-s} + 1_{\R\setminus\Z}(r+s)\frac{1}{2\pi i(r+s)}(e(r+s)-1)}.
%\lim_{N\to\infty}\E_{n\in[N]} e((t - \beta s) a(n)- \{ a(n)\}t +\{\beta a(n)+u\}s)
\end{multline*}
% {\color{red} If we want to simplify the notations, we could define a function 
% $\text{Int}_{e}(\alpha):=\int_{[0,1)}e(\alpha t)\,dt$, which is equal to $\frac{e(\alpha)-1}{2\pi i\alpha}$ if $\alpha\notin\R\backslash\{0\}$ and is equal to 1 otherwise.
% In particular, $\text{Int}_{e}(\alpha)=0\Leftrightarrow \alpha\in\Z\backslash\{0\}$.
% Then
%   the RHS of the above equation becomes
% $$e(-ru)1_\Z(q) \text{Int}_{e}(q-t)
% 1_\Z(r)\text{Int}_{e}(r+s)$$
%   We did something similar in
%   my paper with Andreas.
% } \BK{I'd stick to what we have, the level of complexity of the current expression is still human understandable, and it is easier to read off when the exponential sum vanishes.}
 In particular, this is 0 if and only if one of the followings holds:
    \begin{enumerate}
    \item $q\notin \Z$;
    \item $q-t\in\mathbb{Z}\backslash\{0\}$;
    \item $r\notin \mathbb{Z}$;
    \item $r+s\in\mathbb{Z}\backslash\{0\}$.
\end{enumerate}
\end{lemma}
\begin{proof}
By Theorem \ref{T: Boshernitzan} and Weyl's equidistribution theorem, the condition on $a$ is equivalent to the sequence
\begin{align}\label{E: 2-dim seq 1}
    (\{a(n)/c\}, \{(\beta a(n)+u)/c\})_n    
\end{align}
being equidistributed on $[0,1)^2$ for every nonzero $c\in\Z$. {Indeed, for the sake of equidistribution, we want $l\frac{a(x)}{c}+k\frac{\beta a(x)+u}{c} = \frac{k\beta + l}{c}a(x)+\frac{ku}{c}$ not to lie in $\Q[x]+\R$ for any $k,l\in\Z$ not both 0 (the cases when it equals an element of $\Q[x]+\R$ plus a nonconstant function $g$ such that $g\to 0$ or $1\prec g\ll \log$ are excluded by the assumption $\clubsuit$).}  This happens precisely if $a$ takes the prescribed form.

Let $c\in\Z$ be such that $q, r\in\frac{1}{c}\Z$.
We therefore rephrase our exponential sum as
 \begin{align*}
     %\lim_{N\to\infty}\E_{n\in[N]} e(qa(n)+\beta r a(n)-\{ a(n)\}t +\{\beta a(n)+u\}s) = 
     \lim_{N\to\infty}\E_{n\in[N]} F(\{a(n)/c\}, \{(\beta a(n)+u)/c\})
 \end{align*}
 for $F(x,y) =  e(cq x+cr y-\{cx\}t +\{cy\}s-ru)$. The equidistribution of \eqref{E: 2-dim seq 1}
 %$(\{a(n)/c\}, \{(\beta a(n)+u)/c\})_n$ 
 on $[0,1)^2$ gives that our average equals
 \begin{align*}
     %\E_n F(\{a(n)/c\}, \{\beta a(n)/c\}) 
      \int_{[0,1)^2}F(x,y)\, dx dy
     &= e(-ru)\int_{[0,1)^2}e(cqx+cr y-\{cx\}t +\{cy\}s)\, dx dy.
 \end{align*}
 This expression then equals
 \begin{equation*}
     \begin{split}
       & %\int_{[0,1)^2}e(\{cqx + cry\}-\{cx\}t +\{cy\}s)\, dxdy
       e(-ru)\int_{[0,1)^2}e(cq x+ cry-\{cx\}t +\{cy\}s)\, dxdy
       \\&=e(-ru)\sum_{i,j=0}^{c-1}\int_{[0,\frac{1}{c})^2}e\brac{cq\brac{x+i/c}+cr\brac{y+j/c}-\{c\brac{x+i/c}\}t +\{c\brac{y+j/c}\}s}\, dxdy
       \\&=e(-ru)\sum_{i,j=0}^{c-1}e\brac{qi}e\brac{rj}\int_{[0,\frac{1}{c})^2}e(cqx + cry-\{cx\}t +\{cy\}s)\, dxdy
        \\&=e(-ru)\frac{1}{c}\sum_{i=0}^{c-1}e\brac{qi}\cdot \frac{1}{c}\sum_{j=0}^{c-1}e\brac{rj}\cdot \int_{[0,1)}e((q-t)x)\, dx\cdot\int_{[0,1)}e((r+s)y)\, dy\\
        &= e(-ru)1_\Z(q)  \brac{1_{q=t} + 1_{\R\setminus\Z}(q-t)\frac{1}{2\pi i(q-t)}(e(q-t)-1)}\\
        &\qquad\qquad\qquad\qquad\qquad\qquad 1_\Z(r)\brac{1_{r=-s} + 1_{\R\setminus\Z}(r+s)\frac{1}{2\pi i(r+s)}(e(r+s)-1)}.
     \end{split}
 \end{equation*}
 The four terms in the product are zero if and only if (iii), (iv), (i) and (ii) holds respectively.
\end{proof}




The next few lemmas cover the cases when $a = \frac{p - ku}{k\beta + l}$ for some nonconstant {$p\in\Z[x]+\R$} and various choices of $k,l\in\Z$.
%, , and $t-\beta s, \beta, 1$ are $\Q$-rationally dependent.  

%{\color{yellow} This is LEMMA 4.13 from smoothing file.}
\begin{lemma}[{Case $l=0$}]\label{L: beta a  = p + g}
Let $c\in\N$, $k\in\Z\backslash\{0\}$, $q,r\in\frac{1}{c}\Z$ and suppose that $a = \frac{p - ku}{k\beta}$ for some nonconstant $p\in\Z[x]{+\R}$. Then
\begin{multline*}%\label{E: ml22}
    \lim_{N\to\infty}\E_{n\in[N]} e(qa(n)+\beta r a(n)-\{ a(n)\}t +\{\beta a(n)+u\}s)\\
    =    e\brac{-ru}\E_{j\in[ck]}e\brac{\frac{r p(j)}{k}+ \rem{\frac{p(j)}{k}}s}1_\Z(q)\\
    \brac{1_{q=t} + 1_{\R\setminus\Z}(q-t)\frac{1}{2\pi i(q-t)}(e(q-t)-1)}.
\end{multline*}
 In particular, this is 0 if and only if one of the followings holds:
    \begin{enumerate}
    \item $\E\limits_{j\in[ck]}e\brac{\frac{r p(j)}{k}+ \rem{\frac{p(j)}{k}}s}=0$;
    \item $q\notin\Z$;
    \item $q-t\in\mathbb{Z}\backslash\{0\}$.
\end{enumerate}
\end{lemma}

We remark that the first condition is tantamount to saying that the infinite average $\lim\limits_{N\to\infty}\E\limits_{n\in[N]}e\brac{\frac{rp(n)}{k}+ \rem{\frac{p(n)}{k}}s}$ converges to 0.


\begin{proof}

{If $k<0$, then we can make it positive by letting $k\mapsto - k$, $p\mapsto - p$, so we assume that $k>0$.}
%Let $k\in\N$ be a number such that $k\cdot p\in\Z[n]$. 
On splitting the average into arithmetic progressions of difference $ck$, it equals
\begin{multline*}
    %\lim_{N\to\infty}\E_{n\in[N]}e(qa(n)+\beta r a(n)-\{ a(n)\}t +\{\beta a(n)+u\}s)\\
 e\brac{-ru}\cdot\E_{j\in[ck]}e\brac{\frac{r p(j)}{k}+\rem{\frac{p(j)}{k}}s} \\
 \lim_{N\to\infty}\E_{n\in[N]}e\left(\frac{qp(ckn+j)-qku}{k\beta} -\rem{\frac{p(ckn+j)-ku}{k\beta}}t\right)
\end{multline*}
since
\begin{align*}
\beta a(ckn+j) +u= \frac{p(j)}{k}\!\!\!\!\mod \Z\quad \textrm{and}\quad 
    \beta r a(ckn+j) = \frac{r p(j)}{k}-ru \!\!\!\!\mod \Z.
\end{align*}
By Theorem \ref{T: Boshernitzan}, the sequence $\brac{\rem{\frac{p(ckn+j)-ku}{ck\beta}}}_n$ is equidistributed on $[0,1)$; 
% We now apply Lemma \ref{L: uniform weyl for Riemann integrable} with $\ell=1$ and
% \begin{align}\label{E: p + w}
%     p_1(n) = \frac{1}{ck\beta}(p(ckn+i)+w);
% \end{align}
%$\big(\frac{1}{ c}(\frac{1}{\beta}p(k(cn+i)+l)+w)\big)_n$; 
the reason why we divide by $c$ is for $\frac{qp(ckn+j)-qku}{k\beta}$ 
to be an integer multiple of $\frac{p(ckn+j)-ku}{ck\beta}$. Then our average equals
%The equidistribution of the sequence \eqref{E: p + w} on $\T$ (and the relation $\{mx\}=\{m\{x\}\}$ for $m\in\Z$) allows us to conclude that
\begin{align*}
        %\E_{n}e(qa(n)+\beta r a(n)-\{ a(n)\}t +\{\beta a(n)\}s)\\
         e\brac{-ru}\cdot\E_{j\in[ck]}e\brac{\frac{r p(j)}{k}+ \rem{\frac{p(j)}{k}}s}\cdot \int_0^1 e(cq x- \{cx\}t)\, dx.
\end{align*}
The same integral computation as in the proof of Lemma \ref{L: a, beta a eq} shows that this is
\begin{align*}
    e\brac{-ru}\E_{j\in[ck]}e\brac{\frac{r p(j)}{k}+ \rem{\frac{p(j)}{k}}s} 1_\Z(q)\brac{1_{q=t} + 1_{\R\setminus\Z}(q-t)\frac{1}{2\pi i(q-t)}(e(q-t)-1)},
\end{align*}
as claimed.
\end{proof}


\begin{lemma}[{Case $k=0$}]\label{L: a = p + g}
Let $c\in\N$, $l\in\Z\backslash\{0\}$, $q,r\in\frac{1}{c}\Z$ and suppose that $a = \frac{p}{l}$ for some nonconstant $p\in\Z[x]{+\R}$. Then
\begin{multline*}%\label{E: ml22}
    \lim_{N\to\infty}\E_{n\in[N]} e(qa(n)+\beta r a(n)-\{ a(n)\}t +\{\beta a(n)+u\}s)\\
=e(-ru)\E_{j\in[cl]}e\brac{\frac{q p(j)}{l}-\rem{\frac{p(j)}{l}}t}1_\Z(r)\\
\brac{1_{r=-s} + 1_{\R\setminus\Z}(r+s)\frac{1}{2\pi i(r+s)}(e(r+s)-1)}.
\end{multline*}
 In particular, this is 0 if and only if one of the followings holds:
    \begin{enumerate}
    \item $\E\limits_{j\in[cl]}e\brac{\frac{q p(j)}{l}-\rem{\frac{p(j)}{l}}t}=0$;
    \item $r\notin \Z$;
    \item $r+s\in\mathbb{Z}\backslash\{0\}$.
    %\item $q-t\in\mathbb{Z}\backslash\{0\}$;
    %\item $q\notin\Z$;
\end{enumerate}
\end{lemma}

% {\color{red}
% Suppose that $a = \frac{1}{l}(p + g)$ for  $p\in\Z[n], g\in\CS$. Then $A(t,s)\neq 0$  if and only if  
% $t-\beta s,\beta, 1$ are $\mathbb{Q}$-dependent, and that writing $t-\beta s=q+\beta r$ be such that $q,r\in\mathbb{Q}$, we have that:
%     \begin{enumerate}
%     \item $r+s\notin\mathbb{Z}\backslash\{0\}$;
%     %\item $q-t\in\mathbb{Z}\backslash\{0\}$;
%     %\item $q\notin\Z$;
%     \item $r\in \Z$.
%     \item $\E_{i\in[cl]}e\brac{\frac{q}{l}p(i)-\rem{\frac{1}{l}p(i)}t} \neq 0$.
% \end{enumerate}

% }

\begin{proof}
{Like in the previous lemma, we can assume without loss of generality that $l>0$.}
On splitting the average into arithmetic progressions of difference $cl$, it equals
\begin{align*}
\E_{j\in[cl]}e\brac{\frac{q p(j)}{l}-\rem{\frac{p(j)}{l}}t} \lim_{N\to\infty}\E_{n\in[N]} e\left(\frac{\beta r p(cln+j)}{l} \right. 
    \left.    +\rem{\frac{\beta p(cln+j)}{l} +u}s\right)
%e\left(\beta r a(cln+i)}\right.\\  \left.-\rem{\frac{1}{l}(p(i)+ g(cln + i))}t+\{\beta a(cl n+i)\}s\right)
\end{align*}
since
\begin{align*}
    qa(cl n + j) = \frac{q p(j)}{l}\!\!\!\! \mod \Z \quad \textrm{and}\quad 
    a(cl n + j) = \frac{p(j)}{l}\!\!\!\! \mod\Z.
\end{align*}
By Theorem \ref{T: Boshernitzan}, the sequence $\brac{\rem{\frac{\beta p(cln+j)}{cl}+u}}_n$ is equidistributed on $[0,1)$, and hence the average equals
\begin{multline*}
    e(-ru)\E_{j\in[cl]}e\brac{\frac{q p(j)}{l}-\rem{\frac{p(j)}{l}}t}\int_{0}^{1} e\big(rcx+s\{cx\}\big)dx=\\
     e(-ru)\E_{j\in[cl]}e\brac{\frac{q p(j)}{l}-\rem{\frac{p(j)}{l}}t}1_\Z(r)\brac{1_{r=-s} + 
    1_{\R\setminus\Z}(r+s)\frac{1}{2\pi i(r+s)}(e(r+s)-1)},
\end{multline*}
where the integral is evaluated as in the proof of Lemma \ref{L: a, beta a eq}.
\end{proof}






In the special case where $a(n) = n$ and $u=0$, Lemma \ref{L: a = p + g} gives the following corollary that can be used to describe the zero set of $B$.
\begin{corollary}[Case $a(n)=n$]\label{C: a = n}
%Let $c\in\N$ and $q,r\in\frac{1}{c}\Z$. 
Let $q,r\in\Q$.
Then
\begin{multline*}
    \lim_{N\to\infty}\E_{n\in[N]} e((q+\beta r)n +\{\beta n\}s)\\
    =     1_\Z(q) 1_\Z(r)
\brac{1_{r=-s} + 1_{\R\setminus\Z}(r+s)\frac{1}{2\pi i(r+s)}(e(r+s)-1)}.
\end{multline*}
    In particular, it is 0 if and only if one of the followings holds:
    \begin{enumerate}
        \item $q\notin\Z$;
    \item $r\notin \Z$;
    \item $r+s\in\mathbb{Z}\backslash\{0\}$.
    %\item $q-t\in\mathbb{Z}\backslash\{0\}$;
\end{enumerate}
\end{corollary}























































































































\begin{comment}
    \begin{lemma}\label{L: (k beta + l) a  = p + g}
Let $c, k, l\in\N$ with $k, l\neq 0$, $q,r\in\frac{1}{c}\Z$ and suppose that $a = \frac{1}{k\beta + l}(p + g)$ for  $p\in\Z[n], g\ll \log$. Then, if the averages 
\begin{equation*}
    \E_ne(\{qa(n)+\beta r a(n)\}-\{ a(n)\}t +\{\beta a(n)\}s)\iffalse\\
%=    \E_{j\in[ck]}e(rp(j)+\{p(j)\}s)\cdot 1_\Z(q)\brac{1_{q=t} + 1_{\R\setminus\Z}(q-t)\frac{1}{2\pi i(q-t)}(e(q-t)-1)}.
    =1_{(\R\setminus kl\cdot\Z)\cup\{0\}}(k(q-t)-l(r+s))\cdot 1_{(\mathbb{Z}/c)\backslash\mathbb{Z}}(kt+ls)\cdot
    \E_{j\in[ck]}\E_{i_2\in[0, l-1]}\E_{i_3\in[0, k-1]}\\
    %&\qquad\qquad\qquad\qquad\qquad
    e\left(\frac{{r}+s}{k}p(j) + \floor{\frac{k i_2 + i_3}{l}}t + i_2 s- \floor{\frac{-i_3-1+p(j)}{k}}s\right)\fi
\end{equation*}
 converge to 0, then one of the followings holds:
      \begin{enumerate}
    \item $k(q-t)-l(r+s)\in (\R\setminus kl\cdot\Z)\cup\{0\}$;
    \item $kt+ls\in(\mathbb{Z}/c)\backslash\mathbb{Z}$;
    \item $\E_{j\in[ck]}\E_{i_2\in[0, l-1]}\E_{i_3\in[0, k-1]}
    e\left(\frac{{r}+s}{k}p(j) + \floor{\frac{k i_2 + i_3}{l}}t + i_2 s- \floor{\frac{-i_3-1+p(j)}{k}}s\right)=0$.
\end{enumerate}
\end{lemma}
\begin{proof}

Let $\tilde{q} = qc, \tilde{r} = rc\in\Z$, so that the average in question equals
\begin{align*}
    \E_{n}e\brac{\rem{\frac{\tilde{q} +\tilde{r}\beta}{c(k\beta + l)}(p(n)+g(n))}-\rem{\frac{1}{k\beta + l}(p(n)+g(n))}t +\rem{\frac{\beta}{k\beta + l} (p(n)+g(n))}s}.
\end{align*}
Using the identity
\begin{align*}
    \frac{x\beta + y}{k\beta + l} = \frac{ky - lx}{k(k\beta + l)} + \frac{x}{k},
\end{align*}
we want to bring all the fractions above to the form $\Q + \frac{\Z}{ck(k\beta + l)}$ and then apply Weyl's equidistribution theorem. Indeed, our average equals
\begin{multline*}
    \E_{n}e\left(\rem{\brac{\frac{\tilde{q}k -\tilde{r}l}{ck(k\beta + l)}+\frac{\tilde{r}}{ck}}(p(n)+g(n))}-\rem{\frac{ck}{ck(k\beta + l)}(p(n)+g(n))}t\right.\\
    \left.+\rem{\brac{\frac{-cl}{ck(k\beta + l)}+\frac{1}{k}} (p(n)+g(n))}s\right).
\end{multline*}
To remove the rational shifts, we split the range of $n$ into arithmetic progressions of difference $ck$. Thus our average becomes
\begin{align*}
    \E_{j\in[ck]}&\E_{n}e\left({{\frac{\tilde{q}k -\tilde{r}l}{ck(k\beta + l)}}(p(ckn + j)+g(ckn+j))+\frac{\tilde{r}}{ck}p(j)+\frac{\tilde{r}}{ck}g(ckn+j)}\right.\\
    &\qquad\qquad-\rem{\frac{ck}{ck(k\beta + l)}(p(ckn+j)+g(ckn+j))}t\\
    &\qquad\qquad\left.+\rem{{\frac{-cl}{ck(k\beta + l)}} (p(ckn+j)+g(ckn+j) )+\frac{1}{k}p(j)+\frac{1}{k} g(ckn+j)}s\right).
\end{align*}

$\bullet$ Case 1: $g(t)\to w$ as $t\to\infty$.
Once again, we assume that $g$ is eventually increasing, since the other case is similar.
Furthermore, we write $g(ckn+j)=w+o_n(1)$ and let $\delta>0$ be a small parameter.

We want to replace $\frac{1}{k}g(ckn+j)$ in the last fractional part with $\frac{1}{k}w$, since $g(t)\to w$. 
 \begin{multline*}
    {{\frac{-cl}{ck(k\beta + l)}} (p(ckn+j)+g(ckn+j) )+\frac{1}{k}p(j)+\frac{1}{k} g(ckn+j)}=\\
    {{\frac{-cl}{ck(k\beta + l)}} (p(ckn+j)+g(ckn+j) )+\frac{1}{k}p(j)+\frac{1}{k} w}+o_n(1)
\end{multline*}
Therefore, for $n$ sufficiently large (depending on $\delta$), we get the equality \begin{multline*}
    \rem{{\frac{-cl}{ck(k\beta + l)}} (p(ckn+j)+g(ckn+j) )+\frac{1}{k}p(j)+\frac{1}{k} g(ckn+j)}=\\
    \rem{{\frac{-cl}{ck(k\beta + l)}} (p(ckn+j)+g(ckn+j) )+\frac{1}{k}p(j)+\frac{1}{k} w}+o_n(1)
\end{multline*}whenever \begin{equation*}
    \delta \leq
    \rem{{\frac{-cl}{ck(k\beta + l)}} (p(ckn+j)+g(ckn+j) )+\frac{1}{k}p(j)+\frac{1}{k} w}\leq 1-\delta.
\end{equation*}
Since the sequence inside the last fractional part is equidistributed modulo 1, we conclude that this equality holds for at least $(1-4\delta)N$ values of $n\in [1,N]$, for every $N$ large enough.

Using the last approximation, we conclude that \begin{multline*}
    \E_{j\in[ck]}\E_{n\leq N}e\left({{\frac{\tilde{q}k -\tilde{r}l}{ck(k\beta + l)}}(p(ckn + j)+g(ckn+j))+\frac{\tilde{r}}{ck}(p(j)+g(ckn+j))}\right.\\
    -\rem{\frac{ck}{ck(k\beta + l)}(p(ckn+j)+g(ckn+j))}t\\
    \left.+\rem{{\frac{-cl}{ck(k\beta + l)}} (p(ckn+j)+g(ckn+j) )+\frac{1}{k}(p(j)+ g(ckn+j))}s\right)=\\
     \E_{j\in[ck]}e\left(\frac{\tilde{r}}{ck}(p(j)+w)\right)\E_{n\leq N}e\left({\frac{\tilde{q}k -\tilde{r}l}{ck(k\beta + l)}}(p(ckn + j)+g(ckn+j))\right.\\
    -\rem{\frac{ck}{ck(k\beta + l)}(p(ckn+j)+g(ckn+j))}t\\ \left.+\rem{{\frac{-cl}{ck(k\beta + l)}} (p(ckn+j)+g(ckn+j) )+\frac{1}{k}(p(j)+w)}s\right)+O(\delta).
\end{multline*}

By Weyl's equidistribution theorem for the sequence $\left(\frac{1}{ck(k\beta + l)}(p(ckn+j)+g(ckn+j))\right)_n$, the limit as $N\to\infty$ average above equals
\begin{align*}
    &\E_{j\in[ck]}e\left(\frac{\tilde{r}}{ck}(p(j)+w)\right)\int_0^1 e\brac{(\tilde{q}k - \tilde{r}l)x  - \rem{ck x}t + \rem{-clx + \frac{1}{k}(p(j)+w)}s}\; dx +O(\delta).
\end{align*}

To remove the fractional parts, we split the domain of integration into intervals of length $\frac{1}{|ckl|}$. We note that if $x\in\left(\frac{i}{|ckl|}, \frac{i+1}{|ckl|}\right)$, then 
\begin{align*}
    \floor{ck x} = \floor{\frac{i}{l}}\quad \textrm{and}\quad \floor{-clx + \frac{1}{k}(p(j)+w)} = \floor{\frac{-i-1+(p(j)+w)}{k}},
\end{align*}
hence our expression equals
\begin{align*}
    &\E_{j\in[ck]}e\left(\frac{\tilde{r}}{ck}(p(j)+w)\right)\sum_{i\in[0, |ckl|-1]}\int_{\frac{i}{|ckl|}}^{\frac{i+1}{|ckl|}} e\left((\tilde{q}k - \tilde{r}l)x  - \brac{ck x - \floor{\frac{i}{l}}}t\right.\\
    &\qquad\qquad\qquad\qquad\qquad\left.+ \brac{-clx + \frac{1}{k}(p(j)+w) - \floor{\frac{-i-1+(p(j)+w)}{k}}}s\right)\; dx\\
    &=\E_{j\in[ck]}e\left(\frac{\tilde{r}}{ck}(p(j)+w)\right)\E_{i\in[0, |ckl|-1]}\int_{0}^1 e\left((\tilde{q}k - \tilde{r}l-ckt-cls)\frac{x+i}{|ckl|}  + \floor{\frac{i}{l}}t\right.\\
    &\qquad\qquad\qquad\qquad\qquad\left.+ \brac{\frac{1}{k}(p(j)+w) - \floor{\frac{-i-1+(p(j)+w)}{k}}}s\right)\; dx\\
    &=1_{(\R\setminus ckl\cdot\Z)\cup\{0\}}(\tilde{q}k - \tilde{r}l-ckt-cls)\\
    &\qquad\qquad\qquad\E_{j\in[ck]}\E_{i\in[0, |ckl|-1]}
    e\left(\frac{\tilde{r}+cs}{ck}(p(j)+w) + \floor{\frac{i}{l}}t- \floor{\frac{-i-1+(p(j)+w)}{k}}s\right)\\
    &=1_{(\R\setminus kl\cdot\Z)\cup\{0\}}(k(q-t)-l(r+s))\\
    &\qquad\qquad\qquad\E_{j\in[ck]}\E_{i\in[0, |ckl|-1]}
    e\left(\frac{r+s}{k}(p(j)+w) + \floor{\frac{i}{l}}t- \floor{\frac{-i-1+(p(j)+w)}{k}}s\right).\\
\end{align*}
We then split 
\begin{align*}
    i = kl i_1 + k i_2 + i_3\quad \textrm{for}\quad i_1\in[0,c-1],\; i_2\in[0,l-1],\; i_3\in[0,k-1],
\end{align*}
so that the double exponential sum above becomes
\begin{multline*}
    \E_{i_1\in[0,c-1]}e((kt + ls)i_1)
    \E_{j\in[ck]}\E_{i_2\in[0, l-1]}\E_{i_3\in[0, k-1]}\\
    %&\qquad\qquad\qquad\qquad\qquad
    e\left(\frac{{r}+s}{k}(p(j)+w) + \floor{\frac{k i_2 + i_3}{l}}t + i_2 s- \floor{\frac{-i_3-1+(p(j)+w)}{k}}s\right).
\end{multline*}

The conclusion follows since $ \E_{i_1\in[0,c-1]}e((kt + ls)i_1)=1_{(\mathbb{Z}/c)\backslash\mathbb{Z}}(kt+ls)$.









$\bullet$ Case 2: $1\prec g(t)\ll \log t$.

Assume that the averages in the statement converge to zero. Then, for any $\delta>0$, Lemma \ref{L: dyadic intervals} and the calculations that have been carried out imply that the averages\begin{align*}
    \E_{j\in[ck]}&\E_{N\leq n<(1+\delta)N}e\left({{\frac{\tilde{q}k -\tilde{r}l}{ck(k\beta + l)}}(p(ckn + j)+g(ckn+j))+\frac{\tilde{r}}{ck}p(j)+\frac{\tilde{r}}{ck}g(ckn+j)}\right.\\
    &\qquad\qquad-\rem{\frac{ck}{ck(k\beta + l)}(p(ckn+j)+g(ckn+j))}t\\
    &\qquad\qquad\left.+\rem{{\frac{-cl}{ck(k\beta + l)}} (p(ckn+j)+g(ckn+j) )+\frac{1}{k}p(j)+\frac{1}{k} g(ckn+j)}s\right).
\end{align*}converge to $0$, as $N\to\infty$.

 Let $(N_j)_m$ be a subsequence that will be chosen later.
Since $g(t)=O(\log t)$ and $i,l$ are bounded, the mean value theorem implies that \begin{equation}\label{E: approximation 416}
    g(ckn+j)=g(ckN_j)+O(\log (1+\delta)),\ \text{for all } n\in[N_j,(1+\delta)N_j).
\end{equation}
Now, we use this approximation to replace all instances of $ g(ckn+j)$ in the previous equation by $ g(ckN_j)$.


In view of \eqref{E: approximation 416}, we have the equality \begin{equation*}
    \rem{\frac{ck}{ck(k\beta + l)}(p(ckn+j)+g(ckn+j))}=\rem{\frac{ck}{ck(k\beta + l)}(p(ckn+j)+g(ckN_j))}+O(\delta)
\end{equation*}whenever $$\rem{\frac{ck}{ck(k\beta + l)}(p(ckn+j)+g(ckN_j))}\in [O(\log(1+\delta)),1-O(\log(1+\delta))].$$
However, since the implicit polynomial has at least one irrational coefficient, the last inclusion holds for an at least $(1-O(\delta))$ proportion of $n\in [N_j,(1+\delta)N_j)$, where $N_j$ is assumed to be sufficiently large.

Similarly, we have \begin{multline*}
    \rem{{\frac{-cl}{ck(k\beta + l)}} (p(ckn+j)+g(ckn+j) )+\frac{1}{k}p(j)+\frac{1}{k} g(ckn+j)}=\\
    \rem{{\frac{-cl}{ck(k\beta + l)}} (p(ckn+j)+g(ckN_j) )+\frac{1}{k}p(j)+\frac{1}{k} g(ckN_j)}
\end{multline*}for an at least $(1-O(\delta))$ proportion of $n\in [N_j,(1+\delta)N_j)$.

Combining everything together, we deduce that 
\begin{align*}
    &\E_{j\in[ck]}\E_{N_j\leq n<(1+\delta)N_j} e\left({{\frac{\tilde{q}k -\tilde{r}l}{ck(k\beta + l)}}(p(ckn + j)+g(ckn+j))+\frac{\tilde{r}}{ck}p(j)+\frac{\tilde{r}}{ck}g(ckn+j)}\right.\\
    &\qquad\qquad\qquad\qquad\qquad\qquad\qquad\qquad-\rem{\frac{ck}{ck(k\beta + l)}(p(ckn+j)+g(ckn+j))}t\\
    &\qquad\qquad\left.+\rem{{\frac{-cl}{ck(k\beta + l)}} (p(ckn+j)+g(ckn+j) )+\frac{1}{k}p(j)+\frac{1}{k} g(ckn+j)}s\right)=\\
    &\E_{j\in[ck]}e\left(\frac{\tilde{r}}{ck}(p(j)+g(ckN_j))\right)\E_{N_j\leq n<(1+\delta)N_j}e\left({\frac{\tilde{q}k -\tilde{r}l}{ck(k\beta + l)}}(p(ckn + j)+g(ckN_j)) \right.\\
    &\qquad\qquad\qquad\qquad\qquad\qquad\qquad\qquad-\rem{\frac{ck}{ck(k\beta + l)}(p(ckn+j)+g(ckN_j))}t\\
    &\qquad\qquad\qquad\left.+\rem{{\frac{-cl}{ck(k\beta + l)}} (p(ckn+j)+g(ckN_j))+\frac{1}{k}p(j)+\frac{1}{k} g(ckN_j)}s\right).
\end{align*}

Using Lemma \ref{L: uniform weyl for Riemann integrable} the sequence $\left(\frac{1}{ck(k\beta + l)}p(ckn+j)\right)$, we get the asymptotic
\begin{multline*}
    \E_{N_j\leq n<(1+\delta)N_j}e\left({\frac{\tilde{q}k -\tilde{r}l}{ck(k\beta + l)}}(p(ckn + j)+g(ckN_j)) \right.\\
    -\rem{\frac{ck}{ck(k\beta + l)}(p(ckn+j)+g(ckN_j))}t\\
    \qquad\qquad\qquad\left.+\rem{{\frac{-cl}{ck(k\beta + l)}} (p(ckn+j)+g(ckN_j))+\frac{1}{k}p(j)+\frac{1}{k} g(ckN_j)}s\right)=\\
\int_0^1 e\brac{(\tilde{q}k - \tilde{r}l)x  - \rem{ck x}t + \rem{-clx + \frac{1}{k}(p(j)+g(kcN_j))}s}\; dx +o_j(1).
\end{multline*}
We compute the last integral as previously to deduce that \begin{multline*}
   1_{(\R\setminus kl\cdot\Z)\cup\{0\}}(k(q-t)-l(r+s)) \E_{i_1\in[0,c-1]}e((kt + ls)i_1)
    \E_{j\in[ck]}\E_{i_2\in[0, l-1]}\E_{i_3\in[0, k-1]}\\
    %&\qquad\qquad\qquad\qquad\qquad
    e\left(\frac{{r}+s}{k}(p(j)+g(kcN_j)) + \floor{\frac{k i_2 + i_3}{l}}t + i_2 s- \floor{\frac{-i_3-1+(p(j)+g(kcN_j))}{k}}s\right)
\end{multline*}is $O(\delta)$.

If $1_{(\R\setminus kl\cdot\Z)\cup\{0\}}(k(q-t)-l(r+s))=0$ is zero, then we get condition (i).
Assume that $1_{(\R\setminus kl\cdot\Z)\cup\{0\}}(k(q-t)-l(r+s))\neq 0$.
Now, we choose the sequence $N_j$ so that $\{\frac{g(kcN_j)}{k}\}\in [\e,2\e]$ where $\e$ is small in terms of $1/k$.
This can be done because the sequence $\left(g(kcn)/k\right)_n$ is dense modulo 1.
Then, we have \begin{equation*}
     \floor{\frac{-i_3-1+(p(j)+g(kcN_j))}{k}}= \floor{\frac{-i_3-1+p(j)}{k}}+\floor{\frac{g(kcN_j)}{k}}.
\end{equation*}Indeed, this holds since \begin{equation*}
     \{\frac{-i_3-1+p(j)}{k}\}+\{\frac{g(kcN_j)}{k}\}<
    1,
\end{equation*}which is true because the first fractional part is at most $(k-1)/k$ and the second is at most $2\e$.
We rewrite the last averages as \begin{multline*}
    e\left(\frac{r+s}{k}g(kcN_j)-\floor{\frac{g(kcN_j)}{k}}s\right) \E_{i_1\in[0,c-1]}e((kt + ls)i_1) \E_{j\in[ck]}\E_{i_2\in[0, l-1]}\E_{i_3\in[0, k-1]}\\
    %&\qquad\qquad\qquad\qquad\qquad
    e\left(\frac{{r}+s}{k}p(j)+ + \floor{\frac{k i_2 + i_3}{l}}t + i_2 s- \floor{\frac{-i_3-1+p(j)}{k}}s\right)
\end{multline*}Since this $O(\delta)$, we deduce that \begin{multline*}
    \E_{i_1\in[0,c-1]}e((kt + ls)i_1) \E_{j\in[ck]}\E_{i_2\in[0, l-1]}\E_{i_3\in[0, k-1]}\\
    %&\qquad\qquad\qquad\qquad\qquad
    e\left(\frac{{r}+s}{k}p(j)+ + \floor{\frac{k i_2 + i_3}{l}}t + i_2 s- \floor{\frac{-i_3-1+p(j)}{k}}s\right)=O(\delta)
\end{multline*}
Sending $\delta \to 0$, we deduce that the last expression is equal to zero, which implies that one of conditions (ii) or (iii) holds. The proof is finished.

\end{proof}
\end{comment}

In the last remaining case, the computations become so complicated that the conclusion is less explicit than in the previous results. In what follows, $\sgn(x) = 1_{x\geq 0}-1_{x\leq 0}$ is the usual sign function.

\begin{lemma}[Case $k,l\neq 0$]\label{L: (k beta + l) a  = p + g}
Let $c\in\N$, $k, l\in\Z\backslash\{0\}$ {with $l>0$}, and $q,r\in\frac{1}{c}\Z$, and suppose that $a = \frac{p-ku}{k\beta + l}$ for some nonconstant $p\in\Z[x]{+\R}$. For $j\in[ck]$ and $j'~\in~[0,|ckl|~-~1]$, define
\begin{align*}
  %w &:= c(k(q-t)-l(r+s))\\
  w &:= {\frac{k(q-t)-l(r+s)}{|kl|}}\\
  x_{j,j'} &:= |k|\rem{\frac{p(j)-\sgn(k)j'}{k}}\\
  b_{j, j'} &:=\floor{\frac{\sgn(k)(j' +1_{k<0})}{l}}t-\floor{\frac{p(j)-\sgn(k)j'}{k}}s. 
\end{align*}
Then
%Then there exist $w, b_{j,j'}\in\Z t + \Z s + \R$ (for $j\in[ck], j'\in[ckl]$) such that
\begin{equation*}
    \lim_{N\to\infty}\E_{n\in[N]} e(qa(n)+\beta r a(n)-\{ a(n)\}t +\{\beta a(n)+u\}s) = 0
\end{equation*}
 if and only if one of the followings holds: 
      \begin{enumerate}
    \item $w=0$ and 
    \begin{align*}
    \E_{j\in[ck]}e\left(\frac{{r+s}}{k}p(j)\right)\E_{j'\in[0, |ckl|-1]}e\left(b_{j,j'}\right)(e(s) + \min(x_{j,j'},1)(1-e(s))) = 0;
\end{align*}
    \item $w\neq 0$ and 
\begin{multline*}
\E_{j\in[ck]}e\left(\frac{{r+s}}{k}p(j)\right)\E_{j'\in[0, |ckl|-1]}e\left({j'w}+b_{j,j'}\right)\\
    \brac{e(s)\Bigbrac{e({w})-1}+(1-e(s))\Bigbrac{e({w\min(x_{j,j'},1)})-1}} = 0.
    % \E_{j\in[ck]}e\left(\frac{{r+s}}{k}p(j)\right)\E_{j'\in[0, |ckl|-1]}e\left(\frac{j'}{|ckl|}w+b_{j,j'}\right)\\
    % \brac{e(s)\Bigbrac{e\Bigbrac{\frac{w}{|ckl|}}-1}+(1-e(s))\Bigbrac{e\Bigbrac{w\min(x_{j,j'},\frac{1}{|ckl|})}-1}} = 0.
\end{multline*}    
\end{enumerate}
%Moreover, if $g$ converges, then the converse is also true, i.e. if any of the three conditions holds, then the original average converges to 0. 
\end{lemma}
{We observe that there is no loss of generality in assuming that $l>0$ for if $l<0$, then we can reduce to the case $l>0$ by changing $l\mapsto -l, k\mapsto -k, p\mapsto -p$.} 
\begin{proof}

Let $\tilde{q} = qc, \tilde{r} = rc\in\Z$, and $\tilde{p}(n) = p(n)-ku$. We can recast our average as
\begin{align*}
    \lim_{N\to\infty}\E_{n\in[N]} e\brac{{\frac{(\tilde{q} +\beta\tilde{r})\tilde{p}(n)}{c(k\beta + l)}}-\rem{\frac{\tilde{p}(n)}{k\beta + l}}t +\rem{\frac{\beta \tilde{p}(n)}{k\beta + l}+u}s}.
\end{align*}
Using the identity
\begin{align*}
    \frac{x\beta + y}{k\beta + l} = \frac{ky - lx}{k(k\beta + l)} + \frac{x}{k},
\end{align*}
we want to bring all the fractions above to the form $\Q + \frac{\Z}{ck(k\beta + l)}$ and then apply Theorem~\ref{T: Boshernitzan} and Weyl's equidistribution theorem. Indeed, our average equals
\begin{multline*}
    \lim_{N\to\infty}\E_{n\in[N]} e\left({\brac{\frac{\tilde{q}k -\tilde{r}l}{ck(k\beta + l)}+\frac{\tilde{r}}{ck}}\tilde{p}(n)}-\rem{\frac{ck}{ck(k\beta + l)}\tilde{p}(n)}t\right.\\
    \left.+\rem{\brac{\frac{-cl}{ck(k\beta + l)}+\frac{1}{k}} \tilde{p}(n)+u}s\right).
\end{multline*}
To remove the rational shifts, we split the range of $n$ into arithmetic progressions of difference $ck$. Thus our average becomes
\begin{multline*}
    \E_{j\in[ck]} e\brac{\frac{\tilde{r}\tilde p(j)}{ck}}  \lim_{N\to\infty}\E_{n\in[N]}e\left({{\frac{\tilde{q}k -\tilde{r}l}{ck(k\beta + l)}}\tilde p(ckn + j)}\right.\\
    \left. -\rem{\frac{ck}{ck(k\beta + l)}\tilde p(ckn+j)}t+\rem{{\frac{-cl}{ck(k\beta + l)}} \tilde p(ckn+j)+\frac{\tilde p(j)}{k}+u}s\right).
\end{multline*}
By Theorem \ref{T: Boshernitzan}, the sequence $\brac{\rem{\frac{\tilde p(ckn+j)}{ck(k\beta + l)}}}_n$ is equidistributed on $[0,1)$, and hence our average  equals
\begin{align*}
    \E_{j\in[ck]}e\left(\frac{\tilde{r}\tilde p(j)}{ck}\right)\int_0^1 e\brac{(\tilde{q}k - \tilde{r}l)x  - \rem{ck x}t + \rem{-clx +\frac{\tilde p(j)}{k}+u}s}\; dx.
\end{align*}
Recalling that $\tilde{r} = cr$ and $\tilde{p}(j) = p(j) - ku$, we can slightly simplify this expression to
\begin{align*}
    e(-ru)\E_{j\in[ck]}e\left(\frac{{r}p(j)}{k}\right)\int_0^1 e\brac{(\tilde{q}k - \tilde{r}l)x  - \rem{ck x}t + \rem{-clx +\frac{ p(j)}{k}}s}\; dx.
\end{align*}

To remove the fractional parts, we split the domain of integration into intervals of length $\frac{1}{|ckl|}$. 
%\BK{(Before, this part of the argument implicitly assumed that $k,l,kl>0$, which need not be the case. I had to make changes to account for the missing cases, which unfortunately make the expressions even worse. Please check correctness)} \AK{I checked it; more people should do it too}
%(it is to simplify the following discussion where we use the assumption that $kl > 0$). 
% We note that if $x\in\left(\frac{i'}{|ckl|}, \frac{i'+\sgn(kl)}{|ckl|}\right)$ (where $\sgn(x) = 1_{x>0}-1_{x<0}$ is the sign function), then 
% \begin{align*}
%     \floor{ck x} = \floor{\frac{i' \BK{-1_{l<0}}}{l}}\quad \textrm{and}\quad \floor{-clx + \frac{p(i)}{k}} = \floor{\frac{-i'\BK{-1_{k>0}}+p(i)}{k}},
% \end{align*}
%From now on, we assume that $l>0$ (if $l<0$, then by changing $l\mapsto -l, k\mapsto -k, p\mapsto -p$, we can reduce to the previous case). 
If $x\in\left(\frac{j'}{|ckl|}, \frac{j'+1}{|ckl|}\right)$, then 
\begin{align*}
    \floor{ck x} = \floor{\frac{\sgn(k)(j' +1_{k<0})}{l}}. %\quad \textrm{and}\quad \floor{-clx + \frac{p(i)}{k}} = \floor{\frac{-\sgn(k)(i'+1_{l>0})+p(i)}{k}},
\end{align*}
On the other hand, let $x_{j,j'} := |k|\rem{\frac{p(j)-\sgn(k)j'}{k}}$.  Then an easy computation shows that
\begin{align*}
     \floor{-clx + \frac{p(j)}{k}} = \begin{cases}
         \floor{\frac{p(j)-\sgn(k)j'}{k}},\; &x\in \left(\frac{j'}{|ckl|}, \frac{j'+\min(x_{j,j'},1)}{|ckl|} \right]=:C_{j,j'}\\
         \floor{\frac{p(j)-\sgn(k)j'}{k}}-1,\; &x\in \left(\frac{j'+\min(x_{j,j'},1)}{|ckl|},\frac{j'+1}{|ckl|}\right) =: C'_{j,j'}.
     \end{cases}
\end{align*}
% \AK{Indeed, if $x\in C_{j,j'},$ then 
% \[-clx+\frac{p(j)}{k}\in \Big[\frac{p(j)-\sgn(k)j'}{k}-\min\Big\{clx_{j,j'}, \frac{\sgn(k)}{k}\Big\},\frac{p(j)-\sgn(k)j'}{k}\Big),\]
% so the integer part of the number equals \[\floor{\frac{p(j)-j'\sgn(k)}{k}}.\] To see this take the minimum to be $x_{j,j'}$; the other case will give something larger (but of the same integer part). On the other hand, if $x\in C'_{j,j'},$ then 
% \[-clx+\frac{p(j)}{k}\in \Big(\frac{p(j)-(j'+1)\sgn(k)}{k},\frac{p(j)-\sgn(k)j'}{k}-\min\Big\{clx_{j,j'},\frac{\sgn(k)}{k}\Big\}\Big),\]
% so the integer part is equal to 
% \[\floor{\frac{p(j)-j'\sgn(k)}{k}}-1,\] as only the $x_{j,j'}$ can happen as minimum.
% } \BK{(Shall we keep Andreas' extra explanation in the paper or just use it for our convenience while proofreading?

% I think I'd be inclined to keep the extra explanation out after we all verify that the computation is correct. Since you were able to replicate my reasoning, we can expect this of the reader as well, and it makes the proof quite lengthy.)} \AK{I would definitely remove it. Whatever has to do with computations which can be verified (relatively easily) IMO should be removed; I just wrote them for ourselves}

% On $x\in \left(\frac{i'}{|ckl|}, \frac{i'}{|ckl|} + \min\bigbrac{x_{i,i'}, \frac{1}{|ckl|}}\right]$, we have
% \begin{align*}
%     \floor{-clx + \frac{p(i)}{k}} =    \floor{\frac{-\sgn(k)i'+p(i)}{k}}
% \end{align*}

Now, letting
\begin{align*}
  %w &:= \tilde{q}k - \tilde{r}l-ckt-cls = c(k(q-t)-l(r+s))\\
  w &:= {\frac{\tilde{q}k - \tilde{r}l-ckt-cls}{|ckl|} = \frac{k(q-t)-l(r+s))}{|kl|}}\\
  \textrm{and} \quad b_{j, j'} &:=\floor{\frac{\sgn(k)(j' +1_{k<0})}{l}}t-\floor{\frac{p(j)-\sgn(k)j'}{k}}s  
\end{align*}
 %$w := \tilde{q}k - \tilde{r}l-ckt-cls = c(k(q-t)-l(r+s))$ and $b_{j, j'}:=\floor{\frac{\sgn(k)(j' +1_{k<0})}{l}}t-\floor{\frac{p(j)-\sgn(k)j'}{k}}s$
 to simplify the discussion, our expression equals $e(-ru)$ times  
\begin{align*}
    & \E_{j\in[ck]}e\left(\frac{{r+s}}{k}p(j)\right)\sum_{j'\in[0, |ckl|-1]}\int_{\frac{j'}{|ckl|}}^{\frac{j'+1}{|ckl|}} e\left({|ckl|}wx  +\floor{ckx}t-\floor{-clx + \frac{p(j)}{k}}s\right)\; dx\\
    &=\E_{j\in[ck]}e\left(\frac{{r+s}}{k}p(j)\right)\sum_{j'\in[0, |ckl|-1]}\left(\int_{C_{j,j'}} e({|ckl|}wx  +b_{j,j'})\; dx+\int_{C'_{j,j'}} e\left({|ckl|}wx  +b_{j,j'}+s\right)\; dx\right)\\
    &=\E_{j\in[ck]}e\left(\frac{{r+s}}{k}p(j)\right)\sum_{j'\in[0, |ckl|-1]}\left(e(s)\int_{\frac{j'}{|ckl|}}^{\frac{j'+1}{|ckl|}} e({|ckl|}wx  +b_{j,j'})\; dx\right.\\
    &\qquad\qquad\qquad\qquad\qquad\qquad\qquad\qquad\qquad\qquad\left.+(1-e(s))\int_{C_{j,j'}} e\left({|ckl|}wx  +b_{j,j'}\right)\; dx\right).
    \end{align*}
Performing the change of variables, this equals
\begin{align*}
    &\E_{j\in[ck]}e\left(\frac{{r+s}}{k}p(j)\right)\E_{j'\in[0, |ckl|-1]}e\left({j'w}+b_{j,j'}\right)\\
    &\qquad\qquad\qquad\qquad\left(e(s)\int_{0}^1 e\left({w}x\right)\; dx+(1-e(s))\int_{0}^{\min(x_{j,j'},1)} e\left({w}x\right)\; dx\right).
\end{align*}
If $w=0$, the expression above equals
\begin{align*}
    \E_{j\in[ck]}e\left(\frac{{r+s}}{k}p(j)\right)\E_{j'\in[0, |ckl|-1]}e\left(b_{j,j'}\right)(e(s) + \min(x_{j,j'},1)(1-e(s))),
\end{align*}
otherwise we obtain
\begin{multline*}
    \E_{j\in[ck]}e\left(\frac{{r+s}}{k}p(j)\right)\E_{j'\in[0, |ckl|-1]}e\left({j'w}+b_{j,j'}\right)\\
    \frac{1}{2\pi i  {w}}\brac{e(s)\Bigbrac{e({w})-1}+(1-e(s))\Bigbrac{e({w\min(x_{j,j'},1)})-1}}.
\end{multline*}
This completes the proof.
\end{proof}

\subsection{Structure of $E_3$}


Combining the content of the previous lemmas, we move on to prove Proposition \ref{P: E_3}, i.e. we describe the set $E_3$ consisting of those points $(t,s)\in[0,1)^2$ for which $A(t,s) =0$ while $B(t,s)\neq 0$. 

     \begin{proof}[Proof of Proposition \ref{P: E_3}]
     Our objective can be restated as follows: we want to show that $E_3$ is countable, and that it is contained in $\{\brac{\beta \gamma, \gamma}: \gamma\in\R\}$ mod $\Z^2$.
    
     %and that modulo $\Z^2$, it is contained in $\{\brac{\beta \gamma_i, \gamma_i}: i\in\N \} \cup {E}_3'$ for some nonzero $\gamma_1, \gamma_2, \ldots\in\R$. The last claim $\sigma_f(E_3') = 0$ will then follow from the countability of $E_3$ and Lemma \ref{L: E_3'}.
     
     Let $(t,s)\in E_3$. By Lemma \ref{L: Q-independence} and Corollary \ref{C: a = n}, we have 
     \begin{align}\label{E: admissible (t,s)}
     t-\beta s = q + \beta r\quad \textrm{for\; some}\quad q,r\in\Z\quad \textrm{satisfying}\quad r+s\notin\Z\setminus\{0\}.
     \end{align}
     {In particular, the value $c$ appearing in various lemmas throughout Section \ref{SS: Q-dependent} is $1$.} Then
     \begin{align*}
         (t,s) = (\beta(r+s), r+s) + (q,-r),
     \end{align*}
     from which we deduce that $(t,s)$ lies in
     \begin{align}\label{E: hyperplane}
        \{(\beta \gamma, \gamma):\; \gamma\in\R\}\mod \Z^2.
     \end{align}
     It remains to show that $E_3$ is countable.
     %Our goal then is to show that the set $E_3$ is countable, and that $\{(\beta y, y): y\in\R\}$ can be replaced by countably many elements of the form $(\beta \gamma, \gamma)$.
     To obtain this last claim, we split the proof into the following four cases whose exponential sums we analyzed separately in the previous section, and each of which requires a slightly different treatment now:
     \begin{enumerate}
    \item\label{i: a 1} $a(x)$ does not equal
    \begin{align}\label{E: special form of a_2}
        \frac{p(x)-ku}{k\beta + l} \quad \textrm{for\; any}\;\; (k, l)\in\Z^2\setminus{\{(0,0)\}},\; p\in\Z[x]{+\R}; %\; g \ll \log;
    \end{align}
    \item\label{i: a 2} $a(x)$ equals \eqref{E: special form of a_2} with $k\neq 0, l=0$;
    \item\label{i: a 3} $a(x)$ equals \eqref{E: special form of a_2} with $k=0, l\neq 0$;
    \item\label{i: a 4} $a(x)$ equals \eqref{E: special form of a_2} with $k, l\neq 0$.
\end{enumerate}
     %The gist of these arguments is however the same. First, we want to obtain a fairly precise description of $E_3''$, the ``structured'' part of $E_3$. Second, we want to show that $E_3$ is countable; once this is established, the claim $\sigma_f(E_3') = 0$ will follow from Lemma \ref{L: E_3'}.
     
     \smallskip
     \textbf{Case \eqref{i: a 1}.}
     \smallskip

    Combining the conclusions of Lemma \ref{L: a, beta a eq} and Corollary \ref{C: a = n}, the set $E_3$ consists of those pairs $(t,s)\in[0,1)^2$ satisfying \eqref{E: admissible (t,s)} for which additionally $q-t\in\mathbb{Z}\backslash\{0\}$.
    %, $q, r\in\Z$, and $r+s$ is either 0 or not an integer. 
    Hence $t\in\Z\setminus\{q\}$ and $s  = \frac{1}{\beta}(t-q)-r\in \frac{1}{\beta}\Z\setminus\{0\} + \Z$, and so the collection of pairs $(t,s)\in E_3$ forms a countable subset of \eqref{E: hyperplane}.
    
    %Conversely, suppose that $t\in\Z$ and $s = \frac{1}{\beta}r_1 + r_2$ for $r_1, r_2\in\Z$ with $r_1\neq 0$. Then $t-\beta s = (t-r_1)-\beta r_2$, and since $t\in\Z$, we have $q = t-r_1$, $r = -r_2$. Clearly, $q-t = -r_1\in\Z\setminus\{0\}$, $q\in\Z$ and $r\in\Z$, implying (ii), the inverse of (iii) and the inverse of (iv) respectively. We also get the inverse of (i) since as argued above, (i) and (ii) cannot hold simultaneously.

   %      . If this is not the case, then by Lemma \ref{L: ii holds} and Corollary \ref{C: a = n}, we have $A(t,s) = 0$ while $B(t,s) \neq 0$ if and only if $(t,s)\in\langle (0,1/\beta)\rangle\setminus\{(0,0)\}$, and the claim follows.

    \smallskip
    \textbf{Case \eqref{i: a 2}.}
    \smallskip

    In this case, Lemma \ref{L: beta a  = p + g} says that on top of \eqref{E: admissible (t,s)}, each $(t,s)\in E_3$ satisfies one of the two conditions:
       %By Lemma \ref{L: beta a  = p + g} and Corollary \ref{C: a = n}, we have $q,r\in\Z$, $r+s\notin\Z\setminus\{0\}$, and at least one of the two conditions happens: 
       $q-t\in\Z\setminus\{0\}$ or $\E\limits_{i\in[{k}]}e\brac{\frac{r p(i)}{k}+ \rem{\frac{p(i)}{k}}s}=0$. If $q-t\in\Z\setminus\{0\}$, then  we have countably many choices of $(t,s)$ inside \eqref{E: hyperplane} just like in the previous case.
       %necessarily $t\in\Z$ (since $q-t, q\in\Z$) and $$s = \frac{1}{\beta}(t-q) - r\in \frac{1}{\beta}\Z\setminus\{0\}+\Z,$$ and so $(t,s)\in \langle (0,1/\beta)\rangle_{[0,1)^2}$.
        If instead $\E_{i\in[k]}e\brac{\frac{r p(i)}{k}+ \rem{\frac{p(i)}{k}}s}=0$, then $r\in\Z$ and rudimentary manipulations imply that
        \begin{align}\label{E: reduced exp sum}
            \E_{i\in[k]} e\brac{(r+s)\rem{\frac{p(i)}{k}}} = 0.    
        \end{align}
         Consider therefore the remaining set
        \begin{align*}
            \tilde{E}_3 = \left\{(t,s)\in E_3:\; t = q+ \beta (r+s),\; q,r\in\Z,\; r+s\notin\Z\setminus\{0\},\; \eqref{E: reduced exp sum}\; \textrm{holds}\right\}.
            %\tilde{E}_3 = \left\{(t,s)\in E_3:\; t - \beta s = q + \beta r,\; q,r\in\Z,\; r+s\in\Z\setminus\{0\},\; \E_{i\in[k]} e((r+s)\{p(i)\}) = 0\right\}.
        \end{align*}
        The crucial observation is that the exponential sum \eqref{E: reduced exp sum} can be realized as a polynomial of bounded degree. Indeed, for each $j\in\{0,\ldots, k-1\}$, let
        $$c_j =\frac{1}{k}\cdot \abs{\rem{i\in[k]:\; \rem{\frac{p(i)}{k}} = \frac{j+\{p(0)\}}{k}}}$$ and $Q(z) = \sum\limits_{j=0}^{k-1} c_j z^j$ be the rational polynomial satisfying
        \begin{align*}
            \E_{i\in[k]} e\brac{\alpha\rem{\frac{p(i)}{k}}} = Q(e(\alpha/k)){e(\alpha \{p(0)\}/k)}
        \end{align*}
        for every $\alpha\in\T$. By the fundamental theorem of algebra, there exist at most $k-1$ values $\alpha\in\T$ for which $Q(e(\alpha/k)) = 0$. In our case, $\alpha = r+ s$ and $t = q +\beta (r+s)$, and so we deduce that $\tilde{E}_3$ is countable. Combining both subcases, we deduce that $E_3$ is a countable subset of \eqref{E: hyperplane}.   



\smallskip
\textbf{Case \eqref{i: a 3}.}
\smallskip
         
         %In this case, Lemma \ref{L: a = p + g} and Corollary \ref{C: a = n} imply that $A(t,s) = 0$ while $B(t,s)\neq 0$ if only if $q,r\in\Z$, $r+s\notin\Z\setminus\{0\}$ and 
         In this case, Lemma \ref{L: a = p + g} guarantees that $(t,s)\in E_3$ satisfies 
         \begin{align}\label{E: reduced exp sum 3}
             \E_{i\in[{l}]}e\brac{(q-t)\rem{\frac{p(i)}{l}}} =0
         \end{align}
         on top of \eqref{E: admissible (t,s)}. The countability of $E_3$ then follows from a similar argument as in the previous case.
        %Arguing similarly as in the previous case, we deduce that the set $E_3$ is countable.% and $\sigma_f({E}_3') = 0$. 
        

        \smallskip
        \textbf{Case \eqref{i: a 4}}
        \smallskip

        While in the previous two cases we deduced the countability of $E_3$ from the fundamental theorem of algebra, this time we will rely on basic properties of holomorphic functions.
        Let
        \begin{align*}
            w:= {\frac{k(q-t)-l(r+s)}{|kl|} = \frac{-(k\beta +l)(r+s)}{|kl|},}
        \end{align*}
        and recall that every $(t,s)\in E_3$ must satisfy one of the two conditions from Lemma~\ref{L: (k beta + l) a  = p + g} on top of \eqref{E: admissible (t,s)}.
        %Let $(t,s)\in E_3$; by Corollary \ref{C: a = n}, we have $t-\beta s = q +\beta r$ for $q,r\in\Z$ and $r+s\in(\R\backslash\Z)\cup\{0\}$. 
        If $w=0$, then necessarily $s=-r\in\Z$, which is precluded by \eqref{E: admissible (t,s)}.
        %and there are countably many such values $s$. Since we also have countably many choices of $q\in\Z$, and $t$ is uniquely determined by $q,r,s$, we have countably many choices of $(t,s)$ for which $w=0$.
        Therefore $w\neq 0$, in which case the second condition from Lemma \ref{L: (k beta + l) a  = p + g} holds, i.e. %$w:=-(k\beta + l)(r+s)\notin \Z$ and 
\begin{multline}\label{E: tricky exp sum}
    \E_{j\in[k]}e\left(\frac{{r+s}}{k}p(j)\right)\E_{j'\in[0, |kl|-1]}e\left(-\frac{j'}{|kl|}(k\beta + l)(r+s)+b_{j,j'}\right)\\
    \brac{e(s)\Bigbrac{e({w})-1}+(1-e(s))\Bigbrac{e({w\min(x_{j,j'},1)})-1}}
    %\brac{e(s)\Bigbrac{e\Bigbrac{\frac{w}{|kl|}}-1}+(1-e(s))\Bigbrac{e\Bigbrac{w\min(x_{j,j'},\frac{1}{|kl|})}-1}} = 0
\end{multline}    
is 0
for 
%some $x_{j,j'}\in\R$ (independent of $t, s$) 
\begin{align*}
    x_{j,j'} &=|k|\left\{\frac{p(j)-\sgn(k)j'}{k}\right\}\\
    \textrm{and}\quad b_{j,j'} &= \floor{\frac{\sgn(k)(j' +1_{k<0})}{l}}(q+\beta(r+s)) -\floor{\frac{p(j)-\sgn(k)j'}{k}}s. 
    %b_{1,j,j'}(q+\beta r)+(b_{1,j,j'}\beta + b_{2,j,j'})s
\end{align*}


% and $b_{j,j'}$'s, which by Lemma \ref{L: (k beta + l) a  = p + g} take the form
% $$b_{j,j'} = b_{1,j,j'}(q+\beta r)+(b_{1,j,j'}\beta + b_{2,j,j'})s$$ for 
% given by
% \begin{align}\label{E: b_{i,j,j'}}
%     b_{1,j,j'}:=\floor{\frac{\sgn(k)(j' +1_{k<0})}{l}}\quad \textrm{and}\quad    b_{2,j,j'}:=-\floor{\frac{p(j)-\sgn(k)j'}{k}}.
% \end{align}
%$b_{j,j'}\in (q+\beta r)\Z + s(\Z+\beta\Z)$.
For each fixed $q,r\in\mathbb{Z}$, the average \eqref{E: tricky exp sum} is a finite linear combination of functions of the form $s\mapsto e(v s)$ for $v\in\R$. 
Each such function is holomorphic, and therefore the average in \eqref{E: tricky exp sum} is also holomorphic as a function of $s$. Hence, one of the two things may happen:
\begin{enumerate}
    \item \eqref{E: tricky exp sum} is identically 0 for all $s\in\C$;
    \item \eqref{E: tricky exp sum} has at most countably many zeros.
\end{enumerate}
If we are able to exclude the first case, then the second case immediately implies the countability of $E_3$.

\ 

We show that the former option is not possible, i.e. the average in \eqref{E: tricky exp sum} cannot vanish for all $s\in\C$, by identifying the value $s=-|kl|-r$ for which \eqref{E: tricky exp sum} is nonzero.
For this value $s$, we can simplify various expressions appearing in \eqref{E: tricky exp sum} to
\begin{align*}
        \frac{r+s}{k}p(j) &\equiv \frac{r+s}{k}p(0)\!\!\!\mod 1,\\
        -\frac{j'}{|kl|}(k\beta + l)(r+s) &\equiv \frac{\sgn(k)j'}{l}|kl|\beta\!\!\!\mod 1,\\
        b_{j,j'}&\equiv -\floor{\frac{\sgn(k)(j' +1_{k<0})}{l}}|kl|\beta\!\!\!\mod 1\\
        w &= k\beta + l\notin\Z.
        %\frac{w}{|kl|} &= k\beta + l\notin\Z.
\end{align*}
With these simplifications, \eqref{E: tricky exp sum} reduces to
\begin{align*}
    %\left(e(k\beta + l)-1\right)
    \E_{j'\in [0,|kl|-1]}
    e\brac{\Bigbrac{\frac{\sgn(k)j'}{l}-\floor{\frac{\sgn(k)(j' +1_{k<0})}{l}}}|kl|\beta} = 0.
\end{align*}
The left-hand side then equals  $e\Bigbrac{-\frac{\sgn(k)1_{k<0}}{l}|kl|\beta}$ times  
\begin{align*}
    %\left(e(k\beta + l)-1\right)
    \E_{j'\in [0,|kl|-1]}
    e\brac{\Big\{\frac{\sgn(k)(j'+1_{k<0})}{l}\Big\}|kl|\beta} &= \E_{j'\in [0,l-1]}
    e\brac{\Big\{\frac{j'}{l}\Big\}|kl|\beta}\\ &= \E_{j'\in [0,l-1]}e(j'|k|\beta) =\frac{e(|kl|\beta)-1}{l(e(|k|\beta)-1)},
    %e\brac{\Big\{\frac{\sgn(k)(j'+1_{k<0})}{l}\Big\}|kl|\beta}\\ &= \E_{j'\in [0,l-1]}e(j'k\beta) =\frac{e(kl\beta)-1}{l(e(k\beta)-1)},
\end{align*}
which is nonzero since $\beta\notin\Q$.
 \end{proof}



\subsection{Proof of Proposition \ref{P: beta irrational}}\label{SS: seminorm comparison base case}
We now bring all the pieces together in order to prove Proposition \ref{P: beta irrational}.

For a non-empty subset $H\subseteq[0,1)^2$, set
$$\CK(H) := \langle f\in L^2(\mu):\; T_1 f = e(t)f,\; T_2f = e(s)f\;\; \textrm{for\; some}\;\; (t,s)\in H\rangle,$$
and set $\CK(\emptyset):=\{0\}$. Let $P_H f$ be the orthogonal projection of $f$ onto $\CK(H)$. If $H=\{(t,s)\}$, then we also write
\begin{align*}
    P_{(t,s)}f =P_{\{(t,s)\}}f= \lim_{N\to\infty}\E_{n_1, n_2\in[N]}T_1^{n_1}T_2^{n_2}f\cdot e(-tn_1 - sn_2)
\end{align*}
for simplicity, and note that the single limit over $([N]^2)_N$ could be replaced by an iterated limit. If $H\cap H' = \emptyset$, then clearly $\langle P_H f, P_{H'} f\rangle = 0$ and $P_{H\cup H'} f = P_H f + P_{H'}f$.  It follows in particular that if $H\subseteq H'$, then
\begin{align}
    \nnorm{P_{H'} f}_{\be_1-\beta \be_2}^2 = \nnorm{P_{H} f}_{\be_1-\beta \be_2}^2 + \nnorm{P_{H'\setminus H} f}_{\be_1-\beta \be_2}^2\geq \nnorm{P_{H} f}_{\be_1-\beta \be_2}^2;
\end{align}
to see this, we write $P_{H'}f = P_H f +P_{H'\setminus H} f,$ expand the left-hand side as a Gowers-Cauchy-Schwarz inner product, and then use the orthogonality of $P_H f$, $P_{H'\setminus H} f$ to cancel the cross terms.
%split $P_H f = \sum_{(t,s)\in H} P_{(t,s)} f$, and similarly for $P_{H'\setminus H} f$ (these are countable sums), and then use orthogonality to cancel the cross terms. 
%\AK{If you mean here that you are using an infinite "Pythagorean", I think we have to be more careful.} \BK{No, there are just two parts: projection onto $H$ and $H'\setminus H$. Expand $\nnorm{P_{H'} f}_{\be_1-\beta \be_2}^2$ as a GCS inner product, then the cross terms should cancel because of the orthogonality of $P_H$ and $P_{H'\setminus H}$. Do carry out the computation though and see if it makes sense to you.} \AK{Vale}

The next lemma will be crucial in completing the proof of Proposition \ref{P: beta irrational}. For eigenfunctions with eigenvalues of the form $(\beta \gamma, \gamma)$ mod $\Z^2$ with $\gamma\neq 0$, it allows us to replace the single-parameter action of $(T_1^n T_2^{-\floor{\beta n}})_n$ by the double-parameter action of $(T_1^{\floor{\frac{n}{\beta \gamma}}} T_2^{\floor{-\frac{m}{\gamma}}})_{n,m}$.

\begin{lemma}\label{L: scaling}
Let $\gamma\in\R$ be nonzero and $f\in L^2(\mu)$. Then there exists $c_{\beta, \gamma}\in\C$ such that  %Then the function $$F=\E_n T_1^n T_2^{-\floor{\beta n}} P_{(\beta \gamma, \gamma)} f$$ is either 0, or there exists $c_{\beta, \gamma}\neq 0$ such that
\begin{align}\label{E: 2-parameter action}
\lim_{N\to\infty}\E_{n\in[N]} T_1^n T_2^{-\floor{\beta n}} P_{(\beta \gamma, \gamma)} f = c_{\beta, \gamma} \cdot \lim_{N\to\infty}\E_{n, m\in[N]} T_1^{\floor{\frac{n}{\beta \gamma}}} T_2^{\floor{-\frac{m}{\gamma}}} P_{(\beta \gamma, \gamma)} f.
\end{align}
\end{lemma}

One could alternatively try to compare the left-hand side of \eqref{E: 2-parameter action} with $$\lim\limits_{N\to\infty}\E\limits_{n, m\in[N]} T_1^{\floor{\xi n}} T_2^{\floor{-\xi \beta m}} P_{(\beta \gamma, \gamma)} f$$ for other choices of $\xi\in\R\backslash\{0\}$. However, 
%the exponential sum coming from the right-hand side of \eqref{E: 2-parameter action} 
this average vanishes unless $\xi\in\frac{1}{\beta\gamma}\Z$.
% \begin{align*}
%     \E_n e\brac{\xi \beta \gamma n - \rem{\xi n}\beta\gamma}\cdot \E_m e\brac{-\xi \beta \gamma m - \rem{\xi\beta m}\beta\gamma}.
% \end{align*}
% They vanish unless $\xi\in\frac{1}{\beta\gamma}\Z$, I think (but needs a check)

% \begin{align*}
%     \xi = \frac{p}{q + \beta r + s\beta\gamma}
% \end{align*}
% for some $p,q,r,s\in\Z$, RHS would be zero, so such a comparison would not be possible. For $\xi$ of the form above, the exponential sum coming from the RHS would be
% \begin{align*}
%     \E_{n,m} e\brac{\frac{p\beta\gamma}{q + \beta r + s\beta\gamma}(n-m)}
% \end{align*}
% pretty complicated
% One can also consider a linear combo over various $\xi$'s.}

\begin{proof}
Let 
\begin{align*}
F &:= \lim_{N\to\infty}\E_{n\in[N]} T_1^n T_2^{-\floor{\beta n}} P_{(\beta \gamma, \gamma)} f = P_{(\beta \gamma, \gamma)} f \cdot \lim_{N\to\infty}\E_{n\in[N]} e(\beta\gamma n - \floor{\beta n}\gamma)\\
&= P_{(\beta \gamma, \gamma)} f \cdot \lim_{N\to\infty}\E_{n\in[N]} e(\rem{\beta n}\gamma) = P_{(\beta \gamma, \gamma)} f \cdot \int_0^1 e(\gamma y)\, dy
\end{align*}
and 
\begin{align*}
G &:= \lim_{N\to\infty}\E_{n, m\in[N]} T_1^{\floor{\frac{n}{\beta \gamma}}} T_2^{\floor{-\frac{m}{\gamma}}} P_{(\beta \gamma, \gamma)} f\\
&= P_{(\beta \gamma, \gamma)} f\cdot \lim_{N\to\infty}\E_{n, m\in[N]} e\brac{\floor{\frac{n}{\beta \gamma}}\beta \gamma + \floor{-\frac{m}{\gamma}}\gamma}\\
&= P_{(\beta \gamma, \gamma)} f\cdot \lim_{N\to\infty}\E_{n, m\in[N]} e\brac{-\rem{\frac{n}{\beta \gamma}}\beta \gamma - \rem{-\frac{m}{\gamma}}\gamma}.
\end{align*}

If $\gamma\in \Z\setminus\{0\}$, then $F = 0$, and so the claim follows trivially with $c_{\beta, \gamma} = 0$. 
Otherwise $\int_0^1 e(\gamma y)\, dy\neq 0$, and the goal is to show that
\begin{align*}
    C_{\beta, \gamma} := \lim_{N\to\infty}\E_{n, m\in[N]} e\brac{-\rem{\frac{n}{\beta \gamma}}\beta \gamma - \rem{-\frac{m}{\gamma}}\gamma}
\end{align*}
is nonzero as well so that
\begin{align*}
    c_{\beta, \gamma} := C_{\beta, \gamma}\inv\cdot\int_0^1 e(\gamma y)\, dy
    %\brac{\E_{n, m} e\brac{-\rem{\frac{n}{\beta \gamma}}\beta \gamma + \rem{\frac{m}{\gamma}}\gamma}}
\end{align*}
is well-defined.

We split the proof into several cases.
%        \smallskip
 %       \textbf{Case $\gamma \notin\frac{1}{\beta}\cdot \Q\cup \Q$}
  %      \smallskip
If $\gamma \notin(\frac{1}{\beta} \Q)\cup \Q$, then both $\beta \gamma$ and $\gamma$ are irrational, and so 
\begin{align*}
	C_{\beta, \gamma} = \int_0^1 e(-\beta \gamma x)\, dx\cdot \int_0^1 e(-\gamma y)\, dy.
\end{align*}
Both integrals are nonzero since neither $\gamma$ nor $\beta\gamma$ is an integer.


If $\gamma\in\Q\setminus\Z$, write $\gamma = -\frac{p}{q}$ for coprime integers $p,q$. Then 
\begin{align*}
C_{\beta, \gamma} = \int_0^1 e(-\beta \gamma x)\, dx \cdot \E_{m\in[p]} e\brac{\rem{\frac{qm}{p}}\frac{p}{q}}.
\end{align*}
The integral is nonzero since $\beta\gamma\notin\Z$. To see that the exponential sum is nonzero, let $q^*$ be the multiplicative inverse of $q$ mod $p$. Then the map $m\mapsto q^* m$ is a bijection on $\Z/p\Z$, and so 
\begin{align*}
	\E_{m\in[p]} e\brac{\rem{\frac{qm}{p}}\frac{p}{q}} &= \E_{m\in[p]} e\brac{\rem{\frac{qq^*m}{p}}\frac{p}{q}} = \E_{m\in[p]} e\brac{\frac{m}{q}}\neq 0 % = \frac{1}{p}\cdot\frac{1-e(p/q)}{e(-1/q)-1}\neq 0
\end{align*}
%since $|q|>1$.
due to the coprimality of $p,q$.

Lastly, suppose that $\gamma\in\frac{1}{\beta}\Q$. Setting $\gamma = \frac{p}{q\beta}$ for coprime integers $p,q$, we deduce that
\begin{align*}
	C_{\beta, \gamma} = \int_0^1 e(-\gamma y)\, dy \cdot  \E_{n\in[p]} e\brac{-\rem{\frac{qn}{p}}\frac{p}{q}}.
\end{align*}
The integral is nonzero since $\gamma\notin\Z$, and the exponential sum is nonzero by the same argument as before.
%and hence we obtain the claim with $c_{\beta, \gamma} = \brac{\E_{n\in[p]} e\brac{\rem{\frac{qn}{p}}\frac{p}{q}}}^{-1}.$
\end{proof}
We conclude this section with the proof of Proposition \ref{P: beta irrational}.
\begin{proof}[Proof of Proposition \ref{P: beta irrational}]
If $a$ is constant, then the action can only be ergodic if the system is conjugate to a 1-point system, in which case the result follows trivially. Suppose therefore that $a$ is not constant.

    By the arguments below the statement of Proposition \ref{P: beta irrational}, we have
    \begin{align*}
        \nnorm{{f}}_{\be_1-\beta \be_2}^2 = \int_{E_3\cup\{(0,0)\}} |B(t,s)|^2\, d\sigma_{{f}}(t,s).
    \end{align*}
By Proposition \ref{P: E_3}, every element of the support $\tilde{E}_3:=E_3\cup\{(0,0)\}$ of the preceding integral takes the form $(\beta \gamma, \gamma)$ mod $\Z^2$.
Using this and orthogonality of projections on different eigenspaces, we get
\begin{align*}
    \nnorm{f}_{\be_1-\beta \be_2}^2 = \nnorm{P_{\tilde{E}_3}f}_{\be_1-\beta \be_2}^2 = \sum_{(\beta \gamma, \gamma)\in \Tilde{E}_3}\nnorm{P_{(\beta\gamma,\gamma)}f}_{\be_1-\beta \be_2}^2.
\end{align*}
By Lemma \ref{L: scaling}, for any nonzero $\gamma\in\R$,
  there exists $c_{\beta, \gamma}\geq 0$ such that 
\begin{align*}
\nnorm{P_{(\beta\gamma,\gamma)}f}_{\be_1-\beta \be_2}&=\lim_{N\to\infty}\norm{\E_{n\in[N]} T_1^n T_2^{-\floor{\beta n}} P_{(\beta \gamma, \gamma)} f}_{L^2(\mu)}\\
&= c_{\beta, \gamma} \cdot \lim_{N\to\infty}\norm{\E_{n, m\in[N]} T_1^{\floor{\frac{n}{\beta \gamma}}} T_2^{\floor{-\frac{m}{\gamma}}} P_{(\beta \gamma, \gamma)} f}_{L^2(\mu)}\\
&= c_{\beta, \gamma} \cdot \nnorm{P_{(\beta\gamma,\gamma)}f}_{\langle (1/\beta \gamma)\be_1, (-1/\gamma)\be_2\rangle}\ll_{\beta,\gamma} \nnorm{P_{(\beta\gamma,\gamma)}f}_{\langle \be_1, \be_2\rangle}^+
\end{align*}
by Lemma \ref{L: basic properties of seminorms}. 
% The fact that $P_{(\beta\gamma,\gamma)}f$ is an eigenfunction of $T_1, T_2$ implies that
% \begin{align*}
%     \nnorm{P_{(\beta\gamma,\gamma)}f}_{\langle \be_1, \be_2\rangle}^+ = \norm{P_{(\beta\gamma,\gamma)}f}_{L^2(\mu)}.
% \end{align*}
A usual computation gives
\begin{align*}
    \nnorm{f}_{\langle \be_1, \be_2\rangle}^+ = \nnorm{P_{[0,1)^2}f}_{\langle \be_1, \be_2\rangle}^+
\end{align*}
(where $P_{[0,1)^2}f$ is the projection on the Kronecker factor of the joint action of $T_1, T_2$),
and so combining the computations so far, we get
\begin{align*}
    \nnorm{P_{(\beta\gamma,\gamma)}f}_{\be_1-\beta \be_2}\ll_{\beta,\gamma} \nnorm{P_{(\beta\gamma,\gamma)}f}_{\langle \be_1, \be_2\rangle}^+ \leq  \nnorm{P_{[0,1)^2}f}_{\langle \be_1, \be_2\rangle}^+ =  \nnorm{f}_{\langle \be_1, \be_2\rangle}^+
\end{align*}
for any $\gamma\in\R$ (including $\gamma = 0$, which is trivial). Hence if $\nnorm{f}_{\langle \be_1, \be_2\rangle}^+ = 0$, then $\nnorm{f}_{\be_1-\beta \be_2} = 0$ as well. 
% For general $\alpha, \beta\in\R$, we can combine the argument above with \eqref{E: alpha e_1 - beta e_2 vs. e_1, e_2} to deduce that $\nnorm{f}_{\alpha \be_1 - \beta \be_2} = 0$ whenever $\nnorm{f}_{\langle \be_1, \be_2\rangle^2}^+ = 0$. \BK{This last line unfortunately does not work since it relies on the following chain of implications: $\nnorm{f}_{\langle \be_1, \be_2\rangle^2}^+ = 0$ implies $\nnorm{\Delta_{(m,m')}f}_{\langle \be_1, \be_2\rangle}^+ = 0$ for almost all $(m,m')$ implies $\nnorm{\Delta_{(m,m')}f}_{\be_1 - \beta\be_2} = 0$ for almost all $(m,m')$ implies $\nnorm{f}_{\alpha\be_1 - \beta\be_2} = 0$. Unfortunately, the first of these implications doesn't work because the vanishing of $\nnorm{f}_{\langle \be_1, \be_2\rangle^2}^+$ can happen if $\nnorm{\Delta_{(m,m')}f}_{\langle \be_1, \be_2\rangle}^+$ are close to 0 but positive. To avoid this, we probably need to upgrade the implication ``0 implies 0'' to ``$\delta$ implies $\veps$''. If we cannot do this by uniformly bounding $c_{\beta, \gamma}$ from Lemma \ref{L: scaling}, then we can use \cite[Proposition A.2]{FrKu22a}.}
\end{proof}
A glimpse at the proof shows that we cannot make the dependence between $\nnorm{f}_{\langle \be_1, \be_2\rangle}^+$ and $\nnorm{f}_{\be_1-\beta \be_2}$ quantitative since the constants $c_{\beta, \gamma}$ can be arbitrarily close to 0 for $(\beta\gamma,\gamma)\in \tilde{E}_3$.

\section{Proofs of main theorems}\label{S: proofs of main theorems}
In this section, we finally prove the theorems from the introduction. We start with a discussion on the ergodicity for slowly growing sequences.  In what follows, we say that a sequence of transformations $(T_n)_n$ is \textit{strongly/weakly ergodic} on $(X, \CX, \mu)$ if for any $f\in L^\infty(\mu)$, we have
\begin{align*}
    \lim_{N\to\infty}\E_{n\in[N]}T_n f= \int f\; d\mu
\end{align*}
in the strongly/weakly $L^{2}(\mu)$-sense (thus, the notion of strong ergodicity is the same as ergodicity). 
%\AK{Since ``ergodicity'' is a universal notion, why do we have to define ``strong ergodicity'' too? Just define the weak one.} \BK{The discussion below is easier to carry out if we have a strong and weak version of everything. Let's leave it as it is.}




 




%First, we state one auxiliary lemma which allows us to toss away all slowly growing functions.
\begin{lemma}[Ergodicity of slowly growing sequences]\label{L: slowly growing}
    Let $a\in\CH$ with $a\ll \log$, and let $(X, \CX, \mu, T)$ be a system.  
    \begin{enumerate}
        \item Then $(T^{\floor{a(n)}})_n$ is strongly ergodic for $(X, \CX, \mu)$ if and only if $(X, \CX, \mu, T)$ is conjugate to a one-point system, or equivalently $\mu$ is the Dirac measure $\delta_x$ for~some~$x~\in~X$;
        \item If in addition $a\succ 1$, then $(T^{\floor{a(n)}})_n$ is weakly ergodic for $(X, \CX, \mu)$ if and only if $T$ is mixing.
    \end{enumerate}
\end{lemma}
\begin{proof}
Part (i). The converse direction is trivial, so we only prove the forward direction.

Suppose first that $a$ converges to a constant. Since $a$ is eventually monotone, we can find $c\in\N$ such that $\norm{T^c f - \int f\; d\mu}_{L^2(\mu)} = 0$ for all $f\in L^\infty(\mu)$. This implies that $f$ is constant $\mu$-a.e.; since $f$ was arbitrary, we deduce that the measure must be concentrated on a single point.


Suppose now that $1\prec a\ll\log$. First, we show that $a\inv(x)\gg (1+c)^x$ for some $c>0$. If this is not the case, then $\lim\limits_{x\to\infty}\frac{\log a\inv(x)}{x} = 0$. Taking the limit along $a(x)$ rather than $x$, we get $\lim\limits_{x\to\infty} \frac{\log x}{a(x)}=0$, contradicting our assumption $a\ll \log$.

Then, by Lemma \ref{L: dyadic intervals}, we get 
\begin{align*}
    \lim_{N\to\infty}\norm{\E_{n\in[a\inv(N), a\inv(N+1))}T^{\floor{a(n)}}f-\int f\; d\mu}_{L^2(\mu)} = 0
\end{align*}
for all $f\in L^\infty(\mu)$. Noticing that $\floor{a(n)} = N$ for $n\in [a\inv(N), a\inv(N+1))$, we can get rid of the average, obtaining
\begin{align*}
     \lim_{N\to\infty}\norm{T^N f-\int f\; d\mu}_{L^2(\mu)} = \norm{f-\int f\; d\mu}_{L^2(\mu)} = 0.
\end{align*}
Once again, the last equality implies that $f$ is constant $\mu$-a.e, and hence the measure must be concentrated on a single point.

\

Part (ii). We thank Nikos Frantzikinakis for showing this argument to us. Suppose first that $(T^{\floor{a(n)}})_n$ is weakly ergodic. Arguing as in the proof of Lemma \ref{L: slowly growing}, we deduce that
\begin{align*}
    \lim_{N\to\infty}\int f_1\cdot \E_{n\in[a\inv(N), a\inv(N+1))}T^{\floor{a(n)}}f_2\; d\mu = \int f_1\; d\mu\cdot \int f_2\; d\mu
\end{align*}
for all $f_1,f_2\in L^\infty(\mu)$. By noticing that $\floor{a(n)} = N$ for $n\in [a\inv(N), a\inv(N+1))$, we infer that
\begin{align*}
    \lim_{N\to\infty}\int f_1\cdot T^N f_2\; d\mu = \int f_1\; d\mu\cdot \int f_2\; d\mu,
\end{align*}
hence $T$ is mixing. The converse implication is immediate. 
\end{proof}

We start with establishing the seminorm control provided by Theorem \ref{T: seminorm control}.

\begin{proof}[Proof of Theorem \ref{T: seminorm control}]
Suppose that $a_{1},\dots,a_{\ell}\in\CH$ satisfy the difference ergodicity condition on a system $(X, \CX, \mu, T_1, \ldots, T_\ell)$, and that $T_1, \ldots, T_\ell$ are all ergodic. We will show that they are good for seminorm control. If $a_j\ll \log$ for some $j\in[\ell]$, then Lemma \ref{L: slowly growing} implies that the measure is concentrated on a single point, in which case seminorm control follows immediately. We assume therefore for the rest of the proof that $a_1, \ldots, a_\ell\succ \log$.

%By Corollary 12.5 of paper 1, the product ergodicity condition implies that $a_{1},\dots,a_{\ell}$ are good for equidistribution for the system. By ???, it suffices to show that $a_{1},\dots,a_{\ell}$ are good for seminorm estimate. 

By \cite[Lemma A.3]{R22}, we can find an ordered Hardy family $\CQ = (q_1, \ldots, q_m)\subseteq\CH$ such that $a_{1},\dots,a_{\ell}\in\extspan_\R\CQ$.
%By ???, $\CQ$ is good for smoothing.
By Theorem \ref{T: new estimates}, there exists a positive integer $s=O_{\ell, \CQ}(1)$ and vectors $\bv_{1},\dots,\bv_{s}$ coming from the set of the leading coefficients of 
        \begin{align*}
%            \{a_1 \be_1\}\cup\{a_i \be_i - a_j \be_j:\; i\neq j,\; a_i, a_j\; \textrm{are\; dependent}\},
            \{a_1 \be_1\}\cup\{a_i \be_i - a_j \be_j:\; i\neq j,\; a_i, a_j\; \textrm{are\; dependent}\}
            % and } \deg(a_{i})=\deg(a_{j})\geq \deg(a_{1})\}
        \end{align*}
        (taken with respect to $\CQ$) such that for all $f_{1},\dots,f_{\ell}\in L^{\infty}(\mu)$, we have
 % Let $f_{1},\dots,f_{\ell}\in L^{\infty}(\mu)$ and fix $j\in[\ell]$.
  %If $a_{i}\prec \log$, {\color{red} then add the details.} If $a_{i}\succ \log$, then we have that $\deg(a_{i})>m_{1}$. 
  % By Theorem \ref{T: new estimates general}, 
        %$a_1 \be_1$ and those $a_i \be_i - a_j \be_j$ for which $i\neq j$ and $a_i$, $a_j$ are pairwise dependent,
        %$$a_1 \be_1\qquad \textrm{and}\qquad a_i \be_i - a_j \be_j\quad \textrm{as\; long\; as}\; i\neq j\; \textrm{and}\; a_i,\; a_j\; \textrm{are\; pairwise\; dependent},$$
         % such that 
\begin{align}\label{goal0}
                  \lim_{N\to\infty} %\sup_{|c_n|\leq 1}
         \norm{\E_{n\in [N]} 
         %c_{n}\cdot
         \prod_{j\in[\ell]}  T_{j}^{\sfloor{a_{j}(n)}} f_{j}}_{L^2(\mu)}=0
    \end{align}
    whenever $\nnorm{f_1}_{\bv_{1},\dots,\bv_{s}}=0$. Note that if $a_i, a_j$ are dependent, then the leading coefficient of $a_i \be_i - a_j \be_j$ equals $\beta_i \be_i - \beta_j \be_j$ for nonzero $\beta_i, \beta_j\in\R$ for which $\beta_j a_i - \beta_i a_j\ll\log$. Applying Proposition \ref{P: seminorm comparison} iteratively, we can replace each such vector in our seminorm by groups $\Z^\ell$ and $\langle \be_i, \be_j\rangle$. Combined with Lemma \ref{L: basic properties of seminorms}, the ergodicity of $T_i, T_j$ allows us to replace each such group by $\be_1$, and so we can find an integer $s'=O_{\ell, \CQ}(1)$ such that (\ref{goal0}) holds whenever $\nnorm{f_1}_{\be_1^{\times s'}}=0$. By symmetry, we can get seminorm control of (\ref{goal0}) in terms of respective bounded-degree Host-Kra seminorms of other functions. 
%     or whenever $\nnorm{f_1}^{+}_{\bv_{1},\dots,\bv_{s}}=0$ by ???, or whenever $\nnorm{f_1}^{+}_{\langle\be_{1},\be_{i}\rangle^{s'}, 1\leq i\leq \ell}=0$ for some $s'=O_{\ell, \CQ}(1)$ by repeatedly using the difference ergodicity condition, Lemma \ref{L: basic properties of seminorms}, and the following proposition
% {\color{red} (We need to state the proposition in the following form)
%  \begin{proposition} 
%       Let $a_{1},a_{2}$ be elements of some ordered Hardy family $\CQ$ with $\deg(a_{1}),\deg(a_{2})>m_{1}$. Let $\bv$ be the leading coefficients of $a_{1}\be_{1}-a_{2}\be_{2}$. Suppose that $(T_{1}^{\lfloor a_{1}(n)\rfloor}T_{2}^{-\lfloor a_{2}(n)\rfloor})_{n}$ is ergodic. Then for all $s\in\N\cup\{0\}$, $\bv_{1},\dots,\bv_{s}$ and $f\in L^{\infty}(\mu)$, we have that
%       $$\text{$\nnorm{f}^{+}_{\bv_{1},\dots,\bv_{s},\bv}=0$ whenever $\nnorm{f}^{+}_{\bv_{1},\dots,\bv_{s},\langle\be_{1},\be_{2}\rangle^{O(1)}}=0$.}$$
%  \end{proposition}
%  }
%Note that the difference ergodicity condition implies that $\langle\be_{1},\be_{i}\rangle^{s'}$ is ergodic for all $2\leq i\leq \ell$, whereas the product ergodicity condition implies that $\be_{1}$ is ergodic. So by ???, we have that (\ref{goal0}) holds whenever $\nnorm{f_1}^{+}_{\e_{1}^{s''}}=0$ or whenever $\nnorm{f_1}_{\e_{1}^{s''+1}}=0$ for some $s''=O_{\ell, \CQ}(1)$. By symmetry, (\ref{goal0}) holds whenever  $\nnorm{f_i}_{\e_{i}^{O_{\ell, \CQ}(1)}}=0$ for some $i\in[\ell]$. 
Therefore, $a_{1},\dots,a_{\ell}$ are good for seminorm control for the system and we are done.
\end{proof}

We are now ready to prove Theorem \ref{T: main theorem}, the main result of this paper.

\begin{proof}[Proof of Theorem \ref{T: main theorem}]
Suppose that $a_{1},\dots,a_{\ell}\in\CH$ satisfy the difference and product ergodicity conditions on a system $(X, \CX, \mu, T_1, \ldots, T_\ell)$; we will show that they are jointly ergodic. Again, we can assume that $a_1, \ldots, a_\ell\succ \log$ since otherwise the result is trivial by Lemma \ref{L: slowly growing}.

%Suppose that $a_{1},\dots,a_{\ell}$ satisfy the difference and product ergodicity conditions from Definition \ref{D: good for joint ergodicity}. We will show that $a_{1},\dots,a_{\ell}$ are good for joint ergodicity. 
Using Theorem \ref{T: joint ergodicity criteria}, it suffices to establish that $a_1, \ldots, a_\ell$ are good for  equidistribution and seminorm control. By \cite[Corollary 12.5]{DKKST24}, the condition of being good for equidistribution follows from the product ergodicity condition. Moreover, the product ergodicity condition implies that $T_1, \ldots, T_\ell$ are all ergodic, which jointly with the difference ergodicity condition and Theorem \ref{T: seminorm control} give seminorm control.
\end{proof}

% \begin{proof}[Proof of Theorem \ref{T: main theorem}]
% Suppose that $a_{1},\dots,a_{\ell}\in\CH$ satisfy the difference and product ergodicity conditions on a system $(X, \CX, \mu, T_1, \ldots, T_\ell)$; we will show that they are jointly ergodic. If $a_j\ll \log$ for some $j\in[\ell]$, then Lemma \ref{L: slowly growing} implies that the measure is concentrated on a single point in which case, in which case joint ergodicity follows immediately. We assume therefore for the rest of the proof that $a_1, \ldots, a_\ell\succ \log$.

% %Suppose that $a_{1},\dots,a_{\ell}$ satisfy the difference and product ergodicity conditions from Definition \ref{D: good for joint ergodicity}. We will show that $a_{1},\dots,a_{\ell}$ are good for joint ergodicity. 
% Using Theorem \ref{T: joint ergodicity criteria}, it suffices to establish that $a_1, \ldots, a_\ell$ are good for  equidistribution and seminorm control. By \cite[Corollary 12.5]{DKKST24}, the condition of being good for equidistribution follows from the product ergodicity condition. It remains to prove that $a_1, \ldots, a_\ell$ are good for seminorm control. 


% %By Corollary 12.5 of paper 1, the product ergodicity condition implies that $a_{1},\dots,a_{\ell}$ are good for equidistribution for the system. By ???, it suffices to show that $a_{1},\dots,a_{\ell}$ are good for seminorm estimate. 

% By \cite[Lemma A.3]{R22}, we can find an ordered Hardy family $\CQ = \{q_1, \ldots, q_m\}\subseteq\CH$ such that $a_{1},\dots,a_{\ell}\in\extspan_\R\CQ$.
% %By ???, $\CQ$ is good for smoothing.
% By Theorem \ref{T: new estimates}, there exists a positive integer $s=O_{\ell, \CQ}(1)$ and vectors $\bv_{1},\dots,\bv_{s}$ coming from the set of the leading coefficients of 
%         \begin{align*}
% %            \{a_1 \be_1\}\cup\{a_i \be_i - a_j \be_j:\; i\neq j,\; a_i, a_j\; \textrm{are\; dependent}\},
%             \{a_1 \be_1\}\cup\{a_i \be_i - a_j \be_j:\; i\neq j,\; a_i, a_j\; \textrm{are\; dependent}\}
%             % and } \deg(a_{i})=\deg(a_{j})\geq \deg(a_{1})\}
%         \end{align*}
%         (taken with respect to $\CQ$) such that for all $f_{1},\dots,f_{\ell}\in L^{\infty}(\mu)$, we have
%  % Let $f_{1},\dots,f_{\ell}\in L^{\infty}(\mu)$ and fix $j\in[\ell]$.
%   %If $a_{i}\prec \log$, {\color{red} then add the details.} If $a_{i}\succ \log$, then we have that $\deg(a_{i})>m_{1}$. 
%   % By Theorem \ref{T: new estimates general}, 
%         %$a_1 \be_1$ and those $a_i \be_i - a_j \be_j$ for which $i\neq j$ and $a_i$, $a_j$ are pairwise dependent,
%         %$$a_1 \be_1\qquad \textrm{and}\qquad a_i \be_i - a_j \be_j\quad \textrm{as\; long\; as}\; i\neq j\; \textrm{and}\; a_i,\; a_j\; \textrm{are\; pairwise\; dependent},$$
%          % such that 
% \begin{align}\label{goal0}
%                   \lim_{N\to\infty} %\sup_{|c_n|\leq 1}
%          \norm{\E_{n\in [N]} 
%          %c_{n}\cdot
%          \prod_{j\in[\ell]}  T_{j}^{\sfloor{a_{j}(n)}} f_{j}}_{L^2(\mu)}=0
%     \end{align}
%     whenever $\nnorm{f_1}_{\bv_{1},\dots,\bv_{s}}=0$. Note that if $a_i, a_j$ are dependent, then the leading coefficient of $a_i \be_i - a_j \be_j$ equals $\beta_i \be_i - \beta_j \be_j$ for nonzero $\beta_i, \beta_j\in\R$ for which $\beta_j a_i - \beta_i a_j\ll\log$. Applying Proposition \ref{P: seminorm comparison} iteratively, we can replace each such vector in our seminorm by groups $\Z^\ell$ and $\langle \be_i, \be_j\rangle$. Combined with Lemma \ref{L: basic properties of seminorms}, the ergodicity of $T_i, T_j$ (itself a consequence of the product ergodicity condition) allows us to replace each such group by $\be_1$, and so we can find an integer $s'=O_{\ell, \CQ}(1)$ such that (\ref{goal0}) holds whenever $\nnorm{f_1}_{\be_1^{s'}}=0$. By symmetry, we can get seminorm control of (\ref{goal0}) in terms of respective bounded-degree Host-Kra seminorms of other functions. 
% %     or whenever $\nnorm{f_1}^{+}_{\bv_{1},\dots,\bv_{s}}=0$ by ???, or whenever $\nnorm{f_1}^{+}_{\langle\be_{1},\be_{i}\rangle^{s'}, 1\leq i\leq \ell}=0$ for some $s'=O_{\ell, \CQ}(1)$ by repeatedly using the difference ergodicity condition, Lemma \ref{L: basic properties of seminorms}, and the following proposition
% % {\color{red} (We need to state the proposition in the following form)
% %  \begin{proposition} 
% %       Let $a_{1},a_{2}$ be elements of some ordered Hardy family $\CQ$ with $\deg(a_{1}),\deg(a_{2})>m_{1}$. Let $\bv$ be the leading coefficients of $a_{1}\be_{1}-a_{2}\be_{2}$. Suppose that $(T_{1}^{\lfloor a_{1}(n)\rfloor}T_{2}^{-\lfloor a_{2}(n)\rfloor})_{n}$ is ergodic. Then for all $s\in\N\cup\{0\}$, $\bv_{1},\dots,\bv_{s}$ and $f\in L^{\infty}(\mu)$, we have that
% %       $$\text{$\nnorm{f}^{+}_{\bv_{1},\dots,\bv_{s},\bv}=0$ whenever $\nnorm{f}^{+}_{\bv_{1},\dots,\bv_{s},\langle\be_{1},\be_{2}\rangle^{O(1)}}=0$.}$$
% %  \end{proposition}
% %  }
% %Note that the difference ergodicity condition implies that $\langle\be_{1},\be_{i}\rangle^{s'}$ is ergodic for all $2\leq i\leq \ell$, whereas the product ergodicity condition implies that $\be_{1}$ is ergodic. So by ???, we have that (\ref{goal0}) holds whenever $\nnorm{f_1}^{+}_{\e_{1}^{s''}}=0$ or whenever $\nnorm{f_1}_{\e_{1}^{s''+1}}=0$ for some $s''=O_{\ell, \CQ}(1)$. By symmetry, (\ref{goal0}) holds whenever  $\nnorm{f_i}_{\e_{i}^{O_{\ell, \CQ}(1)}}=0$ for some $i\in[\ell]$. 
% Therefore, $a_{1},\dots,a_{\ell}$ are good for seminorm control for the system and we are done.
% \end{proof}


Next, we establish Theorem \ref{T: easy direction}.
\begin{proof}[Proof of Theorem \ref{T: easy direction}]
    Suppose now that $a_{1},\dots,a_{\ell}\in\CH$ are jointly ergodic for a system $(X, \CX, \mu, T_1, \ldots, T_\ell)$. 
    Once again, we can assume that  $a_1, \ldots, a_\ell\succ \log$ for otherwise the result follows trivially from Lemma \ref{L: slowly growing}. Then the product ergodicity condition follows from \cite[Lemmas 12.6 and 12.13]{DKKST24}.

    To get the difference ergodicity condition, we need to additionally assume that for every $i\neq j$ and $c_i,c_j\in\R$, the sequence $c_i a_i - c_j a_j$ is either an almost rational polynomial or it stays logarithmically away from rational polynomials. We will show that under these conditions, the exponential sum
    \begin{align}\label{E: convergence of exp sum}
        \E_{n\in[N]}e(t\floor{a_1(n)}+s\floor{a_2(n)})
    \end{align}
    converges, and hence by the spectral theorem, $\E_{n\in[N]}T_i^{\floor{a_i(n)}}T_j^{-\floor{a_j(n)}}f$ converges in $L^2(\mu)$ for any $f\in L^2(\mu)$. Therefore the $L^2(\mu)$ limit must equal the weak limit, which equals $\int f\; d\mu$ by joint ergodicity because
    \begin{align*}
        \lim_{N\to\infty}\int g\cdot \E_{n\in[N]}T_i^{\floor{a_i(n)}}T_j^{-\floor{a_j(n)}}f\; d\mu &= \lim_{N\to\infty}\int \E_{n\in[N]} T_j^{\floor{a_j(n)}}g\cdot T_i^{\floor{a_i(n)}}f\; d\mu\\
        &= \int g\; d\mu \cdot \int f\; d\mu.
    \end{align*}

    To show that \eqref{E: convergence of exp sum} converges, we expand it as
    \begin{align*}
        \E_{n\in[N]}e(t a_1(n)+s a_2(n) - t\{a_1(n)\}-s\{a_2(n)\}).
    \end{align*}
    Since the function $x\mapsto e(-t\{x\})$ is Riemann integrable for any $t\in\R$, we can approximate it arbitrarily well by trigonometric polynomials, and hence it suffices to show the convergence of 
    \begin{align*}
        \E_{n\in[N]}e(t' a_1(n)+s' a_2(n))
    \end{align*}
    for all $t'\in t+\Z,\; s'\in s+\Z$. By assumption, either $t' a_1(n)+s' a_2(n)$ equals $p(n)+c + o(1)$ for some $p\in\Q[x]$ and $c\in\R$ as $n\to\infty$, in which case the exponential sum obviously converges, or it stays logarithmically away from rational polynomials, in which case it converges by Theorem \ref{T: Boshernitzan}.    
%If $a_i, a_j$ are pairwise independent, then the ergodicity of $(T_i^{\floor{a_i(n)}}T_j^{-\floor{a_j(n)}})_n$ follows from \cite[Lemma 12.12]{DKKST24}. Suppose therefore that $a_i, a_j$ are pairwise dependent, in which case we can parametrize them as $a_i = a$ and $a_j = \beta a + g$ for some $\beta\in\R$ and $a, g\in\CH$ with $a\succ \log$ and $g\ll \log$. The joint ergodicity of $(T_i^{\floor{a_i(n)}}, T_j^{\floor{a_j(n)}})_n = (T_i^{\floor{a(n)}}, T_j^{\floor{\beta a(n)+g(n)}})_n$ implies that $\E\limits_{n\in[N]} T_i^{\floor{a(n)}}T_j^{-\floor{\beta a(n)+g(n)}}f\to\int f\; d\mu$ weakly for any $f\in L^\infty(\mu)$. It therefore suffices to show that $\E\limits_{n\in[N]} T_i^{\floor{a(n)}}T_j^{-\floor{\beta a(n)+g(n)}}f$ converges in $L^2(\mu)$.
\end{proof}


In order to prove Theorem \ref{T: counterexample}, we make use of the following seminorm estimate.

\begin{proposition}\label{est12}
    Let $(X,\mathcal{X},\mu,T_{1},T_{2})$ be a system and $f_{1},f_{2}\in L^{\infty}(\mu)$. Assume that $T_{2}\vert_{\CI(T^{-1}_{1}T_{2})}$ is mixing. We have that 
    $$\lim_{N\to\infty}\Bigl\Vert\E_{n\in[N]} T_{1}^n f_1\cdot T_{2}^{n+\floor{\log_{2}n}}f_2\Bigr\Vert_{L^2(\mu)}=0$$
    if $\nnorm{f_{j}}_{T_{j},T_{j}}=0$ for any $j=1,2$.
\end{proposition}
\begin{proof}
%We show that \begin{equation*}
 %       \lim_{N\to\infty} \E_{n\in[N]}  T_{1}^n f_1\cdot T_{2}^{n+\floor{\log_{2}n}}f_2=0
  %  \end{equation*}in $L^2(\mu)$ whenever $\min\limits_{j\in[2]}\nnorm{f_j}_{T_j^{\times O(1)}}=0$.
%  % Dyadically decomposing the sum, it suffices to show that \begin{equation*}
 % %       \lim_{N\to\infty}\frac{1}{N}\sum_{k\leq \floor{\log_2 N}} \sum_{2^k\leq n< 2^{k+1}}T_{1}^n f_1\cdot T_{2}^{n+k}f_2=0
%  %   \end{equation*}in $L^2(\mu)$.
        Dyadically decomposing the sum, it suffices to show that 
        %\AK{What am I missing here? $[N]=\{1,\ldots,N\},$ and when $k=\floor{\log_2N},$ then $2^{k+1}>2^{\log_2N}=N,$ so why do we have the remaining terms? Also, for $\floor{\log_2 n}$ to equal $N,$ we have to have $n=2^N,$ which never happens.} \BK{We know that Cesaro convergence implies dyadic convergence, not vice versa. So if we split $[N]$ into dyadic intervals, there is a tail that we have to account for. In general, the averages along $[N]$, $[2^{\sfloor{\log_2 N}}]$, $[2^{\sfloor{\log_2 N}+1}]$ need not be the same, so we can't look purely at dyadic averages.} 
        \begin{equation*}
        \frac{1}{N}\brac{\sum_{k=0}^{\floor{\log_2 N}-1} \sum_{n\in[2^k,2^{k+1})}T_{1}^n f_1\cdot T_{2}^{n+k}f_2 + \sum_{n\in[2^\floor{\log_2 N}, N]}T_{1}^n f_1\cdot T_{2}^{n+\floor{\log_2 N}}f_2}\to 0
    \end{equation*}in $L^2(\mu)$.
By the triangle inequality and the bound $N-2^\floor{\log_2(N)}\leq N$, the $L^2(\mu)$ norm of the previous expression is smaller than \begin{equation*}
    \frac{1}{N}\sum_{k=0}^{\floor{\log_2 N}-1}2^k \norm{\E_{n\in[2^k,2^{k+1})} T_{1}^n f_1\cdot T_{2}^{n+k}f_2}_{L^2(\mu)} + \norm{\E_{n\in[2^\floor{\log_2 N}, N]}T_{1}^n f_1\cdot T_{2}^{n+\floor{\log_2 N}}f_2}_{L^2(\mu)}.
\end{equation*}
    By applying the Cauchy-Schwarz inequality, we infer that the square of the first expression above is bounded by
    \begin{multline*}
        \frac{1}{N^2}\left(\sum_{k=0}^{\floor{\log_2 N}-1}2^k \right)\left(\sum_{k=0}^{\floor{\log_2 N}-1} 2^k\norm{\E_{n\in[2^k,2^{k+1})} T_{1}^n f_1\cdot T_{2}^{n+k}f_2}_{L^2(\mu)}^2\right)\\
        \ll \frac{1}{N}\sum_{k=0}^{\floor{\log_2 N}-1} 2^k\norm{\E_{n\in[2^k,2^{k+1})} T_{1}^n f_1\cdot T_{2}^{n+k}f_2}_{L^2(\mu)}^2.
    \end{multline*}
    Let $S:=T_{1}^{-1}T_{2}$.
    By the van der Corput inequality, this is at most 
    %\BK{(below, where does $\log_2 N$ come from in the error term? I think that for each fixed $k$, we get $O(H/N)$, so summing over $k$, we get $\ll\frac{1}{N}\sum_{k\leq \floor{\log_2 N}}2^k H/N \ll H/N$?)}    
    \begin{multline*}
%\frac{1}{N^2}\left(\sum_{k\leq \floor{\log_2 N}}2^k \right)\left(\sum_{k\leq \floor{\log_2 N}} 2^k\norm{\E_{n\in[2^k,2^{k+1})} T^n(f_1\cdot T^{k}f_2)}_{L^2(\mu)}^2\right)\ll\\
       \ll\frac{1}{N}\sum_{k=0}^{\floor{\log_2 N}-1} 2^k\E_{h\in[H]}\left| \E_{n\in[2^k,2^{k+1})}\int(\overline{f_1}\cdot T_1^h f_1)\cdot S^n T_2^k(\overline{f_2}\cdot T_2^{h} f_2) \, d\mu\right|+ {\frac{1}{H}+\frac{H\log_2 N}{N}}.
         %=\E_{h\in[H]} \left(\frac{1}{N}\sum_{k\leq \floor{\log_2 N}}2^k \left|\int (\overline{f_1}\cdot T_1^h f_1)\cdot S^k(\overline{f_2}\cdot T_2^{h} f_2) \, d\mu\right|\right)+{\frac{1}{H}+\frac{H\log_2 N}{N}}.
    \end{multline*}
    % Since the system is mixing, for any fixed $h\in H$ we have
    % \begin{equation*}
    %     \lim\limits_{k\to\infty} \int (\overline{f_1}\cdot T^h f_1)\cdot (T^k\overline{f_2}\cdot T^{k+h} f_2) \, d\mu=\int \overline{f_1}\cdot T^h f_1\, d\mu \int \overline{f_2}\cdot T^h f_2\, d\mu,
    % \end{equation*}
    % and taking $N\to\infty$, we conclude that the desired expression is bounded by 
    % \begin{multline*}
    % \limsup_{N\to\infty}\brac{\frac{1}{N}\sum_{k\leq \floor{\log_2 N}}2^k \norm{\E_{n\in[2^k,2^{k+1})} T^n(f_1\cdot T^{k}f_2)}_{L^2(\mu)}}^2\\
    % \ll     \E_{h\in[H]}\left|\int \overline{f_1}\cdot T^h f_1\, d\mu \int \overline{f_2}\cdot T^h f_2\, d\mu\right|+\frac{1}{H}.
    % \end{multline*}
    % Taking $H\to\infty$, we conclude that the last quantity tends to zero, since the system is weak-mixing (mixing implies weak-mixing) and one of the functions $f_1,f_2$ has zero integral. Arguing analogously, we deduce that
    % \begin{align*}
    %     \lim_{N\to\infty}\norm{\E_{n\in[\floor{\log_2 N}, N]}T^n(f_1\cdot T^{k}f_2)}_{L^2(\mu)} = 0
    % \end{align*}
    % under the assumption that one of $f_1,f_2$ has zero integral.
    % \BK{STOP}




    
% By the triangle inequality, the $L^2$-norm of the previous expression is smaller than \begin{equation*}
%     \frac{1}{N}\sum_{k\leq \floor{\log_2 N}}2^k \norm{\frac{1}{2^k}\sum_{2^k\leq n< 2^{k+1}}T_{1}^n f_1\cdot T_{2}^{n+k}f_2}_{L^2(\mu)}.
% \end{equation*}
% Let $S:=T_{1}^{-1}T_{2}$.
%     By squaring, Cauchy-Schwarzing and van der Corputing this, we infer that the square of the expression above is bounded by \begin{multline*}
% \frac{1}{N^2}\left(\sum_{k\leq \floor{\log_2 N}}2^k \right)\left(\sum_{k\leq \floor{\log_2 N}} 2^k\norm{\frac{1}{2^k}\sum_{2^k\leq n< 2^{k+1}}T_{1}^n f_1\cdot T_{2}^{n+k}f_2}_{L^2(\mu)}^2\right)\ll\\
%        \frac{1}{N}\sum_{k\leq \floor{\log_2 N}} \frac{2^k}{H}\sum_{h\leq H}\left| \frac{1}{2^k} \sum_{2^k\leq n< 2^{k+1}}\int(\overline{f_1}\cdot T_{1}^h f_1)\cdot S^{n}(T_{2}^k\overline{f_2}\cdot T_{2}^{k+h} f_2) \, d\mu\right|+O(H^{-1})+O(H\log_2 N/N)=\\
%          \frac{1}{H}\sum_{h\leq H} \left(\frac{1}{N}\sum_{k\leq \floor{\log_2 N}}2^k \left|\frac{1}{2^k} \sum_{2^k\leq n< 2^{k+1}}\int (\overline{f_1}\cdot T_{1}^h f_1)\cdot S^{n}(T_{2}^k\overline{f_2}\cdot T_{2}^{k+h} f_2) \, d\mu\right|\right)+O(H^{-1})+O(H\log_2 N/N)=
%          \\
%          \frac{1}{H}\sum_{h\leq H} \left(\frac{1}{N}\sum_{k\leq \floor{\log_2 N}}2^k \left|\frac{1}{2^k} \sum_{2^k\leq n< 2^{k+1}}\int T_{2}^{-k}(\overline{f_1}\cdot T_{1}^h f_1)\cdot S^{n}(\overline{f_2}\cdot T_{2}^{h} f_2) \, d\mu\right|\right)+O(H^{-1})+O(H\log_2 N/N).
%     \end{multline*}    

Note that for each fixed $h$ and $\e>0$, if $k$ is sufficiently large depending only on $h,\e$ and $f_{2}, T_{2}, S$, then 
$$\Bigl\Vert \E_{n\in[2^k,2^{k+1})}\Bigbrac{S^{n}(\overline{f_2}\cdot T_{2}^{h} f_2)-\E(\overline{f_2}\cdot T_{2}^{h} f_2\vert \CI(S))}\Bigr\Vert_{2}<\e$$
by the mean ergodic theorem for F\o lner sequences. So the previous expression can be bounded by 
\begin{multline*}
 \ll\E_{h\in[H]} \left(\frac{1}{N}\sum_{k=0}^{\floor{\log_2 N}-1}2^k \left|\int T_{2}^{-k}(\overline{f_1}\cdot T_{1}^h f_1)\cdot \E(\overline{f_2}\cdot T_{2}^{h} f_2\vert \CI(S)) \, d\mu\right|\right)
         \\+{\frac{1}{H}+\frac{H\log_2 N}{N}}+o_{H;N\to\infty}(1),
    \end{multline*}
    Since $T_{2}\vert_{\CI(S)}$ is mixing by assumption, for any fixed $h\in [H]$ we have
    \begin{multline*}
       \lim_{N\to\infty}\frac{1}{N}\sum_{k=0}^{\floor{\log_2 N}-1}2^k\abs{\int T_{2}^{-k}(\overline{f_1}\cdot T_{1}^h f_1)\cdot \E(\overline{f_2}\cdot T_{2}^{h} f_2\vert \CI(S)) \, d\mu}\\
       =\abs{\int \overline{f_1}\cdot T_{1}^h f_1\, d\mu \cdot\int \overline{f_2}\cdot T_{2}^h f_2\, d\mu},
    \end{multline*} and hence
    %and taking $N\to\infty$, we conclude that the desired expression is bounded by 
    \begin{equation*}
         \abs{\frac{1}{N}\sum_{k=0}^{\floor{\log_2 N}-1} \sum_{n\in[2^k,2^{k+1})}T_{1}^n f_1\cdot T_{2}^{n+k}f_2}^2\ll\E_{h\in[H]}\left|\int \overline{f_1}\cdot T_{1}^h f_1\, d\mu \cdot\int \overline{f_2}\cdot T_{2}^h f_2\, d\mu\right|+\frac{1}{H}.
    \end{equation*}
    The square of this is bounded by
    \begin{equation*}
         \ll\E_{h\in[H]}\left|\int \overline{f_j}\cdot T_{j}^h f_j\, d\mu\right|^{2}+\frac{1}{H}
    \end{equation*}
    for any $j=1,2$.
       Taking $H\to\infty$, we conclude that the last quantity vanishes whenever $\nnorm{f_{j}}_{T_{j},T_{j}}$. An identical proof allows us to control $$\limsup\limits_{N\to\infty}\norm{\E_{n\in[2^\floor{\log_2 N}, N]}T_{1}^n f_1\cdot T_{2}^{n+\floor{\log_2 N}}f_2}_{L^2(\mu)}$$ by degree-2 Host-Kra seminorms.    
\end{proof}

\begin{proof}[Proof of Theorem \ref{T: counterexample}]
    By Lemma \ref{L: slowly growing}, $(T^{\floor{\log_2 n}})_n$ is ergodic if and only if $T$ is conjugate to a one-point system, so it remains to show that $(T^n, T^{n+\floor{\log_2 n}})_n$ is jointly ergodic if and only if $T$ is mixing.  Suppose first that $(T^n, T^{n+\floor{\log_2 n}})_n$ is jointly ergodic. Then in particular
    \begin{align*}
        \lim_{N\to\infty}\int \E_{n\in[N]}T^n f_1\cdot T^{n+\floor{\log_2 n}}f_2\; d\mu = \lim_{N\to\infty}\int \E_{n\in[N]} f_1\cdot T^{\floor{\log_2 n}}f_2\; d\mu,
    \end{align*}
    hence $(T^{\floor{\log_2 n}})_n$ is weakly ergodic. It follows from Lemma \ref{L: slowly growing} that $T$ is mixing. 

    Now we assume that $T$ is mixing. Then Proposition \ref{est12} implies that $(T^n, T^{n+\floor{\log_2 n}})_n$ is good for seminorm estimates. Since $T$ has no nontrivial eigenvalues, $(T^n, T^{n+\floor{\log_2 n}})_n$ is also good for equidistribution. So, by Theorem \ref{T: joint ergodicity criteria}, we have that $(T^n, T^{n+\floor{\log_2 n}})_n$ is jointly ergodic.
\end{proof}


\iffalse

Before we prove Theorem \ref{T: counterexample}, yielding a negative answer to Problem \ref{Pr: joint ergodicity problem}, we need a weak-limit variation of Lemma \ref{L: slowly growing}. We thank Nikos Frantzikinakis for showing this argument to us. In what follows, we say that a sequence of transformations $(T_n)_n$ is \textit{weakly ergodic} on $(X, \CX, \mu)$ if for any $f_1,f_2\in L^\infty(\mu)$, we have
\begin{align*}
    \lim_{N\to\infty}\int f_1\cdot \E_{n\in[N]}T^n f_2\; d\mu = \int f_1\; d\mu\cdot\int f_2\; d\mu.
\end{align*}
\begin{lemma}[Weak ergodicity of slowly growing sequences]\label{L: slowly growing weak}
    Let $a\in\CH$ with $1\prec a\ll \log$, and let $(X, \CX, \mu, T)$ be a system. Then $(T^{\floor{a(n)}})_n$ is weakly ergodic for $(X, \CX, \mu)$ if and only if $T$ is mixing.
\end{lemma}
\begin{proof}
Suppose first that $(T^{\floor{a(n)}})_n$ is weakly ergodic. Arguing as in the proof of Lemma \ref{L: slowly growing}, we deduce that
\begin{align*}
    \lim_{N\to\infty}\int f_1\cdot \E_{n\in[a\inv(N), a\inv(N+1))}T^{\floor{a(n)}}f_2\; d\mu = \int f_1\; d\mu\cdot \int f_2\; d\mu
\end{align*}
for all $f_1,f_2\in L^\infty(\mu)$. By noticing that $\floor{a(n)} = N$ for $n\in [a\inv(N), a\inv(N+1))$, we infer that
\begin{align*}
    \lim_{N\to\infty}\int f_1\cdot T^N f_2\; d\mu = \int f_1\; d\mu\cdot \int f_2\; d\mu,
\end{align*}
hence $T$ is mixing. The converse implication is immediate. 
\end{proof}


We are now in the position to prove Theorem \ref{T: counterexample}. 
\begin{proof}[Proof of Theorem \ref{T: counterexample}]
    The statement that $(T^{\floor{\log_2 n}})_n$ is ergodic if and only if the system is conjugate to a one-point system follows directly from Lemma \ref{L: slowly growing}, so we instead prove that $(T^n, T^{n+\floor{\log_2 n}})_n$ is jointly ergodic if and only if $T$ is mixing. Suppose first that $(T^n, T^{n+\floor{\log_2 n}})_n$ is jointly ergodic. Then in particular
    \begin{align*}
        \lim_{N\to\infty}\int \E_{n\in[N]}T^n f_1\cdot T^{n+\floor{\log_2 n}}f_2\; d\mu = \lim_{N\to\infty}\int \E_{n\in[N]} f_1\cdot T^{\floor{\log_2 n}}f_2\; d\mu,
    \end{align*}
    hence $(T^{\floor{\log_2 n}})_n$ is weakly ergodic. It follows from Lemma \ref{L: slowly growing} that $T$ is mixing. 
    
    Conversely, suppose that $T$ is mixing. It suffices to show that for any $f_1,f_2\in L^{\infty}(\mu)$, at least one of which has a zero integral,  the sequence \begin{equation*}
        \E_{n\in[N]} T^n f_1\cdot  T^{n+\floor{\log_2 n}}f_2 
    \end{equation*}
    converges to zero in $L^2(\mu)$.
    Dyadically decomposing the sum, it suffices to show that \begin{equation*}
        \frac{1}{N}\brac{\sum_{k=0}^{\floor{\log_2 N}-1} \sum_{n\in[2^k,2^{k+1})} + \sum_{n\in[\floor{\log_2 N}, N]}}T^n(f_1\cdot T^{k}f_2)\to 0
    \end{equation*}in $L^2(\mu)$.

By the triangle inequality and the bound $N-\floor{\log_2(N)}\leq N$, the $L^2(\mu)$ norm of the previous expression is smaller than \begin{equation*}
    \frac{1}{N}\sum_{k=0}^{\floor{\log_2 N}-1}2^k \norm{\E_{n\in[2^k,2^{k+1})} T^n(f_1\cdot T^{k}f_2)}_{L^2(\mu)} + \norm{\E_{n\in[\floor{\log_2 N}, N]}T^n(f_1\cdot T^{k}f_2)}_{L^2(\mu)}
\end{equation*}
    By applying the Cauchy-Schwarz inequality, we infer that the square of the first expression above is bounded by
    \begin{multline*}
        \frac{1}{N^2}\left(\sum_{k=0}^{\floor{\log_2 N}-1}2^k \right)\left(\sum_{k=0}^{\floor{\log_2 N}-1} 2^k\norm{\E_{n\in[2^k,2^{k+1})} T^n(f_1\cdot T^{k}f_2)}_{L^2(\mu)}^2\right)\\
        \ll \frac{1}{N}\sum_{k=0}^{\floor{\log_2 N}-1} 2^k\norm{\E_{n\in[2^k,2^{k+1})} T^n(f_1\cdot T^{k}f_2)}_{L^2(\mu)}^2.
    \end{multline*}
    By the van der Corput inequality, this is at most 
    %\BK{(below, where does $\log_2 N$ come from in the error term? I think that for each fixed $k$, we get $O(H/N)$, so summing over $k$, we get $\ll\frac{1}{N}\sum_{k=0}^{\floor{\log_2 N}-1}2^k H/N \ll H/N$?)}    
    \begin{multline*}
%\frac{1}{N^2}\left(\sum_{k=0}^{\floor{\log_2 N}-1}2^k \right)\left(\sum_{k=0}^{\floor{\log_2 N}-1} 2^k\norm{\E_{n\in[2^k,2^{k+1})} T^n(f_1\cdot T^{k}f_2)}_{L^2(\mu)}^2\right)\ll\\
       \ll\frac{1}{N}\sum_{k=0}^{\floor{\log_2 N}-1} 2^k\E_{h\in[H]}\left| \E_{n\in[2^k,2^{k+1})}\int(\overline{f_1}\cdot T^h f_1)\cdot (T^k\overline{f_2}\cdot T^{k+h} f_2) \, d\mu\right|+ {\frac{1}{H}+\frac{H\log_2 N}{N}}\\
         =\E_{h\in[H]} \left(\frac{1}{N}\sum_{k=0}^{\floor{\log_2 N}-1}2^k \left|\int (\overline{f_1}\cdot T^h f_1)\cdot (T^k\overline{f_2}\cdot T^{k+h} f_2) \, d\mu\right|\right)+{\frac{1}{H}+\frac{H\log_2 N}{N}}.
    \end{multline*}
    Since the system is mixing, for any fixed $h\in H$ we have
    \begin{equation*}
        \lim\limits_{k\to\infty} \int (\overline{f_1}\cdot T^h f_1)\cdot (T^k\overline{f_2}\cdot T^{k+h} f_2) \, d\mu=\int \overline{f_1}\cdot T^h f_1\, d\mu \int \overline{f_2}\cdot T^h f_2\, d\mu,
    \end{equation*}
    and taking $N\to\infty$, we conclude that the desired expression is bounded by 
    \begin{multline*}
    \limsup_{N\to\infty}\brac{\frac{1}{N}\sum_{k=0}^{\floor{\log_2 N}-1}2^k \norm{\E_{n\in[2^k,2^{k+1})} T^n(f_1\cdot T^{k}f_2)}_{L^2(\mu)}}^2\\
    \ll     \E_{h\in[H]}\left|\int \overline{f_1}\cdot T^h f_1\, d\mu \int \overline{f_2}\cdot T^h f_2\, d\mu\right|+\frac{1}{H}.
    \end{multline*}
    Taking $H\to\infty$, we conclude that the last quantity tends to zero, since the system is weak-mixing (mixing implies weak-mixing) and one of the functions $f_1,f_2$ has zero integral. Arguing analogously, we deduce that
    \begin{align*}
        \lim_{N\to\infty}\norm{\E_{n\in[\floor{\log_2 N}, N]}T^n(f_1\cdot T^{k}f_2)}_{L^2(\mu)} = 0
    \end{align*}
    under the assumption that one of $f_1,f_2$ has zero integral.
    \end{proof}


% {\color{red} 

% We know that the difference ergodicity condition holds for $(T^{a(n)},T^{a(n)+\floor{\log n}})$ only if the system is trivial. So to prove the necessary direction, we need to show that the joint ergodicity of $(T^{a(n)},T^{a(n)+\floor{\log n}})$ implies that the system is trivial.

% However, when $a(n)=n$, we still don't know how to do it according to the discussion on page the discussion on page 60. (although for $a(n)=\log n$ this implies that  the system is trivial by the discussion on page 61)
% }

% \KT{I think this equivalence is false. Here is a claim that mixing implies that the pair above is jointly ergodic.
% \begin{lemma}
%     Let $(X,\mu,T)$ be a mixing system. Then, the sequence $(T^n, T^{n+\floor{\log_2 n}})$ is jointly ergodic.
% \end{lemma}
% \begin{proof}
%     It suffices to show that for any $f_1,f_2\in L^{\infty}(\mu)$ such that at least one of them has integral equal to zero, we have that the sequence \begin{equation*}
%         \frac{1}{N}\sum_{ n\leq N} T^n f_1\cdot  T^{n+\floor{\log_2 n}}f_2 
%     \end{equation*}converges to zero in $L^2(\mu)$.
%     Dyadically decomposing the sum, it suffices to show that \begin{equation*}
%         \frac{1}{N}\sum_{k\leq \floor{\log_2 N}} \sum_{2^k\leq n< 2^{k+1}}T^n(f_1\cdot T^{k}f_2)\to 0
%     \end{equation*}in $L^2(\mu)$.

%     \BK{(Technically, we also need to include the interval from $2^{\floor{\log_2 N}}$ to $N$)}

% By the triangle inequality, the $L^2$-norm of the previous expression is smaller than \begin{equation*}
%     \frac{1}{N}\sum_{k\leq \floor{\log_2 N}}2^k \norm{\frac{1}{2^k}\sum_{2^k\leq n< 2^{k+1}}T^n(f_1\cdot T^{k}f_2)}_{L^2(\mu)}.
% \end{equation*}
%     By squaring, Cauchy-Schwarzing and van der Corputing this, we infer that the square of the expression above is bounded by \begin{multline*}
% \frac{1}{N^2}\left(\sum_{k\leq \floor{\log_2 N}}2^k \right)\left(\sum_{k\leq \floor{\log_2 N}} 2^k\norm{\frac{1}{2^k}\sum_{2^k\leq n< 2^{k+1}}T^n(f_1\cdot T^{k}f_2)}_{L^2(\mu)}^2\right)\ll\\
%        \frac{1}{N}\sum_{k\leq \floor{\log_2 N}} \frac{2^k}{H}\sum_{h\leq H}\left| \frac{1}{2^k} \sum_{2^k\leq n< 2^{k+1}}\int(\overline{f_1}\cdot T^h f_1)\cdot (T^k\overline{f_2}\cdot T^{k+h} f_2) \, d\mu\right|+O(H^{-1})+O(H\log_2 N/N)=\\
%          \frac{1}{H}\sum_{h\leq H} \left(\frac{1}{N}\sum_{k\leq \floor{\log_2 N}}2^k \left|\int (\overline{f_1}\cdot T^h f_1)\cdot (T^k\overline{f_2}\cdot T^{k+h} f_2) \, d\mu\right|\right)+O(H^{-1})+O(H\log_2 N/N)
%     \end{multline*}Since the system is mixing, we have that for any fixed $h\leq H$, \begin{equation*}
%         \lim\limits_{k\to+\infty} \int (\overline{f_1}\cdot T^h f_1)\cdot (T^k\overline{f_2}\cdot T^{k+h} f_2) \, d\mu=\int \overline{f_1}\cdot T^h f_1\, d\mu \int \overline{f_2}\cdot T^h f_2\, d\mu
%     \end{equation*}and taking $N\to+\infty$, we conclude that the desired expression is bounded by \begin{equation*}
%          \frac{1}{H}\sum_{h\leq H}\left|\int \overline{f_1}\cdot T^h f_1\, d\mu \int \overline{f_2}\cdot T^h f_2\, d\mu\right|+O(H^{-1}).
%     \end{equation*}Taking $H\to+\infty$, we conclude that the last quantity tends to zero, since the system is weak-mixing (mixing implies weak-mixing) and one of the functions $f_1,f_2$ has integral zero.
%     \end{proof}
% }


\fi

\section{Joint ergodicity for pathological Hardy sequences}\label{section: repaired conjecture}

Given Corollary \ref{C:np}, it is natural to ask what can be said for sequences $a_{1},\dots,a_{\ell}\in\mathcal{H}$ for which there exist $i\neq j$, $c_i, c_j\in\R$, and $p\in\Q[x]$ satisfying $1\prec c_i a_i - c_j a_j-p\ll \log$. For convenience, we say that such a family of sequences is \emph{pathological}. Unlike the non-pathological case, the $L^{2}(\mu)$-limit 
$$\lim_{N\to\infty}\E_{n\in[N]}\prod_{j=1}^\ell T_j^{\floor{a_{j}(n)}}f_j$$
may not even exist (for example, by Lemma \ref{L: slowly growing}, $\E\limits_{n\in[N]}T^{\floor{\log_{2}(n)}}f$ converges in $L^2(\mu)$ if and only if $T^{n}f$ converges in $L^2(\mu)$, which happens precisely when $(X, \CX, \mu, T)$ is conjugate to a one-point system).

We propose here a couple of ways to amend the joint ergodicity problem to account for pathological sequences.

\subsection{The weak joint ergodicity problem}
We start with some definitions. 

%For convenience, we say that a sequence of transformations $(T_n)_n$ on a Lebesgue probability space $(X, \CX, \mu)$ is  \textit{strongly/weakly} \emph{ergodic} (abbreviated as \emph{SE/WE}) if  for all $f\in L^\infty(\mu)$, we have that
%\[  \lim_{N\to\infty}\E_{n\in[N]} T_n f = \int f\, d\mu\]     
%    in the strong/weak sense. 




\begin{definition}
      Let $a_1, \ldots, a_\ell:\N\to\R$ be sequences and $(X, \CX, \mu, T_1, \ldots, T_\ell)$ be a system. We say that $a_1, \ldots, a_\ell$ are 
      \begin{enumerate}
      \item \textit{strongly/weakly jointly ergodic} (or \emph{SJ/WJ}) for $(X, \CX, \mu, T_1, \ldots, T_\ell)$  if for all $f_1, \ldots, f_\ell\in L^\infty(\mu)$, we have that
\begin{align}\label{E: joint ergodicity2}
        \lim_{N\to\infty}\E_{n\in[N]}\prod_{j=1}^\ell T_j^{\floor{a_{j}(n)}}f_j=\prod_{j=1}^\ell \int f_j\, d\mu
    \end{align}
    in the strong/weak sense.
    \item   \emph{strongly/weakly difference ergodic} (or \emph{SD/WD}) if $(T_i^{\floor{a_i(n)}}T_j^{-\floor{a_j(n)}})_n$ is strongly/ weakly ergodic on $(X, \CX, \mu)$ for all distinct $i, j\in[\ell]$.
     \item   \emph{strongly/weakly product ergodic} (or \emph{SP/WP}) if $(T_1^{\floor{a_1(n)}}\times \cdots \times T_\ell^{\floor{a_\ell(n)}})_n$ is strongly/weakly ergodic on the product system $(X^\ell, \CX^{\otimes \ell}, \mu^\ell)$.
      \end{enumerate}
\end{definition}

We observe that all the strong notions above correspond to the notions of ergodicity, joint ergodicity, as well as difference and product ergodicity conditions defined in the introduction.

%Let $P$ denote ``jointly ergodic", ``difference ergodic" or ``product ergodic".
Let $Y$ denote $D,J$ or $P$.
It is clear that $SY$ implies $WY$. Conversely, if the relevant limit exists in the strong sense, then $SY$ is equivalent to $WY$.


%For example, Lemma \ref{L: slowly growing}  imply the following.

%\begin{proposition}\label{9.1}
%    Let $(X,\mathcal{X},\mu,T)$ be a system. The sequence $(T^{\floor{\log_{2}n}})_{n}$ is 
%    \begin{enumerate}
 %       \item $SE$ if and only if the systems is conjugate to a one-point system;
%        \item $WE$ if and only if $T$ is mixing.
%    \end{enumerate}
%\end{proposition}

\iffalse

It is clear that Theorem \ref{T: counterexample} follows from Propositions \ref{L: slowly growing} and \ref{9.2}.

\begin{proposition}\label{9.2}
    Let $(X,\mathcal{X},\mu,T)$ be a system. Then $(T^{n},T^{n+\floor{\log_{2}n}})_{n}$ is $SJ$ if and only if $T$ is mixing.
\end{proposition}
\begin{proof}
   Suppose first that $(T^n, T^{n+\floor{\log_2 n}})_n$ is jointly ergodic. Then in particular
    \begin{align*}
        \lim_{N\to\infty}\int \E_{n\in[N]}T^n f_1\cdot T^{n+\floor{\log_2 n}}f_2\; d\mu = \lim_{N\to\infty}\int \E_{n\in[N]} f_1\cdot T^{\floor{\log_2 n}}f_2\; d\mu,
    \end{align*}
    hence $(T^{\floor{\log_2 n}})_n$ is weakly ergodic. It follows from Lemma \ref{L: slowly growing} that $T$ is mixing. 
    
    Conversely, suppose that $T$ is mixing. It suffices to show that for any $f_1,f_2\in L^{\infty}(\mu)$, at least one of which has a zero integral,  the sequence \begin{equation*}
        \E_{n\in[N]} T^n f_1\cdot  T^{n+\floor{\log_2 n}}f_2 
    \end{equation*}
    converges to zero in $L^2(\mu)$.
    Dyadically decomposing the sum, it suffices to show that \begin{equation*}
        \frac{1}{N}\brac{\sum_{k=0}^{\floor{\log_2 N}-1} \sum_{n\in[2^k,2^{k+1})} + \sum_{n\in[\floor{\log_2 N}, N]}}T^n(f_1\cdot T^{k}f_2)\to 0
    \end{equation*}in $L^2(\mu)$.

By the triangle inequality and the bound $N-\floor{\log_2(N)}\leq N$, the $L^2(\mu)$ norm of the previous expression is smaller than \begin{equation*}
    \frac{1}{N}\sum_{k\leq \floor{\log_2 N}}2^k \norm{\E_{n\in[2^k,2^{k+1})} T^n(f_1\cdot T^{k}f_2)}_{L^2(\mu)} + \norm{\E_{n\in[\floor{\log_2 N}, N]}T^n(f_1\cdot T^{k}f_2)}_{L^2(\mu)}
\end{equation*}
    By applying the Cauchy-Schwarz inequality, we infer that the square of the first expression above is bounded by
    \begin{multline*}
        \frac{1}{N^2}\left(\sum_{k\leq \floor{\log_2 N}}2^k \right)\left(\sum_{k\leq \floor{\log_2 N}} 2^k\norm{\E_{n\in[2^k,2^{k+1})} T^n(f_1\cdot T^{k}f_2)}_{L^2(\mu)}^2\right)\\
        \ll \frac{1}{N}\sum_{k\leq \floor{\log_2 N}} 2^k\norm{\E_{n\in[2^k,2^{k+1})} T^n(f_1\cdot T^{k}f_2)}_{L^2(\mu)}^2.
    \end{multline*}
    By the van der Corput inequality, this is at most 
    %\BK{(below, where does $\log_2 N$ come from in the error term? I think that for each fixed $k$, we get $O(H/N)$, so summing over $k$, we get $\ll\frac{1}{N}\sum_{k\leq \floor{\log_2 N}}2^k H/N \ll H/N$?)}    
    \begin{multline*}
%\frac{1}{N^2}\left(\sum_{k\leq \floor{\log_2 N}}2^k \right)\left(\sum_{k\leq \floor{\log_2 N}} 2^k\norm{\E_{n\in[2^k,2^{k+1})} T^n(f_1\cdot T^{k}f_2)}_{L^2(\mu)}^2\right)\ll\\
       \ll\frac{1}{N}\sum_{k\leq \floor{\log_2 N}} 2^k\E_{h\in[H]}\left| \E_{n\in[2^k,2^{k+1})}\int(\overline{f_1}\cdot T^h f_1)\cdot (T^k\overline{f_2}\cdot T^{k+h} f_2) \, d\mu\right|+ {\frac{1}{H}+\frac{H\log_2 N}{N}}\\
         =\E_{h\in[H]} \left(\frac{1}{N}\sum_{k\leq \floor{\log_2 N}}2^k \left|\int (\overline{f_1}\cdot T^h f_1)\cdot (T^k\overline{f_2}\cdot T^{k+h} f_2) \, d\mu\right|\right)+{\frac{1}{H}+\frac{H\log_2 N}{N}}.
    \end{multline*}
    Since the system is mixing, for any fixed $h\in H$ we have
    \begin{equation*}
        \lim\limits_{k\to\infty} \int (\overline{f_1}\cdot T^h f_1)\cdot (T^k\overline{f_2}\cdot T^{k+h} f_2) \, d\mu=\int \overline{f_1}\cdot T^h f_1\, d\mu \int \overline{f_2}\cdot T^h f_2\, d\mu,
    \end{equation*}
    and taking $N\to\infty$, we conclude that the desired expression is bounded by 
    \begin{multline*}
    \limsup_{N\to\infty}\brac{\frac{1}{N}\sum_{k\leq \floor{\log_2 N}}2^k \norm{\E_{n\in[2^k,2^{k+1})} T^n(f_1\cdot T^{k}f_2)}_{L^2(\mu)}}^2\\
    \ll     \E_{h\in[H]}\left|\int \overline{f_1}\cdot T^h f_1\, d\mu \int \overline{f_2}\cdot T^h f_2\, d\mu\right|+\frac{1}{H}.
    \end{multline*}
    Taking $H\to\infty$, we conclude that the last quantity tends to zero, since the system is weak-mixing (mixing implies weak-mixing) and one of the functions $f_1,f_2$ has zero integral. Arguing analogously, we deduce that
    \begin{align*}
        \lim_{N\to\infty}\norm{\E_{n\in[\floor{\log_2 N}, N]}T^n(f_1\cdot T^{k}f_2)}_{L^2(\mu)} = 0
    \end{align*}
    under the assumption that one of $f_1,f_2$ has zero integral.
\end{proof}


\fi



We summarize some basic properties:

\begin{proposition}\label{all}
Let $a_{1},\dots,a_{\ell}\in\mathcal{H}$.
    \begin{enumerate}
    \item For any system, $SD+SP\Rightarrow SJ$;
    \item For any system,  $SJ\Rightarrow SP$;
    \item For any system,  $WJ\Rightarrow WD$;
    \item For some systems,  $SJ\not\Rightarrow SD$.
    \end{enumerate}
\end{proposition}
\begin{proof}
Part (i) follows from Theorem \ref{T: main theorem}. Part (ii) follows from Theorem \ref{T: easy direction}. Part (iv) follows from Theorem \ref{T: counterexample}. To see Part (iii), take $f_{i}=f, f_{j}=g$ and all other functions to be 1. Then $WJ$ implies that 
    $$\int f\; d\mu\cdot\int g\; d\mu=\lim_{N\to\infty}\E_{n\in[N]}\int T_{i}^{a_{i}(n)}f\cdot T_{j}^{a_{j}(n)}g\; d\mu=\lim_{N\to\infty}\E_{n\in[N]}\int T_{i}^{a_{i}(n)}T_{j}^{-a_{j}(n)}f\cdot g\; d\mu,$$
    which implies $WD$.
\end{proof}

Since the jointly ergodic conjecture ($SD+SP\Longleftrightarrow SJ$) fails, it is natural to ask the following relevant questions.

\begin{question}\label{c2} 
    Let $a_{1},\dots,a_{\ell}\in\mathcal{H}$. Is it true that $WD+WP\Leftrightarrow WJ$ for any system?
\end{question}

\begin{question}\label{c3}  
    Let $a_{1},\dots,a_{\ell}\in\mathcal{H}$. Is it true that $WD+SP\Leftrightarrow SJ$ for any system?
\end{question}


%We remark that question \ref{c2} can be very difficult, as it implies that 2-mixing implies 3-mixing (see the following proposition).
We remark that Question \ref{c2} can be very difficult, as the positive answer to it would imply that 2-mixing implies a property reminiscent of 3-mixing (see the following proposition).

\begin{proposition}
     Let $(X,\mathcal{X},\mu,T)$ be a system. The pair $(T^{\floor{\log_{2}n}},T^{2\floor{\log_{2}n}})_{n}$ is 
    \begin{enumerate}
        \item $WJ$ if and only if $T$ has the following property: for all $f_0, f_1, f_2\in L^\infty(\mu)$, we have $$\lim_{n\to\infty}\int f_0\cdot T^n f_1\cdot T^{2n}f_2\; d\mu = \int f_0\; d\mu \cdot \int f_1\; d\mu \cdot \int f_2\; d\mu;$$
        \item $WD$ {and $WP$} if and only if $T$ is (2-)mixing;
        \item $SJ, SD$ or $SP$ if and only if $T$ is conjugate to a one-point system.
    \end{enumerate}
\end{proposition}
\begin{proof}
 The proof of Part (i) is similar to that of Part (ii) of Lemma \ref{L: slowly growing}, and so we leave it to the interested readers.
    Parts (ii) and (iii) follow from Lemma \ref{L: slowly growing}.
\end{proof}

To make Question \ref{c3} more plausible, we address it for a tuple for which the (strong) joint ergodicity problem fails by Theorem \ref{T: counterexample}. 

\begin{proposition}\label{last}
    For any system $(X,\mathcal{X},\mu,T_{1},T_{2})$ and for the pair $(T_{1}^{n},T_{2}^{n+\floor{\log_{2}n}})_{n}$, we have that $WD+SP\Leftrightarrow SJ$.
\end{proposition}

%{\color{red} I think we don't need to include the three term case in this paper}

\begin{proof}
    It follows from Proposition \ref{all} that $SJ\Rightarrow WD+SP$. So now we show the opposite direction. Note that $SP$ implies that $(T_{1}^{n},T_{2}^{n+\floor{\log_{2}n}})_{n}$ is good for equidistribution. By Theorem \ref{T: joint ergodicity criteria}, it suffices to show that $(T_{1}^{n},T_{2}^{n+\floor{\log_{2}n}})_{n}$ is good for seminorm control.
    
     Since $(T_{1}^{n},T_{2}^{n+\floor{\log_{2}n}})_{n}$ is WD, for any $\CI(T_1T_2\inv)$-measurable $f_0, f_1$ we have $$\int f_0 \cdot \E_{n\in[N]} T_2^{\floor{\log n}} f_1\to\int f_0\; d\mu \int f_1\; d\mu.$$ 
This means that $(T_{2}^{\floor{\log_{2}n}})_{n}$ is weakly ergodic for the system $(X,\mathcal{I}(T_1T_2\inv),\mu,T_{2})$. By Lemma \ref{L: slowly growing}, $T_{2}\vert_{\mathcal{I}(T_1T_2\inv)}$ is mixing, and by Proposition \ref{est12}, $(T_{1}^{n},T_{2}^{n+\floor{\log_{2}n}})_{n}$ is good for seminorm control.
\end{proof}







 \subsection{Changing the averaging scheme}
Another approach to revising the joint ergodicity classification problem for Hardy sequences (involving slowly growing functions) is to consider a different averaging scheme. This point of view was adopted in \cite{BMR20}, where some multiple recurrence results were obtained that cannot be established through the study of the corresponding (Ces\`{a}ro) ergodic averages. For instance, it can be easily proven that for any ergodic system $(X,\CX,\mu, T)$, the logarithmic averages $$\frac{1}{\log N}\sum_{n\leq N} \frac{T^{\floor{\log n}}f}{n}$$ 
converge to $\int f\; d\mu$ for all bounded measurable functions $f$. The main input in this result is the spectral theorem and the fact that $a\log n$ is equidistributed with respect to logarithmic averaging for any $a\in (0,1)$. In this sense, logarithmic averaging is more natural for sequences involving $\log$, and so we can ``repair'' Theorem \ref{T: counterexample} by considering logarithmic averages. We call a sequence $a:\N\to \Z$ \emph{ergodic with respect to logarithmic averages} on a system $(X,\CX, \mu,T)$  if \begin{equation*}
    \lim_{N\to\infty} \norm{\frac{1}{\log N}\sum_{n\in [N]} \frac{T^{a(n)}f}{n}-\int f\; d\mu }_{L^2(\mu)}=0.
\end{equation*}We define joint ergodicity of sequences with respect to logarithmic averages similarly to Definition \ref{D: joint ergodicity} by replacing Ces\`{a}ro averages with logarithmic ones. 



\begin{conjecture}\label{Conj: log averaged counterexample}
    For any system $(X,\CX, \mu,T_1,T_2)$, the tuple $(T_1^n,T_2^{n+\floor{\log n}})_{n}$ is jointly ergodic for logarithmic averages if and only if $(T_2^{n+\floor{\log n}}T_1^{-n})_{n}$ and $(T_1^{n}\times T_2^{n+\floor{\log n}})_{n}$ are ergodic for logarithmic averages on $(X, \CX, \mu)$ and $(X\times X, \CX\otimes \CX, \mu\times\mu)$ respectively.
\end{conjecture}
The conjecture is true whenever $T_1=T_2$. Indeed, the difference condition is equivalent to $T$ being ergodic, and, in this case, we can easily see that the product condition holds as well. 
Additionally, the pair we study is good for seminorm control for a single transformation due to the results in \cite{BMR20}. Furthermore, it can be shown that they are good for equidistribution (for logarithmic averages). Using a variant of Theorem \ref{T: joint ergodicity criteria} for logarithmic averages \cite[Theorem 2.19]{FrKu22a}, we conclude that $(T^n, T^{n+\floor{\log n}})_{n}$ is jointly ergodic for all ergodic systems. This argument makes Conjecture \ref{Conj: log averaged counterexample} plausible.

We remark, though, that logarithmic averaging is not enough to establish joint ergodicity characterization when we deal with even slower functions, like $\log\log x$. In this example, one would need to consider an even weaker averaging scheme, namely, double logarithmic averaging. For this reason, we have to modify this line of reasoning.

We formalize the prior discussion.
Let $W(n)$ be a non-decreasing sequence of real numbers satisfying $W(n)\to\infty$ and denote $w(n)=W(n+1)-W(n)$. We define the $W$-averaging operators of a sequence $a_n$ on a normed space by the formula\begin{equation*}
   A^W_{N}\left(\left(a(n)\right)_{n}\right) :=\frac{1}{W(N)}\sum_{n=1}^N w(n)a_n.
\end{equation*}This is well-defined for $N$ sufficiently large.

Let $a_1,\ldots,a_\ell:\N\to \R$ be sequences and let $(X,\CX,\mu,T_1,\ldots, T_\ell)$ be a system.
We call $(T_1^{a_1(n)},\ldots,T_\ell^{a_\ell(n)})$  {\em $W$-jointly ergodic} for $(X,\CX,\mu,T_1,\ldots, T_\ell)$ if for all $f_1,\ldots, f_\ell\in L^{\infty}(\mu)$, we have 
\begin{equation*}
    \lim_{N\to\infty} \norm{A^W_{N}(T_1^{a(n)}f_1\ldots T_k^{a_k(n)}f_k)-\int f_1\; d\mu\cdots\int f_\ell\; d\mu}_{L^2(\mu)}=0.
\end{equation*}
%If $\ell=1$, then we call the sequence $T^{a_1(n)}$ {\em $W$-ergodic} for $(X,\CX, \mu)$. More generally, 
Moreover, we say that a sequence of transformations $(T_n)_n$ is {\em $W$-ergodic} on $(X,\CX,\mu)$ if for all $f\in L^{\infty}(\mu)$, we have
\begin{equation*}
    \lim_{N\to\infty} \norm{A^W_{N}(T_n f)-\int f\; d\mu }_{L^2(\mu)}=0.
\end{equation*}
This brings us to the following possible modification of Problem \ref{Pr: joint ergodicity problem}.

\begin{conjecture}\label{C: weighted conjecture}
    Let $a_1,\ldots, a_\ell\in \mathcal{H}$ with $a_j\succ 1$ for every $j\in[\ell]$. Then there exists a non-decreasing sequence $(W(n))_n$ with $1\prec W(x)\ll x$
    %tending to $\infty$ as $n\to\infty$, 
    such that for any system $(X,\CX, \mu,T_1,\ldots, T_\ell)$, the following are equivalent:\\
\begin{enumerate}
    \item $(T_1^{\floor{a_1(n)}},\ldots, T_\ell^{\floor{a_\ell(n)}})_n$ is $W$-jointly ergodic on $(X,\CX,\mu)$;\\
    \item the sequence $(T_1^{\floor{a_1(n)}}\times\dots\times T_\ell^{\floor{a_\ell(n)}})_n$ is $W$-ergodic on $(X^\ell, \CX^{\otimes \ell},\mu^{\ell})$ while the sequence $(T_i^{\floor{a_i(n)}}T_j^{-\floor{a_j(n)}})_n$ is $W$-ergodic on $(X,\CX, \mu)$ for all distinct $i, j\in[\ell]$.
    \end{enumerate}
\end{conjecture}
    %The most promising method to attack this conjecture should resemble the approach used in the case of integer polynomials. \BK{(Meaning? Why integer polynomials?)} 
    A major missing ingredient needed to resolve Conjecture \ref{C: weighted conjecture} would be seminorm estimates for ergodic averages along Hardy sequences with suitable weights. To obtain such estimates, one should try to merge the argument of the authors for usual (unweighted) ergodic averages along Hardy sequences together with the method of Bergelson, Moreira and Richter \cite{BMR20}.
    %The bulk of the work for resolving such a conjecture would focus on establishing seminorm control for weighted ergodic averages with Hardy iterates and suitably chosen weights. This is known to be the case for $\Z$-actions through the arguments in \cite{BMR20} and it is reasonable to expect that their PET induction argument can be adapted to the $\Z^\ell$ setting. Of course, we anticipate that the multidimensional case will require additional input, such as concatenation theorems and smoothing arguments.

{One potential shortcoming of Conjecture \ref{C: weighted conjecture} vs. the original problem has to do with the choice of a suitable averaging scheme. While such schemes can always be found for Hardy sequences of polynomial growth, it is not clear what would be the ``correct'' averaging scheme if we go beyond Hardy sequences and consider e.g. sequences with oscillations, such as $x\sin{x}$.}

%{\color{red} I think we need to require that $w$ is non-increasing, as otherwise we can choose $w(n)=2^{n}$, in which case the W-erogodicity of $T_{i}^{[a_{i}(n)]}$ (which is implied by both the W-joint and W-product ergodicity conditions) is equivalent of saying that $\lim_{n\to\infty}T_{i}^{[a_{i}(n)]}f=\int f$, which I think means that $(X,T_{i})$ is trivial.}

\medskip


\subsection{More general sequences of polynomial growth} While Corollary \ref{C:np} asserts that Problem \ref{Pr: joint ergodicity problem} admits an affirmative answer for all non-pathological Hardy sequences of polynomial growth, we believe that this result should extend to a more general class of iterates. In this section, we conjecture a natural generalization of Corollary \ref{C:np}, which will involve tempered functions defined below.

\begin{definition}[Tempered functions]
    Let $i\in\mathbb{N}_{0}$ and $x_0\geq 0$. A real-valued function $a\in C^{i+1}(x_0, \infty)$ is called a \emph{tempered function of degree $i$} if the following hold:
\begin{enumerate}
\item $a^{(i+1)}(x)$ tends monotonically to $0$ as $x\to\infty;$
\item  $\lim\limits_{x\to\infty}x|a^{(i+1)}(x)|=\infty.$
\end{enumerate} 
\end{definition}
 

A significant hurdle arising in the study of tempered functions is that in contrast to Hardy field functions, ratios of tempered functions may not converge. In order to avoid various problematic cases, one typically restricts attention to a well-behaved subclass of tempered functions. On letting 
\begin{align*}
    \mathcal{R} &:=\Big\{a\in C^\infty(\mathbb{R}_+):\;\lim_{x\to\infty}\frac{xa^{(i+1)}(x)}{a^{(i)}(x)}\in \mathbb{R}\;\;\text{for all}\;\;i\in\mathbb{N}_{0}\Big\},\\
    \mathcal{T}_i&:=\Big\{a\in\mathcal{R}:\;\exists\;i<\alpha\leq i+1,\;\lim_{x\to\infty}\frac{xa'(x)}{t(x)}=\alpha,\;\lim_{x\to\infty}a^{(i+1)}(x)=0\Big\},
\end{align*}
the aforementioned well-behaved class is given by $\mathcal{T}:=\bigcup_{i=0}^\infty \mathcal{T}_i$. (See \cite{BeHK09, Kouts21} for a detailed discussion of tempered functions.)

It is known that every function $a\in \mathcal{T}_i$ is a tempered function of degree $i$ and satisfies the growth condition $x^i\log x\prec a(x)\prec x^{i+1}$ (see \cite{BeHK09}).

In \cite{DKS23}, the first, second and fourth authors showed that Problem~\ref{Pr: joint ergodicity problem} holds whenever $$a_1=\ldots=a_\ell=a=c_1h+c_2t$$ for some $(c_1,c_2)\in \mathbb{R}^2\setminus\{(0,0)\}$ as well as a Hardy function $h\in \CH$ and a tempered function $t\in \mathcal{T}$ satisfying some growth conditions. 
For example, it holds for the function 
\[x^{\pi}/ \log x+x^{1/2}(2+\cos\sqrt{\log x})\] which is of polynomial growth but is not Hardy.

Although dealing with tempered functions is outside the scope of this article, we expect Problem~\ref{Pr: joint ergodicity problem} to hold in the following general case.

\begin{conjecture}
Let $a_1, \ldots, a_\ell:\N\to\R$ be sequences of the form 
\[a_j=c_{1,j} h_j+c_{2,j}t_j\] for some $(c_{1,j},c_{2,j})\in \mathbb{R}^2\setminus\{(0,0)\},$ $h_j\in \CH,$ and $t_j\in \mathcal{T}$ such that for all distinct $i, j\in[\ell]$ and $c_i, c_j\in\R$, the function $c_i a_i - c_j a_j$ is either an almost rational polynomial or it stays logarithmically away from $\Q[x]$. Then the sequences $a_1, \ldots, a_\ell$ are jointly ergodic for a system $(X, \CX, \mu, T_1, \ldots, T_\ell)$ if and only  if they satisfy the product and difference ergodicity conditions on the system. 
\end{conjecture}














  \bibliography{library}
\bibliographystyle{plain}
\end{document}
